\documentclass[11pt]{article}

\usepackage[a4paper,margin=1in]{geometry}
\usepackage[T1]{fontenc}
\usepackage{lmodern}
\usepackage{microtype}
\usepackage{amsmath,amssymb,amsthm,mathtools}
\usepackage{aliascnt}
\usepackage{bm}
\usepackage{enumitem}
\usepackage{booktabs}
\usepackage{array}
\usepackage{url}
\usepackage[hidelinks]{hyperref}
\usepackage[nameinlink,noabbrev]{cleveref}
\usepackage[backend=biber,style=alphabetic,maxbibnames=20]{biblatex}
\addbibresource{references.bib}

\newcommand{\Fq}{\mathbb{F}_q}
\newcommand{\F}{\mathbb{F}}
\newcommand{\N}{\mathbb{N}}
\newcommand{\B}{\{0,1\}}
\newcommand{\M}{\mathcal{M}}
\newcommand{\RS}{\operatorname{RS}}
\newcommand{\Enc}{\operatorname{Enc}}
\newcommand{\ev}{\operatorname{ev}}
\newcommand{\eq}{\operatorname{eq}}
\newcommand{\wt}{\operatorname{wt}}
\newcommand{\supp}{\operatorname{supp}}
\newcommand{\rev}{\star}
\newcommand{\ip}[2]{\left\langle #1,#2\right\rangle}
\newcommand{\coeff}[2]{[X^{#1}]\,#2}
\newcommand{\polylt}[1]{\F[X]_{<#1}}
\newcommand{\polyltq}[1]{\Fq[X]_{<#1}}
\newcommand{\had}{\circ}
\newcommand{\defeq}{\coloneqq}
\newcommand{\bits}{\{0,1\}^n}

\theoremstyle{plain}
\newtheorem{theorem}{Theorem}[section]

\newaliascnt{lemma}{theorem}
\newtheorem{lemma}[lemma]{Lemma}
\aliascntresetthe{lemma}

\newaliascnt{proposition}{theorem}
\newtheorem{proposition}[proposition]{Proposition}
\aliascntresetthe{proposition}

\newaliascnt{corollary}{theorem}
\newtheorem{corollary}[corollary]{Corollary}
\aliascntresetthe{corollary}

\theoremstyle{definition}
\newaliascnt{definition}{theorem}
\newtheorem{definition}[definition]{Definition}
\aliascntresetthe{definition}

\newaliascnt{remark}{theorem}
\newtheorem{remark}[remark]{Remark}
\aliascntresetthe{remark}

\newaliascnt{example}{theorem}
\newtheorem{example}[example]{Example}
\aliascntresetthe{example}

\crefname{theorem}{Theorem}{Theorems}
\crefname{lemma}{Lemma}{Lemmas}
\crefname{proposition}{Proposition}{Propositions}
\crefname{corollary}{Corollary}{Corollaries}
\crefname{definition}{Definition}{Definitions}
\crefname{remark}{Remark}{Remarks}
\crefname{example}{Example}{Examples}


\title{\textbf{Coefficient-Extraction Arguments for Transparent Post-Quantum Proof Systems}\\
\large From Kronecker Encodings to Inner Products, Hadamard Relations, Lookups, and R1CS}
\author{Ali Mkhida\\Algorizk Labs\\\texttt{ali.mkhida@algorizk.xyz}}
\date{Working manuscript, September 2026}

\begin{document}
\maketitle

\begin{abstract}
We develop a coefficient-extraction framework for transparent arguments over multilinear tables.  Its basic identity is the middle-coefficient formula
\[
  \sum_{j=0}^{N-1}a_jb_j
  =[X^{N-1}]\Bigl(\sum_{j=0}^{N-1}a_jX^j\Bigr)
                  \Bigl(\sum_{j=0}^{N-1}b_jX^{N-1-j}\Bigr),
\]
combined with a split identity that turns the selected coefficient into local relations among degree-$<N$ Reed--Solomon words.  On inversion-closed multiplicative domains, reversal and the high part of the split are virtual words; DEEP quotients similarly bind out-of-domain values without separate commitments.  We use these facts to give self-contained protocols for public linear functionals, committed inner products, and pointwise Hadamard products, and prove that heterogeneous instances of these relations can be combined before a single FRI-style folding test.  Logarithmic-derivative lookup and multiset arguments reduce to reciprocal Hadamard relations and linear/inner-product claims; indexed selection then yields sparse matrix--vector multiplication, from which we obtain R1CS and offline memory-consistency arguments.  The interactive soundness analysis is information theoretic in the unique-decoding regime, apart from a quoted correlated-agreement theorem for Reed--Solomon codes.  We further prove the relaxed round-by-round decoding-knowledge property needed by the salted BCS compiler of Chiesa, Di, Hu, and Zheng, obtaining post-quantum straight-line knowledge soundness in the quantum random-oracle model for the compiled non-interactive framework.  The construction does not claim zero knowledge or performance superiority over optimized Sumcheck systems; we give symbolic cost accounting and isolate list decoding, higher-rate identities, and implementation fusion as the principal remaining efficiency questions.
\end{abstract}

\tableofcontents

\section{Introduction}
\label{sec:introduction}

\subsection{From Sumcheck to coefficient extraction}

The Sumcheck protocol of Lund, Fortnow, Karloff, and Nisan is one of the basic algebraic reductions underlying modern succinct proof systems~\cite{LFKN92}.  Given a low-degree polynomial in several variables, Sumcheck reduces a claim about its sum over the Boolean hypercube to a single evaluation at a verifier-chosen random point.  In the multilinear setting the protocol is especially natural: a polynomial in $n$ variables is represented by $N=2^n$ field elements, each round eliminates one variable, and the prover performs a recursive sequence of folds.  This structure is used directly in general-purpose systems such as Spartan~\cite{Setty20} and HyperPlonk~\cite{HyperPlonk23}, in lookup-oriented systems such as Lasso and Jolt~\cite{Lasso23,Jolt23}, and in recent zkVM-oriented memory arguments~\cite{TwistShout26}.  It is also tightly coupled to several multilinear polynomial-commitment constructions based on error-correcting codes, including BaseFold and DeepFold~\cite{BaseFold24,DeepFold25}.

The cryptographic role of Sumcheck should be stated carefully.  Sumcheck itself is an information-theoretic interactive protocol; it does not rely on elliptic curves, pairings, discrete logarithms, or any other number-theoretic assumption that is known to be vulnerable to quantum computers.  The post-quantum issue arises in the surrounding commitment and compilation layers.  A multilinear polynomial IOP may be compiled with an algebraic polynomial commitment such as KZG, or with a transparent code-based commitment and a hash-based oracle compiler.  The latter route is attractive for post-quantum arguments because Reed--Solomon proximity tests such as FRI use field arithmetic, error-correcting codes, and hash commitments rather than pairings~\cite{FRI18,WHIR24}.  Existing hash-based multilinear systems therefore combine two different mechanisms: Sumcheck carries the multilinear algebraic relation, while folding/proximity testing certifies low degree and commitment consistency.

This paper asks whether that division of labor is always necessary.  Our starting point is a classical coefficient identity.  Let $n\ge 1$, let $N=2^n$, and for vectors
\[
  a=(a_0,\ldots,a_{N-1}),\qquad
  b=(b_0,\ldots,b_{N-1})\in \F^N
\]
define the univariate generating polynomials
\[
  V_a(X)=\sum_{j=0}^{N-1}a_jX^j,
  \qquad
  V_b^{\rev}(X)=\sum_{j=0}^{N-1}b_jX^{N-1-j}.
\]
Then
\begin{equation}
  \ip{a}{b}
  =\sum_{j=0}^{N-1}a_jb_j
  =\coeff{N-1}{V_a(X)V_b^{\rev}(X)}.
  \label{eq:intro-ip-coeff}
\end{equation}
The identity is elementary.  Its relevance to proof systems comes from the observation that a coefficient claim can be turned into a low-degree polynomial identity with exactly the degree bound naturally enforced by a Reed--Solomon code.  Namely, for $U,K\in\polylt{N}$ and $v\in\F$,
\begin{equation}
  v=\coeff{N-1}{U(X)K(X)}
  \quad\Longleftrightarrow\quad
  XU(X)K(X)=A(X)+vX^N+X^{N+1}H(X)
  \label{eq:intro-master-identity}
\end{equation}
for unique polynomials $A,H\in\polylt{N}$ obtained by splitting the product around degree $N$.  The factor $X$ moves the selected coefficient from degree $N-1$ to degree $N$; consequently both auxiliary polynomials retain the power-of-two degree bound $N$.  If $U$, $A$, and the necessary transforms of $K$ are represented as Reed--Solomon words, the value of $H$ at a queried domain point is determined locally by~\eqref{eq:intro-master-identity}.  Thus $H$ need not be committed separately: it is a virtual word.  After a random linear combination of the participating words, one FRI-style proximity test can simultaneously enforce their low-degree structure and the coefficient identity.

The resulting viewpoint is different from a Sumcheck protocol.  Sumcheck proves a hypercube sum by recursively eliminating variables,
\[
  n\text{-variate claim}
  \longrightarrow
  (n-1)\text{-variate claim}
  \longrightarrow\cdots\longrightarrow
  \text{evaluation at a random point}.
\]
Coefficient extraction instead transforms the relation into a static identity among univariate low-degree objects,
\[
  \text{algebraic relation}
  \longrightarrow
  \text{selected coefficient of a product}
  \longrightarrow
  \text{Reed--Solomon proximity}.
\]
The point is not that every use of Sumcheck should be replaced.  Sumcheck is exceptionally efficient, often has a linear-time prover, and remains the right abstraction in many systems.  The question is structural: which relations normally expressed through Sumcheck already have coefficient-extraction forms that can be checked directly by the same FRI machinery used by the commitment layer?  For those relations, the multilinear reduction and the proximity argument need not remain separate protocols.

\subsection{Kronecker encodings and succinct kernels}

Our notation follows the multilinear-polynomial convention used in~\cite{MkhidaIguider26}.  We write $\Fq$ for a finite field, $n\in\N$ for the number of variables, $N=2^n$, and
\[
  \M\subseteq \Fq[x_1,\ldots,x_n]
\]
for the $N$-dimensional $\Fq$-vector space of multilinear polynomials.  For $j\in[0,N-1]$ with binary decomposition
\[
  j=\sum_{i=1}^{n}j_i2^{i-1},\qquad j_i\in\{0,1\},
\]
the canonical monomial is
\[
  m_j=\prod_{i=1}^{n}x_i^{j_i}.
\]
When a multilinear polynomial is represented by its table on $\B^n$, we identify the bit vector $(j_1,\ldots,j_n)$ with the corresponding integer $j$ and write
\begin{equation}
  V_f(X)=\sum_{j=0}^{N-1}f(j)X^j.
  \label{eq:intro-table-encoding}
\end{equation}
This is simply a Kronecker encoding of the length-$N$ table.  It turns table entries into coefficients of one univariate polynomial without a M\"obius transform.

The first important special case of~\eqref{eq:intro-ip-coeff} is multilinear evaluation.  Let $\widetilde f$ be the multilinear extension of a table $f:\B^n\to\F$.  For $z=(z_1,\ldots,z_n)\in\F^n$, define
\begin{equation}
  E_z^{\rev}(X)
  =\prod_{i=1}^{n}\left(z_i+(1-z_i)X^{2^{i-1}}\right).
  \label{eq:intro-eval-kernel}
\end{equation}
The coefficient of $X^{N-1-j}$ in $E_z^{\rev}$ is the Boolean equality polynomial $\eq(j,z)$, and therefore
\begin{equation}
  \widetilde f(z)=\coeff{N-1}{V_f(X)E_z^{\rev}(X)}.
  \label{eq:intro-evaluation-extraction}
\end{equation}
The verifier never materializes the $N$ coefficients of $E_z^{\rev}$: at a point $\xi$, the product in~\eqref{eq:intro-eval-kernel} is evaluated in $O(n)$ field operations.  This motivates a useful abstraction.  A family of vectors $p_\theta\in\F^N$ is a \emph{succinct kernel family} if its reversed generating polynomial
\[
  K_\theta(X)=\sum_{j=0}^{N-1}(p_\theta)_jX^{N-1-j}
\]
can be evaluated at verifier query points in $\operatorname{polylog}(N)$ field operations.  Every such family gives a public linear functional
\[
  \ip{f}{p_\theta}=\coeff{N-1}{V_fK_\theta}
\]
that can be plugged into the same coefficient-extraction protocol.  Multilinear evaluation is the tensor-structured example~\eqref{eq:intro-eval-kernel}; hypercube summation corresponds to the all-ones kernel; selectors, sparse public functionals, and several wiring predicates give further examples.

The second special case is more consequential.  If both vectors are committed, no public kernel is available.  Nevertheless the reversal
\[
  V_b^{\rev}(X)=X^{N-1}V_b(X^{-1})
\]
can be read virtually from the same Reed--Solomon commitment whenever the evaluation domain $L\le \Fq^{\times}$ is a multiplicative subgroup.  Indeed $L^{-1}=L$, and at $\xi\in L$ the verifier obtains
\begin{equation}
  V_b^{\rev}(\xi)=\xi^{N-1}V_b(\xi^{-1}).
  \label{eq:intro-virtual-reversal}
\end{equation}
Thus an arbitrary inner product of two committed tables is reduced to~\eqref{eq:intro-master-identity} without a second commitment to the reversed polynomial.  This is the first indication that the Kronecker formula is not merely an alternative multilinear PCS opening: it is a general proof primitive for bilinear relations.

\subsection{From bilinear sums to pointwise arithmetic}

A canonical degree-two Sumcheck relation is
\begin{equation}
  S=\sum_{j=0}^{N-1}a_jb_j.
  \label{eq:intro-degree-two-sum}
\end{equation}
By~\eqref{eq:intro-ip-coeff}, this is exactly one coefficient-extraction claim.  Consequently the usual $n$-round degree-two Sumcheck can, at least algebraically, be replaced by one committed opening polynomial, one virtual quotient polynomial, and one Reed--Solomon proximity test.  This already covers many terminal reductions in GKR-, Spartan-, and lookup-style protocols.

More general pointwise arithmetic requires an additional ingredient.  Suppose three committed tables satisfy
\begin{equation}
  a\had b=c,
  \qquad\text{i.e.}\qquad
  a_jb_j=c_j\quad\text{for every }j\in[0,N-1].
  \label{eq:intro-hadamard}
\end{equation}
Choose a random $\gamma\in\F$ and consider
\[
  D(\gamma)=\sum_{j=0}^{N-1}\gamma^j(a_jb_j-c_j).
\]
If~\eqref{eq:intro-hadamard} is false then $D$ is a nonzero polynomial of degree at most $N-1$, so a random $\gamma$ hides the error with probability at most $(N-1)/|\F|$.  The weighted product term is again a coefficient:
\begin{equation}
  \sum_{j=0}^{N-1}\gamma^ja_jb_j
  =\coeff{N-1}{V_a(\gamma X)V_b^{\rev}(X)},
  \label{eq:intro-weighted-hadamard}
\end{equation}
while the weighted linear term is $V_c(\gamma)$.  The values $V_a(\gamma\xi)$ and $V_c(\gamma)$ lie outside the original Reed--Solomon domain in general.  We bind them back to the committed words with DEEP quotient relations, represented themselves as virtual low-degree words.  The resulting Hadamard argument therefore uses the same ingredients as the inner-product protocol: coefficient extraction, virtual transformations, random batching, and one FRI-style proximity test.

This pointwise multiplication primitive changes the scope of the framework.  Let
\[
  G(Y_1,\ldots,Y_t)
\]
be an arithmetic expression of bounded size, and let $f_1,\ldots,f_t$ be committed tables.  Instead of applying Sumcheck directly to
\[
  \sum_{j=0}^{N-1}G(f_1(j),\ldots,f_t(j)),
\]
one may materialize the intermediate wire tables of the pointwise arithmetic circuit.  Additions and scalar multiplications are linear relations; multiplication gates are Hadamard relations; the final hypercube sum is an inner product with the all-ones vector.  Thus every fixed arithmetic circuit over table coordinates reduces to a collection of linear functionals and Hadamard checks.  This does not automatically imply a faster prover---materializing intermediate tables may lose the excellent constants and streaming structure of optimized Sumcheck---but it shows that Sumcheck is not algebraically necessary for this class of relations.  The remaining question becomes whether the resulting low-degree objects can be batched and encoded efficiently enough to compete with a dedicated Sumcheck prover.

\subsection{Lookup, permutation, sparse algebra, and R1CS}

The same two primitives, inner product and Hadamard product, also capture relations that initially appear unrelated to coefficient extraction.

\paragraph{Lookup through logarithmic derivatives.}
Let $f=(f_i)$ be a lookup vector, let $t=(t_j)$ be a public or committed table, and let $m_j$ denote the claimed multiplicity of $t_j$ among the entries of $f$.  After sampling a challenge $\beta$ away from the relevant poles, define
\[
  q_i=(\beta-f_i)^{-1},\qquad h_j=(\beta-t_j)^{-1}.
\]
The logarithmic-derivative identity underlying modern lookup arguments is
\begin{equation}
  \sum_i q_i=\sum_jm_jh_j.
  \label{eq:intro-lookup-logderivative}
\end{equation}
The reciprocal definitions are pointwise products,
\[
  q\had(\beta\mathbf 1-f)=\mathbf 1,
  \qquad
  h\had(\beta\mathbf 1-t)=\mathbf 1,
\]
and~\eqref{eq:intro-lookup-logderivative} is an equality of two inner products.  Therefore logarithmic-derivative lookup reduces directly to Hadamard checks plus coefficient-extraction inner products.  In particular, the FRI backend need not receive an intermediate Sumcheck claim merely to aggregate the reciprocal table.

\paragraph{Permutation and multiset equality.}
For two committed vectors $a,b\in\F^N$, multiset equality can be tested by the same logarithmic-derivative principle.  With
\[
  q_i=(\beta-a_i)^{-1},\qquad r_i=(\beta-b_i)^{-1},
\]
we prove the two reciprocal Hadamard relations and the linear equality
\[
  \ip{q}{\mathbf 1}=\ip{r}{\mathbf 1}.
\]
Tagged values then turn multiset equality into a permutation/wiring check.  This gives a route complementary to the grand-product constructions common in Plonkish systems and to the Sumcheck-based permutation checks used in HyperPlonk~\cite{HyperPlonk23}.

\paragraph{Sparse matrix--vector products.}
For R1CS, the central remaining linear operation is multiplication by a sparse public matrix.  Algebraically, if $y=Mz$, a random row-compression challenge $r$ gives
\[
  \widetilde y(r)=\ip{p_r}{z},
  \qquad
  p_r(j)=\sum_iM_{ij}\eq(i,r),
\]
and a second random point can bind $p_r$ to the multilinear extension of $M$.  This form is useful conceptually but does not by itself preserve sparsity: succinctly authenticating a large sparse matrix is precisely one of the roles played by specialized sparse-polynomial arguments in Spartan-like systems.  We therefore also develop a sparse-entry construction.  Write the nonzero entries as triples
\[
  (\mathsf{row}_e,\mathsf{col}_e,\mathsf{val}_e),
  \qquad e\in[0,K-1],
\]
where $K=\operatorname{nnz}(M)$.  The sparse-entry construction authenticates the selected column values $z(c(e))$ by indexed lookup, authenticates the random row weights $\alpha^{\iota(r(e))}$ by a second indexed selection, forms the entry contributions $\mu(e)z(c(e))$ by a Hadamard relation, and compares one committed inner product with a public geometric linear functional of the claimed output.  Thus no dense $N\times N$ matrix is formed, and the meaningful public metadata remains proportional to the sparse representation rather than $N^2$.

\paragraph{R1CS.}
An R1CS instance asks for $z$ satisfying
\begin{equation}
  (Az)\had(Bz)=Cz.
  \label{eq:intro-r1cs}
\end{equation}
Once sparse matrix--vector multiplication is expressed by the primitives above,~\eqref{eq:intro-r1cs} decomposes into three sparse linear maps
\[
  a=Az,\qquad b=Bz,\qquad c=Cz
\]
and one Hadamard relation
\[
  a\had b=c.
\]
This gives a complete algebraic route from R1CS to coefficient extraction and Reed--Solomon proximity without invoking the classical recursive Sumcheck reduction.  Whether the resulting prover is competitive is an implementation question, not an algebraic one; a principal purpose of the accompanying implementation is to measure this tradeoff under controlled engineering.

\subsection{Why FRI is the natural backend}

Coefficient extraction alone is a rewriting, not a proof system.  The cryptographic value comes from matching the rewriting to the closure properties and locality of Reed--Solomon codes.

First, every object introduced above is either a polynomial of degree less than $N$ or a virtual word that should agree with the evaluation of such a polynomial.  Second, the transformations needed by the verifier are local.  Reversal reads the committed word at $\xi^{-1}$ as in~\eqref{eq:intro-virtual-reversal}; a DEEP quotient reads one committed value and performs one division; the virtual $H$ word in~\eqref{eq:intro-master-identity} is determined pointwise by the other words.  Third, random linear combinations preserve the code.  If the relevant words are valid Reed--Solomon codewords, their random batch is a codeword of the same degree bound.  Conversely, correlated-agreement theorems for Reed--Solomon codes can be used to argue that a sufficiently large set of challenges for which the batch is close to the code forces the individual words to agree with a common low-degree explanation~\cite{BSCI23}.

After this first algebraic/batching step, the remainder is an ordinary proximity problem.  Any suitable FRI-style test can in principle serve as the backend.  The present development uses the same smooth multiplicative domains and even--odd folding structure as Kronecker-FRI, because inversion closure and powers-of-two degree bounds make the virtual transformations particularly clean.  Higher folding arity and Merkle-tree optimizations are orthogonal to the coefficient identities.  More generally, the framework suggests a design rule: use coefficient extraction to convert the proof-system relation into a small family of low-degree words, and then use the strongest available Reed--Solomon proximity machinery to certify the batch.  BaseFold and DeepFold interleave multilinear Sumcheck with foldable-code reductions~\cite{BaseFold24,DeepFold25}; WHIR works directly with constrained Reed--Solomon codes and achieves very low query complexity~\cite{WHIR24}.  Our direction is complementary: simplify the algebraic front end so that a plain low-degree proximity argument can enforce relations that would otherwise be carried by a separate Sumcheck transcript.

This separation is also important for post-quantum claims.  The interactive algebraic protocols proved in this paper are information theoretic.  A non-interactive argument additionally needs a commitment to the oracle words, a transcript/hash construction, and a theorem justifying Fiat--Shamir-style compilation in the quantum random-oracle model.  Hash-based IOP compilation is precisely the setting in which FRI-based systems are commonly viewed as post-quantum candidates~\cite{WHIR24}.  We therefore do not say that coefficient extraction ``makes Sumcheck post-quantum.''  Rather, it removes Sumcheck from a broad algebraic layer and exposes that layer directly to a transparent, hash-based Reed--Solomon backend for which post-quantum compilation can be analyzed as a separate theorem.

\subsection{Contributions and scope}
\label{sec:intro-contributions}

The paper develops the following results from first principles.

\begin{enumerate}[leftmargin=*,itemsep=0.5em]
  \item \textbf{A coefficient-extraction calculus.}
  We give a uniform treatment of reversal, public kernels, table-form multilinear evaluation, hypercube sums, and the master split identity~\eqref{eq:intro-master-identity}.  The central abstraction is a succinct kernel family: a public vector whose reversed generating polynomial can be evaluated locally without materializing the vector.

  \item \textbf{One table-form commitment, several proof primitives.}
  The Kronecker table encoding~\eqref{eq:intro-table-encoding} supports multilinear evaluation, public linear functionals, arbitrary committed inner products, and the subsequent compound relations without changing the underlying Reed--Solomon commitment.

  \item \textbf{A sumcheck-free committed inner-product argument.}
  For two committed vectors we use the virtual reversal~\eqref{eq:intro-virtual-reversal}, one opening polynomial, one virtual quotient word, and one FRI-style proximity test.  The protocol proves~\eqref{eq:intro-degree-two-sum} directly rather than reducing it by $n$ Sumcheck rounds.

  \item \textbf{A sumcheck-free Hadamard argument.}
  We prove pointwise multiplication~\eqref{eq:intro-hadamard} by geometric random compression~\eqref{eq:intro-weighted-hadamard} and DEEP quotient consistency.  This supplies the multiplication gate needed to vectorize arbitrary bounded pointwise arithmetic circuits.

  \item \textbf{Lookup and permutation from the same primitives.}
  We derive logarithmic-derivative lookup and multiset checks from reciprocal Hadamard relations plus inner-product equalities.  Tagged variants yield permutation and wiring checks.  The reductions are proved algebraically, including pole and Schwartz--Zippel error terms.

  \item \textbf{Sparse matrix--vector multiplication and R1CS.}
  We give both a direct random-compression formulation and a sparse-entry construction.  The latter reduces a sparse matrix product to two indexed selections, one Hadamard multiplication, one committed inner product, and one public geometric linear functional.  Combining three such products with one Hadamard relation yields an R1CS argument for~\eqref{eq:intro-r1cs}.

  \item \textbf{Heterogeneous batching into one proximity layer.}
  We formulate the algebraic relations so that evaluation, inner-product, Hadamard, quotient, lookup, and sparse-consistency words can be combined by random challenges before the folding stage.  The soundness analysis isolates the algebraic challenge errors from the final proximity error.

  \item \textbf{Post-quantum non-interactive compilation.}
  We prove the relaxed round-by-round decoding-knowledge property required by the salted BCS compiler for interactive oracle reductions and then invoke the quantum-random-oracle theorem of Chiesa, Di, Hu, and Zheng.  This yields post-quantum straight-line knowledge soundness for the compiled framework under the exact transcript and Merkle layout stated in \cref{sec:post-quantum}.
\end{enumerate}

The scope of the claim is deliberately narrower than ``Sumcheck is unnecessary.''  Sumcheck handles general low-degree hypercube sums directly and with excellent prover complexity.  Our framework targets relations that can be decomposed into linear functionals, inner products, pointwise multiplication, reciprocal constraints, and sparse structured access.  These already cover important proof-system components, but the reduction may introduce auxiliary tables and Reed--Solomon encoding costs.  The decisive systems question is therefore quantitative: after batching and shared encoding, when is the coefficient-extraction route faster, smaller, or easier to accelerate than the corresponding Sumcheck pipeline?  We therefore make no benchmark superiority claim in this manuscript.  The theorem statements and symbolic cost model are organized so that a subsequent Lean 4 formalization and controlled Rust implementation can mirror the paper section by section.

\subsection{Relation to prior work}

The middle-coefficient expression for an inner product is classical.  Pairing-based multilinear commitments also exploit tensor-structured representations; Mercury, for example, uses a related inner-product viewpoint with KZG commitments.  The contribution here is not the elementary coefficient identity by itself, but its systematic use as the algebraic front end of a transparent Reed--Solomon proximity argument, together with virtual reversal and quotient words that avoid extra commitments.

The closest code-based multilinear commitments include BaseFold, DeepFold, WHIR, and the more recent TensorSwitch.  BaseFold introduced a general foldable-code perspective and couples a multilinear evaluation reduction with FRI-inspired folding~\cite{BaseFold24}.  DeepFold moves the analysis into the list-decoding regime to reduce query complexity and proof size~\cite{DeepFold25}.  WHIR constructs a proximity test for constrained Reed--Solomon codes and emphasizes very fast verification~\cite{WHIR24}, while TensorSwitch obtains a hash-based multilinear PCS from tensor codes using list decoding and correlated agreement~\cite{TensorSwitch26}.  These works show that multilinear algebra and code-based proximity can coexist efficiently.  We investigate a different interface between them: rather than carrying a Sumcheck claim through the folding process, we first rewrite the target relation as coefficient and pointwise identities among univariate low-degree objects.

Spartan uses Sumcheck to prove R1CS satisfiability and couples it to polynomial commitments~\cite{Setty20}; the 2026 BinarySpartan instantiation shows that an aggressively optimized transparent, hash-based Spartan architecture can retain Sumcheck while moving to binary fields~\cite{BinarySpartan26}.  HyperPlonk develops Sumcheck-based zero and permutation checks over multilinear polynomials~\cite{HyperPlonk23}.  Lasso and Jolt demonstrate the importance of efficient lookup arguments for zkVMs~\cite{Lasso23,Jolt23}; Celer targets the complementary regime of very large query vectors~\cite{Celer26}, while Twist and Shout further emphasize memory checking and lookup as central prover bottlenecks~\cite{TwistShout26}.  Our lookup, permutation, sparse-matrix, and memory-oriented reductions are motivated by these applications, but the proof mechanism is different: logarithmic derivatives and sparse access are reduced to the inner-product/Hadamard instruction set and then to Reed--Solomon proximity.

Finally, this paper builds on two earlier pieces of our work.  The tree-circuit representation of multilinear polynomials~\cite{MkhidaIguider26} studies the recursive structure and cost of Sumcheck itself; we retain its notation for $\Fq$, $n$, $N=2^n$, the Boolean indexing of monomials, and the vector-space view of multilinear polynomials.  Kronecker-FRI~\cite{KroneckerFRI26} introduced the coefficient-extraction opening identity for multilinear evaluation and showed how to enforce it with one FRI-style proximity test.  The present work changes the emphasis: multilinear evaluation becomes only the first instance of a broader coefficient-extraction calculus whose main objects are linear and bilinear relations among committed tables.

\subsection{Organization}

The remainder of the paper is developed in the following order.  We first fix the finite-field, multilinear, Reed--Solomon, and proximity-test notation and state the correlated-agreement result used in soundness.  We then define the table-form Kronecker commitment and develop the coefficient-extraction calculus: reversal, succinct kernels, multilinear evaluation, and the master split identity.  The next sections introduce virtual reversal and DEEP quotient words, followed by the generic coefficient-extraction protocol and its specializations to public linear functionals, committed inner products, and Hadamard products.  After that we prove heterogeneous batching and develop lookup, multiset/permutation, sparse matrix--vector, R1CS, and memory-checking reductions.  The final theory sections collect soundness and knowledge extraction, treat the non-interactive hash-based compilation and its post-quantum conditions, and analyze complexity.  We conclude with a systematic comparison to Sumcheck-based and other FRI-based architectures, followed by limitations and open problems.

\section{Preliminaries and Notation}
\label{sec:preliminaries}

This section fixes the algebraic and coding-theoretic notation used throughout the paper.  We keep the multilinear-polynomial conventions of~\cite{MkhidaIguider26}: $\Fq$ is a finite field, $n$ is the number of variables, $N=2^n$, $\M$ is the vector space of multilinear polynomials, and $m_j$ denotes the canonical monomial indexed by the binary expansion of $j$.  We additionally distinguish a challenge field $\F\supseteq\Fq$, introduce the smooth multiplicative domains used by the Reed--Solomon backend, and state the one external coding-theoretic theorem used later in the soundness analysis.

Unless explicitly stated otherwise, all polynomial equalities are identities in the indicated polynomial ring, all vector spaces are over the indicated field, and all finite probabilities are taken with respect to the uniform distribution.

\subsection{Integers, bits, and indexing}
\label{sec:preliminaries-bits}

For integers $a\le b$ we write
\[
  [a,b]\defeq\{a,a+1,\ldots,b\}.
\]
If $a>b$, then $[a,b]=\varnothing$.  Let $n\in\N$ with $n\ge1$, and put
\[
  N\defeq 2^n.
\]
Every $j\in[0,N-1]$ has a unique binary expansion
\begin{equation}
  j=\sum_{i=1}^{n} j_i2^{i-1},
  \qquad j_i\in\{0,1\}.
  \label{eq:binary-expansion}
\end{equation}
As in~\cite{MkhidaIguider26}, the least significant bit is $j_1$.  We identify $j$ with its bit vector
\[
  (j_1,\ldots,j_n)\in\B^n
\]
whenever no confusion can arise.  Conversely, for $w=(w_1,\ldots,w_n)\in\B^n$ we use the same symbol $w$ for the integer $\sum_iw_i2^{i-1}$.

For $j\in[0,N-1]$, define its Hamming weight by
\[
  \wt(j)\defeq\sum_{i=1}^{n}j_i,
\]
and its bitwise complement by
\[
  \bar j\defeq (N-1)-j.
\]
For $a,b\in[0,N-1]$, write $a\preceq b$ if $a_i\le b_i$ for every $i\in[1,n]$.

\begin{lemma}[Complement bits]
\label{lem:complement-bits}
For $j\in[0,N-1]$, the binary digits of $\bar j$ are $1-j_i$ for $i\in[1,n]$.
\end{lemma}

\begin{proof}
Since
\[
  N-1=\sum_{i=1}^{n}2^{i-1},
\]
we have
\[
  \bar j=(N-1)-j
  =\sum_{i=1}^{n}(1-j_i)2^{i-1}.
\]
Each coefficient $1-j_i$ lies in $\{0,1\}$, so uniqueness of binary expansion gives the claim.
\end{proof}

\subsection{Finite fields and univariate polynomials}
\label{sec:preliminaries-fields}

Let $\Fq$ be a finite field of odd characteristic $p$, with $q=p^e$ for some integer $e\ge1$.  Let $\F\supseteq\Fq$ be a finite extension field.  We call $\Fq$ the \emph{base field} and $\F$ the \emph{challenge field}.  The extension is allowed to be trivial.  The odd-characteristic assumption ensures that $2$ is invertible, which is needed by the even--odd decomposition used by the folding backend.

For a field $K$, write $K[X]$ for the univariate polynomial ring over $K$ and
\[
  K[X]_{<d}\defeq\{P\in K[X]:\deg P<d\}
\]
for $d\in\N$.  We use the convention $\deg0=-\infty$.  For $P\in K[X]$ and $r\in\N$, the notation
\[
  [X^r]P
\]
denotes the coefficient of $X^r$ in $P$.

\begin{lemma}[Root bound]
\label{lem:root-bound}
Let $K$ be a field and let $P\in K[X]$ be nonzero of degree $d\ge0$.  Then $P$ has at most $d$ roots in $K$.  Consequently, if $P,Q\in K[X]_{<d}$ agree at $d$ distinct points of $K$, then $P=Q$.
\end{lemma}

\begin{proof}
We prove the first statement by induction on $d$.  If $d=0$, then $P$ is a nonzero constant and has no root.  Suppose $d\ge1$ and $a\in K$ is a root.  Euclidean division by $X-a$ gives
\[
  P=(X-a)R
\]
with $\deg R=d-1$.  Every root of $P$ other than $a$ is a root of $R$, because $K$ has no zero divisors.  By induction, $R$ has at most $d-1$ roots, so $P$ has at most $d$ roots.

For the second statement, $P-Q$ has degree $<d$.  If it vanishes at $d$ distinct points, the first statement forces $P-Q=0$.
\end{proof}

\begin{lemma}[Even--odd polynomial decomposition]
\label{lem:evenodd-polynomial}
Let $K$ be a field and $d\ge1$.  Every $P\in K[X]_{<2d}$ admits a unique decomposition
\begin{equation}
  P(X)=P_{\mathrm e}(X^2)+XP_{\mathrm o}(X^2),
  \qquad
  P_{\mathrm e},P_{\mathrm o}\in K[X]_{<d}.
  \label{eq:evenodd-polynomial}
\end{equation}
Explicitly,
\[
  [X^r]P_{\mathrm e}=[X^{2r}]P,
  \qquad
  [X^r]P_{\mathrm o}=[X^{2r+1}]P
\]
for $r\in[0,d-1]$.  The map $P\mapsto(P_{\mathrm e},P_{\mathrm o})$ is $K$-linear.
\end{lemma}

\begin{proof}
Every monomial $X^s$ with $0\le s<2d$ occurs exactly once on the right-hand side of~\eqref{eq:evenodd-polynomial}: if $s=2r$ it occurs in $P_{\mathrm e}(X^2)$, and if $s=2r+1$ it occurs in $XP_{\mathrm o}(X^2)$.  Comparing coefficients gives existence, the displayed formulas, and uniqueness.  Linearity follows coefficientwise.
\end{proof}

\subsection{Multilinear polynomials}
\label{sec:preliminaries-multilinear}

Let
\[
  A\defeq\Fq[x_1,\ldots,x_n]
\]
be the polynomial ring in $n$ indeterminates over $\Fq$.  Following~\cite{MkhidaIguider26}, let $\M\subseteq A$ be the set of multilinear polynomials:
\[
  \M\defeq
  \{f\in A:\deg_{x_i}(f)\le1\text{ for every }i\in[1,n]\}.
\]
It is an $\Fq$-vector space.

For $j\in[0,N-1]$ with binary expansion~\eqref{eq:binary-expansion}, define
\begin{equation}
  m_j\defeq\prod_{i=1}^{n}x_i^{j_i}.
  \label{eq:canonical-monomial}
\end{equation}
Thus $m_0=1$, $m_1=x_1$, $m_2=x_2$, $m_3=x_1x_2$, and
\[
  \deg m_j=\wt(j).
\]

\begin{proposition}[Canonical basis]
\label{prop:canonical-basis}
The family
\[
  m=(m_j)_{0\le j\le N-1}
\]
is a basis of the $\Fq$-vector space $\M$.  In particular, $\dim_{\Fq}\M=N$ and every $f\in\M$ has a unique expansion
\begin{equation}
  f=\sum_{j=0}^{N-1}\alpha_jm_j,
  \qquad \alpha_j\in\Fq.
  \label{eq:canonical-expansion}
\end{equation}
We call $(\alpha_0,\ldots,\alpha_{N-1})$ the \emph{canonical coordinate vector} of $f$.
\end{proposition}

\begin{proof}
The monomials of the full polynomial ring $A$ form an $\Fq$-basis of $A$.  A monomial lies in $\M$ exactly when each exponent is either $0$ or $1$.  Such exponent vectors are precisely the bit vectors $(j_1,\ldots,j_n)\in\B^n$, hence precisely the monomials~\eqref{eq:canonical-monomial}.  Therefore the $m_j$ span $\M$ and are linearly independent.  There are $2^n=N$ of them.
\end{proof}

The table representation of a multilinear polynomial is its restriction to the Boolean hypercube.

\begin{definition}[Table form]
\label{def:table-form}
For $f\in\M$, its \emph{table form} is the vector
\[
  \bigl(f(w)\bigr)_{w\in\B^n}\in\Fq^N,
\]
indexed by the integer identification of $w\in\B^n$ from~\cref{sec:preliminaries-bits}.
\end{definition}

To extend a table from the Boolean cube to arbitrary field points, define for $w\in\B^n$ and $z=(z_1,\ldots,z_n)\in\F^n$ the equality polynomial
\begin{equation}
  \eq(w,z)
  \defeq
  \prod_{i=1}^{n}\bigl((1-w_i)(1-z_i)+w_iz_i\bigr).
  \label{eq:equality-polynomial}
\end{equation}
Equivalently, the $i$th factor is $1-z_i$ when $w_i=0$ and $z_i$ when $w_i=1$.

\begin{lemma}[Boolean delta property]
\label{lem:eq-delta}
For $u,w\in\B^n$,
\[
  \eq(w,u)=
  \begin{cases}
    1,&w=u,\\
    0,&w\ne u.
  \end{cases}
\]
\end{lemma}

\begin{proof}
If $w=u$, then for every coordinate $i$ the $i$th factor in~\eqref{eq:equality-polynomial} equals $1$.  If $w\ne u$, choose $i$ with $w_i\ne u_i$.  If $(w_i,u_i)=(0,1)$ the factor is $1-u_i=0$, while if $(w_i,u_i)=(1,0)$ the factor is $u_i=0$.  Hence the whole product vanishes.
\end{proof}

\begin{definition}[Multilinear extension]
\label{def:mle}
For a table $t:\B^n\to\F$, its \emph{multilinear extension} is
\begin{equation}
  \widetilde t(z)
  \defeq
  \sum_{w\in\B^n}t(w)\eq(w,z),
  \qquad z\in\F^n.
  \label{eq:mle-definition}
\end{equation}
\end{definition}

\begin{proposition}[Interpolation on the Boolean cube]
\label{prop:boolean-interpolation}
For every table $t:\B^n\to\F$, the polynomial $\widetilde t$ defined by~\eqref{eq:mle-definition} is multilinear and satisfies
\[
  \widetilde t(u)=t(u)
  \qquad\text{for every }u\in\B^n.
\]
Moreover, it is the unique multilinear polynomial over $\F$ with these Boolean values.
\end{proposition}

\begin{proof}
Each factor in~\eqref{eq:equality-polynomial} has degree at most $1$ in its corresponding variable and is independent of the other variables; hence $\eq(w,z)$ is multilinear in $z$, and so is the sum~\eqref{eq:mle-definition}.  For $u\in\B^n$, \cref{lem:eq-delta} gives
\[
  \widetilde t(u)
  =\sum_{w\in\B^n}t(w)\eq(w,u)
  =t(u).
\]
For uniqueness, let $g$ be multilinear and vanish on all of $\B^n$.  We prove $g=0$ by induction on $n$.  For $n=1$, write $g=a+bx_1$; the equalities $g(0)=g(1)=0$ imply $a=b=0$.  For $n>1$, write uniquely
\[
  g=g_0+x_ng_1
\]
with $g_0,g_1$ multilinear in $x_1,\ldots,x_{n-1}$.  Evaluating at $x_n=0$ shows that $g_0$ vanishes on $\B^{n-1}$, hence $g_0=0$ by induction.  Evaluating at $x_n=1$ then shows that $g_1$ vanishes on $\B^{n-1}$, hence $g_1=0$.  Thus $g=0$.
\end{proof}

The following relation between canonical coordinates and the Boolean table will occasionally be useful.

\begin{lemma}[Canonical coordinates and table values]
\label{lem:mobius}
Let $f\in\M$ have canonical expansion~\eqref{eq:canonical-expansion}.  For $b\in[0,N-1]$,
\begin{equation}
  f(b)=\sum_{j\preceq b}\alpha_j.
  \label{eq:zeta-transform}
\end{equation}
Conversely, for $j\in[0,N-1]$,
\begin{equation}
  \alpha_j
  =\sum_{c\preceq j}(-1)^{\wt(j)-\wt(c)}f(c).
  \label{eq:mobius-transform}
\end{equation}
\end{lemma}

\begin{proof}
For Boolean $b$, the value
\[
  m_j(b)=\prod_{i:j_i=1}b_i
\]
equals $1$ exactly when $j_i\le b_i$ for all $i$, that is, when $j\preceq b$, and equals $0$ otherwise.  Evaluating~\eqref{eq:canonical-expansion} gives~\eqref{eq:zeta-transform}.

For~\eqref{eq:mobius-transform}, substitute~\eqref{eq:zeta-transform} into its right-hand side:
\[
  \sum_{c\preceq j}(-1)^{\wt(j)-\wt(c)}
  \sum_{d\preceq c}\alpha_d
  =
  \sum_{d\preceq j}\alpha_d
  \sum_{d\preceq c\preceq j}
  (-1)^{\wt(j)-\wt(c)}.
\]
For fixed $d\preceq j$, the inner sum ranges over all subsets of the coordinates where $j_i=1$ and $d_i=0$.  If there are $r=\wt(j)-\wt(d)$ such coordinates, the sum is
\[
  \sum_{s=0}^{r}\binom{r}{s}(-1)^{r-s}=(1-1)^r,
\]
which is $1$ when $d=j$ and $0$ otherwise.  Therefore only the term $\alpha_j$ remains.
\end{proof}

\subsection{Smooth multiplicative domains}
\label{sec:preliminaries-domains}

The Reed--Solomon and folding layers use multiplicative subgroups of power-of-two order.

\begin{definition}[Smooth domain]
\label{def:smooth-domain}
A \emph{smooth domain} is a multiplicative subgroup
\[
  L\le\Fq^\times
\]
of order
\[
  M=2^\mu
\]
for some $\mu\in\N$.  For $r\in[0,\mu]$, define
\[
  L^{2^r}\defeq\{\xi^{2^r}:\xi\in L\}.
\]
The notation denotes the image under the power map, not a Cartesian power.
\end{definition}

We use the standard fact that the multiplicative group of a finite field is cyclic~\cite{LidlNiederreiter97}.  From it we derive every structural property of smooth domains needed below.

\begin{lemma}[Power maps on smooth domains]
\label{lem:power-maps}
Let $L\le\Fq^\times$ be a smooth domain of order $M=2^\mu$.
\begin{enumerate}[label=(\arabic*),leftmargin=*]
  \item For $r\in[0,\mu]$, the map $\xi\mapsto\xi^{2^r}$ is a surjective group homomorphism from $L$ onto $L^{2^r}$, and $|L^{2^r}|=M/2^r$.
  \item Every element of $L^{2^r}$ has exactly $2^r$ preimages in $L$.
  \item $(L^{2^r})^2=L^{2^{r+1}}$ for $r<\mu$.
  \item If $M\ge2$, then $-1\in L$, $-1\ne1$, and the two square roots in $L$ of any $\eta\in L^2$ are $\xi$ and $-\xi$ for either choice of square root $\xi$.
  \item $L$ is closed under inversion: $\xi\in L$ implies $\xi^{-1}\in L$.
\end{enumerate}
\end{lemma}

\begin{proof}
Let $\omega$ generate the cyclic group $L$.  The image of $\xi\mapsto\xi^{2^r}$ is generated by $\omega^{2^r}$, whose order is $2^{\mu-r}=M/2^r$.  This proves (1).  The kernel therefore has cardinality $2^r$, and every fibre is a coset of the kernel, proving (2).  Item (3) follows directly from $(\xi^{2^r})^2=\xi^{2^{r+1}}$.

For (4), the kernel of the squaring map has cardinality $2$ by (2).  Its elements are roots of $X^2-1=(X-1)(X+1)$.  Since the characteristic is odd, $1\ne-1$, so the kernel is $\{1,-1\}$.  Hence $-1\in L$, and the two elements in the fibre over $\xi^2$ are $\xi$ and $-\xi$.

Finally, (5) is true in every subgroup of a group: if $\xi\in L$, its group inverse $\xi^{-1}$ also lies in $L$.
\end{proof}

\begin{definition}[Fibres]
\label{def:fibres}
Assume $|L|\ge2$.  For $\eta\in L^2$, the \emph{fibre over $\eta$} is the two-element set
\[
  \{\xi,-\xi\},
  \qquad \xi^2=\eta.
\]
A subset $T\subseteq L$ is \emph{fibre-closed} if, whenever $\xi\in T$, also $-\xi\in T$.
\end{definition}

By \cref{lem:power-maps}, the fibres partition $L$ into $M/2$ disjoint pairs.

\subsection{Reed--Solomon codes}
\label{sec:preliminaries-rs}

For a finite set $D\subseteq\F$ and a polynomial $P\in\F[X]$, write
\[
  \ev_D(P)\defeq(P(\xi))_{\xi\in D}.
\]
We use Reed--Solomon codes on smooth domains $L\subseteq\Fq^\times\subseteq\F^\times$.

\begin{definition}[Reed--Solomon code]
\label{def:rs-code}
Let $L$ be a smooth domain of order $M$, and let $1\le d\le M$.  The Reed--Solomon code of degree bound $d$ on $L$ over $\F$ is
\[
  \RS_{\F}[L,d]
  \defeq
  \ev_L\bigl(\F[X]_{<d}\bigr)
  \subseteq\F^L.
\]
Its rate is
\[
  \rho\defeq\frac dM.
\]
Elements of $\F^L$ are called \emph{words}; elements of $\RS_{\F}[L,d]$ are called \emph{codewords}.
\end{definition}

For $u,v\in\F^L$, define their relative Hamming distance by
\[
  \Delta(u,v)
  \defeq
  \frac{|\{\xi\in L:u(\xi)\ne v(\xi)\}|}{M}.
\]
For a nonempty set $C\subseteq\F^L$, let
\[
  \Delta(u,C)\defeq\min_{c\in C}\Delta(u,c).
\]

\begin{lemma}[Hamming triangle inequality]
\label{lem:hamming-triangle}
For $u,v,w\in\F^L$,
\[
  \Delta(u,w)\le\Delta(u,v)+\Delta(v,w).
\]
Consequently, if $u$ and $v$ agree on at least $\alpha M$ points and $v$ and $w$ agree on at least $\beta M$ points, then $u$ and $w$ agree on at least $(\alpha+\beta-1)M$ points.
\end{lemma}

\begin{proof}
If $u(\xi)\ne w(\xi)$, then at least one of $u(\xi)\ne v(\xi)$ or $v(\xi)\ne w(\xi)$ holds.  Therefore the disagreement set of $(u,w)$ is contained in the union of the disagreement sets of $(u,v)$ and $(v,w)$.  Divide cardinalities by $M$.  The agreement statement is the same inequality written in terms of complements.
\end{proof}

\begin{proposition}[Basic Reed--Solomon parameters]
\label{prop:rs-parameters}
Let $C=\RS_{\F}[L,d]$, with $|L|=M$ and $1\le d\le M$.
\begin{enumerate}[label=(\arabic*),leftmargin=*]
  \item The evaluation map
  \[
    \ev_L:\F[X]_{<d}\longrightarrow C
  \]
  is an $\F$-linear bijection.
  \item Two distinct codewords agree on at most $d-1$ points of $L$.
  \item If
  \[
    2\delta<1-\frac{d-1}{M},
  \]
  then every word $w\in\F^L$ lies within relative distance $\delta$ of at most one codeword in $C$.  In particular, this holds for every
  \[
    \delta\le\frac{1-\rho}{2}.
  \]
\end{enumerate}
\end{proposition}

\begin{proof}
Linearity is immediate.  If $P,Q\in\F[X]_{<d}$ have the same evaluation on $L$, then $P-Q$ has at least $M\ge d$ roots and degree $<d$, so \cref{lem:root-bound} gives $P=Q$.  Thus $\ev_L$ is injective, and it is surjective onto $C$ by definition, proving (1).

For (2), if distinct codewords $\ev_L(P)$ and $\ev_L(Q)$ agreed on $d$ points, then $P-Q$ would have at least $d$ roots while having degree $<d$, contradicting \cref{lem:root-bound}.

For (3), suppose two distinct codewords $c,c'$ both lie within distance $\delta$ of $w$.  By \cref{lem:hamming-triangle},
\[
  \Delta(c,c')\le2\delta<1-\frac{d-1}{M}.
\]
Thus $c$ and $c'$ agree on more than $d-1$ points, contradicting (2).  Finally,
\[
  \frac{1-\rho}{2}
  =\frac12\left(1-\frac dM\right)
  <\frac12\left(1-\frac{d-1}{M}\right).
\]
\end{proof}

When $c\in\RS_{\F}[L,d]$, we write $P_c\in\F[X]_{<d}$ for its unique representing polynomial from \cref{prop:rs-parameters}(1).

\begin{lemma}[Decoding a base-field word]
\label{lem:basefield-decoding}
Let $w\in\Fq^L$ and let $c\in\RS_{\F}[L,d]$ satisfy
\[
  \Delta(w,c)<\frac12\left(1-\frac{d-1}{M}\right).
\]
Then $P_c\in\Fq[X]_{<d}$, and hence $c\in\Fq^L$.
\end{lemma}

\begin{proof}
Let $\varphi:\F\to\F$ be the $q$-power Frobenius map, $x\mapsto x^q$.  It is a field automorphism whose fixed field is exactly $\Fq$.  Apply $\varphi$ coefficientwise to polynomials and entrywise to words.  Since every $\xi\in L$ lies in $\Fq$,
\[
  \varphi(P(\xi))=\varphi(P)(\xi).
\]
Therefore $\varphi(c)=\ev_L(\varphi(P_c))$ is another codeword in $\RS_{\F}[L,d]$.  Since $w\in\Fq^L$, $\varphi(w)=w$, and the Frobenius map is injective, so
\[
  \Delta(w,\varphi(c))=\Delta(w,c).
\]
By the uniqueness radius in \cref{prop:rs-parameters}(3), $\varphi(c)=c$.  Hence
\[
  \ev_L(\varphi(P_c))=\ev_L(P_c).
\]
Injectivity of evaluation on $\F[X]_{<d}$ gives $\varphi(P_c)=P_c$.  Every coefficient of $P_c$ is therefore fixed by Frobenius and lies in $\Fq$.
\end{proof}

For knowledge extraction we will use the standard existence of a polynomial-time unique-decoding algorithm for Reed--Solomon codes inside half the minimum distance~\cite{WelchBerlekamp86,Gao03}.  The algebraic soundness arguments below use only uniqueness, which was proved in \cref{prop:rs-parameters}.

\subsection{Even--odd splitting of words and fibre distance}
\label{sec:preliminaries-fibre-distance}

Let $L$ be a smooth domain of order $M\ge2$.  For a word $w\in\F^L$, define words $w_{\mathrm e},w_{\mathrm o}\in\F^{L^2}$ by
\begin{equation}
  w_{\mathrm e}(\xi^2)
  \defeq
  \frac{w(\xi)+w(-\xi)}{2},
  \qquad
  w_{\mathrm o}(\xi^2)
  \defeq
  \frac{w(\xi)-w(-\xi)}{2\xi}.
  \label{eq:evenodd-word}
\end{equation}
The denominators are nonzero because the characteristic is odd and $L\subseteq\Fq^\times$.

\begin{lemma}[Even--odd split of words]
\label{lem:evenodd-word}
For $w\in\F^L$:
\begin{enumerate}[label=(\arabic*),leftmargin=*]
  \item The definitions in~\eqref{eq:evenodd-word} are independent of the choice of square root $\xi$.
  \item The map $w\mapsto(w_{\mathrm e},w_{\mathrm o})$ is $\F$-linear and depends on each fibre independently.
  \item For every $\xi\in L$,
  \[
    w(\xi)=w_{\mathrm e}(\xi^2)+\xi w_{\mathrm o}(\xi^2),
    \qquad
    w(-\xi)=w_{\mathrm e}(\xi^2)-\xi w_{\mathrm o}(\xi^2).
  \]
  \item If $P\in\F[X]_{<2d}$, $d\le M/2$, and $w=\ev_L(P)$, then
  \[
    w_{\mathrm e}=\ev_{L^2}(P_{\mathrm e}),
    \qquad
    w_{\mathrm o}=\ev_{L^2}(P_{\mathrm o}),
  \]
  where $P_{\mathrm e},P_{\mathrm o}$ are from \cref{lem:evenodd-polynomial}.
\end{enumerate}
\end{lemma}

\begin{proof}
Replacing $\xi$ by $-\xi$ leaves both right-hand sides of~\eqref{eq:evenodd-word} unchanged, and by \cref{lem:power-maps}(4) these are the only two square roots in $L$, proving (1).  Linearity and fibre locality are immediate, proving (2).  Adding and subtracting the two equations in~\eqref{eq:evenodd-word} gives (3).

For (4), use
\[
  P(\xi)=P_{\mathrm e}(\xi^2)+\xi P_{\mathrm o}(\xi^2),
  \qquad
  P(-\xi)=P_{\mathrm e}(\xi^2)-\xi P_{\mathrm o}(\xi^2)
\]
and substitute into~\eqref{eq:evenodd-word}.
\end{proof}

The folding soundness argument is naturally expressed in terms of fibres rather than individual positions.

\begin{definition}[Fibre distance]
\label{def:fibre-distance}
For words $u,v\in\F^L$, define
\[
  \Delta_{\mathrm{fib}}(u,v)
  \defeq
  \frac{
    |\{\eta\in L^2:\text{$u$ and $v$ differ at at least one point of the fibre over $\eta$}\}|
  }{|L^2|}.
\]
For a nonempty code $C\subseteq\F^L$, put
\[
  \Delta_{\mathrm{fib}}(u,C)
  \defeq
  \min_{c\in C}\Delta_{\mathrm{fib}}(u,c).
\]
\end{definition}

\begin{lemma}[Hamming and fibre distance]
\label{lem:hamming-fibre}
For $u,v\in\F^L$,
\[
  \Delta(u,v)\le\Delta_{\mathrm{fib}}(u,v)\le 2\Delta(u,v).
\]
In particular, if $u$ and $v$ agree on a fibre-closed set $T\subseteq L$ with $|T|\ge(1-\delta)M$, then
\[
  \Delta_{\mathrm{fib}}(u,v)\le\delta.
\]
\end{lemma}

\begin{proof}
Let $B$ be the number of bad fibres, that is, fibres on which $u$ and $v$ differ at one or both points, and let $D$ be the number of pointwise disagreements.  Every bad fibre contributes at least one and at most two disagreements, so
\[
  B\le D\le2B.
\]
Since $|L^2|=M/2$,
\[
  \Delta_{\mathrm{fib}}(u,v)=\frac{B}{M/2}=\frac{2B}{M},
  \qquad
  \Delta(u,v)=\frac DM.
\]
Thus $\Delta(u,v)\le\Delta_{\mathrm{fib}}(u,v)\le2\Delta(u,v)$.

For the last statement, the complement $L\setminus T$ is fibre-closed and contains at most $\delta M$ points, hence at most $\delta M/2$ fibres.  Every bad fibre is contained in this complement, so the fibre distance is at most $(\delta M/2)/(M/2)=\delta$.
\end{proof}

\subsection{Correlated agreement for Reed--Solomon curves}
\label{sec:preliminaries-correlated-agreement}

The soundness arguments later batch several words by a random challenge.  The required coding-theoretic input is the correlated-agreement theorem of Ben-Sasson, Carmon, Ishai, Kopparty, and Saraf~\cite{BSCI23}.  We state precisely the form used in this paper.

For $t\ge1$ and words $u_0,\ldots,u_t\in\F^L$, define the degree-$t$ word curve
\begin{equation}
  u_\beta\defeq\sum_{i=0}^{t}\beta^iu_i,
  \qquad \beta\in\F.
  \label{eq:word-curve}
\end{equation}

\begin{theorem}[Correlated agreement for Reed--Solomon curves~\cite{BSCI23}]
\label{thm:correlated-agreement}
Let
\[
  C=\RS_{\F}[L,d],
  \qquad |L|=M,
  \qquad \rho=d/M,
\]
and let
\[
  0<\delta\le\frac{1-\rho}{2}.
\]
Fix $t\ge1$ and $u_0,\ldots,u_t\in\F^L$, and let $u_\beta$ be defined by~\eqref{eq:word-curve}.  Define
\[
  Z\defeq\{\beta\in\F:\Delta(u_\beta,C)\le\delta\}.
\]
If
\[
  |Z|>tM,
\]
then there exist a set $T\subseteq L$ with
\[
  |T|\ge(1-\delta)M
\]
and codewords $c_0,\ldots,c_t\in C$ such that
\[
  u_i(\xi)=c_i(\xi)
  \qquad
  \text{for every }i\in[0,t]\text{ and every }\xi\in T.
\]
\end{theorem}

\begin{remark}[External input]
\label{rem:external-input}
\Cref{thm:correlated-agreement} is the only non-elementary coding-theoretic theorem used in the algebraic batching arguments of this paper.  All consequences drawn from it, including the fibre-closed agreement sets needed by the folding protocol, are proved explicitly in the later soundness sections.  The theorem is applied over the challenge field $\F$ even though the evaluation domain satisfies $L\subseteq\Fq^\times$.
\end{remark}

\subsection{Parameters used throughout}
\label{sec:preliminaries-parameters}

The main constructions use the following common parameter regime.  Let
\[
  N=2^n,
  \qquad
  M=2^{n+R}=N2^R
\]
for an integer redundancy parameter $R\ge1$, and let $L\le\Fq^\times$ be the smooth domain of order $M$.  The base Reed--Solomon code is
\begin{equation}
  C_0\defeq\RS_{\F}[L,N],
  \qquad
  \rho\defeq\frac NM=2^{-R}.
  \label{eq:base-code}
\end{equation}
We write
\begin{equation}
  \delta_*\defeq\frac{1-\rho}{2}
  \label{eq:unique-radius}
\end{equation}
for the unique-decoding radius used by the soundness proofs.  When the folding backend is run for $\ell$ binary levels, we use the nested domains
\[
  L_j\defeq L^{2^j},
  \qquad
  |L_j|=M/2^j,
\]
and degree bounds
\[
  d_j\defeq N/2^j
\]
for those $j$ for which $2^j$ divides $N$.  The corresponding Reed--Solomon codes are
\[
  C_j\defeq\RS_{\F}[L_j,d_j].
\]
The precise folding protocol, its query parameter $\kappa$, and its error term $\varepsilon_{\mathrm{fold}}$ are defined later; this section fixes only the algebraic objects on which that protocol acts.

\section{Table-Form Kronecker Encoding}
\label{sec:table-encoding}

The protocols in this paper commit directly to vectors indexed by the Boolean hypercube.  The corresponding univariate polynomial is obtained by using the integer index of a Boolean point as the exponent.  This section defines that encoding, relates it precisely to the canonical and Lagrange representations of multilinear polynomials from \cref{sec:preliminaries-multilinear}, and records the commitment properties used throughout the paper.

\subsection{The table polynomial}
\label{sec:table-polynomial}

Let $t:\B^n\to\F$ be a table.  As in \cref{sec:preliminaries-bits}, we identify a Boolean point
\[
  w=(w_1,\ldots,w_n)\in\B^n
\]
with the integer
\[
  w=\sum_{i=1}^{n}w_i2^{i-1}\in[0,N-1].
\]
Thus we write $t(w)$ interchangeably for the value at the bit vector and at its integer index.

\begin{definition}[Table polynomial]
\label{def:table-polynomial}
For a table $t:\B^n\to\F$, its \emph{table polynomial} is
\begin{equation}
  V_t(X)
  \defeq
  \sum_{w=0}^{N-1}t(w)X^w
  \in\F[X]_{<N}.
  \label{eq:table-polynomial}
\end{equation}
\end{definition}

The construction in \cref{def:table-polynomial} is the Kronecker encoding of the length-$N$ vector $(t(0),\ldots,t(N-1))$: the entries of the vector are used verbatim as the coefficients of a univariate polynomial.  No interpolation on the Boolean cube is involved in the definition of $V_t$.

\begin{proposition}[Table encoding is a linear isomorphism]
\label{prop:table-isomorphism}
The map
\[
  \mathcal K_T:\F^{\B^n}\longrightarrow\F[X]_{<N},
  \qquad
  t\longmapsto V_t,
\]
is an $\F$-linear isomorphism.  Its inverse maps
\[
  P(X)=\sum_{j=0}^{N-1}p_jX^j
\]
to the table $t_P$ defined by $t_P(j)=p_j$.
\end{proposition}

\begin{proof}
Linearity follows coefficientwise from \eqref{eq:table-polynomial}.  Every polynomial in $\F[X]_{<N}$ has a unique coefficient vector $(p_0,\ldots,p_{N-1})$, and therefore a unique table $t_P$ with $t_P(j)=p_j$.  For this table,
\[
  V_{t_P}(X)=\sum_{j=0}^{N-1}p_jX^j=P(X).
\]
Conversely, the coefficient of $X^j$ in $V_t$ is exactly $t(j)$, so $t_{V_t}=t$.  Thus the two maps are inverse linear maps.
\end{proof}

\begin{remark}[No distinction between a vector and its table polynomial]
\label{rem:vector-table-polynomial}
Whenever a vector $a=(a_0,\ldots,a_{N-1})\in\F^N$ is indexed by $[0,N-1]$, we may regard it as the table $a:\B^n\to\F$ under the fixed binary indexing and write
\[
  V_a(X)=\sum_{j=0}^{N-1}a_jX^j.
\]
This convention will be used for witness vectors, lookup columns, matrix-derived vectors, and intermediate pointwise-arithmetic wires.  The encoding depends only on the ordered vector, not on an a priori multivariate polynomial.
\end{remark}

\subsection{Relation with multilinear polynomial representations}
\label{sec:table-vs-multilinear}

The table polynomial $V_t$ should not be confused with the canonical-coordinate Kronecker polynomial of the multilinear extension $\widetilde t$.  Both are degree-$<N$ univariate polynomials, but their coefficients represent different coordinate systems.

Recall from \cref{prop:boolean-interpolation} that every table $t:\B^n\to\F$ has a unique multilinear extension
\[
  \widetilde t(z)
  =\sum_{w\in\B^n}t(w)\eq(w,z).
\]
The family
\begin{equation}
  e_w(x_1,\ldots,x_n)
  \defeq
  \eq(w,x)
  =\prod_{i=1}^{n}\bigl((1-w_i)(1-x_i)+w_ix_i\bigr)
  \label{eq:boolean-lagrange-basis}
\end{equation}
is the Boolean Lagrange basis of the multilinear space over $\F$.

\begin{proposition}[Table values are Lagrange coordinates]
\label{prop:table-lagrange-coordinates}
For every table $t:\B^n\to\F$,
\begin{equation}
  \widetilde t(x)
  =\sum_{w\in\B^n}t(w)e_w(x).
  \label{eq:mle-lagrange-expansion}
\end{equation}
Consequently, the coordinate vector of $\widetilde t$ in the Boolean Lagrange basis $(e_w)_{w\in\B^n}$ is exactly the table vector $(t(w))_{w\in\B^n}$, and $V_t$ is the univariate Kronecker encoding of those Lagrange coordinates.
\end{proposition}

\begin{proof}
Equation \eqref{eq:mle-lagrange-expansion} is the definition of the multilinear extension in \eqref{eq:mle-definition}.  The Boolean delta property \cref{lem:eq-delta} implies that the polynomials $e_w$ interpolate the standard coordinate vectors on $\B^n$, hence form the Lagrange basis.  Therefore the coefficient of $e_w$ in \eqref{eq:mle-lagrange-expansion} is $t(w)$.  By \cref{def:table-polynomial}, these same coordinates are the coefficients of $V_t$.
\end{proof}

Now let $f\in\M$ be written in the canonical basis of \cref{prop:canonical-basis} as
\[
  f=\sum_{j=0}^{N-1}\alpha_jm_j.
\]
Its canonical-coordinate Kronecker polynomial is
\begin{equation}
  U_f(X)
  \defeq
  \sum_{j=0}^{N-1}\alpha_jX^j.
  \label{eq:canonical-kronecker-polynomial}
\end{equation}
Its table polynomial, by contrast, is
\[
  V_f(X)
  \defeq
  \sum_{b=0}^{N-1}f(b)X^b.
\]
The two coefficient vectors are related by the Boolean zeta and M\"obius transforms from \cref{lem:mobius}.

\begin{proposition}[Canonical and table Kronecker polynomials]
\label{prop:canonical-table-kronecker}
Let $f\in\M$ have canonical coordinates $(\alpha_j)_{j=0}^{N-1}$.  Then
\begin{equation}
  [X^b]V_f
  =f(b)
  =\sum_{j\preceq b}\alpha_j,
  \label{eq:table-from-canonical}
\end{equation}
and
\begin{equation}
  [X^j]U_f
  =\alpha_j
  =\sum_{c\preceq j}(-1)^{\wt(j)-\wt(c)}[X^c]V_f.
  \label{eq:canonical-from-table}
\end{equation}
In particular, $U_f$ and $V_f$ generally differ even though both belong to $\F[X]_{<N}$.
\end{proposition}

\begin{proof}
The first equality in \eqref{eq:table-from-canonical} is the definition of $V_f$, and the second is \eqref{eq:zeta-transform}.  Equation \eqref{eq:canonical-from-table} is \eqref{eq:mobius-transform} after substituting $f(c)=[X^c]V_f$.  No further argument is required.
\end{proof}

\begin{remark}[Why table form is the default in this paper]
\label{rem:why-table-form}
For a protocol whose witness is already represented by values on $\B^n$, the coefficients of $V_t$ are available directly.  In particular, committing to $V_t$ does not require first applying the M\"obius transform \eqref{eq:mobius-transform}.  More importantly, later relations such as pointwise products, lookup columns, permutations, and sparse-matrix streams are naturally relations between table entries.  The table polynomial therefore gives one uniform univariate representation for all of them.
\end{remark}

\subsection{The Reed--Solomon table commitment}
\label{sec:table-commitment}

Fix the common parameters of \cref{sec:preliminaries-parameters}: a smooth multiplicative domain $L\le\Fq^\times$ of order $M\ge N$ and the base code
\[
  C_0=\RS_{\F}[L,N].
\]

\begin{definition}[Table-form Kronecker commitment]
\label{def:table-commitment}
For a table $t:\B^n\to\F$, define its oracle commitment by
\begin{equation}
  \Enc_T(t)
  \defeq
  \ev_L(V_t)
  =\bigl(V_t(\xi)\bigr)_{\xi\in L}
  \in C_0.
  \label{eq:table-commitment}
\end{equation}
In the interactive-oracle model, the word $\Enc_T(t)$ is the commitment oracle.  A later non-interactive compilation will replace the oracle by a hash/Merkle commitment to its entries.
\end{definition}

\begin{proposition}[Commitment is a linear bijection onto the base code]
\label{prop:table-commitment-bijection}
The map
\[
  \Enc_T:\F^{\B^n}\longrightarrow C_0
\]
is an $\F$-linear bijection.
\end{proposition}

\begin{proof}
By \cref{prop:table-isomorphism}, $t\mapsto V_t$ is an $\F$-linear bijection from $\F^{\B^n}$ onto $\F[X]_{<N}$.  By \cref{prop:rs-parameters}(1), evaluation on $L$ is an $\F$-linear bijection from $\F[X]_{<N}$ onto $C_0$.  Their composition is therefore an $\F$-linear bijection.
\end{proof}

The previous proposition gives a canonical decoder from an exact codeword to a table: if $c\in C_0$, let $P_c\in\F[X]_{<N}$ be its unique representing polynomial and define
\begin{equation}
  \operatorname{Tab}(c)(j)
  \defeq
  [X^j]P_c.
  \label{eq:codeword-table-decoder}
\end{equation}
Then $\Enc_T(\operatorname{Tab}(c))=c$.

For soundness we need the stronger statement that a word sufficiently close to the code determines at most one committed table.

\begin{proposition}[Unique table determined by a nearby word]
\label{prop:unique-table-near-word}
Let $0\le\delta\le\delta_*=(1-\rho)/2$, and let $w\in\F^L$.  There is at most one table $t:\B^n\to\F$ satisfying either of the following conditions:
\begin{align}
  \Delta\bigl(w,\Enc_T(t)\bigr)&\le\delta,
  \label{eq:table-near-hamming}\\
  \Delta_{\mathrm{fib}}\bigl(w,\Enc_T(t)\bigr)&\le\delta.
  \label{eq:table-near-fibre}
\end{align}
\end{proposition}

\begin{proof}
For \eqref{eq:table-near-hamming}, \cref{prop:rs-parameters}(3) implies that there is at most one codeword $c\in C_0$ within distance $\delta$ of $w$.  By \cref{prop:table-commitment-bijection}, each codeword in $C_0$ corresponds to exactly one table, so at most one such $t$ exists.

For \eqref{eq:table-near-fibre}, \cref{lem:hamming-fibre} gives
\[
  \Delta\bigl(w,\Enc_T(t)\bigr)
  \le
  \Delta_{\mathrm{fib}}\bigl(w,\Enc_T(t)\bigr)
  \le\delta.
\]
Thus any table satisfying the fibre-distance condition also satisfies the Hamming-distance condition, and uniqueness follows from the first part.
\end{proof}

\begin{proposition}[Base-field consistency]
\label{prop:table-basefield-consistency}
If $t:\B^n\to\Fq$, then
\[
  \Enc_T(t)\in\Fq^L.
\]
Conversely, let $w\in\Fq^L$ and suppose that for some table $t:\B^n\to\F$,
\begin{equation}
  \Delta\bigl(w,\Enc_T(t)\bigr)
  <\frac12\left(1-\frac{N-1}{M}\right).
  \label{eq:table-basefield-radius}
\end{equation}
Then $t$ takes values in $\Fq$.
\end{proposition}

\begin{proof}
If $t(w)\in\Fq$ for every $w$, then all coefficients of $V_t$ lie in $\Fq$.  Since $L\subseteq\Fq^\times$, every evaluation $V_t(\xi)$ also lies in $\Fq$, proving the first statement.

For the converse, put $c=\Enc_T(t)\in C_0$.  Condition \eqref{eq:table-basefield-radius} and \cref{lem:basefield-decoding} imply that the representing polynomial $P_c$ lies in $\Fq[X]_{<N}$.  By injectivity of evaluation and the definition of the commitment,
\[
  P_c=V_t.
\]
Hence every coefficient $[X^j]V_t=t(j)$ lies in $\Fq$.
\end{proof}

\subsection{Commit once, interpret many relations}
\label{sec:table-commitment-interpretation}

A central reason for using \cref{def:table-commitment} is that later protocols do not change the commitment format.  The word
\[
  w_t=\Enc_T(t)=\ev_L(V_t)
\]
commits simultaneously to the ordered vector $t$, to the Boolean table whose multilinear extension is $\widetilde t$, and to the degree-$<N$ univariate polynomial $V_t$.  Different proof statements merely apply different algebraic functionals to the same coefficients.

The next sections develop these functionals systematically.  In particular:
\begin{enumerate}[label=(\arabic*),leftmargin=*]
  \item multilinear evaluation $\widetilde t(z)$ is a coefficient of $V_t$ multiplied by a public product-form kernel;
  \item an arbitrary inner product $\ip{a}{b}$ is a coefficient of $V_a$ multiplied by the reversal of $V_b$;
  \item pointwise products are reduced to randomly weighted inner products;
  \item lookup, permutation, sparse linear algebra, and R1CS are then assembled from the same table commitments and these basic relations.
\end{enumerate}
No additional encoding is introduced for these applications.  This common representation is what allows heterogeneous relations to be batched later into a single Reed--Solomon proximity test.

\begin{remark}[Commitment versus proof of a relation]
\label{rem:commitment-not-proof}
The map $\Enc_T$ by itself only binds a table to a Reed--Solomon codeword.  It does not prove any evaluation, inner-product, Hadamard, lookup, or R1CS relation.  Those statements require the coefficient-extraction identities and soundness arguments developed in the following sections.  Keeping this distinction explicit will be important when the interactive protocols are later compiled with hash commitments.
\end{remark}

\section{Coefficient-Extraction Calculus}
\label{sec:coefficient-calculus}

This section isolates the algebraic identities that drive every protocol in the paper.  The basic mechanism is elementary: an inner product of two length-$N$ vectors is the middle coefficient of the product of one generating polynomial with the reversal of the other.  What makes this useful for proof systems is that several verifier-defined vectors admit compact product-form generating polynomials, and that the selected coefficient can be isolated by one bounded-degree polynomial identity.

We first define reversal and the general coefficient functional, then derive the canonical-coordinate Kronecker kernel and the table-form multilinear-evaluation kernel, and finally prove the universal opening identity used later by the evaluation, inner-product, and Hadamard protocols.

\subsection{Reversal and the middle coefficient}
\label{sec:coefficient-reversal}

\begin{definition}[Reversal]
\label{def:reversal}
For $P\in\F[X]_{<N}$, define its \emph{reversal} by
\begin{equation}
  P^{\star}(X)
  \defeq
  X^{N-1}P(X^{-1}).
  \label{eq:reversal}
\end{equation}
Equivalently, if
\[
  P(X)=\sum_{j=0}^{N-1}p_jX^j,
\]
then
\begin{equation}
  P^{\star}(X)
  =\sum_{j=0}^{N-1}p_jX^{N-1-j}.
  \label{eq:reversal-coefficients}
\end{equation}
\end{definition}

Although \eqref{eq:reversal} is written using $X^{-1}$, the right-hand side is a polynomial because $\deg P<N$.

\begin{lemma}[Elementary properties of reversal]
\label{lem:reversal-properties}
For $P,Q\in\F[X]_{<N}$ and $a,b\in\F$:
\begin{enumerate}[label=(\arabic*),leftmargin=*]
  \item $[X^j]P^{\star}=[X^{N-1-j}]P$ for every $j\in[0,N-1]$;
  \item $(P^{\star})^{\star}=P$;
  \item $(aP+bQ)^{\star}=aP^{\star}+bQ^{\star}$.
\end{enumerate}
\end{lemma}

\begin{proof}
Write $P(X)=\sum_{j=0}^{N-1}p_jX^j$.  Then \eqref{eq:reversal-coefficients} gives
\[
  P^{\star}(X)=\sum_{j=0}^{N-1}p_jX^{N-1-j}.
\]
The coefficient of $X^j$ is therefore $p_{N-1-j}=[X^{N-1-j}]P$, proving (1).  Applying the same identity twice sends $p_j$ first to position $N-1-j$ and then back to $j$, proving (2).  Statement (3) follows coefficientwise from the definition.
\end{proof}

\begin{theorem}[Middle-coefficient inner-product identity]
\label{thm:middle-coefficient-inner-product}
Let $a=(a_0,\ldots,a_{N-1})$ and $b=(b_0,\ldots,b_{N-1})$ be vectors in $\F^N$, and let
\[
  V_a(X)=\sum_{j=0}^{N-1}a_jX^j,
  \qquad
  V_b(X)=\sum_{j=0}^{N-1}b_jX^j.
\]
Then
\begin{equation}
  \ip{a}{b}
  =\sum_{j=0}^{N-1}a_jb_j
  =\coeff{N-1}{V_a(X)V_b^{\star}(X)}.
  \label{eq:middle-coefficient-inner-product}
\end{equation}
\end{theorem}

\begin{proof}
By \cref{def:reversal},
\[
  V_b^{\star}(X)=\sum_{k=0}^{N-1}b_kX^{N-1-k}.
\]
Hence the coefficient of $X^{N-1}$ in $V_aV_b^{\star}$ is
\[
  \sum_{j+k'=N-1} [X^j]V_a\,[X^{k'}]V_b^{\star}
  =\sum_{j=0}^{N-1}a_jb_j,
\]
because $[X^{N-1-j}]V_b^{\star}=b_j$ by \cref{lem:reversal-properties}(1).
\end{proof}

\begin{definition}[Public kernel]
\label{def:public-kernel}
For a public vector $p=(p_0,\ldots,p_{N-1})\in\F^N$, define its \emph{reversed generating kernel}
\begin{equation}
  K_p(X)
  \defeq
  V_p^{\star}(X)
  =\sum_{j=0}^{N-1}p_jX^{N-1-j}.
  \label{eq:public-kernel}
\end{equation}
\end{definition}

The previous theorem immediately gives the following linear-functional form.

\begin{corollary}[Public linear functionals]
\label{cor:public-linear-functional}
For every $a,p\in\F^N$,
\begin{equation}
  \sum_{j=0}^{N-1}a_jp_j
  =\coeff{N-1}{V_aK_p}.
  \label{eq:public-linear-functional}
\end{equation}
\end{corollary}

\begin{proof}
Apply \cref{thm:middle-coefficient-inner-product} with $b=p$ and use \cref{def:public-kernel}.
\end{proof}

\begin{definition}[Succinct kernel family]
\label{def:succinct-kernel-family}
Let $\Theta$ be a parameter set and let $p_\theta\in\F^N$ for $\theta\in\Theta$.  We call $(p_\theta)_{\theta\in\Theta}$ a \emph{succinct kernel family} if the value
\[
  K_{p_\theta}(\xi)
\]
can be computed from $(\theta,\xi)$ using $\operatorname{poly}(n)$ field operations, without materialising all $N=2^n$ coordinates of $p_\theta$.
\end{definition}

This definition is descriptive rather than cryptographic: later protocols will still have to prove low degree and consistency of the committed words.  Its role is to identify those public linear functionals whose local verifier work remains polynomial in $n$.

\subsection{The canonical-coordinate Kronecker kernel}
\label{sec:canonical-kernel}

We now recover the coefficient-extraction identity for a multilinear polynomial written in the canonical basis of \cref{sec:preliminaries-multilinear}.  Let
\[
  f=\sum_{j=0}^{N-1}\alpha_jm_j\in\M,
  \qquad
  U_f(X)=\sum_{j=0}^{N-1}\alpha_jX^j.
\]
For $z=(z_1,\ldots,z_n)\in\F^n$, define
\begin{equation}
  K_z(X)
  \defeq
  \prod_{i=1}^{n}\left(X^{2^{i-1}}+z_i\right).
  \label{eq:canonical-kernel}
\end{equation}

\begin{lemma}[Coefficients of the canonical kernel]
\label{lem:canonical-kernel-coefficients}
For every $j\in[0,N-1]$ with bits $(j_1,\ldots,j_n)$,
\begin{equation}
  [X^{N-1-j}]K_z(X)
  =\prod_{i=1}^{n}z_i^{j_i}
  =m_j(z).
  \label{eq:canonical-kernel-coefficient}
\end{equation}
\end{lemma}

\begin{proof}
Expanding \eqref{eq:canonical-kernel}, from each factor indexed by $i$ we choose either $X^{2^{i-1}}$ or $z_i$.  To obtain the monomial $X^{N-1-j}$, whose $i$-th binary digit is $1-j_i$ by \cref{lem:complement-bits}, we must choose $X^{2^{i-1}}$ exactly when $j_i=0$, and choose $z_i$ exactly when $j_i=1$.  The resulting coefficient is therefore
\[
  \prod_{i:j_i=1}z_i
  =\prod_{i=1}^{n}z_i^{j_i}
  =m_j(z).
\]
\end{proof}

\begin{theorem}[Canonical Kronecker extraction]
\label{thm:canonical-kronecker-extraction}
Let $f\in\M$ and $z\in\F^n$.  Then
\begin{equation}
  f(z)
  =\coeff{N-1}{U_f(X)K_z(X)}.
  \label{eq:canonical-kronecker-extraction}
\end{equation}
\end{theorem}

\begin{proof}
Using \cref{lem:canonical-kernel-coefficients},
\begin{align*}
  \coeff{N-1}{U_fK_z}
  &=\sum_{j=0}^{N-1}[X^j]U_f\,[X^{N-1-j}]K_z\\
  &=\sum_{j=0}^{N-1}\alpha_jm_j(z)
   =f(z).
\end{align*}
\end{proof}

\begin{remark}[Relation with the middle-coefficient identity]
\label{rem:canonical-kernel-inner-product}
Equation \eqref{eq:canonical-kronecker-extraction} is exactly \cref{thm:middle-coefficient-inner-product} applied to the canonical-coordinate vector $(\alpha_j)_j$ and the vector of monomial values $(m_j(z))_j$.  The special feature is that the reversed generating polynomial of $(m_j(z))_j$ factors as the product \eqref{eq:canonical-kernel}.  Thus an exponentially long verifier vector is represented by $n$ simple factors.
\end{remark}

\begin{proposition}[Local evaluation cost of $K_z$]
\label{prop:canonical-kernel-cost}
Given $\xi\in\F$, the value $K_z(\xi)$ can be computed using $n-1$ squarings to obtain
\[
  \xi,\xi^2,\xi^4,\ldots,\xi^{2^{n-1}},
\]
then $n$ additions and $n-1$ multiplications to multiply the $n$ factors in \eqref{eq:canonical-kernel}.
\end{proposition}

\begin{proof}
Starting with $\xi$, repeated squaring computes $\xi^{2^{i-1}}$ from $\xi^{2^{i-2}}$ for $i=2,\ldots,n$, which uses $n-1$ squarings.  Forming each $\xi^{2^{i-1}}+z_i$ uses one addition, hence $n$ additions.  Multiplying $n$ field elements uses $n-1$ multiplications.
\end{proof}

\subsection{The table-form multilinear-evaluation kernel}
\label{sec:table-evaluation-kernel}

The table-form commitment of \cref{sec:table-encoding} stores the Boolean values of a multilinear polynomial rather than its canonical coefficients.  Its evaluation kernel is therefore the generating polynomial of the Boolean equality weights.

For $z\in\F^n$, define
\begin{equation}
  E_z(X)
  \defeq
  \sum_{w\in\B^n}\eq(w,z)X^w.
  \label{eq:eq-generating-polynomial}
\end{equation}

\begin{lemma}[Product form of the equality generating polynomial]
\label{lem:eq-kernel-product}
For every $z\in\F^n$,
\begin{equation}
  E_z(X)
  =\prod_{i=1}^{n}\left((1-z_i)+z_iX^{2^{i-1}}\right).
  \label{eq:eq-kernel-product}
\end{equation}
Its reversal is
\begin{equation}
  E_z^{\star}(X)
  =\prod_{i=1}^{n}\left(z_i+(1-z_i)X^{2^{i-1}}\right).
  \label{eq:eq-kernel-reversal-product}
\end{equation}
\end{lemma}

\begin{proof}
Expanding the product in \eqref{eq:eq-kernel-product}, the coefficient of $X^w$ is obtained by choosing $z_iX^{2^{i-1}}$ when $w_i=1$ and $(1-z_i)$ when $w_i=0$.  Hence
\[
  [X^w]E_z
  =\prod_{i=1}^{n}z_i^{w_i}(1-z_i)^{1-w_i}
  =\eq(w,z),
\]
which proves \eqref{eq:eq-kernel-product} by \eqref{eq:eq-generating-polynomial}.

For the reversal, use $N-1=\sum_i2^{i-1}$:
\begin{align*}
  E_z^{\star}(X)
  &=X^{N-1}E_z(X^{-1})\\
  &=\prod_{i=1}^{n}X^{2^{i-1}}
      \left((1-z_i)+z_iX^{-2^{i-1}}\right)\\
  &=\prod_{i=1}^{n}\left(z_i+(1-z_i)X^{2^{i-1}}\right).
\end{align*}
\end{proof}

\begin{theorem}[Table-form multilinear evaluation by coefficient extraction]
\label{thm:table-evaluation-extraction}
Let $t:\B^n\to\F$ be a table and let $\widetilde t$ be its multilinear extension.  Then for every $z\in\F^n$,
\begin{equation}
  \widetilde t(z)
  =\coeff{N-1}{V_t(X)E_z^{\star}(X)}.
  \label{eq:table-evaluation-extraction}
\end{equation}
\end{theorem}

\begin{proof}
By \cref{thm:middle-coefficient-inner-product}, applied to the vectors $(t(w))_w$ and $(\eq(w,z))_w$,
\[
  \coeff{N-1}{V_tE_z^{\star}}
  =\sum_{w\in\B^n}t(w)\eq(w,z).
\]
The right-hand side is $\widetilde t(z)$ by \eqref{eq:mle-definition}.
\end{proof}

\begin{proposition}[Local evaluation cost of $E_z^{\star}$]
\label{prop:table-kernel-cost}
Given $\xi\in\F$, the value $E_z^{\star}(\xi)$ can be computed using $n-1$ squarings, at most $2n$ additions/subtractions, $n$ scalar multiplications by the public values $1-z_i$, and $n-1$ multiplications to combine the factors.  In particular, its cost is $O(n)$ field operations.
\end{proposition}

\begin{proof}
The powers $\xi^{2^{i-1}}$ are obtained by repeated squaring as in \cref{prop:canonical-kernel-cost}.  Each factor
\[
  z_i+(1-z_i)\xi^{2^{i-1}}
\]
requires computing $1-z_i$, one multiplication by $\xi^{2^{i-1}}$, and one addition.  Multiplying the $n$ factors uses $n-1$ further multiplications.  The stated bound follows.
\end{proof}

\begin{corollary}[Hypercube sums]
\label{cor:hypercube-sum-extraction}
For every table $t:\B^n\to\F$,
\begin{equation}
  \sum_{w\in\B^n}t(w)
  =\coeff{N-1}{V_t(X)K_{\mathbf 1}(X)},
  \qquad
  K_{\mathbf 1}(X)=\sum_{j=0}^{N-1}X^j
  =\prod_{i=1}^{n}\left(1+X^{2^{i-1}}\right).
  \label{eq:hypercube-sum-extraction}
\end{equation}
\end{corollary}

\begin{proof}
Apply \cref{cor:public-linear-functional} with the all-ones vector $\mathbf1$.  Its reversed generating polynomial equals its ordinary generating polynomial because all coefficients are one.  The product identity follows by expanding $\prod_i(1+X^{2^{i-1}})$: each exponent in $[0,N-1]$ has a unique binary expansion, so every monomial $X^j$ appears exactly once.
\end{proof}

\begin{remark}[One table commitment, two evaluation viewpoints]
\label{rem:two-evaluation-viewpoints}
The canonical identity \eqref{eq:canonical-kronecker-extraction} and the table identity \eqref{eq:table-evaluation-extraction} are the same middle-coefficient mechanism applied in two coordinate systems.  In canonical coordinates the verifier vector is $(m_j(z))_j$ and its reversed generating polynomial is $K_z$.  In table/Lagrange coordinates the verifier vector is $(\eq(w,z))_w$ and its reversed generating polynomial is $E_z^{\star}$.  The rest of this paper uses table form because pointwise arithmetic, lookup columns, permutation streams, and sparse linear-algebra witnesses are naturally represented by their values.
\end{remark}

\subsection{The universal opening identity}
\label{sec:universal-opening-identity}

The preceding identities reduce the desired claim to a selected coefficient of a product.  We now isolate that coefficient by a polynomial identity whose auxiliary polynomials obey the same degree bound $N$ as the committed Reed--Solomon code.

\begin{definition}[Coefficient functional]
\label{def:coefficient-functional}
For $U,K\in\F[X]_{<N}$, define
\begin{equation}
  \mathsf{CE}_N(U,K)
  \defeq
  \coeff{N-1}{U(X)K(X)}.
  \label{eq:coefficient-functional}
\end{equation}
\end{definition}

\begin{theorem}[Universal split identity]
\label{thm:universal-split}
Let $U,K\in\F[X]_{<N}$ and put
\[
  v=\mathsf{CE}_N(U,K).
\]
Then there exist unique polynomials $A\in\F[X]_{<N}$ and $H\in\F[X]_{<N}$ such that
\begin{equation}
  XU(X)K(X)
  =A(X)+vX^N+X^{N+1}H(X).
  \label{eq:universal-split}
\end{equation}
Moreover, the honest polynomial $H$ actually satisfies $\deg H<N-1$ unless $H=0$.
\end{theorem}

\begin{proof}
Because $\deg U,\deg K\le N-1$,
\[
  \deg(XUK)\le2N-1.
\]
Write
\[
  XUK=\sum_{r=0}^{2N-1}c_rX^r.
\]
The coefficient at degree $N$ is
\[
  c_N=[X^N](XUK)=[X^{N-1}]UK=v.
\]
Define
\[
  A(X)=\sum_{r=0}^{N-1}c_rX^r
\]
and
\[
  H(X)=\sum_{r=N+1}^{2N-1}c_rX^{r-(N+1)}.
\]
Then $\deg A<N$, while the largest possible exponent of $H$ is
\[
  (2N-1)-(N+1)=N-2,
\]
so $H\in\F[X]_{<N}$, and \eqref{eq:universal-split} holds.

For uniqueness, suppose also
\[
  XUK=A'+vX^N+X^{N+1}H'
\]
with $A',H'\in\F[X]_{<N}$.  The degrees occupied by $A'$, $vX^N$, and $X^{N+1}H'$ are respectively contained in
\[
  [0,N-1],\qquad \{N\},\qquad [N+1,2N].
\]
These sets are disjoint.  Comparing coefficients in the first and third ranges gives $A'=A$ and $H'=H$.
\end{proof}

The converse direction is what later soundness proofs use.

\begin{lemma}[Universal identity lemma]
\label{lem:universal-identity}
Let $U,K,A,H\in\F[X]_{<N}$ and $v\in\F$, and define
\begin{equation}
  \Phi(X)
  \defeq
  XU(X)K(X)-A(X)-vX^N-X^{N+1}H(X).
  \label{eq:identity-residual}
\end{equation}
Then
\begin{equation}
  \deg\Phi\le2N.
  \label{eq:identity-residual-degree}
\end{equation}
If $\Phi=0$, then
\begin{equation}
  v=\mathsf{CE}_N(U,K)=\coeff{N-1}{UK}.
  \label{eq:identity-implies-coefficient}
\end{equation}
\end{lemma}

\begin{proof}
The four terms in \eqref{eq:identity-residual} have degrees at most
\[
  2N-1,\qquad N-1,\qquad N,\qquad 2N,
\]
respectively, proving \eqref{eq:identity-residual-degree}.  If $\Phi=0$, compare the coefficient of $X^N$.  The term $A$ contributes nothing because $\deg A<N$, and $X^{N+1}H$ has no monomial of degree $N$.  Therefore
\[
  [X^N](XUK)=v.
\]
Since $[X^N](XUK)=[X^{N-1}]UK$, \eqref{eq:identity-implies-coefficient} follows.
\end{proof}

\begin{corollary}[General split criterion]
\label{cor:general-split-criterion}
For fixed $U,K\in\F[X]_{<N}$ and $v\in\F$, the following are equivalent:
\begin{enumerate}[label=(\arabic*),leftmargin=*]
  \item $v=\mathsf{CE}_N(U,K)$;
  \item there exist $A,H\in\F[X]_{<N}$ satisfying \eqref{eq:universal-split}.
\end{enumerate}
When these conditions hold, $A$ and $H$ are unique.
\end{corollary}

\begin{proof}
If (1) holds, existence and uniqueness follow from \cref{thm:universal-split}.  If (2) holds, the residual polynomial \eqref{eq:identity-residual} is zero, so \cref{lem:universal-identity} gives (1).  Uniqueness follows again from \cref{thm:universal-split}.
\end{proof}

\begin{remark}[Why the factor $X$ is essential]
\label{rem:factor-x}
Without the factor $X$, one might try to certify
\[
  UK=A_0+vX^{N-1}+X^NH_0.
\]
To prevent $A_0$ from absorbing an incorrect value of $v$, one would need the stricter bound $\deg A_0<N-1$.  The Reed--Solomon backend used in this paper naturally enforces the power-of-two degree bound $N$, not $N-1$.  Multiplication by $X$ moves the selected coefficient from degree $N-1$ to degree $N$, so the ordinary bound $A\in\F[X]_{<N}$ suffices.
\end{remark}

\begin{corollary}[Pointwise reconstruction of the high part]
\label{cor:virtual-high-part}
Suppose $U,K,A,H\in\F[X]_{<N}$ and $v\in\F$ satisfy \eqref{eq:universal-split}.  Then for every $\xi\in\F^{\times}$,
\begin{equation}
  H(\xi)
  =\frac{\xi U(\xi)K(\xi)-A(\xi)-v\xi^N}{\xi^{N+1}}.
  \label{eq:virtual-high-part}
\end{equation}
\end{corollary}

\begin{proof}
Evaluate \eqref{eq:universal-split} at $\xi$ and divide by the nonzero value $\xi^{N+1}$.
\end{proof}

Equation \eqref{eq:virtual-high-part} is the bridge from coefficient extraction to the proximity backend: on the multiplicative evaluation domain $L\subseteq\Fq^{\times}$, the verifier can compute the would-be evaluation of $H$ locally from the evaluations of $U$, the kernel $K$, and the single auxiliary polynomial $A$.  The next section formalises this and other transformations as virtual Reed--Solomon words.

\section{Virtual Reed--Solomon Words}
\label{sec:virtual-oracles}

The coefficient identities of \cref{sec:coefficient-calculus} introduce low-degree polynomials that need not be committed independently.  Their evaluations can often be reconstructed locally from words that are already committed and from public challenges.  This section isolates the three transformations used throughout the rest of the paper: reversal on an inversion-closed domain, reconstruction of the high polynomial in the universal split identity, and DEEP quotient words for binding out-of-domain evaluations.

The word ``virtual'' is used only to describe how the verifier obtains oracle values.  A virtual word is not assumed to be low degree.  Low degree is a conclusion that must be enforced by the later batching and proximity arguments.

\subsection{Virtual words and locality}
\label{sec:virtual-locality}

\begin{definition}[Virtual word]
\label{def:virtual-word}
Let $L$ be the base evaluation domain.  Suppose the verifier has oracle access to words
\[
  w_1,\ldots,w_s\in\F^L.
\]
A word $v\in\F^L$ is called a \emph{virtual word derived from $(w_1,\ldots,w_s)$} if, for every $\xi\in L$, the verifier can compute $v(\xi)$ from public data and from a prescribed finite set of oracle values of the $w_i$.

The virtual word is \emph{$q$-local} if at most $q$ oracle positions in total are required to compute one value $v(\xi)$.  If the required positions lie in one fibre $\{\zeta,-\zeta\}$ whenever the folding verifier queries one fibre, we call the transformation \emph{fibre local}.
\end{definition}

\begin{remark}[Virtuality does not imply low degree]
\label{rem:virtual-not-low-degree}
\Cref{def:virtual-word} is purely operational.  An arbitrary pointwise function of low-degree codewords need not itself be a Reed--Solomon codeword.  Every later protocol therefore batches the virtual words with the committed words and applies a proximity test.  Soundness is obtained only after showing that a batch which is close to the code for sufficiently many verifier challenges forces the individual words to admit compatible low-degree explanations.
\end{remark}

\subsection{Reversal on a multiplicative domain}
\label{sec:virtual-reversal}

The first virtual transformation is the reversal of \cref{def:reversal}.  It is available without a new commitment because the smooth domain $L\le\Fq^\times$ is closed under inversion.

\begin{definition}[Virtual reversal]
\label{def:virtual-reversal}
For a word $w\in\F^L$, define $w^{\star}\in\F^L$ by
\begin{equation}
  w^{\star}(\xi)
  \defeq
  \xi^{N-1}w(\xi^{-1}),
  \qquad \xi\in L.
  \label{eq:virtual-reversal}
\end{equation}
This is well defined because $L$ is a multiplicative group and therefore $\xi^{-1}\in L$.
\end{definition}

\begin{lemma}[Reversal commutes with evaluation]
\label{lem:reversal-commutes-evaluation}
Let $P\in\F[X]_{<N}$ and $w=\ev_L(P)$.  Then
\begin{equation}
  w^{\star}=\ev_L(P^{\star}).
  \label{eq:reversal-commutes-evaluation}
\end{equation}
Consequently, the map $w\mapsto w^{\star}$ maps the base code $C_0$ bijectively onto itself.
\end{lemma}

\begin{proof}
For every $\xi\in L$,
\[
  w^{\star}(\xi)
  =\xi^{N-1}w(\xi^{-1})
  =\xi^{N-1}P(\xi^{-1})
  =P^{\star}(\xi)
\]
by \cref{def:reversal}.  This proves \eqref{eq:reversal-commutes-evaluation}.  Since reversal is an involution on $\F[X]_{<N}$ by \cref{lem:reversal-properties}, it is a bijection on the representing polynomials of $C_0$, hence $w\mapsto w^{\star}$ is a bijection of $C_0$.
\end{proof}

\begin{lemma}[Algebraic properties of virtual reversal]
\label{lem:virtual-reversal-properties}
The map $\mathcal R:\F^L\to\F^L$, $\mathcal R(w)=w^{\star}$, satisfies:
\begin{enumerate}[label=(\arabic*),leftmargin=*]
  \item $\mathcal R$ is $\F$-linear;
  \item $\mathcal R^2=\operatorname{id}$;
  \item for $u,v\in\F^L$ and $\xi\in L$,
  \[
    u^{\star}(\xi)=v^{\star}(\xi)
    \quad\Longleftrightarrow\quad
    u(\xi^{-1})=v(\xi^{-1});
  \]
  \item $\Delta(u^{\star},v^{\star})=\Delta(u,v)$.
\end{enumerate}
\end{lemma}

\begin{proof}
Linearity follows immediately from \eqref{eq:virtual-reversal}.  For the involution property,
\begin{align*}
  (w^{\star})^{\star}(\xi)
  &=\xi^{N-1}w^{\star}(\xi^{-1})\\
  &=\xi^{N-1}(\xi^{-1})^{N-1}w((\xi^{-1})^{-1})\\
  &=w(\xi).
\end{align*}
For (3), the scalar $\xi^{N-1}$ is nonzero, so equality of the two reversed values is equivalent to equality of the two original values at $\xi^{-1}$.  Finally, inversion is a bijection of $L$, so the number of positions on which $u^{\star}$ and $v^{\star}$ differ equals the number of positions on which $u$ and $v$ differ.  Dividing by $|L|$ proves (4).
\end{proof}

The folding backend works with fibre-closed agreement sets, so we also need inversion to preserve fibres.

\begin{lemma}[Inversion preserves fibres]
\label{lem:inversion-preserves-fibres}
For $T\subseteq L$, define
\[
  T^{-1}\defeq\{\xi^{-1}:\xi\in T\}.
\]
Then:
\begin{enumerate}[label=(\arabic*),leftmargin=*]
  \item $|T^{-1}|=|T|$;
  \item $T$ is fibre-closed if and only if $T^{-1}$ is fibre-closed;
  \item for every $u,v\in\F^L$,
  \begin{equation}
    \Delta_{\mathrm{fib}}(u^{\star},v^{\star})
    =\Delta_{\mathrm{fib}}(u,v).
    \label{eq:reversal-fibre-distance}
  \end{equation}
\end{enumerate}
\end{lemma}

\begin{proof}
Inversion is a bijection on the finite group $L$, proving (1).  Moreover
\[
  (-\xi)^{-1}=-\xi^{-1},
\]
so inversion maps the fibre $\{\xi,-\xi\}$ bijectively to the fibre $\{\xi^{-1},-\xi^{-1}\}$.  This proves (2).  By \cref{lem:virtual-reversal-properties}(3), a fibre is bad for $(u^{\star},v^{\star})$ exactly when its inverse fibre is bad for $(u,v)$.  Inversion permutes the fibres, so the number of bad fibres is unchanged, proving \eqref{eq:reversal-fibre-distance}.
\end{proof}

\begin{corollary}[Fibre-local queries for reversal]
\label{cor:reversal-fibre-local}
If a folding check requires $w^{\star}(\xi)$ and $w^{\star}(-\xi)$, the verifier obtains both values by querying $w$ only on the single inverse fibre
\[
  \{\xi^{-1},-\xi^{-1}\}.
\]
Thus virtual reversal is fibre local and requires no additional oracle commitment.
\end{corollary}

\begin{proof}
Apply \eqref{eq:virtual-reversal} at $\xi$ and $-\xi$ and use $(-\xi)^{-1}=-\xi^{-1}$.
\end{proof}

\subsection{The virtual high-part word}
\label{sec:virtual-high-word}

We now turn the pointwise reconstruction formula of \cref{cor:virtual-high-part} into a word transformation.  We formulate it first for three arbitrary words; later specializations will take the second word either from a public kernel or from the virtual reversal of another commitment.

\begin{definition}[Virtual high-part word]
\label{def:virtual-high-word}
Let $u,k,a\in\F^L$ and let $v\in\F$.  Define
\begin{equation}
  \mathsf H[u,k,a;v](\xi)
  \defeq
  \frac{\xi\,u(\xi)k(\xi)-a(\xi)-v\xi^N}{\xi^{N+1}},
  \qquad \xi\in L.
  \label{eq:virtual-high-word}
\end{equation}
The denominator is nonzero because $L\subseteq\Fq^{\times}$.
\end{definition}

\begin{lemma}[Pointwise identity and locality]
\label{lem:virtual-high-locality}
For all $u,k,a\in\F^L$, $v\in\F$, and $\xi\in L$, if
\[
  h=\mathsf H[u,k,a;v],
\]
then
\begin{equation}
  \xi\,u(\xi)k(\xi)
  =a(\xi)+v\xi^N+\xi^{N+1}h(\xi).
  \label{eq:virtual-high-pointwise-identity}
\end{equation}
Moreover $h(\xi)$ depends only on the three values $u(\xi)$, $k(\xi)$, $a(\xi)$ and public field data.  In particular, if $k(\xi)$ is itself public or virtually computable from a bounded number of oracle reads, then so is $h(\xi)$.
\end{lemma}

\begin{proof}
Equation \eqref{eq:virtual-high-pointwise-identity} is obtained by multiplying \eqref{eq:virtual-high-word} by the nonzero scalar $\xi^{N+1}$.  The locality statement is immediate from the same formula.
\end{proof}

\begin{lemma}[Honest low-degree realization]
\label{lem:virtual-high-honest}
Let $U,K,A,H\in\F[X]_{<N}$ and $v\in\F$ satisfy
\begin{equation}
  XUK=A+vX^N+X^{N+1}H.
  \label{eq:virtual-section-split}
\end{equation}
Put
\[
  u=\ev_L(U),\qquad
  k=\ev_L(K),\qquad
  a=\ev_L(A).
\]
Then
\begin{equation}
  \mathsf H[u,k,a;v]=\ev_L(H)\in C_0.
  \label{eq:virtual-high-honest}
\end{equation}
\end{lemma}

\begin{proof}
Evaluate \eqref{eq:virtual-section-split} at $\xi\in L$ and solve for $H(\xi)$.  The resulting expression is exactly \eqref{eq:virtual-high-word}.  Since $H\in\F[X]_{<N}$, its evaluation word belongs to $C_0$.
\end{proof}

\begin{lemma}[Linearity for a fixed multiplier]
\label{lem:virtual-high-linearity}
Fix $k\in\F^L$.  For $u,u',a,a'\in\F^L$ and $v,v'\in\F$,
\begin{equation}
  \mathsf H[u+u',k,a+a';v+v']
  =\mathsf H[u,k,a;v]+\mathsf H[u',k,a';v'].
  \label{eq:virtual-high-linearity}
\end{equation}
More generally, for $\lambda\in\F$,
\[
  \mathsf H[\lambda u,k,\lambda a;\lambda v]
  =\lambda\mathsf H[u,k,a;v].
\]
\end{lemma}

\begin{proof}
Both identities follow pointwise from distributivity in the numerator of \eqref{eq:virtual-high-word}.  The multiplier word $k$ is held fixed.
\end{proof}

\begin{remark}[Committed second factors]
\label{rem:committed-second-factor}
For an arbitrary committed inner product, the multiplier word $k$ will be the virtual reversal $w_b^{\star}$ of a second committed word.  The map $(u,k)\mapsto \mathsf H[u,k,a;v]$ is bilinear in its product term rather than jointly linear.  This causes no difficulty: the verifier computes the product locally, while the later soundness argument uses correlated agreement to obtain low-degree explanations for $u$, $k$, $a$, and $h$ separately, after which the polynomial identity of \cref{lem:universal-identity} applies.
\end{remark}

\subsection{DEEP quotient words}
\label{sec:deep-quotients}

The Hadamard, lookup, and sparse-relation protocols will need to certify values of a committed polynomial at points outside the Reed--Solomon domain.  We use the standard quotient identity
\[
  P(X)-P(x)=(X-x)Q(X)
\]
to turn such a claim into another low-degree word.

\begin{definition}[DEEP quotient word]
\label{def:deep-quotient-word}
Let $w\in\F^L$, let $x\in\F\setminus L$, and let $y\in\F$.  Define
\begin{equation}
  \mathsf Q[w;y,x](\xi)
  \defeq
  \frac{w(\xi)-y}{\xi-x},
  \qquad \xi\in L.
  \label{eq:deep-quotient-word}
\end{equation}
The denominator is nonzero because $x\notin L$.
\end{definition}

\begin{lemma}[Honest DEEP quotient]
\label{lem:deep-honest}
Let $P\in\F[X]_{<N}$, let $x\in\F\setminus L$, and set
\[
  w=\ev_L(P),\qquad y=P(x).
\]
Then there exists a unique polynomial
\[
  R(X)=\frac{P(X)-y}{X-x}\in\F[X]_{<N-1}
\]
and
\begin{equation}
  \mathsf Q[w;y,x]=\ev_L(R)\in C_0.
  \label{eq:deep-honest}
\end{equation}
\end{lemma}

\begin{proof}
Since $P(x)-y=0$, the factor theorem gives a unique $R\in\F[X]$ satisfying
\[
  P(X)-y=(X-x)R(X).
\]
If $P$ is nonconstant, then $\deg R=\deg P-1\le N-2$; if $P$ is constant, then $R=0$.  Thus $R\in\F[X]_{<N-1}$.  For every $\xi\in L$,
\[
  R(\xi)=\frac{P(\xi)-y}{\xi-x}
  =\frac{w(\xi)-y}{\xi-x}
  =\mathsf Q[w;y,x](\xi),
\]
which proves \eqref{eq:deep-honest}.
\end{proof}

The converse is the soundness property needed later: if a word and its quotient word both agree with low-degree polynomials on sufficiently many positions, then the claimed out-of-domain value is forced.

\begin{lemma}[DEEP consistency on an agreement set]
\label{lem:deep-consistency}
Let $x\in\F\setminus L$, $y\in\F$, and let $P,R\in\F[X]_{<N}$.  Let $w,q\in\F^L$ satisfy
\[
  w(\xi)=P(\xi),
  \qquad
  q(\xi)=R(\xi)
  \qquad
  \text{for every }\xi\in T,
\]
for a set $T\subseteq L$.  Suppose also that
\[
  q=\mathsf Q[w;y,x]
\]
pointwise on $L$.  If
\begin{equation}
  |T|>N,
  \label{eq:deep-agreement-size}
\end{equation}
then
\begin{equation}
  P(x)=y
  \label{eq:deep-value-forced}
\end{equation}
and
\begin{equation}
  P(X)-y=(X-x)R(X).
  \label{eq:deep-polynomial-identity}
\end{equation}
In particular $\deg R<N-1$ unless $R=0$.
\end{lemma}

\begin{proof}
For every $\xi\in T$, the definition of $q$ gives
\[
  (\xi-x)q(\xi)-w(\xi)+y=0.
\]
Using the agreements with $P$ and $R$ on $T$,
\[
  (\xi-x)R(\xi)-P(\xi)+y=0.
\]
Therefore the polynomial
\[
  S(X)\defeq (X-x)R(X)-P(X)+y
\]
vanishes on all points of $T$.  Since $P,R\in\F[X]_{<N}$,
\[
  \deg S\le N.
\]
The set $T$ contains more than $N$ distinct field elements by \eqref{eq:deep-agreement-size}, so root counting forces $S=0$.  Evaluating at $X=x$ gives
\[
  -P(x)+y=0,
\]
which proves \eqref{eq:deep-value-forced}; the identity $S=0$ is exactly \eqref{eq:deep-polynomial-identity}.  The final degree statement follows because the right-hand side divides the degree of the nonzero polynomial $P-y$ by one.
\end{proof}

\begin{corollary}[Uniqueness of a DEEP opening]
\label{cor:deep-opening-unique}
Let $P\in\F[X]_{<N}$ and $x\in\F\setminus L$.  There is at most one $y\in\F$ for which a polynomial $R\in\F[X]_{<N}$ can satisfy
\[
  P(X)-y=(X-x)R(X).
\]
That value is $y=P(x)$.
\end{corollary}

\begin{proof}
Evaluating the displayed identity at $X=x$ gives $P(x)-y=0$.
\end{proof}

\begin{lemma}[Locality of DEEP quotients]
\label{lem:deep-locality}
For fixed public $(x,y)$ with $x\notin L$, the value $\mathsf Q[w;y,x](\xi)$ requires exactly one oracle value $w(\xi)$ and $O(1)$ field operations.  If a folding check queries the fibre $\{\xi,-\xi\}$, the two quotient values require only the two values $w(\xi)$ and $w(-\xi)$ from that same fibre.
\end{lemma}

\begin{proof}
This is immediate from \eqref{eq:deep-quotient-word}.
\end{proof}

\begin{remark}[Sampling points outside the domain]
\label{rem:deep-poles}
We formulate DEEP quotients with $x\in\F\setminus L$ so that completeness is exact and no pole convention is needed.  Later protocols may enforce this by sampling from a prescribed subset of the challenge field disjoint from $L$; alternatively, one may sample from all of $\F$ and account explicitly for the event $x\in L$ in the completeness and soundness bounds.  The algebraic statements of this section require only $x\notin L$.
\end{remark}

\subsection{Composition of virtual transformations}
\label{sec:virtual-composition}

The protocols below combine the preceding transformations.  The most important example is the committed inner product: if $w_a=\ev_L(V_a)$ and $w_b=\ev_L(V_b)$, then the verifier first forms $w_b^{\star}$ virtually and then uses it as the multiplier word in \cref{def:virtual-high-word}.  In the honest case, \cref{lem:reversal-commutes-evaluation,lem:virtual-high-honest} give
\[
  w_b^{\star}=\ev_L(V_b^{\star})
\]
and, for the unique $A,H\in\F[X]_{<N}$ satisfying
\[
  XV_aV_b^{\star}=A+SX^N+X^{N+1}H,
\]
\[
  \mathsf H[w_a,w_b^{\star},\ev_L(A);S]=\ev_L(H).
\]
Thus the verifier can query all four low-degree objects
\[
  V_a,\qquad V_b^{\star},\qquad A,\qquad H
\]
while only $V_a$, $V_b$, and $A$ require actual commitments.

Similarly, an out-of-domain value $P(x)=y$ can be tied to the commitment of $P$ by placing the virtual quotient $\mathsf Q[w;y,x]$ in the same random batch as the original word.  The local verifier sees only ordinary oracle values and field operations; global consistency is delegated to the Reed--Solomon proximity argument.

\begin{proposition}[Closure of honest constructions]
\label{prop:honest-virtual-closure}
Every virtual transformation introduced in this section maps the corresponding honest polynomial relation to a word in $C_0$:
\begin{enumerate}[label=(\arabic*),leftmargin=*]
  \item if $w=\ev_L(P)$ with $P\in\F[X]_{<N}$, then $w^{\star}=\ev_L(P^{\star})\in C_0$;
  \item if $u,k,a$ are evaluations of $U,K,A\in\F[X]_{<N}$ satisfying the universal split identity with $H\in\F[X]_{<N}$, then $\mathsf H[u,k,a;v]=\ev_L(H)\in C_0$;
  \item if $w=\ev_L(P)$ and $y=P(x)$ for $x\notin L$, then $\mathsf Q[w;y,x]=\ev_L((P-y)/(X-x))\in C_0$.
\end{enumerate}
\end{proposition}

\begin{proof}
The three statements are exactly \cref{lem:reversal-commutes-evaluation,lem:virtual-high-honest,lem:deep-honest}.
\end{proof}

The next section turns this closure property into a generic coefficient-extraction protocol.  The prover commits only to those auxiliary polynomials whose values cannot be reconstructed locally; all remaining low-degree objects are virtual.  A random linear combination then reduces simultaneous low-degree enforcement to one Reed--Solomon proximity test.

\section{The Folding Backend and the Generic Coefficient-Extraction Protocol}
\label{sec:generic-protocol}

The previous sections reduced coefficient claims to the existence of several words that should all be Reed--Solomon codewords of the same degree bound.  This section supplies the common enforcement mechanism.  We first give a self-contained FRI-style folding proximity test for the base code $C_0=\RS_{\F}[L,N]$, including the soundness statement used later.  We then combine it with the universal split identity of \cref{thm:universal-split} to obtain a generic protocol for claims of the form
\[
  v=[X^{N-1}]U(X)K(X),
  \qquad U,K\in\F[X]_{<N}.
\]
The protocol is deliberately stated for two potentially committed factors.  Public structured kernels are obtained by omitting the second factor from the random batch, and arbitrary committed inner products are obtained by taking the second factor to be a virtual reversal.

\subsection{Two elementary probability lemmas}
\label{sec:generic-probability}

We use two counting facts repeatedly.  They are recorded here so that the soundness proof does not depend on an external probability formalism.

\begin{lemma}[Sequential bad challenges]
\label{lem:sequential-bad-challenges}
Let $S_1,\ldots,S_t$ be finite nonempty sets.  For each $j\in[1,t]$, suppose that after every fixed prefix $(s_1,\ldots,s_{j-1})$ there are at most $b_j$ values $s_j\in S_j$ declared bad.  If $s_1,\ldots,s_t$ are sampled independently and uniformly, then
\[
  \Pr[\text{some round chooses a bad value}]
  \le
  \sum_{j=1}^{t}\frac{b_j}{|S_j|}.
\]
\end{lemma}

\begin{proof}
For a fixed round $j$, condition on the prefix $(s_1,\ldots,s_{j-1})$.  By hypothesis the conditional probability that $s_j$ is bad is at most $b_j/|S_j|$.  Averaging over prefixes gives the same unconditional bound.  A union bound over the $t$ rounds proves the claim.
\end{proof}

\begin{lemma}[Independent uniform queries]
\label{lem:independent-uniform-queries}
Let $\Omega$ be a finite nonempty set, let $D$ be a finite nonempty set, and let $W(\omega)\subseteq D$ be assigned to every $\omega\in\Omega$.  If $\omega$ is uniform in $\Omega$ and $\xi^{(1)},\ldots,\xi^{(\kappa)}$ are independent uniform points of $D$, independent of $\omega$, then
\[
  \Pr\bigl[\xi^{(1)},\ldots,\xi^{(\kappa)}\in W(\omega)\bigr]
  =
  \mathbb E_{\omega}\left[
    \left(\frac{|W(\omega)|}{|D|}\right)^{\kappa}
  \right].
\]
In particular, if $|W(\omega)|\le\pi|D|$ for every $\omega$ in an event $E\subseteq\Omega$, then
\[
  \Pr\bigl[E\text{ and }\xi^{(1)},\ldots,\xi^{(\kappa)}\in W(\omega)\bigr]
  \le \pi^{\kappa}.
\]
\end{lemma}

\begin{proof}
For fixed $\omega$, exactly $|W(\omega)|^{\kappa}$ of the $|D|^{\kappa}$ ordered query tuples lie in $W(\omega)^{\kappa}$.  This gives the conditional probability.  Averaging over $\omega$ proves the first identity.  The second statement follows by restricting the expectation to $E$ and using $|W(\omega)|/|D|\le\pi$ there.
\end{proof}

\subsection{The kernel fold}
\label{sec:generic-kernel-fold}

For this subsection let $L'\le\Fq^\times$ be any smooth domain of order at least $2$.  For $w\in\F^{L'}$, write $w_{\mathrm e},w_{\mathrm o}\in\F^{(L')^2}$ for the even and odd parts of \cref{eq:evenodd-word}.

\begin{definition}[Kernel fold]
\label{def:kernel-fold}
For $r\in\F$, define
\begin{equation}
  \operatorname{fold}_r(w)
  \defeq
  (2r-1)w_{\mathrm e}+(1-r)w_{\mathrm o}.
  \label{eq:kernel-word-fold}
\end{equation}
If
\[
  P(X)=P_{\mathrm e}(X^2)+XP_{\mathrm o}(X^2),
\]
define
\begin{equation}
  \operatorname{pfold}_r(P)
  \defeq
  (2r-1)P_{\mathrm e}+(1-r)P_{\mathrm o}.
  \label{eq:kernel-poly-fold}
\end{equation}
\end{definition}

\begin{lemma}[Local formula and line form]
\label{lem:fold-local-line}
Let $w\in\F^{L'}$ and $r\in\F$.
\begin{enumerate}[label=(\arabic*),leftmargin=*]
  \item For $\xi\in L'$,
  \begin{equation}
  \operatorname{fold}_r(w)(\xi^2)
  =
  \frac{((2r-1)\xi+(1-r))w(\xi)
       +((2r-1)\xi-(1-r))w(-\xi)}{2\xi}.
  \label{eq:fold-local-formula}
  \end{equation}
  \item The map $w\mapsto\operatorname{fold}_r(w)$ is $\F$-linear and fibre local.
  \item Put
  \[
    u_0=w_{\mathrm o}-w_{\mathrm e},
    \qquad
    u_1=2w_{\mathrm e}-w_{\mathrm o}.
  \]
  Then
  \begin{equation}
    \operatorname{fold}_r(w)=u_0+ru_1,
    \label{eq:fold-line-form}
  \end{equation}
  and conversely
  \[
    w_{\mathrm e}=u_0+u_1,
    \qquad
    w_{\mathrm o}=2u_0+u_1.
  \]
\end{enumerate}
\end{lemma}

\begin{proof}
Substitute the definitions
\[
  w_{\mathrm e}(\xi^2)=\frac{w(\xi)+w(-\xi)}2,
  \qquad
  w_{\mathrm o}(\xi^2)=\frac{w(\xi)-w(-\xi)}{2\xi}
\]
into \eqref{eq:kernel-word-fold} and collect the coefficients of $w(\xi)$ and $w(-\xi)$.  This gives \eqref{eq:fold-local-formula}.  Linearity and fibre locality follow immediately.  Finally,
\[
  (2r-1)w_{\mathrm e}+(1-r)w_{\mathrm o}
  =(w_{\mathrm o}-w_{\mathrm e})
   +r(2w_{\mathrm e}-w_{\mathrm o}),
\]
which is \eqref{eq:fold-line-form}; adding the two displayed definitions of $u_0,u_1$ gives $w_{\mathrm e}=u_0+u_1$, and $2u_0+u_1=w_{\mathrm o}$.
\end{proof}

\begin{proposition}[Consistency of word and polynomial folding]
\label{prop:fold-consistency}
Let $|L'|=M'$, let $1\le d\le M'/2$, let $P\in\F[X]_{<2d}$, and put $w=\ev_{L'}(P)$.  Then
\begin{equation}
  \operatorname{fold}_r(w)
  =
  \ev_{(L')^2}\bigl(\operatorname{pfold}_r(P)\bigr).
  \label{eq:fold-consistency}
\end{equation}
In particular,
\[
  \operatorname{fold}_r\bigl(\RS_{\F}[L',2d]\bigr)
  \subseteq
  \RS_{\F}[(L')^2,d].
\]
\end{proposition}

\begin{proof}
By \cref{lem:evenodd-word}(4),
\[
  w_{\mathrm e}=\ev_{(L')^2}(P_{\mathrm e}),
  \qquad
  w_{\mathrm o}=\ev_{(L')^2}(P_{\mathrm o}).
\]
Linearity of evaluation and \cref{def:kernel-fold} give \eqref{eq:fold-consistency}.  Since $P_{\mathrm e},P_{\mathrm o}\in\F[X]_{<d}$, also $\operatorname{pfold}_r(P)\in\F[X]_{<d}$.
\end{proof}

For $r=(r_1,\ldots,r_j)\in\F^j$, write
\[
  \operatorname{pfold}^{(j)}_r
  =
  \operatorname{pfold}_{r_j}\circ\cdots\circ\operatorname{pfold}_{r_1},
\]
with $\operatorname{pfold}^{(0)}$ the identity.

\subsection{The folding proximity test}
\label{sec:folding-test}

Fix integers $1\le\ell\le n$ and $\kappa\ge1$.  For $j\in[0,\ell]$, let
\[
  L_j=L^{2^j},
  \qquad
  M_j=M/2^j,
  \qquad
  N_j=N/2^j,
  \qquad
  C_j=\RS_{\F}[L_j,N_j].
\]
Every $C_j$ has rate $\rho=N/M$.

\begin{definition}[Folding proximity test]
\label{def:folding-test}
The protocol $\Pi_{\mathrm{FT}}[\ell,\kappa]$ has as instance an oracle word $w_0\in\F^L$.
\begin{enumerate}[label=\textnormal{\arabic*.},leftmargin=*]
  \item The verifier samples $r_1\in\F$.
  \item For $j=2,\ldots,\ell$, after seeing $r_1,\ldots,r_{j-1}$ the prover sends an oracle word $w_{j-1}\in\F^{L_{j-1}}$, and the verifier samples $r_j\in\F$.
  \item After $r_1,\ldots,r_\ell$ are fixed, the prover sends a polynomial
  \[
    P\in\F[X]_{<N_\ell}
  \]
  in coefficient form.  Put $w_\ell=\ev_{L_\ell}(P)$.
  \item The verifier samples independent uniform points
  \[
    \xi^{(1)},\ldots,\xi^{(\kappa)}\in L.
  \]
  For each query point $\xi=\xi^{(i)}$ and each level $j\in[1,\ell]$, it checks
  \begin{equation}
    \operatorname{fold}_{r_j}(w_{j-1})(\xi^{2^j})
    =
    w_j(\xi^{2^j}).
    \label{eq:fold-check}
  \end{equation}
\end{enumerate}
For $j\in[1,\ell]$, write
\[
  \widehat w_j=\operatorname{fold}_{r_j}(w_{j-1}).
\]
\end{definition}

The left side of \eqref{eq:fold-check} is computed by reading $w_{j-1}$ on the fibre above $\xi^{2^j}$, namely at $\pm\xi^{2^{j-1}}$.  For $j<\ell$, the right side is among the values read at the next level; at the last level it is obtained by evaluating the explicit polynomial $P$.

\begin{proposition}[Completeness of folding]
\label{prop:folding-completeness}
If $w_0=\ev_L(U)$ for some $U\in\F[X]_{<N}$, then the honest prover defined recursively by
\[
  U_j=\operatorname{pfold}_{r_j}(U_{j-1}),
  \qquad U_0=U,
\]
with $w_j=\ev_{L_j}(U_j)$ and $P=U_\ell$, is accepted with probability $1$.
\end{proposition}

\begin{proof}
By \cref{prop:fold-consistency},
\[
  \operatorname{fold}_{r_j}(w_{j-1})
  =\ev_{L_j}(U_j)=w_j
\]
for every level.  Hence every check \eqref{eq:fold-check} holds at every point.
\end{proof}

\subsection{Good folding challenges}
\label{sec:good-folding-challenges}

For a word $w\in\F^{L_j}$ satisfying $\Delta_{\mathrm{fib}}(w,C_j)\le\delta$, let $\operatorname{dec}_j(w)$ denote the unique codeword in $C_j$ within fibre distance $\delta$.  Uniqueness follows from \cref{lem:hamming-fibre,prop:rs-parameters}.  We use this notation only when such a codeword exists.

\begin{lemma}[Folding a fibre-far word]
\label{lem:folding-far-word}
Let $j\in[0,\ell-1]$ and let $w\in\F^{L_j}$ satisfy
\[
  \Delta_{\mathrm{fib}}(w,C_j)>\delta.
\]
Then
\begin{equation}
  \left|
    \left\{r\in\F:
      \Delta(\operatorname{fold}_r(w),C_{j+1})\le\delta
    \right\}
  \right|
  \le M_{j+1}.
  \label{eq:folding-far-word-bound}
\end{equation}
\end{lemma}

\begin{proof}
Suppose instead that the set in \eqref{eq:folding-far-word-bound} has more than $M_{j+1}$ elements.  By \cref{lem:fold-local-line},
\[
  \operatorname{fold}_r(w)=u_0+ru_1,
\]
where $u_0=w_{\mathrm o}-w_{\mathrm e}$ and $u_1=2w_{\mathrm e}-w_{\mathrm o}$ are words on $L_{j+1}$.  Apply \cref{thm:correlated-agreement} with $t=1$ to the code $C_{j+1}$.  There exist a set $T\subseteq L_{j+1}$ with $|T|\ge(1-\delta)M_{j+1}$ and codewords $c_0,c_1\in C_{j+1}$ such that $u_0=c_0$ and $u_1=c_1$ on $T$.

Let $P_0,P_1\in\F[X]_{<N_{j+1}}$ represent $c_0,c_1$, and define
\[
  E=P_0+P_1,
  \qquad
  O=2P_0+P_1,
  \qquad
  Q(X)=E(X^2)+XO(X^2).
\]
Then $Q\in\F[X]_{<N_j}$.  By the inverse relations in \cref{lem:fold-local-line}, $w_{\mathrm e}=\ev(E)$ and $w_{\mathrm o}=\ev(O)$ on $T$.  Therefore for every $\eta\in T$ and every $\xi\in L_j$ with $\xi^2=\eta$, \cref{lem:evenodd-word}(3) gives
\[
  w(\pm\xi)=E(\eta)\pm\xi O(\eta)=Q(\pm\xi).
\]
Thus $w$ agrees with the codeword $\ev_{L_j}(Q)\in C_j$ on all fibres above $T$, so \cref{lem:hamming-fibre} gives $\Delta_{\mathrm{fib}}(w,C_j)\le\delta$, a contradiction.
\end{proof}

\begin{lemma}[A fibre error survives all but one challenges]
\label{lem:fibre-error-survives}
Let $j\in[0,\ell-1]$, let $w,c\in\F^{L_j}$, and let $\eta\in L_{j+1}$ be a fibre on which $w$ and $c$ differ.  Then there is at most one $r\in\F$ satisfying
\[
  \operatorname{fold}_r(w)(\eta)
  =
  \operatorname{fold}_r(c)(\eta).
\]
\end{lemma}

\begin{proof}
Put $y=w-c$.  By linearity and \cref{lem:fold-local-line},
\[
  \operatorname{fold}_r(y)(\eta)
  =
  (y_{\mathrm o}(\eta)-y_{\mathrm e}(\eta))
  +r(2y_{\mathrm e}(\eta)-y_{\mathrm o}(\eta)),
\]
a polynomial of degree at most $1$ in $r$.  If it were the zero polynomial, the inverse formulas of \cref{lem:fold-local-line} would give $y_{\mathrm e}(\eta)=y_{\mathrm o}(\eta)=0$.  By \cref{lem:evenodd-word}(3), $y$ would vanish on the fibre above $\eta$, contrary to hypothesis.  Hence the polynomial is nonzero and has at most one root.
\end{proof}

\begin{definition}[Good folding challenge]
\label{def:good-folding-challenge}
Let $j\in[1,\ell]$, let $w\in\F^{L_{j-1}}$, and let $r\in\F$.  We call $r$ \emph{good for $w$ at level $j$} if the following holds.
\begin{enumerate}[label=(\arabic*),leftmargin=*]
  \item If $\Delta_{\mathrm{fib}}(w,C_{j-1})>\delta$, then
  \[
    \Delta(\operatorname{fold}_r(w),C_j)>\delta.
  \]
  \item If $\Delta_{\mathrm{fib}}(w,C_{j-1})\le\delta$, put $c=\operatorname{dec}_{j-1}(w)$.  For every fibre $\eta\in L_j$ on which $w$ and $c$ differ,
  \[
    \operatorname{fold}_r(w)(\eta)
    \ne
    \operatorname{fold}_r(c)(\eta).
  \]
\end{enumerate}
\end{definition}

\begin{lemma}[Few bad folding challenges]
\label{lem:few-bad-folding-challenges}
For fixed $j\in[1,\ell]$ and $w\in\F^{L_{j-1}}$, at most $M_j$ values $r\in\F$ are not good for $w$ at level $j$.
\end{lemma}

\begin{proof}
In the first case of \cref{def:good-folding-challenge}, apply \cref{lem:folding-far-word}.  In the second case, at most $M_j$ fibres are bad, and by \cref{lem:fibre-error-survives} each bad fibre excludes at most one value of $r$.  A union bound on excluded values gives the result.
\end{proof}

\begin{lemma}[One sound folding step]
\label{lem:one-sound-folding-step}
Let $j\in[1,\ell]$, let $w\in\F^{L_{j-1}}$, and let $r$ be good for $w$ at level $j$.  Suppose $c'\in C_j$ and
\[
  Y=\{\eta\in L_j:\operatorname{fold}_r(w)(\eta)=c'(\eta)\}
\]
satisfies $|Y|\ge(1-\delta)M_j$.  Then
\[
  \Delta_{\mathrm{fib}}(w,C_{j-1})\le\delta.
\]
If $c=\operatorname{dec}_{j-1}(w)$, then
\[
  \operatorname{fold}_r(c)=c',
\]
and $w=c$ on the complete fibre above every $\eta\in Y$.
\end{lemma}

\begin{proof}
Since $c'$ agrees with $\operatorname{fold}_r(w)$ on at least $(1-\delta)M_j$ positions,
\[
  \Delta(\operatorname{fold}_r(w),C_j)\le\delta.
\]
Goodness therefore excludes case (1) of \cref{def:good-folding-challenge}, so $c=\operatorname{dec}_{j-1}(w)$ exists.  By \cref{prop:fold-consistency}, $\operatorname{fold}_r(c)\in C_j$.

The words $w$ and $c$ agree on all fibres above at least $(1-\delta)M_j$ points of $L_j$.  By fibre locality, $\operatorname{fold}_r(w)$ and $\operatorname{fold}_r(c)$ agree on those points.  They also have the common intermediary $\operatorname{fold}_r(w)$ with $c'$ on $Y$.  Hence \cref{lem:hamming-triangle} implies that $\operatorname{fold}_r(c)$ and $c'$ agree on at least
\[
  (1-2\delta)M_j
  \ge
  \rho M_j
  =N_j
\]
points.  Two codewords of $C_j$ that agree on $N_j$ points are equal by \cref{prop:rs-parameters}(2), so $\operatorname{fold}_r(c)=c'$.

Finally, if $\eta\in Y$, then
\[
  \operatorname{fold}_r(w)(\eta)
  =c'(\eta)
  =\operatorname{fold}_r(c)(\eta).
\]
Condition (2) in \cref{def:good-folding-challenge} therefore implies that $w=c$ on the fibre above $\eta$.
\end{proof}

\subsection{Witness sets and the codeword chain}
\label{sec:folding-witness-sets}

Fix a deterministic prover for \cref{def:folding-test}.  The words $w_j$ are then deterministic functions of the preceding folding challenges.  For $j\in[1,\ell]$, define the failed-check set
\[
  B_j
  \defeq
  \{\eta\in L_j:\widehat w_j(\eta)\ne w_j(\eta)\}.
\]

\begin{definition}[Witness sets]
\label{def:folding-witness-sets}
For $c\in C_0$, let
\[
  W_0(c)
  \defeq
  \{\xi\in L:w_0(\xi)=c(\xi)\text{ and }w_0(-\xi)=c(-\xi)\}.
\]
For $j\in[1,\ell]$ and $c\in C_j$, define
\[
  W_j(c)
  \defeq
  \left\{
    \xi\in L:
    \begin{array}{l}
      \xi^{2^i}\notin B_i\text{ for every }i\in[1,j-1],\\
      \widehat w_j(\xi^{2^j})=c(\xi^{2^j})
    \end{array}
  \right\}.
\]
Put
\[
  \operatorname{dens}_j(c)=\frac{|W_j(c)|}{M}.
\]
\end{definition}

\begin{lemma}[Acceptance is membership in the final witness set]
\label{lem:acceptance-witness-set}
For fixed folding challenges $r_1,\ldots,r_\ell$, a point $\xi\in L$ passes all fold checks if and only if
\[
  \xi\in W_\ell(w_\ell).
\]
Consequently, conditioned on the folding challenges, the probability that all $\kappa$ independent query points pass is
\[
  \operatorname{dens}_\ell(w_\ell)^\kappa.
\]
\end{lemma}

\begin{proof}
By definition, $\xi\in W_\ell(w_\ell)$ means that the checks at levels $1,\ldots,\ell-1$ do not fail and that
\[
  \widehat w_\ell(\xi^{2^\ell})=w_\ell(\xi^{2^\ell}),
\]
which is precisely the last check.  The probability statement follows because the $\kappa$ query points are independent and uniform in $L$.
\end{proof}

\begin{lemma}[One round of the codeword chain]
\label{lem:one-round-codeword-chain}
Let $j\in[1,\ell]$ and suppose $r_j$ is good for $w_{j-1}$ at level $j$.  If $c'\in C_j$ satisfies
\[
  \operatorname{dens}_j(c')\ge1-\delta,
\]
then $c=\operatorname{dec}_{j-1}(w_{j-1})$ exists,
\[
  \operatorname{fold}_{r_j}(c)=c',
  \qquad
  W_j(c')\subseteq W_{j-1}(c),
\]
and therefore $\operatorname{dens}_{j-1}(c)\ge1-\delta$.
\end{lemma}

\begin{proof}
Let
\[
  Y=\{\eta\in L_j:\widehat w_j(\eta)=c'(\eta)\}.
\]
Every $\xi\in W_j(c')$ maps to a point $\xi^{2^j}\in Y$.  The map $\xi\mapsto\xi^{2^j}$ is $2^j$-to-one from $L$ onto $L_j$, hence
\[
  |Y|\ge\frac{|W_j(c')|}{2^j}
  \ge(1-\delta)M_j.
\]
Apply \cref{lem:one-sound-folding-step}; this gives existence of $c$, the equality $\operatorname{fold}_{r_j}(c)=c'$, and equality $w_{j-1}=c$ on every fibre above $Y$.

Now let $\xi\in W_j(c')$.  At level $j-1$, the point $\xi^{2^{j-1}}$ lies in the fibre above $\xi^{2^j}\in Y$, so $w_{j-1}$ and $c$ agree there and at its negative.  If $j=1$, this is exactly $\xi\in W_0(c)$.  If $j>1$, the membership conditions defining $W_j(c')$ already guarantee all earlier fold checks, while equality $w_{j-1}=c$ on the relevant fibre converts the level-$(j-1)$ check into the defining equality for $W_{j-1}(c)$.  Thus $W_j(c')\subseteq W_{j-1}(c)$.
\end{proof}

\begin{proposition}[Codeword chain]
\label{prop:codeword-chain}
Suppose every $r_j$ is good for $w_{j-1}$ and
\[
  \operatorname{dens}_\ell(w_\ell)\ge1-\delta.
\]
Then there exist codewords $c_j\in C_j$, $j\in[0,\ell]$, with $c_\ell=w_\ell$, such that
\[
  c_j=\operatorname{dec}_j(w_j)
  \quad (j<\ell),
  \qquad
  \operatorname{fold}_{r_{j+1}}(c_j)=c_{j+1},
\]
and
\[
  W_\ell(w_\ell)\subseteq W_j(c_j)
  \qquad\text{for every }j.
\]
In particular,
\begin{equation}
  \Delta_{\mathrm{fib}}(w_0,C_0)\le\delta.
  \label{eq:codeword-chain-close}
\end{equation}
\end{proposition}

\begin{proof}
Apply \cref{lem:one-round-codeword-chain} successively for $j=\ell,\ell-1,\ldots,1$.  The density hypothesis needed at each step follows from the containment of witness sets obtained at the preceding step.  The final codeword $c_0$ agrees with $w_0$ on every complete fibre represented in $W_\ell(w_\ell)$, whose density is at least $1-\delta$, proving \eqref{eq:codeword-chain-close}.
\end{proof}

\subsection{Soundness of the folding backend}
\label{sec:folding-soundness}

Define
\begin{equation}
  \varepsilon_{\mathrm{fold}}
  \defeq
  \sum_{j=1}^{\ell}\frac{M_j}{|\F|}
  =
  \frac{M(1-2^{-\ell})}{|\F|}
  <\frac{M}{|\F|}.
  \label{eq:epsilon-fold}
\end{equation}

\begin{theorem}[Folding soundness]
\label{thm:folding-soundness}
Fix a deterministic prover for $\Pi_{\mathrm{FT}}[\ell,\kappa]$ and an instance $w_0\in\F^L$.  Let $\mathsf{Good}$ be the event that every $r_j$ is good for $w_{j-1}$ at level $j$.  Then:
\begin{enumerate}[label=(\arabic*),leftmargin=*]
  \item $\Pr[\neg\mathsf{Good}]\le\varepsilon_{\mathrm{fold}}$;
  \item conditioned on $\mathsf{Good}$, either $\operatorname{dens}_\ell(w_\ell)<1-\delta$, or the conclusions of \cref{prop:codeword-chain} hold;
  \item if $\Delta_{\mathrm{fib}}(w_0,C_0)>\delta$, then the verifier accepts with probability at most
  \begin{equation}
    \varepsilon_{\mathrm{fold}}+(1-\delta)^\kappa.
    \label{eq:folding-soundness-bound}
  \end{equation}
\end{enumerate}
The same bound holds for randomized provers whose internal randomness is independent of the verifier challenges.
\end{theorem}

\begin{proof}
For (1), once $r_1,\ldots,r_{j-1}$ are fixed, the word $w_{j-1}$ is fixed.  By \cref{lem:few-bad-folding-challenges}, at most $M_j$ values of $r_j$ are bad.  Apply \cref{lem:sequential-bad-challenges}.

Item (2) is exactly \cref{prop:codeword-chain} when the final density is at least $1-\delta$.

For (3), suppose $\Delta_{\mathrm{fib}}(w_0,C_0)>\delta$.  On $\mathsf{Good}$, item (2) forces
\[
  \operatorname{dens}_\ell(w_\ell)<1-\delta,
\]
otherwise \cref{prop:codeword-chain} would contradict the hypothesis.  By \cref{lem:acceptance-witness-set,lem:independent-uniform-queries}, acceptance together with $\mathsf{Good}$ has probability at most $(1-\delta)^\kappa$.  Add the probability of $\neg\mathsf{Good}$ from (1).

For a randomized prover, condition on its internal random tape.  The deterministic bound applies to every fixed tape; averaging preserves the same upper bound.
\end{proof}

\subsection{Virtual folding levels and higher arity}
\label{sec:virtual-folding-levels}

The prover need not commit to every intermediate folded word.  This standard optimization is useful later because all application protocols already open several source words per query.

\begin{definition}[Committed folding levels]
\label{def:committed-folding-levels}
Let $\mathcal C\subseteq[1,\ell-1]$.  The test $\Pi^{\mathcal C}_{\mathrm{FT}}[\ell,\kappa]$ is obtained from \cref{def:folding-test} by requiring an oracle $w_j$ only when $j\in\mathcal C$.  If $j\notin\mathcal C$, define the virtual level
\[
  w_j=\widehat w_j=\operatorname{fold}_{r_j}(w_{j-1}).
\]
The verifier checks only the levels in $\mathcal C\cup\{\ell\}$, computing skipped folds locally from the previous committed level.
\end{definition}

\begin{proposition}[Virtual levels preserve soundness]
\label{prop:virtual-levels}
For every $\mathcal C\subseteq[1,\ell-1]$, \cref{thm:folding-soundness} holds for $\Pi^{\mathcal C}_{\mathrm{FT}}[\ell,\kappa]$ with exactly the same error bound.
\end{proposition}

\begin{proof}
Every prover for the sparse-level protocol canonically defines a prover for the full protocol: at an omitted level it sends the deterministic virtual word $\operatorname{fold}_{r_j}(w_{j-1})$.  The omitted fold check then holds identically.  At a committed level, the full verifier and the sparse-level verifier compute the same iterated fold from the same previously committed word.  Hence their acceptance predicates coincide for every challenge sequence.  Apply \cref{thm:folding-soundness} to the induced full-protocol prover.
\end{proof}

For an integer $k\ge1$, taking
\[
  \mathcal C_k=\{k,2k,3k,\ldots\}\cap[1,\ell-1]
\]
corresponds to folding arity $2^k$.  No soundness term changes; only the query geometry and Merkle-authentication cost change.

\subsection{The generic coefficient-extraction relation}
\label{sec:generic-ce-relation}

We now attach the folding backend to the universal coefficient identity.

\begin{definition}[Coefficient-extraction relation]
\label{def:generic-ce-relation}
For $0<\delta\le\delta_*$, define $\mathcal R^{\delta}_{\mathrm{CE}}$ as the set of pairs
\[
  \bigl((w_U,w_K;v),(U,K)\bigr)
\]
with
\[
  w_U,w_K\in\F^L,
  \qquad
  U,K\in\F[X]_{<N},
  \qquad
  v\in\F,
\]
such that
\begin{equation}
  \Delta_{\mathrm{fib}}(w_U,\ev_L(U))\le\delta,
  \qquad
  \Delta_{\mathrm{fib}}(w_K,\ev_L(K))\le\delta,
  \label{eq:generic-ce-distance}
\end{equation}
and
\begin{equation}
  v=[X^{N-1}]U(X)K(X).
  \label{eq:generic-ce-value}
\end{equation}
\end{definition}

\begin{definition}[Generic coefficient-extraction protocol $\Pi_{\mathrm{CE}}$]
\label{def:generic-ce-protocol}
The instance is $(w_U,w_K;v)$ with oracle words $w_U,w_K\in\F^L$ and $v\in\F$.
\begin{enumerate}[label=\textnormal{\arabic*.},leftmargin=*]
  \item The prover sends an oracle $w_A\in\F^L$.  The verifier samples $\beta\in\F$.
  \item Define the virtual high-part word
  \begin{equation}
    h=\mathsf H[w_U,w_K,w_A;v]
    \label{eq:generic-ce-high-word}
  \end{equation}
  as in \cref{def:virtual-high-word}, and form the virtual batch
  \begin{equation}
    w_0^{(\beta)}
    =
    w_U+\beta w_K+\beta^2w_A+\beta^3h.
    \label{eq:generic-ce-batch}
  \end{equation}
  \item Run $\Pi_{\mathrm{FT}}[\ell,\kappa]$, or any sparse-level variant of \cref{def:committed-folding-levels}, on $w_0^{(\beta)}$.
\end{enumerate}
The verifier never receives $h$ as an oracle: whenever the folding verifier asks for $w_0^{(\beta)}(\xi)$, it queries the three source oracles at $\xi$, computes $h(\xi)$ by \eqref{eq:virtual-high-word}, and evaluates \eqref{eq:generic-ce-batch}.
\end{definition}

\begin{theorem}[Completeness of the generic protocol]
\label{thm:generic-ce-completeness}
Suppose
\[
  w_U=\ev_L(U),
  \qquad
  w_K=\ev_L(K),
  \qquad
  U,K\in\F[X]_{<N},
\]
and
\[
  v=[X^{N-1}]UK.
\]
Then the honest prover is accepted by $\Pi_{\mathrm{CE}}$ with probability $1$.
\end{theorem}

\begin{proof}
By \cref{thm:universal-split}, there exist unique $A,H\in\F[X]_{<N}$ such that
\[
  XUK=A+vX^N+X^{N+1}H.
\]
The honest prover sends $w_A=\ev_L(A)$.  By \cref{lem:virtual-high-honest}, the virtual word in \eqref{eq:generic-ce-high-word} is $h=\ev_L(H)$.  Therefore
\[
  w_0^{(\beta)}
  =
  \ev_L(U+\beta K+\beta^2A+\beta^3H)
  \in C_0
\]
for every $\beta\in\F$.  Completeness now follows from \cref{prop:folding-completeness}.
\end{proof}

\subsection{Batching soundness for coefficient extraction}
\label{sec:generic-ce-batching}

The batch \eqref{eq:generic-ce-batch} is a degree-$3$ curve in the challenge $\beta$.  Correlated agreement therefore gives simultaneous low-degree explanations for the four component words whenever too many values of $\beta$ make the batch close to the code.

\begin{lemma}[Coefficient-extraction batching lemma]
\label{lem:generic-ce-batching}
Let $w_U,w_K,w_A\in\F^L$, $v\in\F$, and let
\[
  h=\mathsf H[w_U,w_K,w_A;v].
\]
For $\beta\in\F$, define $w_0^{(\beta)}$ by \eqref{eq:generic-ce-batch}.  Let $G$ be the set of $\beta$ for which there exist a codeword $c_\beta\in C_0$ and a fibre-closed set $T_\beta\subseteq L$ with
\[
  |T_\beta|\ge(1-\delta)M
\]
such that $w_0^{(\beta)}=c_\beta$ on $T_\beta$.  Assume
\begin{equation}
  2N<(1-\delta)M.
  \label{eq:generic-degree-margin}
\end{equation}
If
\[
  |G|>3M,
\]
then the instance $(w_U,w_K;v)$ has a witness in $\mathcal R^{\delta}_{\mathrm{CE}}$.
\end{lemma}

\begin{proof}
For every $\beta\in G$, the batch is within Hamming distance $\delta$ of $C_0$.  Apply \cref{thm:correlated-agreement} with $t=3$ to the four words
\[
  u_0=w_U,
  \quad
  u_1=w_K,
  \quad
  u_2=w_A,
  \quad
  u_3=h.
\]
There exist a set $T\subseteq L$, $|T|\ge(1-\delta)M$, and polynomials $P_0,P_1,P_2,P_3\in\F[X]_{<N}$ such that
\[
  u_i(\xi)=P_i(\xi)
  \qquad
  (i=0,1,2,3;\ \xi\in T).
\]

Define
\[
  \Phi(X)
  =
  XP_0(X)P_1(X)-P_2(X)-vX^N-X^{N+1}P_3(X).
\]
By \cref{lem:universal-identity}, $\deg\Phi\le2N$.  For $\xi\in T$, the defining identity of the virtual high-part word gives
\[
  \Phi(\xi)
  =
  \xi w_U(\xi)w_K(\xi)-w_A(\xi)-v\xi^N-\xi^{N+1}h(\xi)
  =0.
\]
By \eqref{eq:generic-degree-margin}, $|T|>2N\ge\deg\Phi$, so \cref{lem:root-bound} implies $\Phi=0$.  Hence \cref{lem:universal-identity} gives
\[
  v=[X^{N-1}]P_0P_1.
\]

It remains to show fibre distance at most $\delta$ for the first two source words.  Correlated agreement gives only the common point set $T$, which need not be fibre-closed.  Choose any $\beta\in G$ whose pointwise discrepancy polynomial against the correlated-agreement codewords does not vanish outside $T$; we show that such a $\beta$ exists and that its fibre-closed agreement set is contained in $T$.

For $\xi\notin T$, define
\[
  d_\xi(Z)
  =
  \sum_{i=0}^{3}Z^i\bigl(u_i(\xi)-P_i(\xi)\bigr).
\]
By maximality of $T$ after enlarging it to contain every point where all four equalities hold, $d_\xi$ is a nonzero polynomial of degree at most $3$.  The union of its roots over $\xi\notin T$ has cardinality at most
\[
  3|L\setminus T|\le3\delta M<3M<|G|.
\]
Choose $\beta\in G$ outside this union, and let $T_\beta$ be its fibre-closed agreement set with $c_\beta$.  Put
\[
  \widehat c_\beta
  =
  \ev_L(P_0+\beta P_1+\beta^2P_2+\beta^3P_3)\in C_0.
\]
On $T\cap T_\beta$, both $c_\beta$ and $\widehat c_\beta$ equal $w_0^{(\beta)}$.  Since
\[
  |T\cap T_\beta|
  \ge(1-2\delta)M
  \ge \rho M
  =N,
\]
the two degree-$<N$ codewords are equal by \cref{prop:rs-parameters}(2).  If $\xi\in T_\beta$, then
\[
  d_\xi(\beta)
  =w_0^{(\beta)}(\xi)-\widehat c_\beta(\xi)
  =c_\beta(\xi)-\widehat c_\beta(\xi)
  =0.
\]
Our choice of $\beta$ therefore forces $\xi\in T$.  Thus $T_\beta\subseteq T$.

On the fibre-closed set $T_\beta$ of size at least $(1-\delta)M$, we have
\[
  w_U=\ev_L(P_0),
  \qquad
  w_K=\ev_L(P_1).
\]
By \cref{lem:hamming-fibre}, both fibre distances are at most $\delta$.  Taking $U=P_0$ and $K=P_1$ gives a witness in $\mathcal R^{\delta}_{\mathrm{CE}}$.
\end{proof}

\begin{remark}[Degree margin]
\label{rem:generic-degree-margin}
The only place where the Reed--Solomon rate enters the algebraic batching proof is \eqref{eq:generic-degree-margin}.  Since $N=\rho M$, it is enough that
\[
  2\rho<1-\delta.
\]
Under the unique-decoding choice $\delta\le(1-\rho)/2$, the worst case is satisfied whenever $3\rho<1$.  In particular, the standard choice $\rho\le1/4$ gives a strict margin.  Later protocols may impose additional Schwartz--Zippel terms, but they reuse the same degree margin.
\end{remark}

\begin{theorem}[Soundness of the generic coefficient-extraction protocol]
\label{thm:generic-ce-soundness}
Assume \eqref{eq:generic-degree-margin}.  If the instance $(w_U,w_K;v)$ has no witness in $\mathcal R^{\delta}_{\mathrm{CE}}$, then every prover is accepted by $\Pi_{\mathrm{CE}}$ with probability at most
\begin{equation}
  \varepsilon_{\mathrm{CE}}
  \defeq
  \frac{3M}{|\F|}
  +\varepsilon_{\mathrm{fold}}
  +(1-\delta)^\kappa.
  \label{eq:generic-ce-soundness}
\end{equation}
The same bound holds when sparse committed folding levels are used.
\end{theorem}

\begin{proof}
Fix first a deterministic prover, hence a fixed first message $w_A$.  Let $G$ be the set in \cref{lem:generic-ce-batching}.  Since the instance has no witness, that lemma gives
\[
  |G|\le3M.
\]
Thus a uniform $\beta\in\F$ lies outside $G$ except with probability at most $3M/|\F|$.

Fix $\beta\notin G$.  If
\[
  \Delta_{\mathrm{fib}}(w_0^{(\beta)},C_0)\le\delta,
\]
then the decoded codeword agrees with $w_0^{(\beta)}$ on a fibre-closed set of size at least $(1-\delta)M$, placing $\beta$ in $G$, a contradiction.  Hence
\[
  \Delta_{\mathrm{fib}}(w_0^{(\beta)},C_0)>\delta.
\]
By \cref{thm:folding-soundness}, conditioned on such a $\beta$ the folding test accepts with probability at most
\[
  \varepsilon_{\mathrm{fold}}+(1-\delta)^\kappa.
\]
Adding the probability $3M/|\F|$ of $\beta\in G$ proves \eqref{eq:generic-ce-soundness}.  The randomized-prover statement follows by conditioning on the prover's random tape.  The sparse-level case follows from \cref{prop:virtual-levels}.
\end{proof}

\subsection{Public kernels and omitted batch components}
\label{sec:generic-public-kernels}

The four-word batch is the universal form needed when both factors are committed.  If the second factor is public and known to be the evaluation of a degree-$<N$ polynomial, it need not be protected by correlated agreement.

\begin{corollary}[Public-kernel specialization]
\label{cor:public-kernel-specialization}
Let $K\in\F[X]_{<N}$ be public and suppose the verifier can evaluate $K(\xi)$ at every queried $\xi\in L$.  Consider the protocol obtained from \cref{def:generic-ce-protocol} by replacing $w_K$ by the public word $\ev_L(K)$ and batching only
\begin{equation}
  w_U+\beta w_A+\beta^2h.
  \label{eq:public-kernel-batch}
\end{equation}
Then, under the same degree-margin assumption, soundness error is at most
\begin{equation}
  \frac{2M}{|\F|}
  +\varepsilon_{\mathrm{fold}}
  +(1-\delta)^\kappa.
  \label{eq:public-kernel-soundness}
\end{equation}
\end{corollary}

\begin{proof}
The proof of \cref{lem:generic-ce-batching} applies with the degree-$2$ curve generated by $(w_U,w_A,h)$ and the fixed polynomial $K$.  Correlated agreement therefore needs more than $2M$ good values of $\beta$ rather than $3M$.  The identity polynomial becomes
\[
  XP_0K-P_1-vX^N-X^{N+1}P_2,
\]
and the same root-counting argument proves the coefficient claim.  The remainder of the proof of \cref{thm:generic-ce-soundness} is unchanged.
\end{proof}

\begin{remark}[What the generic theorem buys us]
\label{rem:generic-ce-roadmap}
The rest of the paper no longer needs a new proximity argument for each algebraic application.  It only needs to exhibit the appropriate second factor $K$ and to show how its values are obtained:
\begin{itemize}[leftmargin=*]
  \item for multilinear evaluation, $K$ is the product-form evaluation kernel of \cref{thm:table-evaluation-extraction};
  \item for a public linear functional, $K$ is the reversed generating polynomial of the public weight vector;
  \item for an arbitrary committed inner product, $K$ is the reversal of the second committed table and is supplied virtually by \cref{def:virtual-reversal};
  \item for Hadamard and later lookup/permutation arguments, coefficient extraction is combined with the DEEP quotient words of \cref{sec:deep-quotients}.
\end{itemize}
Thus the algebraic work of the later sections is separated cleanly from the Reed--Solomon proximity layer proved here.
\end{remark}

\section{Multilinear Evaluation and Public Linear Functionals}
\label{sec:linear-functionals}

The generic protocol of \cref{sec:generic-protocol} turns a coefficient claim
\(
  v=[X^{N-1}]UK
\)
into one Reed--Solomon proximity test.  In this section the second factor is public.  This gives the first family of applications of the framework: arbitrary public linear functionals of a committed table, multilinear evaluation, and summation over the Boolean hypercube.  No new proximity argument is required; the only work is to identify the public kernel and to account for its evaluation cost.

The point of treating these claims first is structural.  Many proof-system verifiers repeatedly ask for linear functionals of the same committed witness table.  In the multilinear language these appear as evaluations of a multilinear extension, Boolean-hypercube sums, or random linear combinations.  In the present framework all of them are instances of one coefficient functional with a public reversed generating polynomial.

\subsection{The public-linear-functional relation}
\label{sec:public-linear-functional-relation}

Let \(p:\B^n\to\F\) be a public table.  Write
\[
  K_p(X)
  =\sum_{j=0}^{N-1}p(j)X^{N-1-j}
\]
for its reversed generating polynomial, as in \cref{def:public-kernel}.  For a committed table \(t\), the desired value is
\[
  \Lambda_p(t)
  \defeq
  \sum_{w\in\B^n} t(w)p(w).
\]
By \cref{cor:public-linear-functional},
\begin{equation}
  \Lambda_p(t)
  = [X^{N-1}]V_t(X)K_p(X).
  \label{eq:linear-functional-coefficient}
\end{equation}

\begin{definition}[Public-linear-functional relation]
\label{def:public-linear-functional-relation}
Fix a public table \(p:\B^n\to\F\) with \(\deg K_p<N\).  The relation \(\mathcal R^{\delta}_{\mathrm{Lin}}[p]\) consists of pairs
\[
  ((w;v),t)
\]
such that
\begin{enumerate}[label=(\roman*),leftmargin=*]
  \item \(w\in\F^L\) and \(t:\B^n\to\F\);
  \item \(\Delta_{\mathrm{fib}}(w,\Enc_T(t))\le\delta\);
  \item \(v=\Lambda_p(t)\).
\end{enumerate}
The instance is \((w;v)\), the witness is \(t\), and the weight table \(p\) is public.
\end{definition}

\begin{definition}[Public-linear-functional protocol]
\label{def:public-linear-functional-protocol}
The protocol \(\Pi_{\mathrm{Lin}}[p]\) is the public-kernel specialization of \(\Pi_{\mathrm{CE}}\) from \cref{cor:public-kernel-specialization}, with
\[
  U=V_t,
  \qquad
  K=K_p.
\]
Thus the honest prover chooses \(A,H\in\F[X]_{<N}\) satisfying
\begin{equation}
  XV_tK_p=A+vX^N+X^{N+1}H,
  \label{eq:linear-functional-opening-identity}
\end{equation}
sends \(w_A=\ev_L(A)\), and after the verifier samples \(\beta\in\F\) the verifier defines the virtual word
\begin{equation}
  h(\xi)
  =\frac{\xi w(\xi)K_p(\xi)-w_A(\xi)-v\xi^N}{\xi^{N+1}}
  \qquad(\xi\in L).
  \label{eq:linear-functional-virtual-h}
\end{equation}
It then runs the folding test on
\begin{equation}
  w_0^{(\beta)}=w+\beta w_A+\beta^2h.
  \label{eq:linear-functional-batch}
\end{equation}
\end{definition}

\begin{theorem}[Completeness and soundness for public linear functionals]
\label{thm:public-linear-functional-protocol}
Assume
\begin{equation}
  2N<(1-\delta)M.
  \label{eq:linear-functional-degree-margin}
\end{equation}
Then \(\Pi_{\mathrm{Lin}}[p]\) is complete for \(\mathcal R^{\delta}_{\mathrm{Lin}}[p]\).  If an instance \((w;v)\) has no witness in \(\mathcal R^{\delta}_{\mathrm{Lin}}[p]\), every prover is accepted with probability at most
\begin{equation}
  \varepsilon_{\mathrm{Lin}}
  \defeq
  \frac{2M}{|\F|}
  +\varepsilon_{\mathrm{fold}}
  +(1-\delta)^\kappa.
  \label{eq:linear-functional-soundness}
\end{equation}
The same statement holds with sparse committed folding levels.
\end{theorem}

\begin{proof}
For an honest witness \(t\), \eqref{eq:linear-functional-coefficient} gives
\(
  v=[X^{N-1}]V_tK_p
\).
Hence \cref{thm:universal-split} gives unique polynomials \(A,H\in\F[X]_{<N}\) satisfying \eqref{eq:linear-functional-opening-identity}.  Evaluating that identity on \(L\) shows that the virtual word in \eqref{eq:linear-functional-virtual-h} is exactly \(\ev_L(H)\).  Therefore \eqref{eq:linear-functional-batch} is the evaluation of
\[
  V_t+\beta A+\beta^2H\in\F[X]_{<N},
\]
and completeness follows from completeness of the folding test.

For soundness, this is precisely \cref{cor:public-kernel-specialization}, with the public polynomial \(K=K_p\).  The degree hypothesis in that corollary is \eqref{eq:linear-functional-degree-margin}.  The stated error bound is therefore \eqref{eq:public-kernel-soundness}, and the sparse-level statement follows from the corresponding part of \cref{cor:public-kernel-specialization}.
\end{proof}

\begin{remark}[Where the verifier cost sits]
\label{rem:linear-functional-verifier-cost}
The folding verifier performs exactly the same authentication and folding checks for every public linear functional.  The only application-specific local work is evaluation of \(K_p(\xi)\) at the queried points.  Consequently, a family of weights is useful when the corresponding reversed generating polynomials form a succinct kernel family in the sense of \cref{def:succinct-kernel-family}.
\end{remark}

\subsection{Multilinear evaluation from a table commitment}
\label{sec:table-mle-opening}

Let \(f\in\M\), and let \(t=f|_{\B^n}\) be its table form.  Recall that the table commitment is
\[
  \Enc_T(t)=\ev_L(V_t),
\]
and that its multilinear extension is
\[
  \widetilde t(z)
  =\sum_{w\in\B^n}t(w)\eq(w,z).
\]
Since a multilinear polynomial is uniquely determined by its Boolean table, \(\widetilde t=f\).  The public weight table for evaluation at \(z\in\F^n\) is therefore
\[
  p_z(w)=\eq(w,z).
\]
Its public kernel is the product-form polynomial
\begin{equation}
  E_z^{\star}(X)
  =\prod_{i=1}^{n}
      \left(z_i+(1-z_i)X^{2^{i-1}}\right),
  \label{eq:section7-evaluation-kernel}
\end{equation}
and \cref{thm:table-evaluation-extraction} gives
\begin{equation}
  f(z)=\widetilde t(z)
  =[X^{N-1}]V_t(X)E_z^{\star}(X).
  \label{eq:section7-mle-coefficient}
\end{equation}

\begin{definition}[Evaluation relation]
\label{def:table-evaluation-relation}
The relation \(\mathcal R^{\delta}_{\mathrm{Eval}}\) consists of pairs
\[
  ((w;z,v),t)
\]
with \(w\in\F^L\), \(z\in\F^n\), \(v\in\F\), and \(t:\B^n\to\F\), such that
\[
  \Delta_{\mathrm{fib}}(w,\Enc_T(t))\le\delta
  \qquad\text{and}\qquad
  \widetilde t(z)=v.
\]
\end{definition}

\begin{definition}[Table-form evaluation protocol]
\label{def:table-evaluation-protocol}
The protocol \(\Pi_{\mathrm{Eval}}\) is \(\Pi_{\mathrm{Lin}}[p_z]\) with the kernel \(E_z^{\star}\) from \eqref{eq:section7-evaluation-kernel}.  Equivalently, the honest prover sends \(w_A=\ev_L(A)\) for
\begin{equation}
  XV_tE_z^{\star}=A+vX^N+X^{N+1}H,
  \label{eq:section7-evaluation-opening-identity}
\end{equation}
the verifier computes
\begin{equation}
  h(\xi)
  =\frac{\xi w(\xi)E_z^{\star}(\xi)-w_A(\xi)-v\xi^N}{\xi^{N+1}},
  \label{eq:section7-evaluation-virtual}
\end{equation}
and it folds the batch \(w+\beta w_A+\beta^2h\).
\end{definition}

\begin{corollary}[Sound table-form multilinear evaluation]
\label{cor:table-evaluation-soundness}
Under \eqref{eq:linear-functional-degree-margin}, \(\Pi_{\mathrm{Eval}}\) is complete for \(\mathcal R^{\delta}_{\mathrm{Eval}}\), and a false instance is accepted with probability at most
\begin{equation}
  \frac{2M}{|\F|}
  +\varepsilon_{\mathrm{fold}}
  +(1-\delta)^\kappa.
  \label{eq:table-evaluation-soundness}
\end{equation}
\end{corollary}

\begin{proof}
By \eqref{eq:section7-mle-coefficient}, evaluation at \(z\) is the public linear functional with weights \(p_z(w)=\eq(w,z)\).  Apply \cref{thm:public-linear-functional-protocol} with \(K_{p_z}=E_z^{\star}\).
\end{proof}

\begin{proposition}[Local verifier cost for multilinear evaluation]
\label{prop:table-evaluation-verifier-cost}
At one queried \(\xi\in L\), after \(\xi^{-1}\) and \(\xi^N\) are available, the application-specific work of \(\Pi_{\mathrm{Eval}}\), excluding the folding and authentication checks, consists of computing \(E_z^{\star}(\xi)\) and then \eqref{eq:section7-evaluation-virtual}.  In particular it uses \(O(n)\) field operations.  More explicitly, \(E_z^{\star}(\xi)\) can be evaluated with the operation count of \cref{prop:table-kernel-cost}.
\end{proposition}

\begin{proof}
This follows directly from the product formula \eqref{eq:section7-evaluation-kernel}, \cref{prop:table-kernel-cost}, and the single rational expression \eqref{eq:section7-evaluation-virtual}.  The denominator is nonzero because \(L\subseteq\F_q^\times\).
\end{proof}

\begin{remark}[No recursive variable elimination]
\label{rem:evaluation-no-sumcheck}
The protocol proves the final multilinear evaluation directly from the committed table.  It does not recursively eliminate the variables \(x_1,\ldots,x_n\) as Sumcheck does.  The exponential vector of Boolean equality weights is represented instead by the succinct product kernel \(E_z^{\star}\), while low degree of the committed and auxiliary words is enforced by the Reed--Solomon proximity test.  This is the first concrete instance of the transformation studied throughout the paper:
\[
  \text{multilinear relation}
  \longrightarrow
  \text{coefficient identity}
  \longrightarrow
  \text{Reed--Solomon proximity}.
\]
\end{remark}

\subsection{Hypercube sums}
\label{sec:hypercube-sum-protocol}

The sum of a committed table is the public linear functional with the all-ones weight vector.  Its kernel
\begin{equation}
  K_{\mathbf1}(X)
  =\sum_{j=0}^{N-1}X^j
  =\prod_{i=1}^{n}(1+X^{2^{i-1}})
  \label{eq:section7-sum-kernel}
\end{equation}
has the same binary product structure as the evaluation kernels.

\begin{definition}[Hypercube-sum relation]
\label{def:hypercube-sum-relation}
The relation \(\mathcal R^{\delta}_{\mathrm{Sum}}\) consists of pairs
\[
  ((w;s),t)
\]
such that
\[
  \Delta_{\mathrm{fib}}(w,\Enc_T(t))\le\delta
  \qquad\text{and}\qquad
  s=\sum_{u\in\B^n}t(u).
\]
\end{definition}

\begin{corollary}[Sumcheck-free hypercube sum]
\label{cor:hypercube-sum-protocol}
Let \(\Pi_{\mathrm{Sum}}\) be the public-linear-functional protocol with kernel \(K_{\mathbf1}\) from \eqref{eq:section7-sum-kernel}.  Under \eqref{eq:linear-functional-degree-margin}, the protocol is complete for \(\mathcal R^{\delta}_{\mathrm{Sum}}\) and has soundness error at most
\begin{equation}
  \frac{2M}{|\F|}
  +\varepsilon_{\mathrm{fold}}
  +(1-\delta)^\kappa.
  \label{eq:hypercube-sum-protocol-soundness}
\end{equation}
At each queried point, the kernel value can be computed in \(O(n)\) field operations by the product in \eqref{eq:section7-sum-kernel}.
\end{corollary}

\begin{proof}
The coefficient identity is \cref{cor:hypercube-sum-extraction}.  Apply \cref{thm:public-linear-functional-protocol}.  The local complexity follows by repeated squaring to obtain \(\xi^{2^{i-1}}\), followed by multiplication of the \(n\) factors \(1+\xi^{2^{i-1}}\).
\end{proof}

\begin{remark}[Relation with evaluation at the barycentre]
\label{rem:sum-as-evaluation}
If the characteristic is odd, then
\[
  \eq\!\left(w,\left(\tfrac12,\ldots,\tfrac12\right)\right)=2^{-n}
\]
for every \(w\in\B^n\).  Hence
\[
  \sum_{w\in\B^n}t(w)
  =2^n\widetilde t\!\left(\tfrac12,\ldots,\tfrac12\right).
\]
Thus \(\Pi_{\mathrm{Sum}}\) could also be obtained from \(\Pi_{\mathrm{Eval}}\).  The direct all-ones kernel is preferable conceptually because it displays the sum as an ordinary public linear functional and avoids the unnecessary scale factor.
\end{remark}

\subsection{Families of succinct public kernels}
\label{sec:succinct-public-kernel-applications}

The preceding examples use tensor-product weight vectors, but the protocol itself does not require tensor structure.  The exact condition needed for succinct verification is only that the public reversed generating polynomial can be evaluated cheaply.

\begin{proposition}[Public kernel family principle]
\label{prop:public-kernel-family-principle}
Let \(\Theta\) be a parameter set and let
\[
  p_\theta:\B^n\to\F,
  \qquad
  \theta\in\Theta,
\]
be a family of public tables.  Suppose there is an algorithm that, on input \((\theta,\xi)\in\Theta\times L\), computes
\[
  K_\theta(\xi)
  =\sum_{j=0}^{N-1}p_\theta(j)\xi^{N-1-j}
\]
in at most \(T_K(n)\) field operations, while \(K_\theta\in\F[X]_{<N}\) for every \(\theta\).  Then the claim
\[
  v=\sum_{w\in\B^n}t(w)p_\theta(w)
\]
for a table committed by \(w=\Enc_T(t)\) has a coefficient-extraction protocol with the soundness bound \eqref{eq:linear-functional-soundness} and application-specific verifier work \(O(T_K(n))\) per queried point.
\end{proposition}

\begin{proof}
For fixed \(\theta\), apply \cref{thm:public-linear-functional-protocol} with public kernel \(K_p=K_\theta\).  The only extra verifier computation relative to the generic public-kernel protocol is evaluation of \(K_\theta\) at queried points, which costs at most \(T_K(n)\) by hypothesis.
\end{proof}

\begin{example}[Selector weights]
\label{ex:selector-kernel}
For a public subset \(S\subseteq\{0,\ldots,N-1\}\), the selector functional
\[
  \sum_{j\in S}t(j)
\]
has kernel
\[
  K_S(X)=\sum_{j\in S}X^{N-1-j}.
\]
If \(S\) has size \(s\), direct evaluation costs \(O(s\log N)\) field operations by binary exponentiation, and less when the exponents or the set \(S\) have additional structure.  Thus sparse public selectors fit the same protocol even without tensor factorization.
\end{example}

\begin{example}[Affine combinations of public kernels]
\label{ex:affine-kernel-combination}
Suppose \(p^{(1)},\ldots,p^{(r)}\) are public weight tables with kernels \(K_1,\ldots,K_r\).  For public scalars \(\lambda_1,\ldots,\lambda_r\), the combined functional
\[
  t\longmapsto
  \sum_{a=1}^{r}\lambda_a\sum_w t(w)p^{(a)}(w)
\]
has kernel
\[
  K(X)=\sum_{a=1}^{r}\lambda_aK_a(X).
\]
Therefore public random linear combinations of already succinct kernels remain succinct.  This closure property will be useful when several public linear checks are batched later in the paper.
\end{example}

\subsection{What this section establishes}
\label{sec:linear-functionals-summary}

A single table-form Reed--Solomon commitment now supports all of the following without an algebraic Sumcheck reduction:
\begin{enumerate}[leftmargin=*]
  \item evaluation of the multilinear extension at any \(z\in\F^n\);
  \item summation over the Boolean hypercube;
  \item any public weighted sum whose reversed generating polynomial is efficiently evaluable.
\end{enumerate}
The proofs use no application-specific coding theorem beyond the generic public-kernel protocol.  The next step removes the assumption that the weight vector is public: the second factor will itself be committed, and inversion of the smooth domain will make its reversed generating polynomial available as a virtual word.  This yields the arbitrary committed inner-product argument.

\section{Committed Inner Products}
\label{sec:inner-product}

The public linear functionals of \cref{sec:linear-functionals} use a public second factor.  We now pass to the first genuinely bilinear relation: both vectors are committed.  Let
\[
  a,b:\B^n\to\F
\]
be tables and let
\[
  S=\sum_{w\in\B^n}a(w)b(w)
\]
be their inner product.  By \cref{thm:middle-coefficient-inner-product},
\begin{equation}
  S=[X^{N-1}]V_a(X)V_b^{\star}(X).
  \label{eq:ip-middle-coefficient}
\end{equation}
The second factor is not public, but \cref{sec:virtual-reversal} shows that its Reed--Solomon word is already available virtually from the commitment to $b$:
\begin{equation}
  \ev_L(V_b^{\star})(\xi)
  =\xi^{N-1}\ev_L(V_b)(\xi^{-1}).
  \label{eq:ip-reversal-value}
\end{equation}
Thus the arbitrary committed inner product is an instance of the generic coefficient-extraction protocol with no additional commitment for the reversed second factor.

This section states the relation, gives the protocol explicitly, and derives its completeness and soundness from the general machinery of \cref{sec:generic-protocol}.  We also record the decoding consequence needed later when inner products are used as subrelations in lookup, sparse-matrix, and R1CS arguments.

\subsection{The inner-product relation}
\label{sec:ip-relation}

\begin{definition}[Hypercube inner product]
\label{def:hypercube-inner-product}
For tables $a,b:\B^n\to\F$, define
\begin{equation}
  \langle a,b\rangle_{\B^n}
  \defeq
  \sum_{w\in\B^n}a(w)b(w).
  \label{eq:hypercube-inner-product}
\end{equation}
When no confusion is possible we write simply $\langle a,b\rangle$.
\end{definition}

\begin{definition}[Committed-inner-product relation]
\label{def:ip-relation}
Fix $0<\delta\le(1-\rho)/2$.  The relation $\mathcal R^{\delta}_{\mathrm{IP}}$ consists of pairs
\[
  ((w_a,w_b;S),(a,b))
\]
with $w_a,w_b\in\F^L$, $S\in\F$, and tables $a,b:\B^n\to\F$ such that
\begin{align}
  \Delta_{\mathrm{fib}}(w_a,\Enc_T(a))&\le\delta,
  \label{eq:ip-relation-distance-a}\\
  \Delta_{\mathrm{fib}}(w_b,\Enc_T(b))&\le\delta,
  \label{eq:ip-relation-distance-b}\\
  S&=\langle a,b\rangle_{\B^n}.
  \label{eq:ip-relation-value}
\end{align}
The instance is $(w_a,w_b;S)$ and the witness is the pair $(a,b)$.
\end{definition}

The relation is meaningful because, in the unique-decoding radius, each source word determines at most one nearby table by \cref{prop:unique-table-near-word}.  In particular, once $w_a$ and $w_b$ are fixed, the value in \eqref{eq:ip-relation-value} is determined whenever both decodings exist.

\begin{lemma}[Inner products are generic coefficient claims]
\label{lem:ip-as-generic-ce}
Let $w_a,w_b\in\F^L$ and define the virtual reversal
\[
  w_b^{\star}(\xi)=\xi^{N-1}w_b(\xi^{-1}).
\]
Then $(w_a,w_b;S)$ has a witness in $\mathcal R^{\delta}_{\mathrm{IP}}$ if and only if $(w_a,w_b^{\star};S)$ has a witness in the generic coefficient-extraction relation $\mathcal R^{\delta}_{\mathrm{CE}}$ of \cref{def:generic-ce-relation}.
\end{lemma}

\begin{proof}
Suppose first that $(a,b)$ witnesses $\mathcal R^{\delta}_{\mathrm{IP}}$.  Put
\[
  U=V_a,
  \qquad
  K=V_b^{\star}.
\]
By \cref{lem:reversal-commutes-evaluation},
\[
  \ev_L(K)=\Enc_T(b)^{\star}.
\]
By the fibre-distance isometry \eqref{eq:reversal-fibre-distance},
\[
  \Delta_{\mathrm{fib}}(w_b^{\star},\ev_L(K))
  =\Delta_{\mathrm{fib}}(w_b,\Enc_T(b))
  \le\delta.
\]
The first distance condition is
\[
  \Delta_{\mathrm{fib}}(w_a,\ev_L(U))
  =\Delta_{\mathrm{fib}}(w_a,\Enc_T(a))
  \le\delta.
\]
Finally, \cref{thm:middle-coefficient-inner-product} gives
\[
  [X^{N-1}]UK
  =[X^{N-1}]V_aV_b^{\star}
  =\langle a,b\rangle_{\B^n}
  =S.
\]
Hence $(U,K)$ witnesses $\mathcal R^{\delta}_{\mathrm{CE}}$.

Conversely, suppose $(U,K)$ witnesses $\mathcal R^{\delta}_{\mathrm{CE}}$ for $(w_a,w_b^{\star};S)$.  By \cref{prop:table-isomorphism}, there is a unique table $a$ with $V_a=U$, and a unique table $b$ with
\[
  V_b=K^{\star};
\]
equivalently $V_b^{\star}=K$, since reversal is an involution by \cref{lem:reversal-properties}.  The first generic distance condition is exactly \eqref{eq:ip-relation-distance-a}.  For the second, use \cref{lem:reversal-commutes-evaluation} and \eqref{eq:reversal-fibre-distance}:
\[
  \Delta_{\mathrm{fib}}(w_b,\Enc_T(b))
  =\Delta_{\mathrm{fib}}(w_b^{\star},\ev_L(V_b^{\star}))
  =\Delta_{\mathrm{fib}}(w_b^{\star},\ev_L(K))
  \le\delta.
\]
The generic coefficient condition and \cref{thm:middle-coefficient-inner-product} give
\[
  S=[X^{N-1}]UK
   =[X^{N-1}]V_aV_b^{\star}
   =\langle a,b\rangle_{\B^n}.
\]
Thus $(a,b)$ witnesses $\mathcal R^{\delta}_{\mathrm{IP}}$.
\end{proof}

\subsection{The protocol}
\label{sec:ip-protocol}

For honest tables $a,b$, \cref{thm:universal-split} applied to $V_a$ and $V_b^{\star}$ gives unique $A,H\in\F[X]_{<N}$ satisfying
\begin{equation}
  XV_a(X)V_b^{\star}(X)
  =A(X)+SX^N+X^{N+1}H(X).
  \label{eq:ip-opening-identity}
\end{equation}
Only $A$ is sent.  The reversed second factor and $H$ are both virtual.

\begin{definition}[Inner-product protocol $\Pi_{\mathrm{IP}}$]
\label{def:ip-protocol}
The instance is $(w_a,w_b;S)$.
\begin{enumerate}[label=\textnormal{\arabic*.},leftmargin=*]
  \item The prover sends an oracle word $w_A\in\F^L$.  The verifier samples $\beta\in\F$.
  \item The verifier defines the virtual reversed word
  \begin{equation}
    w_b^{\star}(\xi)
    =\xi^{N-1}w_b(\xi^{-1})
    \qquad(\xi\in L),
    \label{eq:ip-virtual-reversal}
  \end{equation}
  and the virtual high-part word
  \begin{equation}
    h(\xi)
    =\frac{\xi w_a(\xi)w_b^{\star}(\xi)-w_A(\xi)-S\xi^N}{\xi^{N+1}}.
    \label{eq:ip-virtual-high}
  \end{equation}
  It forms
  \begin{equation}
    w_0^{(\beta)}
    =w_a+\beta w_b^{\star}+\beta^2w_A+\beta^3h.
    \label{eq:ip-batch}
  \end{equation}
  \item The parties run $\Pi_{\mathrm{FT}}[\ell,\kappa]$, or any sparse-level variant of \cref{def:committed-folding-levels}, on $w_0^{(\beta)}$.
\end{enumerate}
The verifier never receives either $w_b^{\star}$ or $h$ as a committed oracle.  At a query point $\xi$ it obtains $w_b^{\star}(\xi)$ by reading $w_b(\xi^{-1})$, and then computes $h(\xi)$ from \eqref{eq:ip-virtual-high}.
\end{definition}

\begin{remark}[Fibre-local query pattern]
\label{rem:ip-fibre-locality}
A folding query needs the batch at $\xi$ and $-\xi$.  The source values of $w_a$ and $w_A$ are read on the fibre $\{\xi,-\xi\}$, while the reversed source values are obtained from $w_b$ on the inverse fibre
\[
  \{\xi^{-1},-\xi^{-1}\}.
\]
By \cref{lem:inversion-preserves-fibres}, this is again one fibre of $L$.  Thus reversal changes the location of the second commitment query but does not destroy the fibre-local structure required by the folding backend.
\end{remark}

\begin{theorem}[Completeness of the inner-product protocol]
\label{thm:ip-completeness}
If
\[
  w_a=\Enc_T(a),
  \qquad
  w_b=\Enc_T(b),
  \qquad
  S=\langle a,b\rangle_{\B^n},
\]
then the honest prover is accepted by $\Pi_{\mathrm{IP}}$ with probability $1$.
\end{theorem}

\begin{proof}
By \eqref{eq:ip-middle-coefficient},
\[
  S=[X^{N-1}]V_aV_b^{\star}.
\]
Hence \cref{thm:universal-split} gives unique $A,H\in\F[X]_{<N}$ satisfying \eqref{eq:ip-opening-identity}, and the prover sends $w_A=\ev_L(A)$.  By \cref{lem:reversal-commutes-evaluation}, the virtual word \eqref{eq:ip-virtual-reversal} equals $\ev_L(V_b^{\star})$.  Evaluating \eqref{eq:ip-opening-identity} at $\xi\in L$ gives
\[
  H(\xi)
  =\frac{\xi V_a(\xi)V_b^{\star}(\xi)-A(\xi)-S\xi^N}{\xi^{N+1}},
\]
so \eqref{eq:ip-virtual-high} equals $\ev_L(H)$.  Therefore, for every $\beta\in\F$,
\[
  w_0^{(\beta)}
  =\ev_L\!\left(V_a+\beta V_b^{\star}+\beta^2A+\beta^3H\right)
  \in C_0.
\]
Completeness follows from \cref{prop:folding-completeness}.
\end{proof}

\subsection{Soundness}
\label{sec:ip-soundness}

The protocol is exactly the generic coefficient-extraction protocol applied to the pair $(w_a,w_b^{\star})$.  The only point that must be transported back to the original instance is the decoding of the second commitment; the reversal isometry of \cref{lem:inversion-preserves-fibres} does precisely this.

\begin{theorem}[Soundness of committed inner products]
\label{thm:ip-soundness}
Assume
\begin{equation}
  2N<(1-\delta)M.
  \label{eq:ip-degree-margin}
\end{equation}
If $(w_a,w_b;S)$ has no witness in $\mathcal R^{\delta}_{\mathrm{IP}}$, then every prover is accepted by $\Pi_{\mathrm{IP}}$ with probability at most
\begin{equation}
  \varepsilon_{\mathrm{IP}}
  \defeq
  \frac{3M}{|\F|}
  +\varepsilon_{\mathrm{fold}}
  +(1-\delta)^\kappa.
  \label{eq:ip-soundness-bound}
\end{equation}
The same bound holds for every sparse-level folding variant covered by \cref{prop:virtual-levels}.
\end{theorem}

\begin{proof}
Suppose the original instance has no witness in $\mathcal R^{\delta}_{\mathrm{IP}}$.  By \cref{lem:ip-as-generic-ce}, the transformed instance
\[
  (w_a,w_b^{\star};S)
\]
has no witness in $\mathcal R^{\delta}_{\mathrm{CE}}$.  The protocol of \cref{def:ip-protocol} is exactly $\Pi_{\mathrm{CE}}$ of \cref{def:generic-ce-protocol} on this transformed instance: the second source word is evaluated virtually using \eqref{eq:ip-virtual-reversal}, while the generic high-part word is exactly \eqref{eq:ip-virtual-high}.  Therefore \cref{thm:generic-ce-soundness} gives \eqref{eq:ip-soundness-bound}.  The sparse-level statement follows from the same theorem.
\end{proof}

\begin{remark}[Why the field term is $3M/|\F|$]
\label{rem:ip-three-M}
Unlike a public linear functional, the second factor is itself committed and must be protected by correlated agreement.  The batch \eqref{eq:ip-batch} contains four degree-$<N$ words,
\[
  w_a,\quad w_b^{\star},\quad w_A,\quad h,
\]
and is a degree-$3$ curve in $\beta$.  Thus the correlated-agreement step uses $t=3$, which is the origin of the term $3M/|\F|$.  The reversal is virtual, but cryptographically it remains an independently decoded source word.
\end{remark}

\subsection{Decoding, uniqueness, and binding consequences}
\label{sec:ip-decoding}

The soundness theorem is stated as a relation statement.  For later composition it is useful to record explicitly what a witness means when the two commitment words are individually decodable.

\begin{proposition}[Decoded value is unique]
\label{prop:ip-decoded-value-unique}
Assume $0<\delta\le(1-\rho)/2$.  Let $w_a,w_b\in\F^L$.  There is at most one pair of tables $(a,b)$ satisfying
\[
  \Delta_{\mathrm{fib}}(w_a,\Enc_T(a))\le\delta,
  \qquad
  \Delta_{\mathrm{fib}}(w_b,\Enc_T(b))\le\delta.
\]
Consequently, there is at most one $S\in\F$ for which $(w_a,w_b;S)$ has a witness in $\mathcal R^{\delta}_{\mathrm{IP}}$.
\end{proposition}

\begin{proof}
Uniqueness of $a$ and $b$ follows independently from \cref{prop:unique-table-near-word}.  Once $a$ and $b$ are fixed, \eqref{eq:ip-relation-value} forces
\[
  S=\langle a,b\rangle_{\B^n},
\]
so $S$ is unique.
\end{proof}

\begin{corollary}[Evaluation binding for an inner-product value]
\label{cor:ip-binding}
Under the assumptions of \cref{thm:ip-soundness}, fix $w_a,w_b$ and two distinct values $S\ne S'$.  It is impossible that both instances $(w_a,w_b;S)$ and $(w_a,w_b;S')$ have witnesses in $\mathcal R^{\delta}_{\mathrm{IP}}$.  Moreover, if one of the two instances is false, every prover for that instance is accepted with probability at most $\varepsilon_{\mathrm{IP}}$.
\end{corollary}

\begin{proof}
The first statement is \cref{prop:ip-decoded-value-unique}.  The second is \cref{thm:ip-soundness}.
\end{proof}

\begin{proposition}[Polynomial-time witness extraction from decoded source words]
\label{prop:ip-decoding-extractor}
Assume $0<\delta\le(1-\rho)/2$.  Suppose a polynomial-time Reed--Solomon unique decoder is available for radius $\delta$.  Given $w_a,w_b$, one can in time polynomial in $M$ either reject because one source word has no codeword within distance $\delta$, or recover the unique tables $a,b$ satisfying the two distance conditions of $\mathcal R^{\delta}_{\mathrm{IP}}$.  The claimed value is then checked by computing $\langle a,b\rangle_{\B^n}$.
\end{proposition}

\begin{proof}
Apply the Reed--Solomon decoder separately to $w_a$ and $w_b$.  If it returns codewords $c_a,c_b\in C_0$, recover their unique degree-$<N$ polynomials using the inverse evaluation map of \cref{prop:rs-parameters}.  By \cref{prop:table-isomorphism}, their coefficients define unique tables $a,b$.  Check the two fibre distances explicitly and then compute the sum in \eqref{eq:hypercube-inner-product}.  All operations require time polynomial in $M$.
\end{proof}

The proposition is the algebraic extractor used later.  The round-by-round knowledge state needed for non-interactive compilation will be developed globally in the later soundness section; no additional algebraic extraction argument is required for inner products.

\subsection{Costs and relation to degree-two Sumcheck}
\label{sec:ip-costs}

The identity \eqref{eq:ip-middle-coefficient} replaces the standard degree-two Boolean-hypercube sum
\begin{equation}
  S=\sum_{w\in\B^n}a(w)b(w)
  \label{eq:ip-degree-two-sum}
\end{equation}
by one coefficient-extraction relation.  A conventional Sumcheck proof of \eqref{eq:ip-degree-two-sum} recursively eliminates the $n$ Boolean variables and terminates in evaluations of the multilinear extensions of $a$ and $b$ at one random point.  Here there are no Sumcheck round polynomials.  Instead, the prover computes the product $V_aV_b^{\star}$, splits it around degree $N$, commits to the low part $A$, and proves one Reed--Solomon proximity statement.

\begin{proposition}[Local verifier work]
\label{prop:ip-local-verifier-work}
Apart from the folding backend, one queried fibre for $\Pi_{\mathrm{IP}}$ requires:
\begin{enumerate}[label=(\roman*),leftmargin=*]
  \item the values of $w_a$ and $w_A$ on one fibre of $L$;
  \item the values of $w_b$ on the corresponding inverse fibre;
  \item $O(1)$ field operations per queried point to compute the reversal \eqref{eq:ip-virtual-reversal}, the high-part value \eqref{eq:ip-virtual-high}, and the batch \eqref{eq:ip-batch}.
\end{enumerate}
In particular, unlike multilinear evaluation in \cref{sec:table-mle-opening}, no $O(n)$ product-form kernel evaluation is needed.
\end{proposition}

\begin{proof}
Items (i) and (ii) follow from \cref{rem:ip-fibre-locality}.  Given the source values, \eqref{eq:ip-virtual-reversal} uses one multiplication by the known scalar $\xi^{N-1}$, and \eqref{eq:ip-virtual-high} and \eqref{eq:ip-batch} use a constant number of field operations.  Powers of the queried point are determined by the domain point and may be precomputed or updated along the folding query path.  No evaluation of a length-$n$ structured kernel appears.
\end{proof}

\begin{remark}[Prover arithmetic]
\label{rem:ip-prover-arithmetic}
The additional first-round algebra beyond the table commitments is the product $V_aV_b^{\star}$ and the evaluation/commitment of its low part $A$.  With FFT/NTT multiplication, the product is quasi-linear in $N$ over a suitable smooth domain.  Exact implementation costs are deliberately postponed to the benchmark section: the mathematical comparison with Sumcheck depends on whether product computation, low-part encoding, Merkle commitments, and the common folding backend are fused or shared with surrounding relations.
\end{remark}

\begin{remark}[Why this primitive matters later]
\label{rem:ip-later-applications}
The inner-product protocol is the committed analogue of the public-linear-functional protocol.  Later sections use it as a component rather than as an isolated end goal.  In particular:
\begin{itemize}[leftmargin=*]
  \item degree-two Sumcheck relations of the form $\sum_w a(w)b(w)$ are immediate instances;
  \item logarithmic-derivative lookup and multiset arguments reduce their global consistency equation to inner products;
  \item sparse matrix--vector products reduce, after random row compression, to an inner product between a compressed matrix row and the committed vector;
  \item bounded-degree pointwise arithmetic expressions are reduced to Hadamard relations followed by a final inner product or public linear functional.
\end{itemize}
Thus \cref{sec:inner-product} is the first step from coefficient extraction as an opening technique to coefficient extraction as a general proof-system primitive.
\end{remark}

\section{Committed Hadamard Relations}
\label{sec:hadamard}

The inner-product protocol of \cref{sec:inner-product} proves one bilinear scalar relation.  Many proof systems need the stronger coordinatewise statement
\[
  a(w)b(w)=c(w)
  \qquad\text{for every }w\in\B^n,
\]
for three committed tables $a,b,c:\B^n\to\F$.  Equivalently,
\[
  a\had b=c,
\]
where $(a\had b)(w)=a(w)b(w)$.  A direct proof of all $N$ coordinate equations would destroy succinctness.  Instead we compress them with one random geometric weighting and then reduce the resulting scalar bilinear claim to the coefficient-extraction primitive of the preceding sections.

The construction uses one additional idea.  After the verifier chooses the geometric challenge $\gamma$, the first factor becomes the polynomial $V_a(\gamma X)$.  Its values are not, in general, values of $V_a$ on the original Reed--Solomon domain $L$.  We therefore let the prover commit to a word for this scaled polynomial and bind it back to the original commitment by two DEEP quotient relations.  A third DEEP quotient binds the compressed output value $V_c(\gamma)$ to the commitment to $c$.  All low-degree words are then placed in one random curve and tested by one folding proximity test.

\subsection{Random geometric compression}
\label{sec:hadamard-geometric-compression}

\begin{definition}[Hadamard product]
\label{def:hadamard-product}
For tables $a,b:\B^n\to\F$, their \emph{Hadamard product} is the table
\[
  a\had b:\B^n\to\F,
  \qquad
  (a\had b)(w)=a(w)b(w).
\]
\end{definition}

For a table $d:\B^n\to\F$, recall that
\[
  V_d(X)=\sum_{w=0}^{N-1}d(w)X^w.
\]
Thus $d=0$ if and only if $V_d=0$.

\begin{lemma}[Geometric fingerprint]
\label{lem:hadamard-geometric-fingerprint}
Let $a,b,c:\B^n\to\F$ and put
\[
  d=a\had b-c.
\]
Then
\begin{equation}
  V_d(\gamma)
  =\sum_{w=0}^{N-1}\gamma^w\bigl(a(w)b(w)-c(w)\bigr)
  \qquad(\gamma\in\F).
  \label{eq:hadamard-fingerprint}
\end{equation}
If $a\had b\ne c$, then $V_d$ is a nonzero polynomial of degree at most $N-1$, and therefore
\begin{equation}
  \Pr_{\gamma\leftarrow\F}
  \left[
    \sum_{w=0}^{N-1}\gamma^wa(w)b(w)
    =\sum_{w=0}^{N-1}\gamma^wc(w)
  \right]
  \le \frac{N-1}{|\F|}.
  \label{eq:hadamard-fingerprint-bound}
\end{equation}
\end{lemma}

\begin{proof}
Equation \eqref{eq:hadamard-fingerprint} is the definition of $V_d$.  If $a\had b\ne c$, then $d(w)\ne0$ for at least one $w$, so $V_d\ne0$ because its coefficients are exactly the values $d(w)$.  Its degree is at most $N-1$.  The event in \eqref{eq:hadamard-fingerprint-bound} is exactly $V_d(\gamma)=0$, so the root bound of \cref{lem:root-bound} gives the result.
\end{proof}

The weighted bilinear sum in \eqref{eq:hadamard-fingerprint} is again a middle coefficient.

\begin{lemma}[Weighted Hadamard coefficient identity]
\label{lem:weighted-hadamard-coefficient}
For tables $a,b:\B^n\to\F$ and every $\gamma\in\F$,
\begin{equation}
  \sum_{w=0}^{N-1}\gamma^wa(w)b(w)
  =
  [X^{N-1}]\,V_a(\gamma X)V_b^{\star}(X).
  \label{eq:weighted-hadamard-coefficient}
\end{equation}
Moreover,
\begin{equation}
  \sum_{w=0}^{N-1}\gamma^wc(w)=V_c(\gamma).
  \label{eq:weighted-c-evaluation}
\end{equation}
\end{lemma}

\begin{proof}
The coefficient of $X^w$ in $V_a(\gamma X)$ is $\gamma^wa(w)$.  Apply \cref{thm:middle-coefficient-inner-product} to the two coefficient vectors of $V_a(\gamma X)$ and $V_b$.  This gives \eqref{eq:weighted-hadamard-coefficient}.  Equation \eqref{eq:weighted-c-evaluation} is the definition of $V_c$.
\end{proof}

Consequently, if we write
\begin{equation}
  Q_\gamma(X)\defeq V_a(\gamma X),
  \label{eq:hadamard-Q-honest}
\end{equation}
then the compressed Hadamard relation is
\begin{equation}
  V_c(\gamma)
  = [X^{N-1}]Q_\gamma(X)V_b^{\star}(X).
  \label{eq:hadamard-compressed-relation}
\end{equation}
The right-hand side can be proved by the committed inner-product/coefficient-extraction mechanism.  It remains to bind $Q_\gamma$ to $V_a$ and the scalar $V_c(\gamma)$ to $V_c$.

\subsection{The relation and protocol}
\label{sec:hadamard-protocol}

\begin{definition}[Committed Hadamard relation]
\label{def:hadamard-relation}
Fix $0<\delta\le(1-\rho)/2$.  The relation $\mathcal R^{\delta}_{\mathrm{Had}}$ consists of pairs
\[
  ((w_a,w_b,w_c),(a,b,c))
\]
with $w_a,w_b,w_c\in\F^L$ and tables $a,b,c:\B^n\to\F$ such that
\begin{align}
  \Delta_{\mathrm{fib}}(w_a,\Enc_T(a))&\le\delta,
  \label{eq:had-relation-a}\\
  \Delta_{\mathrm{fib}}(w_b,\Enc_T(b))&\le\delta,
  \label{eq:had-relation-b}\\
  \Delta_{\mathrm{fib}}(w_c,\Enc_T(c))&\le\delta,
  \label{eq:had-relation-c}\\
  a\had b&=c.
  \label{eq:had-relation-product}
\end{align}
The instance is $(w_a,w_b,w_c)$ and the witness is $(a,b,c)$.
\end{definition}

For the protocol below, the verifier samples $\gamma,\theta$ uniformly from $\F$.  If any of
\begin{equation}
  \gamma,\qquad \theta,\qquad \gamma\theta
  \label{eq:had-poles}
\end{equation}
lies in $L$, the verifier rejects.  This convention keeps the DEEP quotients of \cref{sec:deep-quotients} pole-free.  It costs only a small completeness term and can be removed by rejection-sampling challenges from the allowed set; see \cref{rem:had-perfect-completeness}.

\begin{definition}[Hadamard protocol $\Pi_{\mathrm{Had}}$]
\label{def:hadamard-protocol}
The instance is $(w_a,w_b,w_c)$.
\begin{enumerate}[label=\textnormal{\arabic*.},leftmargin=*]
  \item The verifier samples $\gamma\leftarrow\F$.  The prover sends an oracle word $w_Q\in\F^L$; honestly,
  \[
    w_Q=\ev_L(Q_\gamma),
    \qquad
    Q_\gamma(X)=V_a(\gamma X).
  \]
  \item The verifier samples $\theta\leftarrow\F$.  The prover sends scalars $y_1,y_3\in\F$ and an oracle word $w_A\in\F^L$; honestly,
  \begin{equation}
    y_1=V_a(\gamma\theta)=Q_\gamma(\theta),
    \qquad
    y_3=V_c(\gamma),
    \label{eq:had-honest-values}
  \end{equation}
  and $w_A=\ev_L(A)$, where
  \begin{equation}
    XQ_\gamma(X)V_b^{\star}(X)
    =A(X)+y_3X^N+X^{N+1}H(X)
    \label{eq:had-opening-identity}
  \end{equation}
  is the universal split of \cref{thm:universal-split}.
  \item If one of the three points in \eqref{eq:had-poles} lies in $L$, the verifier rejects.  Otherwise it samples $\beta\leftarrow\F$ and defines the following virtual words:
  \begin{align}
    w_b^{\star}(\xi)
      &\defeq \xi^{N-1}w_b(\xi^{-1}),
      \label{eq:had-wrev}\\
    h(\xi)
      &\defeq
      \frac{\xi w_Q(\xi)w_b^{\star}(\xi)-w_A(\xi)-y_3\xi^N}{\xi^{N+1}},
      \label{eq:had-high}\\
    q_a(\xi)
      &\defeq \mathsf Q[w_a;y_1,\gamma\theta](\xi),
      \label{eq:had-qa}\\
    q_Q(\xi)
      &\defeq \mathsf Q[w_Q;y_1,\theta](\xi),
      \label{eq:had-qQ}\\
    q_c(\xi)
      &\defeq \mathsf Q[w_c;y_3,\gamma](\xi).
      \label{eq:had-qc}
  \end{align}
  It batches the nine words
  \begin{equation}
  \begin{split}
    w_0^{(\beta)}={}&
      w_a+\beta w_c+\beta^2w_Q+\beta^3w_b^{\star}
      +\beta^4w_A+\beta^5h\\
      &+\beta^6q_a+\beta^7q_Q+\beta^8q_c.
  \end{split}
  \label{eq:had-batch}
  \end{equation}
  \item The parties run $\Pi_{\mathrm{FT}}[\ell,\kappa]$, or a sparse-level variant covered by \cref{prop:virtual-levels}, on $w_0^{(\beta)}$.
\end{enumerate}
\end{definition}

The scalar $y_1$ is deliberately reused in two DEEP relations.  If both quotient words decode correctly, they force
\[
  V_a(\gamma\theta)=y_1=Q(\theta),
\]
which links the newly committed scaled polynomial $Q$ to the original commitment to $a$ at a random point.  The third quotient similarly forces $V_c(\gamma)=y_3$.

\begin{theorem}[Completeness]
\label{thm:had-completeness}
Suppose
\[
  w_a=\Enc_T(a),
  \qquad
  w_b=\Enc_T(b),
  \qquad
  w_c=\Enc_T(c),
  \qquad
  a\had b=c.
\]
Conditioned on $\gamma,\theta,\gamma\theta\notin L$, the honest prover is accepted by $\Pi_{\mathrm{Had}}$ with probability $1$.  With uniform unrestricted $\gamma,\theta\leftarrow\F$ as in \cref{def:hadamard-protocol}, the honest acceptance probability is at least
\begin{equation}
  1-\frac{3M}{|\F|}.
  \label{eq:had-completeness-prob}
\end{equation}
\end{theorem}

\begin{proof}
Let $Q_\gamma=V_a(\gamma X)$.  By \cref{lem:weighted-hadamard-coefficient} and $a\had b=c$,
\[
  [X^{N-1}]Q_\gamma V_b^{\star}
  =\sum_w\gamma^wa(w)b(w)
  =\sum_w\gamma^wc(w)
  =V_c(\gamma)=y_3.
\]
Hence \cref{thm:universal-split} gives $A,H\in\F[X]_{<N}$ satisfying \eqref{eq:had-opening-identity}.  The virtual reversal \eqref{eq:had-wrev} is $\ev_L(V_b^{\star})$ by \cref{lem:reversal-commutes-evaluation}, and evaluating \eqref{eq:had-opening-identity} on $L$ shows that \eqref{eq:had-high} is $\ev_L(H)$.

Because the three poles lie outside $L$, \cref{lem:deep-honest} gives
\begin{align*}
  q_a&=\ev_L\!\left(\frac{V_a-y_1}{X-\gamma\theta}\right),\\
  q_Q&=\ev_L\!\left(\frac{Q_\gamma-y_1}{X-\theta}\right),\\
  q_c&=\ev_L\!\left(\frac{V_c-y_3}{X-\gamma}\right),
\end{align*}
and each quotient polynomial has degree at most $N-2$.  Therefore every one of the nine component words in \eqref{eq:had-batch} is a codeword in $C_0$, and so is their linear combination for every $\beta$.  Folding completeness, \cref{prop:folding-completeness}, proves conditional acceptance with probability $1$.

It remains to bound the pole event.  Since $|L|=M$,
\[
  \Pr[\gamma\in L]=\frac{M}{|\F|},
  \qquad
  \Pr[\theta\in L]=\frac{M}{|\F|}.
\]
For fixed $\gamma\ne0$, multiplication by $\gamma$ is a bijection of $\F$, so
\[
  \Pr[\gamma\theta\in L\mid\gamma]=\frac{M}{|\F|}.
\]
If $\gamma=0$, then $\gamma\theta=0\notin L$ because $L\subseteq\F^\times$.  A union bound gives \eqref{eq:had-completeness-prob}.
\end{proof}

\begin{remark}[Perfect completeness by rejection sampling]
\label{rem:had-perfect-completeness}
The completeness loss in \eqref{eq:had-completeness-prob} is not intrinsic.  One may sample $\gamma$ outside $L$ and then sample $\theta$ outside $L\cup\gamma^{-1}L$.  The protocol is then perfectly complete.  Soundness is obtained by replacing full-field root probabilities with the corresponding challenge-set denominators.  We keep unrestricted challenges in the main analysis because they give the cleanest algebraic error terms and make the Schwartz--Zippel contributions explicit.
\end{remark}

\subsection{A degree-eight batching lemma}
\label{sec:hadamard-batching}

The random word \eqref{eq:had-batch} is a degree-$8$ curve in $\beta$.  Correlated agreement therefore recovers simultaneous low-degree explanations for all nine words whenever too many $\beta$ values make the batch close to $C_0$.  The DEEP consistency lemma then turns three of those explanations into out-of-domain evaluation identities.

\begin{lemma}[Hadamard batching lemma]
\label{lem:had-batching}
Fix $\gamma,\theta\in\F$ with
\[
  \gamma,\theta,\gamma\theta\notin L,
\]
and fix the prover's messages $w_Q,y_1,y_3,w_A$.  Let the nine words $u_0,\ldots,u_8$ be
\[
  (u_0,\ldots,u_8)
  =
  (w_a,w_c,w_Q,w_b^{\star},w_A,h,q_a,q_Q,q_c),
\]
and let $w_0^{(\beta)}=\sum_{i=0}^8\beta^iu_i$.  Let $G$ be the set of $\beta\in\F$ for which there exist a codeword $c_\beta\in C_0$ and a fibre-closed set $T_\beta\subseteq L$ with
\[
  |T_\beta|\ge(1-\delta)M
\]
such that $w_0^{(\beta)}=c_\beta$ on $T_\beta$.  Assume
\begin{equation}
  2N<(1-\delta)M.
  \label{eq:had-degree-margin}
\end{equation}
If $|G|>8M$, then there exist tables $a,b,c:\B^n\to\F$ and a polynomial $\widehat Q\in\F[X]_{<N}$ such that
\begin{align}
  \Delta_{\mathrm{fib}}(w_a,\Enc_T(a))&\le\delta,
  \label{eq:had-batch-decode-a}\\
  \Delta_{\mathrm{fib}}(w_b,\Enc_T(b))&\le\delta,
  \label{eq:had-batch-decode-b}\\
  \Delta_{\mathrm{fib}}(w_c,\Enc_T(c))&\le\delta,
  \label{eq:had-batch-decode-c}\\
  \Delta_{\mathrm{fib}}(w_Q,\ev_L(\widehat Q))&\le\delta,
  \label{eq:had-batch-decode-Q}\\
  V_a(\gamma\theta)&=y_1=\widehat Q(\theta),
  \label{eq:had-batch-link-Q}\\
  V_c(\gamma)&=y_3=[X^{N-1}]\widehat Q(X)V_b^{\star}(X).
  \label{eq:had-batch-coefficient}
\end{align}
\end{lemma}

\begin{proof}
For every $\beta\in G$, the batch lies within Hamming distance $\delta$ of $C_0$.  Apply \cref{thm:correlated-agreement} with $t=8$ to $u_0,\ldots,u_8$.  After enlarging the common agreement set maximally, there exist a set $T\subseteq L$, $|T|\ge(1-\delta)M$, and polynomials
\[
  P_0,\ldots,P_8\in\F[X]_{<N}
\]
such that
\begin{equation}
  u_i(\xi)=P_i(\xi)
  \qquad
  (0\le i\le8,\ \xi\in T),
  \label{eq:had-correlated-agreement}
\end{equation}
and for every $\xi\notin T$ at least one of these nine equalities fails.

First recover the coefficient identity.  Define
\[
  \Phi(X)
  =XP_2(X)P_3(X)-P_4(X)-y_3X^N-X^{N+1}P_5(X).
\]
By \cref{lem:universal-identity}, $\deg\Phi\le2N$.  For every $\xi\in T$, the definition \eqref{eq:had-high} gives
\[
  \Phi(\xi)
  =\xi w_Q(\xi)w_b^{\star}(\xi)-w_A(\xi)-y_3\xi^N-\xi^{N+1}h(\xi)
  =0.
\]
Since $|T|\ge(1-\delta)M>2N$, root counting gives $\Phi=0$.  Hence
\begin{equation}
  y_3=[X^{N-1}]P_2P_3.
  \label{eq:had-batch-middle}
\end{equation}

Next use the quotient words.  The degree margin implies
\[
  |T|>2N\ge N+2,
\]
because $N\ge2$.  Apply \cref{lem:deep-consistency} on the common set $T$ to the triples
\[
  (P_0,P_6;\gamma\theta,y_1),
  \qquad
  (P_2,P_7;\theta,y_1),
  \qquad
  (P_1,P_8;\gamma,y_3).
\]
The three poles are outside $L$ by assumption.  We obtain
\begin{equation}
  P_0(\gamma\theta)=y_1,
  \qquad
  P_2(\theta)=y_1,
  \qquad
  P_1(\gamma)=y_3.
  \label{eq:had-batch-deep-values}
\end{equation}

It remains to turn pointwise agreement on $T$ into fibre-distance bounds.  For $\xi\notin T$, define the discrepancy polynomial
\[
  d_\xi(Z)
  =\sum_{i=0}^8 Z^i\bigl(u_i(\xi)-P_i(\xi)\bigr).
\]
By maximality of $T$, $d_\xi$ is nonzero and has degree at most $8$.  The union of all its roots over $\xi\notin T$ has size at most
\[
  8|L\setminus T|\le8\delta M<8M<|G|.
\]
Choose $\beta\in G$ outside this union and let $T_\beta$ be its fibre-closed agreement set with $c_\beta$.  Put
\[
  \widehat c_\beta
  =\ev_L\!\left(\sum_{i=0}^8\beta^iP_i\right)\in C_0.
\]
On $T\cap T_\beta$, both $c_\beta$ and $\widehat c_\beta$ equal $w_0^{(\beta)}$.  Moreover,
\[
  |T\cap T_\beta|
  \ge(1-2\delta)M
  \ge\rho M
  =N,
\]
where the middle inequality uses $\delta\le(1-\rho)/2$.  Two degree-$<N$ codewords agreeing on at least $N$ points are equal by \cref{prop:rs-parameters}(2).  Hence $c_\beta=\widehat c_\beta$.  If $\xi\in T_\beta$, then
\[
  d_\xi(\beta)
  =w_0^{(\beta)}(\xi)-\widehat c_\beta(\xi)
  =0.
\]
Our choice of $\beta$ forces $\xi\in T$, so $T_\beta\subseteq T$.

Therefore all nine component words agree with their polynomial explanations on the same fibre-closed set $T_\beta$ of size at least $(1-\delta)M$.  Define the tables and auxiliary polynomial by
\[
  V_a=P_0,
  \qquad
  V_c=P_1,
  \qquad
  \widehat Q=P_2,
  \qquad
  V_b=P_3^{\star}.
\]
The table isomorphism of \cref{prop:table-isomorphism} makes these definitions unique.  Since $P_3=V_b^{\star}$ and $w_b^{\star}=\ev_L(P_3)$ on $T_\beta$, reversal and inversion transport this agreement to $w_b=\Enc_T(b)$ on $T_\beta^{-1}$, a fibre-closed set of the same size by \cref{lem:inversion-preserves-fibres}.  The fibre-distance bounds \eqref{eq:had-batch-decode-a}--\eqref{eq:had-batch-decode-Q} now follow from \cref{lem:hamming-fibre}.

Finally, \eqref{eq:had-batch-deep-values} becomes
\[
  V_a(\gamma\theta)=y_1=\widehat Q(\theta),
  \qquad
  V_c(\gamma)=y_3,
\]
and \eqref{eq:had-batch-middle} becomes
\[
  y_3=[X^{N-1}]\widehat QV_b^{\star}.
\]
This proves \eqref{eq:had-batch-link-Q} and \eqref{eq:had-batch-coefficient}.
\end{proof}

\subsection{Soundness}
\label{sec:hadamard-soundness}

The batching lemma says that a successful $\beta$-curve determines a decoded auxiliary polynomial $\widehat Q$.  The earlier challenge $\theta$ then certifies that this polynomial is the intended scaling $V_a(\gamma X)$, except with a root-counting error.  Finally the challenge $\gamma$ tests the original coordinatewise Hadamard relation.

\begin{theorem}[Soundness of the Hadamard protocol]
\label{thm:had-soundness}
Assume
\[
  0<\delta\le\frac{1-\rho}{2},
  \qquad
  2N<(1-\delta)M.
\]
If $(w_a,w_b,w_c)$ has no witness in $\mathcal R^{\delta}_{\mathrm{Had}}$, then every prover is accepted by $\Pi_{\mathrm{Had}}$ with probability at most
\begin{equation}
  \varepsilon_{\mathrm{Had}}
  \defeq
  \frac{2(N-1)}{|\F|}
  +\frac{8M}{|\F|}
  +\varepsilon_{\mathrm{fold}}
  +(1-\delta)^\kappa.
  \label{eq:had-soundness-bound}
\end{equation}
The same bound holds for every sparse-level folding variant covered by \cref{prop:virtual-levels}.
\end{theorem}

\begin{proof}
It is enough to consider a deterministic prover; the acceptance probability of a randomized prover is an average over its internal randomness.

If one of $w_a,w_b,w_c$ has no table within fibre distance $\delta$, then \cref{lem:had-batching} cannot hold for any pole-free pair $(\gamma,\theta)$ and any round-two messages.  Hence the only way to pass is through the bounded set of $\beta$ values or through the folding/query error already accounted for below.  We may therefore assume that the three source words uniquely decode, by \cref{prop:unique-table-near-word}, to tables $a,b,c$.

Since the instance is false, these decoded tables satisfy
\[
  d\defeq a\had b-c\ne0.
\]
Let
\[
  D(X)=V_d(X)=\sum_{w=0}^{N-1}\bigl(a(w)b(w)-c(w)\bigr)X^w.
\]
By \cref{lem:hadamard-geometric-fingerprint}, $D$ is nonzero of degree at most $N-1$.  Thus at most $N-1$ choices of $\gamma$ satisfy $D(\gamma)=0$.

Fix now a value of $\gamma$ with $D(\gamma)\ne0$, and fix the prover's first-round word $w_Q$, which may depend on $\gamma$.  If $w_Q$ has no degree-$<N$ codeword within fibre distance $\delta$, then \cref{lem:had-batching} cannot hold for any pole-free $\theta$.  Otherwise let $\widehat Q\in\F[X]_{<N}$ be its unique decoded polynomial.

Suppose first that
\begin{equation}
  \widehat Q(X)=V_a(\gamma X).
  \label{eq:had-Q-correct}
\end{equation}
If, for some pole-free $\theta$ and the prover's corresponding round-two messages, the set $G$ of \cref{lem:had-batching} had size greater than $8M$, that lemma would give
\[
  V_c(\gamma)
  =[X^{N-1}]\widehat QV_b^{\star}.
\]
Using \eqref{eq:had-Q-correct} and \cref{lem:weighted-hadamard-coefficient}, this says
\[
  \sum_w\gamma^wc(w)
  =\sum_w\gamma^wa(w)b(w),
\]
that is, $D(\gamma)=0$, a contradiction.  Hence for a correctly decoded $\widehat Q$, every pole-free $\theta$ has $|G|\le8M$.

Suppose instead that
\[
  \widehat Q(X)\ne V_a(\gamma X).
\]
Then
\[
  R_\gamma(Y)=\widehat Q(Y)-V_a(\gamma Y)
\]
is a nonzero polynomial of degree at most $N-1$.  If a pole-free $\theta$ and the corresponding messages gave $|G|>8M$, \cref{lem:had-batching} would force
\[
  V_a(\gamma\theta)=\widehat Q(\theta),
\]
so $R_\gamma(\theta)=0$.  There are at most $N-1$ such values of $\theta$ by \cref{lem:root-bound}.

We have proved that, for a false decoded Hadamard relation, the probability over the first two unrestricted field challenges of reaching a pair $(\gamma,\theta)$ for which the actual prover messages yield $|G|>8M$ is at most
\begin{equation}
  \frac{N-1}{|\F|}+\frac{N-1}{|\F|}
  =\frac{2(N-1)}{|\F|}.
  \label{eq:had-bad-gamma-theta}
\end{equation}
The protocol rejects pole pairs, so they cannot increase this acceptance probability.

Condition on a remaining pair $(\gamma,\theta)$ and its fixed prover messages.  Then $|G|\le8M$, hence
\[
  \Pr_{\beta\leftarrow\F}[\beta\in G]\le\frac{8M}{|\F|}.
\]
For $\beta\notin G$, the codeword-chain argument of \cref{thm:generic-ce-soundness}, applied to the family $w_0^{(\beta)}$, gives at most
\[
  \varepsilon_{\mathrm{fold}}+(1-\delta)^\kappa
\]
additional acceptance probability: if the folding challenges are good and the final witness density were at least $1-\delta$, the accepted batch would agree with a codeword on a fibre-closed set of size at least $(1-\delta)M$, placing $\beta$ in $G$.  Combining this with \eqref{eq:had-bad-gamma-theta} gives \eqref{eq:had-soundness-bound}.  The sparse-level statement follows from \cref{prop:virtual-levels}.
\end{proof}

\begin{remark}[Origin of the error terms]
\label{rem:had-error-terms}
The four contributions in \eqref{eq:had-soundness-bound} have different meanings:
\begin{itemize}[leftmargin=*]
  \item $(N-1)/|\F|$ fingerprints the false coordinatewise relation $a\had b\ne c$ at $\gamma$;
  \item a second $(N-1)/|\F|$ binds the prover's auxiliary polynomial $\widehat Q$ to the intended scaling $V_a(\gamma X)$ at $\theta$;
  \item $8M/|\F|$ is the correlated-agreement field term for a degree-$8$ batch curve;
  \item $\varepsilon_{\mathrm{fold}}+(1-\delta)^\kappa$ is the common Reed--Solomon proximity error.
\end{itemize}
The protocol therefore separates the algebraic fingerprints from the code-proximity test cleanly.
\end{remark}

\subsection{Decoding and extraction consequences}
\label{sec:hadamard-extraction}

\begin{proposition}[Uniqueness of the decoded Hadamard witness]
\label{prop:had-witness-unique}
Assume $0<\delta\le(1-\rho)/2$.  For fixed words $w_a,w_b,w_c\in\F^L$, there is at most one triple of tables $(a,b,c)$ satisfying the three distance conditions \eqref{eq:had-relation-a}--\eqref{eq:had-relation-c}.  Hence the truth of $a\had b=c$ is determined by the three source words whenever all three unique decodings exist.
\end{proposition}

\begin{proof}
Apply \cref{prop:unique-table-near-word} independently to $w_a,w_b,w_c$.
\end{proof}

\begin{proposition}[Polynomial-time decoded witness extraction]
\label{prop:had-decoding-extractor}
Assume $0<\delta\le(1-\rho)/2$ and a polynomial-time Reed--Solomon unique decoder for radius $\delta$.  Given $(w_a,w_b,w_c)$, one can in time polynomial in $M$ either reject because one of the three words has no codeword within distance $\delta$, or recover the unique decoded tables $(a,b,c)$ and check the coordinatewise relation $a\had b=c$ in $O(N)$ additional field multiplications/comparisons.
\end{proposition}

\begin{proof}
Decode the three source words independently, recover their unique degree-$<N$ polynomials by \cref{prop:rs-parameters}, and interpret their coefficient vectors as tables using \cref{prop:table-isomorphism}.  Check the three fibre distances explicitly.  Finally compare $a(w)b(w)$ with $c(w)$ for every $w\in\B^n$.  All steps are polynomial in $M$ because $N\le M$.
\end{proof}

As in \cref{sec:ip-decoding}, the round-by-round knowledge state for the final Fiat--Shamir compilation will be treated globally later.  The algebraic extractor needed for the Hadamard relation is simply the unique decoding of the three source commitments.

\subsection{Costs and role in larger relations}
\label{sec:hadamard-costs}

The Hadamard protocol is more expensive than a single coefficient opening because it must authenticate the scaled polynomial $Q_\gamma$ and three out-of-domain values.  Nevertheless, all checks share one folding backend.

\begin{proposition}[Local verifier work]
\label{prop:had-local-verifier-work}
Apart from the folding backend, one queried point $\xi\in L$ requires source values from
\[
  w_a,\quad w_c,\quad w_Q,\quad w_A
\]
at $\xi$, and one source value $w_b(\xi^{-1})$ for the reversed factor.  From these values the verifier computes
\[
  w_b^{\star}(\xi),\quad h(\xi),\quad q_a(\xi),\quad q_Q(\xi),\quad q_c(\xi)
\]
with $O(1)$ field operations and three divisions.  The three divisions may be evaluated with one batched inversion per collection of queried points.
\end{proposition}

\begin{proof}
The reversal and high-part words are local by \cref{cor:reversal-fibre-local,lem:virtual-high-locality}.  Each DEEP quotient is local by \cref{lem:deep-locality}.  Their denominators depend only on public challenges and the queried point.  Standard batched inversion computes the inverses of all nonzero denominators using one field inversion and linear many multiplications.
\end{proof}

\begin{remark}[Why Hadamard is the multiplication gate of the framework]
\label{rem:had-multiplication-gate}
Linear relations among tables are handled directly by linearity of the encoding and by public linear functionals.  Coordinatewise multiplication is the only nonlinear gate needed to evaluate an ordinary arithmetic circuit pointwise over the Boolean table.  Thus, once $a\had b=c$ can be proved, any fixed pointwise arithmetic expression may be represented by auxiliary wire tables: addition and scalar multiplication are linear, multiplication gates are Hadamard relations, and a final scalar claim is an inner product or public linear functional.  This observation is the bridge from isolated coefficient identities to lookup, permutation, sparse-matrix, and R1CS relations developed below.
\end{remark}

\begin{remark}[Comparison with degree-two Sumcheck]
\label{rem:had-vs-sumcheck}
A conventional proof of the randomly compressed identity
\[
  \sum_w\gamma^w a(w)b(w)=\sum_w\gamma^w c(w)
\]
can be expressed through a degree-two Sumcheck followed by multilinear openings.  The present construction instead materializes the scaled coefficient polynomial $Q_\gamma$, binds it by DEEP quotients, reduces the bilinear scalar to one middle coefficient, and delegates global low-degree consistency to one Reed--Solomon folding test.  The mathematical tradeoff is therefore explicit: Sumcheck's $n$ recursive round polynomials are replaced by auxiliary low-degree words that can be batched into the same FRI-style proximity argument.  Whether this tradeoff is faster in practice is an implementation question addressed only by controlled benchmarks later in the paper.
\end{remark}

\section{One Folding Test for Heterogeneous Relations}
\label{sec:heterogeneous-batching}

The preceding sections deliberately presented public linear functionals, committed inner products, and Hadamard relations as separate protocols.  Their algebraic cores are different, but their proximity layer is the same: before the folding test starts, the verifier has a finite family of words that should all be Reed--Solomon codewords of degree bound $N$, together with local identities that tie those words to the claimed relation.  Running an independent folding test for every relation would repeat the most expensive verifier work and would obscure this common structure.

This section shows that the folding stage is paid only once.  After all relation-specific challenges and prover messages have been fixed, the verifier collects every low-degree word required by every claim, removes exact repetitions if desired, and places the resulting words on one random polynomial curve
\[
  W_\beta=\sum_{i=0}^{s-1}\beta^iu_i.
\]
One execution of the folding test on $W_\beta$ then protects the entire family.  Correlated agreement recovers simultaneous degree-$<N$ explanations for all $u_i$ on one large set whenever too many values of $\beta$ make the master batch close to the code.  The local coefficient, reversal, and DEEP identities from the preceding sections can therefore be applied simultaneously.

The statement is important for two reasons.  First, it makes the framework compositional: later lookup, permutation, sparse matrix, and R1CS arguments may be assembled from the existing primitives without introducing a new proximity proof for every component.  Second, the proximity and query errors are paid once.  The price of the master batch is only the correlated-agreement field term, which grows with the number of protected component words rather than with the number of coordinates represented by those words.

Throughout the section all claims use the same parameters $n,N,L,M,\rho,\delta,\ell,\kappa$ and the same base code
\[
  C_0=\RS_{\F}[L,N].
\]
We assume
\begin{equation}
  0<\delta\le\frac{1-\rho}{2}
  \qquad\text{and}\qquad
  2N<(1-\delta)M.
  \label{eq:heterogeneous-parameters}
\end{equation}
The second inequality is the degree margin already used in the generic coefficient-extraction and Hadamard arguments.

\subsection{Batch instances and protected words}
\label{sec:batch-protected-words}

We batch three classes of claims already proved in this manuscript.

\begin{definition}[Heterogeneous batch instance]
\label{def:heterogeneous-batch-instance}
A \emph{heterogeneous batch instance} $\mathfrak B$ consists of finite families of the following claims.
\begin{enumerate}[label=(\roman*),leftmargin=*]
  \item \emph{Public linear functionals.}  For every $j\in J_{\mathrm{Lin}}$, a public table $p_j:\B^n\to\F$ and an instance
  \[
    (w_j;v_j)
  \]
  of the relation $\mathcal R^{\delta}_{\mathrm{Lin}}[p_j]$ from \cref{def:public-linear-functional-relation}.

  \item \emph{Committed inner products.}  For every $j\in J_{\mathrm{IP}}$, an instance
  \[
    (w_{a,j},w_{b,j};S_j)
  \]
  of the relation $\mathcal R^{\delta}_{\mathrm{IP}}$ from \cref{def:ip-relation}.

  \item \emph{Hadamard relations.}  For every $j\in J_{\mathrm{Had}}$, an instance
  \[
    (w_{a,j},w_{b,j},w_{c,j})
  \]
  of the relation $\mathcal R^{\delta}_{\mathrm{Had}}$ from \cref{def:hadamard-relation}.
\end{enumerate}
The same oracle word may occur in several claims.  A \emph{joint witness} assigns one table to each distinct original commitment word appearing in $\mathfrak B$, consistently across all of its occurrences, and makes every individual relation true.
\end{definition}

The consistency clause is essential.  For example, if the same commitment word is used in two claimed evaluations, the batch does not allow the prover to decode it to two different tables.  In the unique-decoding regime of \eqref{eq:heterogeneous-parameters}, at most one such table can be within fibre distance $\delta$ of a fixed word.

We next describe the low-degree words that the master fold must protect.  These are exactly the words already used by the individual protocols, but without their individual final batching challenges.

\begin{definition}[Protected component list]
\label{def:protected-component-list}
Fix all prover messages and all relation-specific challenges that occur before the final folding stage.

For a public-linear-functional claim $j\in J_{\mathrm{Lin}}$, let $w_{A,j}$ be the prover's opening word and let
\[
  h_j(\xi)
  =\frac{\xi w_j(\xi)K_{p_j}(\xi)-w_{A,j}(\xi)-v_j\xi^N}{\xi^{N+1}}
\]
be its virtual high-part word.  Its protected words are
\begin{equation}
  w_j,\qquad w_{A,j},\qquad h_j.
  \label{eq:lin-protected-words}
\end{equation}

For an inner-product claim $j\in J_{\mathrm{IP}}$, let $w_{A,j}$ be the prover's opening word and define
\begin{align}
  w_{b,j}^{\star}(\xi)
    &\defeq \xi^{N-1}w_{b,j}(\xi^{-1}),
  \\
  h_j(\xi)
    &\defeq
    \frac{\xi w_{a,j}(\xi)w_{b,j}^{\star}(\xi)-w_{A,j}(\xi)-S_j\xi^N}{\xi^{N+1}}.
\end{align}
Its protected words are
\begin{equation}
  w_{a,j},\qquad w_{b,j}^{\star},\qquad w_{A,j},\qquad h_j.
  \label{eq:ip-protected-words}
\end{equation}

If $J_{\mathrm{Had}}\ne\varnothing$, all Hadamard claims use the same two verifier challenges $\gamma,\theta\in\F$.  For every $j\in J_{\mathrm{Had}}$, let $w_{Q,j},y_{1,j},y_{3,j},w_{A,j}$ be the prover's messages and define the five virtual words
\[
  w_{b,j}^{\star},\qquad
  h_j,\qquad
  q_{a,j},\qquad
  q_{Q,j},\qquad
  q_{c,j}
\]
exactly as in \eqref{eq:had-wrev}--\eqref{eq:had-qc}, using the common $\gamma,\theta$.  Its protected words are
\begin{equation}
\begin{split}
  &w_{a,j},\quad w_{c,j},\quad w_{Q,j},\quad w_{b,j}^{\star},\quad
  w_{A,j},\quad h_j,\\
  &q_{a,j},\quad q_{Q,j},\quad q_{c,j}.
\end{split}
\label{eq:had-protected-words}
\end{equation}

Concatenate all protected words and, optionally, delete exact repetitions.  Denote the resulting ordered list by
\begin{equation}
  \mathcal U=(u_0,\ldots,u_{s-1}).
  \label{eq:protected-list}
\end{equation}
The integer $s=|\mathcal U|$ is the \emph{protected-word count} of the batch.
\end{definition}

Without deduplication,
\begin{equation}
  s
  =3|J_{\mathrm{Lin}}|
   +4|J_{\mathrm{IP}}|
   +9|J_{\mathrm{Had}}|.
  \label{eq:naive-protected-count}
\end{equation}
Repeated original commitments and repeated virtual transformations may reduce this number.  The theorem below depends only on the actual value of $s$.

\subsection{The master batching protocol}
\label{sec:master-batching-protocol}

The timing matters.  Every oracle that appears in $\mathcal U$ must be fixed before the master challenge $\beta$ is sampled.  Hadamard's geometric challenges must in turn be fixed before their corresponding prover messages.

\begin{definition}[Heterogeneous batching protocol $\Pi_{\mathrm{Het}}$]
\label{def:heterogeneous-batching-protocol}
On a batch instance $\mathfrak B$ the parties proceed as follows.
\begin{enumerate}[label=\textnormal{\arabic*.},leftmargin=*]
  \item For every $j\in J_{\mathrm{Lin}}\cup J_{\mathrm{IP}}$, the prover sends the opening word $w_{A,j}$ required by the corresponding coefficient-extraction identity.

  \item If $J_{\mathrm{Had}}\ne\varnothing$, the verifier samples one shared challenge $\gamma\leftarrow\F$.  For every $j\in J_{\mathrm{Had}}$, the prover sends $w_{Q,j}$.  The verifier then samples one shared challenge $\theta\leftarrow\F$, and the prover sends $y_{1,j},y_{3,j},w_{A,j}$ for every Hadamard claim.  If
  \[
    \gamma\in L,
    \qquad
    \theta\in L,
    \qquad\text{or}\qquad
    \gamma\theta\in L,
  \]
  the verifier rejects.

  \item The verifier constructs all virtual words in \cref{def:protected-component-list}, obtaining the fixed list
  \[
    \mathcal U=(u_0,\ldots,u_{s-1}).
  \]
  It samples one master challenge $\beta\leftarrow\F$ and defines
  \begin{equation}
    W_\beta
    \defeq
    \sum_{i=0}^{s-1}\beta^iu_i.
    \label{eq:master-batch-word}
  \end{equation}

  \item The parties run exactly one instance of $\Pi_{\mathrm{FT}}[\ell,\kappa]$, or one sparse-level variant covered by \cref{prop:virtual-levels}, on $W_\beta$.
\end{enumerate}
\end{definition}

The individual challenges called $\beta$ in \cref{def:public-linear-functional-protocol,def:ip-protocol,def:hadamard-protocol} disappear in the heterogeneous protocol.  Their sole purpose was to protect the low-degree component words before a folding test.  The single challenge in \eqref{eq:master-batch-word} performs that role for all claims simultaneously.

\begin{theorem}[Completeness]
\label{thm:heterogeneous-completeness}
Suppose $\mathfrak B$ has a joint witness and every prover message is generated honestly from that witness.  If $J_{\mathrm{Had}}=\varnothing$, the verifier accepts $\Pi_{\mathrm{Het}}$ with probability $1$.  If $J_{\mathrm{Had}}\ne\varnothing$, then conditioned on
\[
  \gamma,\theta,\gamma\theta\notin L,
\]
the verifier accepts with probability $1$, and under unrestricted uniform challenges the acceptance probability is at least
\begin{equation}
  1-\frac{3M}{|\F|}.
  \label{eq:heterogeneous-completeness}
\end{equation}
\end{theorem}

\begin{proof}
For every public-linear-functional claim, the argument in \cref{thm:public-linear-functional-protocol} shows that the three words in \eqref{eq:lin-protected-words} are evaluations of polynomials of degree $<N$.  For every inner-product claim, \cref{thm:ip-completeness} shows the same for the four words in \eqref{eq:ip-protected-words}.  For every Hadamard claim, conditioned on the three poles lying outside $L$, the proof of \cref{thm:had-completeness} shows that all nine words in \eqref{eq:had-protected-words} lie in $C_0$.

Therefore every $u_i$ in \eqref{eq:protected-list} lies in $C_0$.  Since $C_0$ is linear, $W_\beta\in C_0$ for every value of the master challenge.  Folding completeness, \cref{prop:folding-completeness}, gives conditional acceptance with probability $1$.  When Hadamard claims are present, the same union bound as in \cref{thm:had-completeness} gives \eqref{eq:heterogeneous-completeness}.  With rejection sampling outside the three pole sets, the batched protocol is perfectly complete.
\end{proof}

\subsection{Simultaneous correlated agreement}
\label{sec:simultaneous-correlated-agreement}

The key lemma below is the many-word form of the batching arguments already used for one coefficient claim and one Hadamard relation.  It is stated independently of those applications because later sections will reuse it for lookup, permutation, and sparse algebra.

\begin{lemma}[Master correlated-agreement lemma]
\label{lem:master-correlated-agreement}
Let $u_0,\ldots,u_{s-1}\in\F^L$ with $s\ge2$, and let
\[
  W_\beta=\sum_{i=0}^{s-1}\beta^iu_i.
\]
Let $G\subseteq\F$ be the set of $\beta$ for which there exist a codeword $c_\beta\in C_0$ and a fibre-closed set $T_\beta\subseteq L$ satisfying
\begin{equation}
  |T_\beta|\ge(1-\delta)M
  \qquad\text{and}\qquad
  W_\beta=c_\beta\quad\text{on }T_\beta.
  \label{eq:master-good-beta}
\end{equation}
If
\begin{equation}
  |G|>(s-1)M,
  \label{eq:master-many-good}
\end{equation}
then there exist polynomials
\[
  P_0,\ldots,P_{s-1}\in\F[X]_{<N}
\]
and a fibre-closed set $T_\circ\subseteq L$ such that
\begin{equation}
  |T_\circ|\ge(1-\delta)M
  \label{eq:master-common-fibre-size}
\end{equation}
and
\begin{equation}
  u_i(\xi)=P_i(\xi)
  \qquad
  (0\le i<s,\ \xi\in T_\circ).
  \label{eq:master-common-fibre-agreement}
\end{equation}
\end{lemma}

\begin{proof}
For every $\beta\in G$, \eqref{eq:master-good-beta} implies
\[
  \Delta(W_\beta,C_0)\le\delta.
\]
Apply correlated agreement, \cref{thm:correlated-agreement}, with curve degree $t=s-1$.  Since \eqref{eq:master-many-good} is exactly the hypothesis $|G|>tM$, there exist polynomials $P_0,\ldots,P_{s-1}\in\F[X]_{<N}$ and a set $T\subseteq L$ with
\[
  |T|\ge(1-\delta)M
\]
such that
\begin{equation}
  u_i(\xi)=P_i(\xi)
  \qquad
  (0\le i<s,\ \xi\in T).
  \label{eq:master-ca-point-set}
\end{equation}
Enlarge $T$ maximally so that \eqref{eq:master-ca-point-set} remains true for every component.

For $\xi\notin T$, define the discrepancy polynomial
\begin{equation}
  d_\xi(Z)
  =\sum_{i=0}^{s-1}Z^i\bigl(u_i(\xi)-P_i(\xi)\bigr).
  \label{eq:master-discrepancy-polynomial}
\end{equation}
By maximality of $T$, at least one coefficient of $d_\xi$ is nonzero, so $d_\xi$ is a nonzero polynomial of degree at most $s-1$.  Hence the union of all roots of all $d_\xi$ with $\xi\notin T$ has cardinality at most
\begin{equation}
  (s-1)|L\setminus T|
  \le(s-1)\delta M
  <(s-1)M
  <|G|.
  \label{eq:master-root-union}
\end{equation}
Choose $\beta\in G$ outside this union, and let $T_\beta$ and $c_\beta$ satisfy \eqref{eq:master-good-beta}.  Define
\[
  \widehat c_\beta
  =\ev_L\!\left(\sum_{i=0}^{s-1}\beta^iP_i\right)\in C_0.
\]
On $T\cap T_\beta$, both $c_\beta$ and $\widehat c_\beta$ equal $W_\beta$.  Moreover,
\[
  |T\cap T_\beta|
  \ge(1-2\delta)M
  \ge\rho M
  =N,
\]
where the middle inequality is equivalent to $\delta\le(1-\rho)/2$.  Two degree-$<N$ Reed--Solomon codewords agreeing on at least $N$ points are equal by \cref{prop:rs-parameters}(2), so
\[
  c_\beta=\widehat c_\beta.
\]
If $\xi\in T_\beta$, then
\[
  d_\xi(\beta)
  =W_\beta(\xi)-\widehat c_\beta(\xi)
  =c_\beta(\xi)-\widehat c_\beta(\xi)
  =0.
\]
The choice of $\beta$ outside the root union forces $\xi\in T$.  Thus
\[
  T_\beta\subseteq T.
\]
Taking $T_\circ=T_\beta$ gives a fibre-closed set of size at least $(1-\delta)M$ on which every equality in \eqref{eq:master-common-fibre-agreement} holds.
\end{proof}

\begin{remark}[Why a fibre-closed common set is recovered]
\label{rem:master-fibre-closed}
The correlated-agreement theorem itself supplies only an ordinary point set $T$.  Fibre distance, however, is the metric used by the folding proof.  The second half of \cref{lem:master-correlated-agreement} is therefore necessary: one carefully chosen good master challenge transports simultaneous agreement to the fibre-closed set $T_\circ$.  This is the same mechanism used in \cref{lem:generic-ce-batching,lem:had-batching}, now isolated once for an arbitrary number of protected words.
\end{remark}

\subsection{Recovering all relations from one common set}
\label{sec:recover-all-relations}

The master lemma only says that all protected words have simultaneous low-degree explanations.  The next proposition states that the algebra developed earlier is sufficient to recover every individual relation from those explanations.  For Hadamard claims the two outer challenges $\gamma,\theta$ still play their fingerprinting role.

\begin{proposition}[Deterministic recovery on a common agreement set]
\label{prop:heterogeneous-deterministic-recovery}
Fix a transcript of \cref{def:heterogeneous-batching-protocol} through the messages immediately before the master challenge $\beta$, and assume the three Hadamard poles lie outside $L$.  Suppose there exist polynomials $P_i\in\F[X]_{<N}$ and a fibre-closed set $T_\circ$ satisfying \eqref{eq:master-common-fibre-size}--\eqref{eq:master-common-fibre-agreement} for the protected list $\mathcal U$.

Then the following hold.
\begin{enumerate}[label=(\roman*),leftmargin=*]
  \item Every public-linear-functional source word decodes within fibre distance $\delta$ to a table, and its claimed scalar equals the required public linear functional of that table.

  \item Every inner-product source pair decodes within fibre distance $\delta$ to tables $a,b$, and its claimed scalar equals $\langle a,b\rangle_{\B^n}$.

  \item For every Hadamard claim there exist decoded tables $a,b,c$ and a polynomial $\widehat Q\in\F[X]_{<N}$ satisfying the conclusions \eqref{eq:had-batch-decode-a}--\eqref{eq:had-batch-coefficient} of \cref{lem:had-batching}.
\end{enumerate}
All occurrences of the same original commitment word decode to the same table.
\end{proposition}

\begin{proof}
Consider first a public-linear-functional claim.  Let $P_U,P_A,P_H$ be the polynomial explanations of its three protected words.  On $T_\circ$, the defining formula for the virtual high-part word gives
\[
  \xi P_U(\xi)K_p(\xi)
  -P_A(\xi)-v\xi^N-\xi^{N+1}P_H(\xi)=0.
\]
Hence the polynomial
\[
  \Phi(X)=XP_U(X)K_p(X)-P_A(X)-vX^N-X^{N+1}P_H(X)
\]
has degree at most $2N$ and vanishes on more than $2N$ points by \eqref{eq:heterogeneous-parameters}.  Root counting gives $\Phi=0$, so \cref{lem:universal-identity} yields
\[
  v=[X^{N-1}]P_UK_p.
\]
The table is the unique $t$ with $V_t=P_U$, and agreement on the fibre-closed set $T_\circ$ gives
\[
  \Delta_{\mathrm{fib}}(w,\Enc_T(t))\le\delta.
\]
By \cref{cor:public-linear-functional}, the displayed coefficient is $\Lambda_p(t)$.  This proves (i).

For an inner-product claim, let $P_a,P_b^{\star},P_A,P_H$ explain its four protected words.  The virtual-high identity gives a degree-$\le2N$ polynomial
\[
  XP_aP_b^{\star}-P_A-SX^N-X^{N+1}P_H
\]
vanishing on $T_\circ$, hence identically zero.  Therefore
\[
  S=[X^{N-1}]P_aP_b^{\star}.
\]
Let $V_a=P_a$ and $V_b=(P_b^{\star})^{\star}$.  The middle-coefficient identity gives
\[
  S=\langle a,b\rangle_{\B^n}.
\]
Agreement of $w_a$ with $\ev_L(P_a)$ on $T_\circ$ gives the first fibre-distance bound.  For the second source word, the protected word is the virtual reversal of $w_b$.  By \cref{lem:inversion-preserves-fibres}, agreement of $w_b^{\star}$ with $\ev_L(P_b^{\star})$ on $T_\circ$ transports to agreement of $w_b$ with $\Enc_T(b)$ on the fibre-closed set $T_\circ^{-1}$ of the same cardinality.  This proves (ii).

For a Hadamard claim, restrict the common polynomial explanations to the nine words listed in \eqref{eq:had-protected-words}.  The proof of \cref{lem:had-batching} after its correlated-agreement step uses only simultaneous degree-$<N$ explanations on a set of more than $2N$ points.  We already have the stronger fibre-closed set $T_\circ$.  Repeating its coefficient-identity and three DEEP-consistency arguments gives exactly \eqref{eq:had-batch-decode-a}--\eqref{eq:had-batch-coefficient}.  This proves (iii).

Finally, if one original commitment word occurs in several claims, every occurrence agrees with its polynomial explanation on a fibre-closed set of size at least $(1-\delta)M$.  Equivalently, every resulting table commitment is within fibre distance $\delta$ of that same word.  By unique decoding, \cref{prop:unique-table-near-word}, there is at most one such table.  Thus the decoded witnesses are consistent across all occurrences.
\end{proof}

\subsection{Soundness of the heterogeneous batch}
\label{sec:heterogeneous-soundness}

The remaining issue is Hadamard's geometric fingerprint.  A common pair $(\gamma,\theta)$ suffices for arbitrarily many Hadamard claims: if the joint batch relation is false because at least one decoded Hadamard relation is false, fix one such relation.  The root-counting argument need only detect that one relation, so the algebraic error does not grow with $|J_{\mathrm{Had}}|$.

\begin{lemma}[One shared Hadamard fingerprint]
\label{lem:shared-hadamard-fingerprint}
Assume every distinct original commitment word in the batch has a unique decoded table within fibre distance $\delta$, and suppose at least one Hadamard relation among those decoded tables is false.  Fix a deterministic prover.  With shared independent uniform challenges $\gamma,\theta\leftarrow\F$, except with probability at most
\begin{equation}
  \frac{2(N-1)}{|\F|},
  \label{eq:shared-had-bad-pair}
\end{equation}
there cannot exist simultaneous low-degree explanations satisfying conclusion (iii) of \cref{prop:heterogeneous-deterministic-recovery} for every Hadamard claim.
\end{lemma}

\begin{proof}
Choose one false Hadamard claim deterministically, for example the least one in a fixed ordering, and let its decoded tables be $\widehat a,\widehat b,\widehat c$.  Put
\[
  D(X)=V_{\widehat a\had\widehat b-\widehat c}(X).
\]
Since the relation is false, $D$ is nonzero and has degree at most $N-1$.  Thus
\[
  \Pr[D(\gamma)=0]\le\frac{N-1}{|\F|}.
\]
Fix $\gamma$ with $D(\gamma)\ne0$.  The prover's corresponding word $w_Q$ is now fixed because it is sent after $\gamma$ and before $\theta$.  If it has no degree-$<N$ decoding within fibre distance $\delta$, simultaneous recovery is already impossible.  Otherwise let $\widehat Q$ denote its unique decoding.

If
\[
  \widehat Q(X)=V_{\widehat a}(\gamma X),
\]
then the two recovered identities
\[
  V_{\widehat c}(\gamma)
  =[X^{N-1}]\widehat QV_{\widehat b}^{\star}
  \qquad\text{and}\qquad
  \widehat Q=V_{\widehat a}(\gamma X)
\]
would imply
\[
  V_{\widehat c}(\gamma)
  =\sum_w\gamma^w\widehat a(w)\widehat b(w),
\]
that is, $D(\gamma)=0$, a contradiction.  Hence any simultaneous recovery must have
\[
  \widehat Q\ne V_{\widehat a}(\gamma X).
\]
Their difference is then a nonzero polynomial of degree at most $N-1$, fixed before $\theta$ is sampled.  The recovered DEEP equalities require
\[
  \widehat Q(\theta)=V_{\widehat a}(\gamma\theta),
\]
which occurs for at most $(N-1)/|\F|$ of the possible $\theta$.  A union bound with the bad-$\gamma$ event proves \eqref{eq:shared-had-bad-pair}.
\end{proof}

\begin{theorem}[Soundness of one-fold heterogeneous batching]
\label{thm:heterogeneous-soundness}
Let $\mathfrak B$ be a heterogeneous batch instance with protected-word count $s\ge2$, and assume \eqref{eq:heterogeneous-parameters}.  Suppose $\mathfrak B$ has no joint witness.  Then every prover is accepted by $\Pi_{\mathrm{Het}}$ with probability at most
\begin{equation}
  \varepsilon_{\mathrm{Het}}
  \defeq
  \mathbf 1_{J_{\mathrm{Had}}\ne\varnothing}
  \frac{2(N-1)}{|\F|}
  +\frac{(s-1)M}{|\F|}
  +\varepsilon_{\mathrm{fold}}
  +(1-\delta)^\kappa.
  \label{eq:heterogeneous-soundness}
\end{equation}
The same bound holds for every sparse-level folding variant covered by \cref{prop:virtual-levels}.
\end{theorem}

\begin{proof}
It suffices to fix a deterministic prover; a randomized prover is an average over deterministic ones.

First consider the transcript up to the master challenge $\beta$.  If some original commitment word has no table within fibre distance $\delta$, then \cref{lem:master-correlated-agreement,prop:heterogeneous-deterministic-recovery} imply that the set $G$ of master challenges satisfying \eqref{eq:master-good-beta} has size at most $(s-1)M$: otherwise simultaneous agreement would decode that commitment, a contradiction.

Assume therefore that every original commitment word has a unique decoded table.  If $J_{\mathrm{Had}}\ne\varnothing$ and at least one decoded Hadamard relation is false, \cref{lem:shared-hadamard-fingerprint} shows that, except with probability at most $2(N-1)/|\F|$ over the shared pair $(\gamma,\theta)$, simultaneous recovery of all Hadamard claims is impossible.  If no decoded Hadamard relation is false, this outer error is unnecessary.  In either case, outside the event charged by the first term of \eqref{eq:heterogeneous-soundness}, the absence of a joint witness means that at least one linear-functional or inner-product claim is false, or one Hadamard claim cannot satisfy its recovered identities.  By \cref{prop:heterogeneous-deterministic-recovery}, the master good set must again satisfy
\begin{equation}
  |G|\le(s-1)M.
  \label{eq:heterogeneous-good-set-bound}
\end{equation}
Consequently the probability that the uniform master challenge lies in $G$ is at most
\[
  \frac{(s-1)M}{|\F|}.
\]

Condition on a master challenge $\beta\notin G$.  The definition of $G$ says that $W_\beta$ has no fibre-closed agreement set of relative size $1-\delta$ with a codeword in $C_0$.  The folding soundness theorem, \cref{thm:folding-soundness}, applies exactly as in the proof of \cref{thm:generic-ce-soundness}: bad folding challenges contribute at most $\varepsilon_{\mathrm{fold}}$, and if the final witness density remains below $1-\delta$, all $\kappa$ independent queries land in the accepting set with probability at most $(1-\delta)^\kappa$.  Adding the outer fingerprint term and the master-challenge term proves \eqref{eq:heterogeneous-soundness}.  The sparse-level statement follows from \cref{prop:virtual-levels}.
\end{proof}

\begin{remark}[The field term and the number of claims]
\label{rem:heterogeneous-field-term}
The master field term is
\[
  \frac{(s-1)M}{|\F|},
\]
where $s$ counts protected low-degree words, not coordinates and not arithmetic gates.  One Hadamard relation already certifies $N$ coordinatewise multiplications at once.  Likewise one inner-product relation certifies an $N$-term bilinear sum.  Later sections aggregate lookup, sparse-matrix, and R1CS constraints into a small number of such table relations before invoking the master fold.

The bound is useful when $(s-1)M\ll|\F|$.  If a concrete application creates too many protected words, one may enlarge the challenge field, reduce $s$ by deduplicating shared words, or introduce a hierarchy of algebraic reductions before the final proximity batch.  No such optimization is needed for the correctness theorem.
\end{remark}

\begin{remark}[Why shared $\gamma$ and $\theta$ matter]
\label{rem:shared-had-challenges}
Using one pair $(\gamma,\theta)$ for all Hadamard claims prevents the root-counting term from scaling linearly with the number of Hadamard relations.  If the conjunction is false, choose one false decoded relation.  Detecting that relation costs at most two degree-$(N-1)$ root events, irrespective of how many other Hadamard claims are simultaneously present.  The proximity cost is then paid once through the master batch.
\end{remark}

\subsection{Extraction and verifier cost}
\label{sec:heterogeneous-extraction-cost}

The master batch also gives the correct witness structure for later knowledge arguments.

\begin{corollary}[Decoding extraction from an accepting batch]
\label{cor:heterogeneous-extraction}
Assume the hypotheses of \cref{thm:heterogeneous-soundness}.  If the acceptance probability of a prover exceeds the right-hand side of \eqref{eq:heterogeneous-soundness}, then there exists a transcript prefix and simultaneous degree-$<N$ explanations for all protected words from which one obtains a joint witness for $\mathfrak B$ by uniquely decoding each distinct original commitment word.  In particular, every repeated occurrence of one commitment is assigned the same table.
\end{corollary}

\begin{proof}
The contrapositive is \cref{thm:heterogeneous-soundness}.  More concretely, acceptance above the stated error forces the outer Hadamard challenges outside their bad root events and forces more than $(s-1)M$ master challenges to admit the agreement structure required by the folding witness argument.  Then \cref{lem:master-correlated-agreement} gives simultaneous polynomial explanations on one fibre-closed set, and \cref{prop:heterogeneous-deterministic-recovery} recovers all scalar and Hadamard identities.  Unique decoding of the original commitments gives the joint witness.
\end{proof}

\begin{remark}[One folding verifier, many local reads]
\label{rem:heterogeneous-verifier-cost}
The protocol executes one folding chain and one set of $\kappa$ final query locations.  At a queried fibre, however, the verifier must locally evaluate every component of $W_\beta$.  Thus the arithmetic work for forming the master word is $O(s)$ field operations per queried point, in addition to the application-specific local formulas such as public-kernel evaluation, reversal, and DEEP quotients.

In a Merkle compilation, the number of authentication paths depends on the commitment layout rather than directly on $s$.  Words committed in the same round and queried at the same positions can be stored as vector-valued leaves and share authentication paths.  Words fixed in different transcript rounds require separate roots.  We defer the exact hash and proof-size accounting to the implementation and complexity sections; the theorem here is an IOP-level statement and does not assume a particular Merkle layout.
\end{remark}

\subsection{Consequence for the remainder of the paper}
\label{sec:batching-consequence}

The role of \cref{thm:heterogeneous-soundness} is to separate algebra from proximity.  To add a new relation to the framework, it is enough to do three things:
\begin{enumerate}[label=(\arabic*),leftmargin=*]
  \item reduce the relation to a finite family of degree-$<N$ table words and local algebraic identities;
  \item prove that simultaneous low-degree explanations on a set of more than the relevant identity degree force the claimed relation, up to explicit low-degree fingerprint events;
  \item add those words to the protected list $\mathcal U$.
\end{enumerate}
The folding argument itself is unchanged.  In particular, the lookup and permutation constructions of the next sections will reduce their inverse relations to Hadamard constraints and their global equality to linear or inner-product constraints.  Their protected words can then be inserted into the same master list and verified by one folding test.

\section{Logarithmic-Derivative Lookup Arguments}
\label{sec:lookup}

Lookup arguments prove that the entries of one committed table occur in another table.  The logarithmic-derivative method rewrites this set-theoretic statement as an identity between sums of reciprocals.  In the present framework those reciprocals are committed tables, their defining inverse equations are Hadamard relations, and the remaining equality is an inner-product identity.  Consequently a lookup can be reduced to the primitives of \cref{sec:linear-functionals,sec:inner-product,sec:hadamard} and protected by the single folding test of \cref{sec:heterogeneous-batching}.

The use of logarithmic derivatives for lookup arguments goes back to the multivariate construction of Hab\"ock~\cite{Habock23} and subsequent GKR-based refinements~\cite{PapiniHabock23}.  Our point here is not the rational identity itself, but its realization on the table-form Kronecker commitments of this paper without an algebraic Sumcheck: the two inverse relations are Hadamard checks, while the logarithmic-derivative equality is a pair of coefficient-extraction claims.

Throughout this section the query table, lookup table, and multiplicity table all have length $N=2^n$.  Unequal lengths can be embedded into a common power-of-two length by the usual tagging and padding techniques; the equal-length formulation keeps the algebra and soundness statement transparent.

\subsection{The algebraic lookup relation}
\label{sec:lookup-relation}

Let $f,t,m:\B^n\to\F$.  The table $f$ contains the values to be looked up, $t$ is the lookup table, and $m$ is a multiplicity witness attached to the rows of $t$.

\begin{definition}[Residue multiplicities]
\label{def:lookup-residue-multiplicities}
For $a\in\F$, define
\begin{align}
  \nu_f(a)
  &\defeq
  \sum_{w\in\B^n:\,f(w)=a}1_{\F},
  \label{eq:lookup-query-multiplicity}\\
  \mu_{t,m}(a)
  &\defeq
  \sum_{u\in\B^n:\,t(u)=a}m(u).
  \label{eq:lookup-table-multiplicity}
\end{align}
The sums are taken in $\F$.  We say that $(f,t,m)$ satisfies the \emph{algebraic lookup relation} if
\begin{equation}
  \nu_f(a)=\mu_{t,m}(a)
  \qquad\text{for every }a\in\F.
  \label{eq:algebraic-lookup-relation}
\end{equation}
\end{definition}

The formulation in \eqref{eq:algebraic-lookup-relation} permits repeated values in $t$: only the sum of the row multiplicities attached to all occurrences of the same value matters.  It also separates the algebraic relation from the integer interpretation of multiplicities.  The latter follows under the natural characteristic condition recorded below.

\begin{proposition}[Ordinary lookup as a consequence]
\label{prop:ordinary-lookup-consequence}
Assume that the characteristic $p=\operatorname{char}(\F)$ satisfies $p>N$.  If $(f,t,m)$ satisfies \eqref{eq:algebraic-lookup-relation}, then every value occurring in $f$ occurs in $t$:
\begin{equation}
  \{f(w):w\in\B^n\}
  \subseteq
  \{t(u):u\in\B^n\}.
  \label{eq:lookup-support-inclusion}
\end{equation}
\end{proposition}

\begin{proof}
Let $a=f(w_0)$ for some $w_0\in\B^n$.  Then the integer number of occurrences of $a$ in $f$ lies in $[1,N]$.  Since $p>N$, its image in $\F$ is nonzero, so $\nu_f(a)\ne0$.  By \eqref{eq:algebraic-lookup-relation}, $\mu_{t,m}(a)\ne0$.  If no row of $t$ contained $a$, the defining sum \eqref{eq:lookup-table-multiplicity} would be empty and hence zero, a contradiction.  Thus $a$ occurs in $t$.
\end{proof}

We next express \eqref{eq:algebraic-lookup-relation} as a rational-function identity.  This is the logarithmic-derivative step.

\begin{definition}[Lookup rational function]
\label{def:lookup-rational-function}
For tables $f,t,m:\B^n\to\F$, define
\begin{equation}
  R_{f,t,m}(Z)
  \defeq
  \sum_{w\in\B^n}\frac{1}{Z-f(w)}
  -
  \sum_{u\in\B^n}\frac{m(u)}{Z-t(u)}
  \in \F(Z).
  \label{eq:lookup-rational-function}
\end{equation}
\end{definition}

\begin{lemma}[Residue form]
\label{lem:lookup-residue-form}
Let
\begin{equation}
  S_{f,t}
  \defeq
  \operatorname{im}(f)\cup\operatorname{im}(t)
  \subseteq\F.
  \label{eq:lookup-value-set}
\end{equation}
Then
\begin{equation}
  R_{f,t,m}(Z)
  =
  \sum_{a\in S_{f,t}}
  \frac{\nu_f(a)-\mu_{t,m}(a)}{Z-a}.
  \label{eq:lookup-residue-expansion}
\end{equation}
Consequently, $(f,t,m)$ satisfies the algebraic lookup relation if and only if
\begin{equation}
  R_{f,t,m}(Z)=0
  \qquad\text{in }\F(Z).
  \label{eq:lookup-rational-zero}
\end{equation}
\end{lemma}

\begin{proof}
Group the terms of the first sum in \eqref{eq:lookup-rational-function} according to the common value $a=f(w)$, and group the second sum according to $a=t(u)$.  The coefficient of $(Z-a)^{-1}$ is exactly \eqref{eq:lookup-query-multiplicity} minus \eqref{eq:lookup-table-multiplicity}, giving \eqref{eq:lookup-residue-expansion}.

If \eqref{eq:algebraic-lookup-relation} holds, every numerator in \eqref{eq:lookup-residue-expansion} is zero, so \eqref{eq:lookup-rational-zero} follows.  Conversely, suppose \eqref{eq:lookup-rational-zero} holds.  Fix $a_0\in S_{f,t}$ and multiply \eqref{eq:lookup-residue-expansion} by $Z-a_0$.  Evaluating the resulting rational function at $Z=a_0$ leaves only the $a=a_0$ term and gives
\[
  \nu_f(a_0)-\mu_{t,m}(a_0)=0.
\]
This holds for every $a_0\in S_{f,t}$, and outside $S_{f,t}$ both multiplicities are zero.  Hence \eqref{eq:algebraic-lookup-relation} holds.
\end{proof}

For soundness we need a quantitative form of the preceding equivalence.

\begin{lemma}[Bad lookup challenges]
\label{lem:lookup-bad-challenges}
Assume that $(f,t,m)$ does not satisfy \eqref{eq:algebraic-lookup-relation}.  Then there exists a nonzero polynomial
\begin{equation}
  P_{f,t,m}(Z)\in\F[Z]
  \label{eq:lookup-numerator-polynomial}
\end{equation}
with
\begin{equation}
  \deg P_{f,t,m}\le |S_{f,t}|-1\le2N-1
  \label{eq:lookup-numerator-degree}
\end{equation}
such that for every $\zeta\in\F\setminus S_{f,t}$,
\begin{equation}
  R_{f,t,m}(\zeta)=0
  \quad\Longleftrightarrow\quad
  P_{f,t,m}(\zeta)=0.
  \label{eq:lookup-root-equivalence}
\end{equation}
In particular,
\begin{equation}
  \left|
    \left\{
      \zeta\in\F\setminus S_{f,t}:
      R_{f,t,m}(\zeta)=0
    \right\}
  \right|
  \le2N-1.
  \label{eq:lookup-bad-shift-count}
\end{equation}
\end{lemma}

\begin{proof}
Put $S=S_{f,t}$ and
\begin{equation}
  D_S(Z)=\prod_{a\in S}(Z-a).
  \label{eq:lookup-squarefree-denominator}
\end{equation}
By \cref{lem:lookup-residue-form},
\begin{equation}
  P_{f,t,m}(Z)
  \defeq
  D_S(Z)R_{f,t,m}(Z)
  =
  \sum_{a\in S}
  \bigl(\nu_f(a)-\mu_{t,m}(a)\bigr)
  \frac{D_S(Z)}{Z-a}
  \label{eq:lookup-numerator-explicit}
\end{equation}
is a polynomial of degree at most $|S|-1\le2N-1$.

Because the lookup relation is false, choose $a_0\in S$ with
\[
  \nu_f(a_0)-\mu_{t,m}(a_0)\ne0.
\]
Evaluating \eqref{eq:lookup-numerator-explicit} at $a_0$ gives
\[
  P_{f,t,m}(a_0)
  =
  \bigl(\nu_f(a_0)-\mu_{t,m}(a_0)\bigr)
  \prod_{a\in S\setminus\{a_0\}}(a_0-a),
\]
which is nonzero because the elements of $S$ are distinct.  Thus $P_{f,t,m}$ is nonzero.

For $\zeta\notin S$, $D_S(\zeta)\ne0$, so \eqref{eq:lookup-root-equivalence} follows from \eqref{eq:lookup-numerator-explicit}.  The root bound \cref{lem:root-bound} and \eqref{eq:lookup-numerator-degree} imply \eqref{eq:lookup-bad-shift-count}.
\end{proof}

\subsection{Committed lookup relation}
\label{sec:committed-lookup-relation}

\begin{definition}[Committed lookup relation]
\label{def:lookup-relation}
Fix $0<\delta\le(1-\rho)/2$.  The relation $\mathcal R^{\delta}_{\mathrm{Look}}$ consists of pairs
\[
  ((w_f,w_t,w_m),(f,t,m))
\]
with $w_f,w_t,w_m\in\F^L$ and tables $f,t,m:\B^n\to\F$ satisfying
\begin{align}
  \Delta_{\mathrm{fib}}(w_f,\Enc_T(f))&\le\delta,
  \label{eq:lookup-relation-f}\\
  \Delta_{\mathrm{fib}}(w_t,\Enc_T(t))&\le\delta,
  \label{eq:lookup-relation-t}\\
  \Delta_{\mathrm{fib}}(w_m,\Enc_T(m))&\le\delta,
  \label{eq:lookup-relation-m}
\end{align}
and the algebraic lookup relation \eqref{eq:algebraic-lookup-relation}.  The instance is $(w_f,w_t,w_m)$ and the witness is $(f,t,m)$.
\end{definition}

Let $\mathbf 1:\B^n\to\F$ denote the constant-one table, and let
\begin{equation}
  w_{\mathbf1}\defeq\Enc_T(\mathbf1)\in C_0.
  \label{eq:lookup-one-word}
\end{equation}
This codeword is public.  After the lookup challenge $\zeta$ is sampled, define the two affine virtual words
\begin{equation}
  d_f\defeq \zeta w_{\mathbf1}-w_f,
  \qquad
  d_t\defeq \zeta w_{\mathbf1}-w_t.
  \label{eq:lookup-affine-words}
\end{equation}
They represent the denominator tables $\zeta\mathbf1-f$ and $\zeta\mathbf1-t$ whenever the original commitment words represent $f$ and $t$.

\begin{lemma}[Affine virtual words preserve decoding]
\label{lem:lookup-affine-decoding}
Let $w\in\F^L$, let $a:\B^n\to\F$, and let $\zeta\in\F$.  Put
\[
  d=\zeta w_{\mathbf1}-w.
\]
Then
\begin{equation}
  \Delta_{\mathrm{fib}}
  \bigl(d,\Enc_T(\zeta\mathbf1-a)\bigr)
  =
  \Delta_{\mathrm{fib}}\bigl(w,\Enc_T(a)\bigr).
  \label{eq:lookup-affine-distance}
\end{equation}
Consequently, in the unique-decoding regime, if $w$ decodes to $a$ within fibre distance $\delta$, then $d$ decodes uniquely to $\zeta\mathbf1-a$ within the same distance.
\end{lemma}

\begin{proof}
Linearity of the table commitment, \cref{prop:table-commitment-bijection}, gives
\[
  \Enc_T(\zeta\mathbf1-a)
  =\zeta w_{\mathbf1}-\Enc_T(a).
\]
For every $\xi\in L$,
\[
  d(\xi)-\Enc_T(\zeta\mathbf1-a)(\xi)
  =-\bigl(w(\xi)-\Enc_T(a)(\xi)\bigr).
\]
Hence the two words disagree at exactly the same positions and, in particular, on exactly the same fibres.  This proves \eqref{eq:lookup-affine-distance}.  Uniqueness follows from \cref{prop:unique-table-near-word}.
\end{proof}

\subsection{The lookup protocol}
\label{sec:lookup-protocol}

The protocol uses two inverse tables
\begin{equation}
  q(w)=\frac{1}{\zeta-f(w)},
  \qquad
  h(u)=\frac{1}{\zeta-t(u)}.
  \label{eq:lookup-inverse-tables}
\end{equation}
They are defined whenever $\zeta\notin S_{f,t}$.  Their defining equations are the pointwise relations
\begin{equation}
  q\had(\zeta\mathbf1-f)=\mathbf1,
  \qquad
  h\had(\zeta\mathbf1-t)=\mathbf1.
  \label{eq:lookup-inverse-hadamards}
\end{equation}
The logarithmic-derivative equality becomes
\begin{equation}
  \ip{q}{\mathbf1}=\ip{m}{h}.
  \label{eq:lookup-inner-product-equality}
\end{equation}
Thus lookup is reduced to two Hadamard relations, one public linear functional, and one committed inner product.

\begin{definition}[Lookup sub-batch]
\label{def:lookup-sub-batch}
Fix a lookup shift $\zeta\in\F$, inverse commitment words $w_q,w_h\in\F^L$, and a scalar $S\in\F$.  Define the heterogeneous batch
\begin{equation}
  \mathfrak B_{\mathrm{Look}}(\zeta,w_q,w_h,S)
  \label{eq:lookup-sub-batch}
\end{equation}
to contain the following four claims:
\begin{enumerate}[label=(\roman*),leftmargin=*]
  \item the Hadamard relation
  \begin{equation}
    q\had(\zeta\mathbf1-f)=\mathbf1
    \label{eq:lookup-batch-had-f}
  \end{equation}
  instantiated on the oracle words $(w_q,d_f,w_{\mathbf1})$;

  \item the Hadamard relation
  \begin{equation}
    h\had(\zeta\mathbf1-t)=\mathbf1
    \label{eq:lookup-batch-had-t}
  \end{equation}
  instantiated on $(w_h,d_t,w_{\mathbf1})$;

  \item the public linear-functional claim
  \begin{equation}
    \ip{q}{\mathbf1}=S
    \label{eq:lookup-batch-linear}
  \end{equation}
  instantiated on $w_q$ with public table $\mathbf1$;

  \item the committed inner-product claim
  \begin{equation}
    \ip{m}{h}=S
    \label{eq:lookup-batch-ip}
  \end{equation}
  instantiated on $(w_m,w_h)$.
\end{enumerate}
All four claims are protected by one execution of $\Pi_{\mathrm{Het}}$ from \cref{def:heterogeneous-batching-protocol}.
\end{definition}

The affine words $d_f,d_t$ in \eqref{eq:lookup-affine-words} are virtual: querying them requires one query to $w_f$ or $w_t$ and one field operation.  The public word $w_{\mathbf1}$ requires no commitment.  Therefore the lookup reduction introduces only the two new committed inverse tables $w_q,w_h$ before the heterogeneous batch creates its usual opening oracles.

\begin{definition}[Lookup protocol $\Pi_{\mathrm{Look}}$]
\label{def:lookup-protocol}
The instance is $(w_f,w_t,w_m)$.
\begin{enumerate}[label=\textnormal{\arabic*.},leftmargin=*]
  \item The verifier samples a lookup shift $\zeta\leftarrow\F$.

  \item The prover sends two oracle words $w_q,w_h\in\F^L$ and a scalar $S\in\F$.  Honestly, if $\zeta\notin S_{f,t}$, these are
  \begin{equation}
    w_q=\Enc_T(q),
    \qquad
    w_h=\Enc_T(h),
    \qquad
    S=\sum_{w\in\B^n}q(w),
    \label{eq:lookup-honest-messages}
  \end{equation}
  with $q,h$ as in \eqref{eq:lookup-inverse-tables}.  If $\zeta\in S_{f,t}$, the honest prover aborts and the verifier rejects.

  \item The verifier constructs $d_f,d_t$ by \eqref{eq:lookup-affine-words} and the parties execute one heterogeneous batch protocol
  \begin{equation}
    \Pi_{\mathrm{Het}}
    \bigl(\mathfrak B_{\mathrm{Look}}(\zeta,w_q,w_h,S)\bigr).
    \label{eq:lookup-run-heterogeneous}
  \end{equation}
\end{enumerate}
\end{definition}

\begin{theorem}[Lookup completeness]
\label{thm:lookup-completeness}
Suppose $(w_f,w_t,w_m)$ has a witness $(f,t,m)$ in $\mathcal R^{\delta}_{\mathrm{Look}}$, the three instance words are the honest commitments to that witness, and the prover follows \cref{def:lookup-protocol}.  Conditioned on
\begin{equation}
  \zeta\notin S_{f,t}
  \label{eq:lookup-nonpole-shift}
\end{equation}
and on the Hadamard DEEP challenges of the heterogeneous batch avoiding their three public pole sets, the verifier accepts with probability $1$.

With all challenges sampled uniformly from $\F$ and rejection on poles, the acceptance probability is at least
\begin{equation}
  1-\frac{2N}{|\F|}-\frac{3M}{|\F|}.
  \label{eq:lookup-completeness-bound}
\end{equation}
The $3M/|\F|$ term disappears if the public Hadamard DEEP challenges are sampled by rejection outside their pole sets.
\end{theorem}

\begin{proof}
Since $|S_{f,t}|\le2N$, the event \eqref{eq:lookup-nonpole-shift} fails with probability at most $2N/|\F|$.  Fix a shift satisfying \eqref{eq:lookup-nonpole-shift}.  Then the inverse tables \eqref{eq:lookup-inverse-tables} are well defined and satisfy the two Hadamard identities \eqref{eq:lookup-inverse-hadamards} pointwise.

Because $(f,t,m)$ satisfies the algebraic lookup relation, \cref{lem:lookup-residue-form} gives $R_{f,t,m}=0$.  Evaluating at $Z=\zeta$ yields
\[
  \sum_{w\in\B^n}\frac{1}{\zeta-f(w)}
  =
  \sum_{u\in\B^n}\frac{m(u)}{\zeta-t(u)},
\]
which is precisely \eqref{eq:lookup-inner-product-equality}.  Hence the four claims of \cref{def:lookup-sub-batch} have a joint witness $(f,t,m,q,h)$, with the affine denominator tables supplied by \cref{lem:lookup-affine-decoding}.  The completeness theorem for the heterogeneous batch, \cref{thm:heterogeneous-completeness}, gives conditional acceptance with probability $1$.  Its unrestricted Hadamard completeness loss is at most $3M/|\F|$, giving \eqref{eq:lookup-completeness-bound} by the union bound.
\end{proof}

\begin{remark}[Pole-free extension-field challenge]
\label{rem:lookup-extension-shift}
Suppose $f$ and $t$ take values in the base field $\Fq$ and the challenge field is a strict extension $\F\supsetneq\Fq$.  Sampling
\begin{equation}
  \zeta\leftarrow \F\setminus\Fq
  \label{eq:lookup-extension-challenge}
\end{equation}
ensures automatically that $\zeta\notin S_{f,t}$.  Thus the lookup-shift completeness loss vanishes.  The soundness root bound below becomes $(2N-1)/|\F\setminus\Fq|$.  This is often the natural configuration when witness columns live in a base field and verifier challenges already use an extension field.
\end{remark}

\subsection{Soundness}
\label{sec:lookup-soundness}

The next lemma is the bridge between a joint witness for the four batched subclaims and the original lookup relation.

\begin{lemma}[A joint sub-batch witness implies a logarithmic-derivative equality]
\label{lem:lookup-joint-witness}
Fix $\zeta\in\F$.  Suppose the batch $\mathfrak B_{\mathrm{Look}}(\zeta,w_q,w_h,S)$ has a joint witness in the sense of \cref{def:heterogeneous-batch-instance}.  Then there exist tables $f,t,m$ satisfying \eqref{eq:lookup-relation-f}--\eqref{eq:lookup-relation-m} such that
\begin{equation}
  \zeta\notin S_{f,t}
  \label{eq:lookup-joint-nonpole}
\end{equation}
and
\begin{equation}
  R_{f,t,m}(\zeta)=0.
  \label{eq:lookup-joint-rational-zero}
\end{equation}
\end{lemma}

\begin{proof}
A joint witness provides decoded tables $d_f',d_t',m,q,h$ for the denominator, multiplicity, and inverse words that occur in the four subclaims.  Define
\begin{equation}
  f\defeq\zeta\mathbf1-d_f',
  \qquad
  t\defeq\zeta\mathbf1-d_t'.
  \label{eq:lookup-joint-denominators}
\end{equation}
Because $d_f=\zeta w_{\mathbf1}-w_f$ and $d_t=\zeta w_{\mathbf1}-w_t$, the pointwise difference calculation in \cref{lem:lookup-affine-decoding} applied in reverse shows
\[
  \Delta_{\mathrm{fib}}(w_f,\Enc_T(f))
  =\Delta_{\mathrm{fib}}(d_f,\Enc_T(d_f'))\le\delta
\]
and similarly $\Delta_{\mathrm{fib}}(w_t,\Enc_T(t))\le\delta$.  The inner-product subclaim already supplies $m$ with $\Delta_{\mathrm{fib}}(w_m,\Enc_T(m))\le\delta$.  The two Hadamard claims in the joint witness therefore give, for every $w,u\in\B^n$,
\begin{equation}
  q(w)(\zeta-f(w))=1,
  \qquad
  h(u)(\zeta-t(u))=1.
  \label{eq:lookup-joint-inverses}
\end{equation}
Hence no denominator in \eqref{eq:lookup-joint-inverses} is zero, proving \eqref{eq:lookup-joint-nonpole}, and
\[
  q(w)=\frac1{\zeta-f(w)},
  \qquad
  h(u)=\frac1{\zeta-t(u)}.
\]
The public-linear and committed-inner-product claims share the same scalar $S$, so
\[
  \sum_w q(w)
  =S
  =\sum_u m(u)h(u).
\]
Substituting the inverse formulas gives exactly \eqref{eq:lookup-joint-rational-zero}.
\end{proof}

Let $s_{\mathrm{Look}}$ denote the protected-word count of the heterogeneous batch in \cref{def:lookup-sub-batch}, after any exact deduplication.  The direct instantiation of \cref{def:protected-component-list} gives the simple bound
\begin{equation}
  s_{\mathrm{Look}}\le
  2\cdot9+3+4=25.
  \label{eq:lookup-protected-count}
\end{equation}
The actual count can be smaller because $w_q$, $w_{\mathbf1}$, and other public or repeated virtual words may be shared.

\begin{theorem}[Lookup soundness]
\label{thm:lookup-soundness}
Assume
\begin{equation}
  0<\delta\le\frac{1-\rho}{2},
  \qquad
  2N<(1-\delta)M.
  \label{eq:lookup-parameters}
\end{equation}
If the instance $(w_f,w_t,w_m)$ has no witness in $\mathcal R^{\delta}_{\mathrm{Look}}$, then every deterministic prover is accepted by $\Pi_{\mathrm{Look}}$ with probability at most
\begin{equation}
\begin{split}
  \varepsilon_{\mathrm{Look}}
  \le{}&
  \frac{2N-1}{|\F|}
  +\frac{2(N-1)}{|\F|}
  +\frac{(s_{\mathrm{Look}}-1)M}{|\F|}
  \\
  &+\varepsilon_{\mathrm{fold}}
  +(1-\delta)^\kappa.
\end{split}
\label{eq:lookup-soundness-bound}
\end{equation}
The same bound holds for randomized provers by averaging.  If the shift is sampled uniformly from a pole-free set $\Gamma\subseteq\F$, replace the first term by $(2N-1)/|\Gamma|$.
\end{theorem}

\begin{proof}
By \cref{prop:unique-table-near-word}, each of $w_f,w_t,w_m$ has at most one table within fibre distance $\delta$.  If one of the three has none, \cref{lem:lookup-joint-witness} shows that no value of $\zeta$ can make the lookup sub-batch have a joint witness.

Otherwise let $f,t,m$ denote these unique decoded tables.  Because the original lookup instance is false, $(f,t,m)$ does not satisfy \eqref{eq:algebraic-lookup-relation}.  By \cref{lem:lookup-bad-challenges}, there are at most $2N-1$ values $\zeta\in\F\setminus S_{f,t}$ for which $R_{f,t,m}(\zeta)=0$.  By \cref{lem:lookup-joint-witness}, these are the only shifts for which the heterogeneous lookup sub-batch can have a joint witness.  Therefore
\begin{equation}
  \Pr_{\zeta\leftarrow\F}
  [\mathfrak B_{\mathrm{Look}}\text{ has a joint witness}]
  \le\frac{2N-1}{|\F|}.
  \label{eq:lookup-joint-witness-probability}
\end{equation}

Fix any other shift.  The heterogeneous batch is a false instance containing two Hadamard claims, one public linear-functional claim, and one committed inner-product claim.  Apply \cref{thm:heterogeneous-soundness}.  Since Hadamard claims are present, its bound is
\begin{equation}
  \frac{2(N-1)}{|\F|}
  +\frac{(s_{\mathrm{Look}}-1)M}{|\F|}
  +\varepsilon_{\mathrm{fold}}
  +(1-\delta)^\kappa.
  \label{eq:lookup-conditional-batch-soundness}
\end{equation}
Combining \eqref{eq:lookup-joint-witness-probability} and \eqref{eq:lookup-conditional-batch-soundness} gives \eqref{eq:lookup-soundness-bound}.  For a randomized prover, condition on its internal random tape; the deterministic bound applies to every fixed tape, and averaging preserves the same upper bound.

If $\zeta$ is sampled uniformly from a pole-free set $\Gamma$, the polynomial $P_{f,t,m}$ of \cref{lem:lookup-bad-challenges} has at most $2N-1$ roots in $\Gamma$, so the first probability is at most $(2N-1)/|\Gamma|$.
\end{proof}

\begin{corollary}[Lookup knowledge by decoding]
\label{cor:lookup-knowledge}
Under the hypotheses of \cref{thm:lookup-soundness}, if a prover is accepted with probability larger than the right-hand side of \eqref{eq:lookup-soundness-bound}, then the unique decodings of $w_f,w_t,w_m$ within radius $\delta$ exist and form a witness $(f,t,m)\in\mathcal R^{\delta}_{\mathrm{Look}}$.  They can be recovered by a polynomial-time Reed--Solomon unique-decoding algorithm inside half the minimum distance~\cite{WelchBerlekamp86,Gao03}.
\end{corollary}

\begin{proof}
The contrapositive is \cref{thm:lookup-soundness}.  Uniqueness of the three decoded tables is \cref{prop:unique-table-near-word}; polynomial-time recovery follows from standard Reed--Solomon unique decoding~\cite{WelchBerlekamp86,Gao03} and the coefficient/table bijection \cref{prop:table-isomorphism}.
\end{proof}

\subsection{What the reduction buys}
\label{sec:lookup-consequences}

The lookup protocol contains no algebraic Sumcheck.  Its relation-specific algebra is exactly
\begin{equation}
\begin{aligned}
  \text{lookup}
  &={}2\text{ inverse Hadamard relations}\\
  &\quad+1\text{ public linear functional}\\
  &\quad+1\text{ committed inner product}.
\end{aligned}
\label{eq:lookup-decomposition}
\end{equation}
After the two inverse tables have been committed, every low-degree obligation produced by those four relations is protected by the single folding test of \cref{sec:heterogeneous-batching}.  In particular, adding a lookup to a larger proof need not add another independent proximity proof: its protected words can be appended to the same master batch used by evaluations, inner products, and other Hadamard constraints.

The reduction is deliberately stated for an unindexed lookup.  Indexed lookups can be obtained by tagging each value with its index before applying the multiset machinery developed in the next section.  The present theorem isolates the part relevant to the coefficient-extraction framework: logarithmic derivatives turn lookup membership into inverse pointwise products and one scalar inner-product equality, all of which are already native relations of the framework.

\section{Multiset and Permutation Arguments}
\label{sec:permutation}

Permutation checks are a basic consistency primitive in proof systems: two committed columns should contain the same values, possibly in a different order, and tagged variants enforce a prescribed wiring permutation.  Grand-product arguments are a standard way to prove such statements; HyperPlonk, for example, reduces permutation constraints to polynomial identities that are ultimately checked through Sumcheck~\cite{HyperPlonk23}.  In the present framework the logarithmic-derivative lookup of \cref{sec:lookup} has a simpler specialization.  Multiset equality is proved by two inverse Hadamard relations and two public linear functionals, all protected by the single folding test of \cref{sec:heterogeneous-batching}.  A prescribed permutation is then reduced to multiset equality by random compression of index--value pairs.

Throughout this section, tables have length $N=2^n$.  We first formulate multiset equality algebraically over the field and then record the characteristic condition under which it coincides with ordinary equality of finite multisets.

\subsection{Algebraic multiset equality}
\label{sec:algebraic-multiset}

Let $a,b:\B^n\to\F$.  For $v\in\F$, define the field-valued multiplicities
\begin{equation}
  \nu_a(v)\defeq
  \sum_{w\in\B^n:\,a(w)=v}1_{\F},
  \qquad
  \nu_b(v)\defeq
  \sum_{w\in\B^n:\,b(w)=v}1_{\F}.
  \label{eq:multiset-field-multiplicities}
\end{equation}

\begin{definition}[Algebraic multiset equality]
\label{def:algebraic-multiset-equality}
We write
\begin{equation}
  a\equiv_{\mathrm{ms}} b
  \label{eq:algebraic-multiset-equality}
\end{equation}
if
\begin{equation}
  \nu_a(v)=\nu_b(v)
  \qquad\text{for every }v\in\F.
  \label{eq:multiset-multiplicity-equality}
\end{equation}
\end{definition}

\begin{proposition}[Relation with ordinary multiset equality]
\label{prop:ordinary-multiset-equality}
Assume $p=\operatorname{char}(\F)>N$.  Then $a\equiv_{\mathrm{ms}}b$ if and only if the ordinary multisets
\begin{equation}
  \{\!\{a(w):w\in\B^n\}\!\}
  \quad\text{and}\quad
  \{\!\{b(w):w\in\B^n\}\!\}
  \label{eq:ordinary-multisets}
\end{equation}
are equal, including multiplicities.
\end{proposition}

\begin{proof}
Ordinary multiset equality clearly implies \eqref{eq:multiset-multiplicity-equality}.  Conversely, for every $v\in\F$, the two ordinary multiplicities are integers in $[0,N]$.  Since $p>N$, their images in $\F$ are equal only when the integers themselves are equal.  Hence every value has the same ordinary multiplicity in the two tables.
\end{proof}

The logarithmic derivative of a multiset is the rational function
\begin{equation}
  R_{a,b}(Z)
  \defeq
  \sum_{w\in\B^n}\frac{1}{Z-a(w)}
  -
  \sum_{u\in\B^n}\frac{1}{Z-b(u)}.
  \label{eq:multiset-log-derivative}
\end{equation}
Let
\begin{equation}
  S_{a,b}\defeq
  \{a(w):w\in\B^n\}\cup
  \{b(u):u\in\B^n\}.
  \label{eq:multiset-poles}
\end{equation}

\begin{lemma}[Residue characterization]
\label{lem:multiset-residue-characterization}
The rational function in \eqref{eq:multiset-log-derivative} satisfies
\begin{equation}
  R_{a,b}(Z)
  =
  \sum_{v\in S_{a,b}}
  \frac{\nu_a(v)-\nu_b(v)}{Z-v}.
  \label{eq:multiset-residue-form}
\end{equation}
Consequently,
\begin{equation}
  R_{a,b}=0
  \quad\Longleftrightarrow\quad
  a\equiv_{\mathrm{ms}}b.
  \label{eq:multiset-rational-characterization}
\end{equation}
\end{lemma}

\begin{proof}
Group the terms of each sum in \eqref{eq:multiset-log-derivative} by their common value.  This gives \eqref{eq:multiset-residue-form}.  If the multiplicities agree, every residue is zero, so $R_{a,b}=0$.

Conversely, suppose $R_{a,b}=0$.  Fix $v_0\in S_{a,b}$ and multiply \eqref{eq:multiset-residue-form} by $Z-v_0$.  Evaluating at $Z=v_0$ leaves only the term indexed by $v_0$, hence
\[
  \nu_a(v_0)-\nu_b(v_0)=0.
\]
Since $v_0$ was arbitrary, \eqref{eq:multiset-multiplicity-equality} holds.
\end{proof}

For a false multiset relation, a random evaluation of $R_{a,b}$ detects the failure except at few field points.

\begin{lemma}[Bad multiset shifts]
\label{lem:multiset-bad-shifts}
Assume $a\not\equiv_{\mathrm{ms}}b$.  Then there exists a nonzero polynomial $P_{a,b}\in\F[Z]$ of degree at most
\begin{equation}
  |S_{a,b}|-2\le2N-2
  \label{eq:multiset-numerator-degree}
\end{equation}
such that, for every $\zeta\in\F\setminus S_{a,b}$,
\begin{equation}
  R_{a,b}(\zeta)=0
  \quad\Longleftrightarrow\quad
  P_{a,b}(\zeta)=0.
  \label{eq:multiset-root-equivalence}
\end{equation}
Therefore
\begin{equation}
  \left|
    \{\zeta\in\F\setminus S_{a,b}:R_{a,b}(\zeta)=0\}
  \right|
  \le2N-2.
  \label{eq:multiset-bad-shift-count}
\end{equation}
\end{lemma}

\begin{proof}
Put $S=S_{a,b}$ and
\[
  D_S(Z)=\prod_{v\in S}(Z-v).
\]
By \cref{lem:multiset-residue-characterization},
\begin{equation}
  P_{a,b}(Z)
  \defeq D_S(Z)R_{a,b}(Z)
  =
  \sum_{v\in S}
  \bigl(\nu_a(v)-\nu_b(v)\bigr)
  \frac{D_S(Z)}{Z-v}
  \label{eq:multiset-numerator}
\end{equation}
is a polynomial of degree at most $|S|-1$.  The coefficient of $Z^{|S|-1}$ is
\[
  \sum_{v\in S}\bigl(\nu_a(v)-\nu_b(v)\bigr)
  =N-N=0
\]
in $\F$, so $\deg P_{a,b}\le |S|-2$.  Since the multiset relation is false, some residue is nonzero.  Evaluating \eqref{eq:multiset-numerator} at the corresponding pole, exactly as in \cref{lem:lookup-bad-challenges}, shows that $P_{a,b}$ is nonzero.

For $\zeta\notin S$, $D_S(\zeta)\ne0$, giving \eqref{eq:multiset-root-equivalence}.  The root bound \cref{lem:root-bound} and $|S|\le2N$ give \eqref{eq:multiset-bad-shift-count}.
\end{proof}

\subsection{Committed multiset relation}
\label{sec:committed-multiset}

\begin{definition}[Committed multiset relation]
\label{def:committed-multiset-relation}
Fix $0<\delta\le(1-\rho)/2$.  The relation $\mathcal R^{\delta}_{\mathrm{MS}}$ consists of pairs
\[
  ((w_a,w_b),(a,b))
\]
with $w_a,w_b\in\F^L$ and tables $a,b:\B^n\to\F$ satisfying
\begin{equation}
  \Delta_{\mathrm{fib}}(w_a,\Enc_T(a))\le\delta,
  \qquad
  \Delta_{\mathrm{fib}}(w_b,\Enc_T(b))\le\delta,
  \label{eq:multiset-commitment-distances}
\end{equation}
and $a\equiv_{\mathrm{ms}}b$.
\end{definition}

After a challenge $\zeta\in\F$ is sampled, define the virtual denominator words
\begin{equation}
  d_a\defeq\zeta w_{\mathbf1}-w_a,
  \qquad
  d_b\defeq\zeta w_{\mathbf1}-w_b,
  \label{eq:multiset-denominator-words}
\end{equation}
where $w_{\mathbf1}=\Enc_T(\mathbf1)$ is the public codeword from \eqref{eq:lookup-one-word}.  As in \cref{lem:lookup-affine-decoding}, these words preserve fibre distance exactly.

The prover supplies inverse tables
\begin{equation}
  q(w)=\frac1{\zeta-a(w)},
  \qquad
  r(u)=\frac1{\zeta-b(u)}.
  \label{eq:multiset-inverse-tables}
\end{equation}
They satisfy the pointwise inverse relations
\begin{equation}
  q\had(\zeta\mathbf1-a)=\mathbf1,
  \qquad
  r\had(\zeta\mathbf1-b)=\mathbf1,
  \label{eq:multiset-inverse-hadamards}
\end{equation}
and the logarithmic-derivative equality is simply
\begin{equation}
  \ip{q}{\mathbf1}=\ip{r}{\mathbf1}.
  \label{eq:multiset-sum-equality}
\end{equation}
Thus, unlike the general lookup relation of \cref{sec:lookup}, no multiplicity commitment and no committed inner product are needed.

\begin{definition}[Multiset sub-batch]
\label{def:multiset-sub-batch}
Fix $\zeta\in\F$, inverse commitment words $w_q,w_r\in\F^L$, and a scalar $S\in\F$.  Define $\mathfrak B_{\mathrm{MS}}(\zeta,w_q,w_r,S)$ to contain:
\begin{enumerate}[label=(\roman*),leftmargin=*]
  \item the Hadamard claim
  \[
    q\had(\zeta\mathbf1-a)=\mathbf1
  \]
  on $(w_q,d_a,w_{\mathbf1})$;
  \item the Hadamard claim
  \[
    r\had(\zeta\mathbf1-b)=\mathbf1
  \]
  on $(w_r,d_b,w_{\mathbf1})$;
  \item the public linear-functional claim
  \[
    \ip{q}{\mathbf1}=S
  \]
  on $w_q$;
  \item the public linear-functional claim
  \[
    \ip{r}{\mathbf1}=S
  \]
  on $w_r$.
\end{enumerate}
All four claims are protected by one execution of $\Pi_{\mathrm{Het}}$.
\end{definition}

\begin{definition}[Multiset protocol $\Pi_{\mathrm{MS}}$]
\label{def:multiset-protocol}
The instance is $(w_a,w_b)$.
\begin{enumerate}[label=\textnormal{\arabic*.},leftmargin=*]
  \item The verifier samples $\zeta\leftarrow\F$.
  \item The prover sends $w_q,w_r\in\F^L$ and $S\in\F$.  Honestly, if $\zeta\notin S_{a,b}$,
  \begin{equation}
    w_q=\Enc_T(q),
    \qquad
    w_r=\Enc_T(r),
    \qquad
    S=\sum_{w\in\B^n}q(w),
    \label{eq:multiset-honest-messages}
  \end{equation}
  with $q,r$ as in \eqref{eq:multiset-inverse-tables}.  If $\zeta\in S_{a,b}$, the honest prover aborts and the verifier rejects.
  \item The verifier constructs $d_a,d_b$ from \eqref{eq:multiset-denominator-words}, and the parties execute
  \begin{equation}
    \Pi_{\mathrm{Het}}
    \bigl(\mathfrak B_{\mathrm{MS}}(\zeta,w_q,w_r,S)\bigr).
    \label{eq:multiset-run-batch}
  \end{equation}
\end{enumerate}
\end{definition}

\begin{theorem}[Multiset completeness]
\label{thm:multiset-completeness}
Suppose $(w_a,w_b)$ are honest commitments to tables $a,b$ with $a\equiv_{\mathrm{ms}}b$.  Conditioned on $\zeta\notin S_{a,b}$ and on the shared Hadamard DEEP challenges avoiding their pole sets, the honest prover is accepted with probability $1$.

With all challenges uniform in $\F$ and rejection on poles, the acceptance probability is at least
\begin{equation}
  1-\frac{2N}{|\F|}-\frac{3M}{|\F|}.
  \label{eq:multiset-completeness-bound}
\end{equation}
The $3M/|\F|$ term disappears if the public DEEP challenges are sampled by rejection outside the corresponding pole sets.
\end{theorem}

\begin{proof}
The pole set has size at most $2N$.  Fix $\zeta\notin S_{a,b}$.  The inverse tables \eqref{eq:multiset-inverse-tables} are well defined and satisfy \eqref{eq:multiset-inverse-hadamards}.  Since $a\equiv_{\mathrm{ms}}b$, \cref{lem:multiset-residue-characterization} gives $R_{a,b}=0$, so evaluating at $\zeta$ yields
\[
  \sum_w q(w)=\sum_u r(u).
\]
Thus all four claims in \cref{def:multiset-sub-batch} are true on one joint decoding, and \cref{thm:heterogeneous-completeness} gives conditional acceptance with probability $1$.  The stated unconditional lower bound follows by the union bound, exactly as in \cref{thm:lookup-completeness}.
\end{proof}

\begin{remark}[Pole-free extension-field challenge]
\label{rem:multiset-extension-challenge}
If $a$ and $b$ take values in $\Fq$ and $\F\supsetneq\Fq$, sampling
\[
  \zeta\leftarrow\F\setminus\Fq
\]
removes the lookup-shift completeness loss entirely.  In the soundness theorem below, the term $(2N-2)/|\F|$ is then replaced by $(2N-2)/|\F\setminus\Fq|$.
\end{remark}

\subsection{Soundness and extraction}
\label{sec:multiset-soundness}

\begin{lemma}[A joint multiset sub-batch witness fixes the logarithmic derivative]
\label{lem:multiset-joint-witness}
Fix $\zeta\in\F$.  If $\mathfrak B_{\mathrm{MS}}(\zeta,w_q,w_r,S)$ has a joint witness, then there exist tables $a,b$ satisfying \eqref{eq:multiset-commitment-distances} such that $\zeta\notin S_{a,b}$ and
\begin{equation}
  R_{a,b}(\zeta)=0.
  \label{eq:multiset-joint-rational-zero}
\end{equation}
\end{lemma}

\begin{proof}
Let $d_a',d_b',q,r$ be the decoded denominator and inverse tables supplied by the joint witness, and define
\[
  a=\zeta\mathbf1-d_a',
  \qquad
  b=\zeta\mathbf1-d_b'.
\]
The reverse direction of \cref{lem:lookup-affine-decoding} gives the distance conditions \eqref{eq:multiset-commitment-distances}.  The two Hadamard relations give
\[
  q(w)(\zeta-a(w))=1,
  \qquad
  r(u)(\zeta-b(u))=1
\]
for all coordinates.  Hence $\zeta$ is not a pole and $q(w)=(\zeta-a(w))^{-1}$, $r(u)=(\zeta-b(u))^{-1}$.  The two public linear-functional claims share the same scalar $S$, hence
\[
  \sum_wq(w)=S=\sum_ur(u),
\]
which is exactly \eqref{eq:multiset-joint-rational-zero}.
\end{proof}

Let $s_{\mathrm{MS}}$ be the number of protected low-degree words in the heterogeneous batch after exact deduplication.  The direct component count gives
\begin{equation}
  s_{\mathrm{MS}}\le2\cdot9+2\cdot3=24.
  \label{eq:multiset-protected-count}
\end{equation}
As before, sharing $w_q,w_r,w_{\mathbf1}$ and repeated virtual words can reduce this count.

\begin{theorem}[Multiset soundness]
\label{thm:multiset-soundness}
Assume
\begin{equation}
  0<\delta\le\frac{1-\rho}{2},
  \qquad
  2N<(1-\delta)M.
  \label{eq:multiset-parameters}
\end{equation}
If $(w_a,w_b)$ has no witness in $\mathcal R^{\delta}_{\mathrm{MS}}$, every deterministic prover is accepted by $\Pi_{\mathrm{MS}}$ with probability at most
\begin{equation}
\begin{split}
  \varepsilon_{\mathrm{MS}}
  \le{}&
  \frac{2N-2}{|\F|}
  +\frac{2(N-1)}{|\F|}
  +\frac{(s_{\mathrm{MS}}-1)M}{|\F|}
  \\
  &+\varepsilon_{\mathrm{fold}}
  +(1-\delta)^\kappa.
\end{split}
\label{eq:multiset-soundness-bound}
\end{equation}
The same bound holds for randomized provers.  If $\zeta$ is sampled uniformly from a pole-free set $\Gamma\subseteq\F$, replace the first term by $(2N-2)/|\Gamma|$.
\end{theorem}

\begin{proof}
By \cref{prop:unique-table-near-word}, each of $w_a,w_b$ has at most one table within fibre distance $\delta$.  If one has none, \cref{lem:multiset-joint-witness} implies that no shift can make the multiset sub-batch have a joint witness.

Otherwise let $a,b$ be the unique decoded tables.  Since the original instance is false, $a\not\equiv_{\mathrm{ms}}b$.  By \cref{lem:multiset-bad-shifts}, at most $2N-2$ nonpole values of $\zeta$ satisfy $R_{a,b}(\zeta)=0$.  By \cref{lem:multiset-joint-witness}, only those shifts can make the heterogeneous sub-batch have a joint witness.  This contributes at most $(2N-2)/|\F|$.

For every other shift, the heterogeneous batch is false and contains two Hadamard claims and two public linear-functional claims.  Applying \cref{thm:heterogeneous-soundness} contributes
\[
  \frac{2(N-1)}{|\F|}
  +\frac{(s_{\mathrm{MS}}-1)M}{|\F|}
  +\varepsilon_{\mathrm{fold}}
  +(1-\delta)^\kappa.
\]
Adding the two bounds proves \eqref{eq:multiset-soundness-bound}.  Randomized provers follow by conditioning on their internal randomness.  Restricting $\zeta$ to $\Gamma$ changes only the root-count denominator.
\end{proof}

\begin{corollary}[Multiset knowledge by decoding]
\label{cor:multiset-knowledge}
Under the hypotheses of \cref{thm:multiset-soundness}, acceptance with probability larger than the right-hand side of \eqref{eq:multiset-soundness-bound} implies that the unique decodings of $w_a,w_b$ exist and satisfy $a\equiv_{\mathrm{ms}}b$.  They are recoverable in polynomial time by Reed--Solomon unique decoding.
\end{corollary}

\begin{proof}
This is the contrapositive of \cref{thm:multiset-soundness}, together with uniqueness from \cref{prop:unique-table-near-word} and the decoding algorithms cited in \cref{cor:lookup-knowledge}.
\end{proof}

\subsection{Prescribed permutations by tagged compression}
\label{sec:tagged-permutation}

Multiset equality proves that two columns contain the same values but does not specify which occurrence should move to which position.  Let
\begin{equation}
  \sigma:\B^n\longrightarrow\B^n
  \label{eq:known-permutation}
\end{equation}
be a public permutation.  The prescribed permutation relation is
\begin{equation}
  b(\sigma(w))=a(w)
  \qquad\text{for every }w\in\B^n.
  \label{eq:prescribed-permutation-relation}
\end{equation}

Let $\iota:\B^n\to\F$ be the canonical index embedding
\begin{equation}
  \iota(w)=\sum_{k=1}^n w_k2^{k-1}\,1_{\F}.
  \label{eq:index-embedding}
\end{equation}
When $p>N$, this map is injective.  For a compression challenge $\tau\in\F$, define
\begin{equation}
  A_\tau(w)
  \defeq a(w)+\tau\iota(\sigma(w)),
  \qquad
  B_\tau(u)
  \defeq b(u)+\tau\iota(u).
  \label{eq:tagged-compressed-columns}
\end{equation}
If \eqref{eq:prescribed-permutation-relation} holds, then $A_\tau\equiv_{\mathrm{ms}}B_\tau$ for every $\tau$, because the substitution $u=\sigma(w)$ identifies the tagged pairs exactly.

\begin{lemma}[Compression of unequal tagged multisets]
\label{lem:tagged-compression}
Assume $p>N$ and that \eqref{eq:prescribed-permutation-relation} is false.  Let
\begin{align}
  \mathcal A
  &\defeq
  \{\!\{(\iota(\sigma(w)),a(w)):w\in\B^n\}\!\},
  \label{eq:tagged-multiset-a}\\
  \mathcal B
  &\defeq
  \{\!\{(\iota(u),b(u)):u\in\B^n\}\!\}.
  \label{eq:tagged-multiset-b}
\end{align}
Then $\mathcal A\ne\mathcal B$.  Moreover, there is a set $T_{\mathrm{coll}}\subseteq\F$ with
\begin{equation}
  |T_{\mathrm{coll}}|
  \le \binom{2N}{2}
  \label{eq:tag-collision-count}
\end{equation}
such that for every $\tau\notin T_{\mathrm{coll}}$,
\begin{equation}
  A_\tau\not\equiv_{\mathrm{ms}}B_\tau.
  \label{eq:tagged-compression-detects}
\end{equation}
\end{lemma}

\begin{proof}
Because $\sigma$ is a permutation, for each $u\in\B^n$ there is a unique $w=\sigma^{-1}(u)$.  Equality $\mathcal A=\mathcal B$ would therefore force equality of the values attached to every unique first coordinate:
\[
  a(\sigma^{-1}(u))=b(u)
  \qquad\text{for all }u,
\]
which is equivalent to \eqref{eq:prescribed-permutation-relation}.  Hence the tagged multisets are unequal.

Let $\mathcal U$ be the union of the supports of $\mathcal A$ and $\mathcal B$ as a set of distinct pairs; $|\mathcal U|\le2N$.  For two distinct pairs $(i,x),(j,y)\in\mathcal U$, their compressed values collide when
\begin{equation}
  x+\tau i=y+\tau j.
  \label{eq:tag-collision-equation}
\end{equation}
If $i=j$, distinctness of the pairs gives $x\ne y$, so \eqref{eq:tag-collision-equation} has no solution.  If $i\ne j$, injectivity of the index embedding gives $i-j\ne0$ in $\F$, hence there is exactly one solution
\[
  \tau=(y-x)(i-j)^{-1}.
\]
Let $T_{\mathrm{coll}}$ contain all such solutions over unordered pairs of elements of $\mathcal U$.  The union bound gives \eqref{eq:tag-collision-count}.

For $\tau\notin T_{\mathrm{coll}}$, the map $(i,x)\mapsto x+\tau i$ is injective on $\mathcal U$.  Therefore equality of the compressed multisets would imply equality of the multiplicity of every pair in $\mathcal U$, contradicting $\mathcal A\ne\mathcal B$.  This proves \eqref{eq:tagged-compression-detects}.
\end{proof}

The lemma deliberately uses a simple collision argument rather than a tighter protocol-specific fingerprint.  Its purpose is to isolate a fully proved reduction from a prescribed permutation to the multiset primitive; tighter tagging bounds can be substituted without changing the coefficient-extraction backend.

\begin{definition}[Tagged permutation protocol $\Pi_{\mathrm{Perm}}$]
\label{def:tagged-permutation-protocol}
Assume $p>N$ and let $\sigma$ be public.  On instance $(w_a,w_b)$:
\begin{enumerate}[label=\textnormal{\arabic*.},leftmargin=*]
  \item The verifier samples $\tau\leftarrow\F$.
  \item Using linearity of $\Enc_T$, the verifier forms the virtual words
  \begin{equation}
    w_{A,\tau}
    =w_a+\tau w_{\iota\circ\sigma},
    \qquad
    w_{B,\tau}
    =w_b+\tau w_{\iota},
    \label{eq:tagged-virtual-words}
  \end{equation}
  where $w_{\iota\circ\sigma}=\Enc_T(\iota\circ\sigma)$ and $w_{\iota}=\Enc_T(\iota)$ are public codewords.
  \item The parties run $\Pi_{\mathrm{MS}}$ on $(w_{A,\tau},w_{B,\tau})$.
\end{enumerate}
\end{definition}

\begin{theorem}[Tagged permutation soundness]
\label{thm:tagged-permutation-soundness}
Assume $p>N$ and the parameter conditions of \cref{thm:multiset-soundness}.  If the unique decoded tables $a,b$ fail \eqref{eq:prescribed-permutation-relation}, then every prover is accepted by $\Pi_{\mathrm{Perm}}$ with probability at most
\begin{equation}
  \varepsilon_{\mathrm{Perm}}
  \le
  \frac{\binom{2N}{2}}{|\F|}
  +\varepsilon_{\mathrm{MS}}.
  \label{eq:tagged-permutation-soundness}
\end{equation}
If either input word has no decoding within fibre distance $\delta$, the first term is unnecessary and the multiset soundness bound alone applies.
\end{theorem}

\begin{proof}
If one input word has no valid decoding, linearity and the same pointwise-difference argument as \cref{lem:lookup-affine-decoding} show that the corresponding tagged virtual word has no decoding compatible with the prescribed affine shift, so \cref{thm:multiset-soundness} applies directly.

Otherwise let $a,b$ be the unique decoded tables.  By \cref{lem:tagged-compression}, outside a set of at most $\binom{2N}{2}$ challenges $\tau$, the compressed tables $A_\tau,B_\tau$ do not satisfy algebraic multiset equality.  For every such $\tau$, the multiset protocol accepts with probability at most $\varepsilon_{\mathrm{MS}}$ by \cref{thm:multiset-soundness}.  Taking the union bound over the bad compression event and the conditional multiset soundness event gives \eqref{eq:tagged-permutation-soundness}.
\end{proof}

\begin{remark}[Tuple tags and indexed lookups]
\label{rem:tuple-tags-indexed-lookup}
The same construction compresses tuples rather than pairs.  For public tags $t_1,\ldots,t_d$ and random challenges $\tau_1,\ldots,\tau_d$, replace a tuple $(x_0,x_1,\ldots,x_d)$ by
\[
  x_0+\tau_1x_1+\cdots+\tau_dx_d.
\]
A collision analysis analogous to \cref{lem:tagged-compression} reduces equality of tagged tuple multisets to the scalar multiset protocol.  In particular, an indexed lookup can tag a table value by its row index before applying \cref{sec:lookup}, preventing two equal values in different rows from being interchangeable when the application requires row identity.
\end{remark}

\subsection{Consequence for the coefficient-extraction framework}
\label{sec:permutation-consequence}

The scalar multiset check has the decomposition
\begin{equation}
\begin{aligned}
  \text{multiset equality}
  &={}2\text{ inverse Hadamard relations}\\
  &\quad+2\text{ public linear functionals},
\end{aligned}
\label{eq:multiset-decomposition}
\end{equation}
and therefore uses no algebraic Sumcheck.  A prescribed permutation adds only a public random compression of tagged values before the same primitive.  All low-degree obligations from the two inverse relations and the two sums enter the single folding batch of \cref{sec:heterogeneous-batching}.

This closes the second logarithmic-derivative primitive needed later for sparse matrix and memory arguments: \cref{sec:lookup} proves membership with multiplicities, while the present section proves equality of committed multisets and public wiring permutations.  The next step is to use these primitives together with Hadamard products to authenticate sparse matrix--vector multiplication without introducing a separate Sumcheck layer.

\section{Sparse Matrix--Vector Products}
\label{sec:sparse-matvec}

Sparse matrix--vector multiplication is the remaining linear-algebra primitive needed to pass from the table relations of the preceding sections to R1CS.  In Spartan-type systems the matrices are public and sparse, while the vector is witness-dependent; the proof must therefore authenticate the products without expanding a sparse matrix into its dense $N\times N$ table~\cite{Setty20}.  A direct random compression already reduces $y=Mz$ to one weighted inner product, but evaluating the resulting matrix-dependent weight vector naively would still require the verifier to scan all nonzero entries.  This section gives a sparse-entry reduction that avoids that scan after preprocessing.

The construction uses only primitives already proved in this paper.  A public sparse matrix is represented by a list of nonzero entries.  Two indexed lookups authenticate the column reads and the row weights, one Hadamard relation multiplies the fetched vector entries by the public matrix coefficients, and one committed inner product is compared with a public geometric linear functional of the claimed output.  All low-degree obligations enter the single heterogeneous folding test of \cref{sec:heterogeneous-batching}.

\subsection{Ambient sparse representation}
\label{sec:sparse-representation}

The protocols of the preceding sections use tables of one common power-of-two length.  We therefore choose an ambient
\begin{equation}
  N=2^n
  \label{eq:sparse-ambient-N}
\end{equation}
large enough to contain both the vector coordinates of interest and all sparse-entry slots.  A smaller vector or a shorter sparse list is padded to length $N$.  Dummy sparse slots have coefficient zero, so their row and column indices are irrelevant; for definiteness they may both be set to $0$.

Let
\begin{equation}
  r,c:\B^n\longrightarrow\B^n,
  \qquad
  \mu:\B^n\longrightarrow\F
  \label{eq:sparse-entry-tables}
\end{equation}
be public tables.  The slot $e\in\B^n$ represents an entry of value $\mu(e)$ in row $r(e)$ and column $c(e)$.  Repeated row--column pairs are allowed.  They define the matrix
\begin{equation}
  M_{i,j}
  \defeq
  \sum_{e\in\B^n:\,r(e)=i,\,c(e)=j}\mu(e),
  \qquad i,j\in\B^n.
  \label{eq:sparse-matrix-definition}
\end{equation}
Thus the padded zero slots contribute nothing, and combining duplicate entries is automatic.

For a table $z:\B^n\to\F$, matrix--vector multiplication means
\begin{equation}
  (Mz)(i)
  =
  \sum_{j\in\B^n}M_{i,j}z(j)
  =
  \sum_{e:\,r(e)=i}\mu(e)z(c(e)).
  \label{eq:sparse-matvec-coordinate}
\end{equation}
The second equality follows directly from \eqref{eq:sparse-matrix-definition} by exchanging the finite sums.

\begin{definition}[Sparse matrix--vector relation]
\label{def:sparse-matvec-relation}
Fix $0<\delta\le(1-\rho)/2$.  For the public sparse metadata $(r,c,\mu)$, the relation $\mathcal R^{\delta}_{\mathrm{SMV}}[r,c,\mu]$ consists of pairs
\[
  ((w_z,w_y),(z,y))
\]
with $w_z,w_y\in\F^L$ and tables $z,y:\B^n\to\F$ satisfying
\begin{equation}
  \Delta_{\mathrm{fib}}(w_z,\Enc_T(z))\le\delta,
  \qquad
  \Delta_{\mathrm{fib}}(w_y,\Enc_T(y))\le\delta,
  \label{eq:sparse-matvec-commitment-distances}
\end{equation}
and
\begin{equation}
  y(i)=\sum_{e:\,r(e)=i}\mu(e)z(c(e))
  \qquad\text{for every }i\in\B^n.
  \label{eq:sparse-matvec-relation}
\end{equation}
\end{definition}

The public preprocessing associated with $(r,c,\mu)$ consists only of length-$N$ field-valued tables: the encodings of $\iota\circ r$, $\iota\circ c$, and $\mu$, together with the index table of \cref{sec:tagged-permutation} and the public multiplicity tables introduced below.  No dense $N^2$ matrix is formed.

\subsection{Indexed selection from lookup}
\label{sec:indexed-selection}

The first task is to authenticate a table of selected coordinates.  Let
\begin{equation}
  a:\B^n\longrightarrow\B^n
  \label{eq:index-map-a}
\end{equation}
be any public index map.  For a source table $z$ and a claimed selected table $u$, the indexed-selection relation is
\begin{equation}
  u(e)=z(a(e))
  \qquad\text{for every }e\in\B^n.
  \label{eq:indexed-selection-relation}
\end{equation}
Let $\iota:\B^n\to\F$ be the injective index embedding from \eqref{eq:index-embedding}; throughout this subsection assume
\begin{equation}
  p=\operatorname{char}(\F)>N.
  \label{eq:sparse-index-characteristic}
\end{equation}
Define the public multiplicity table
\begin{equation}
  m_a(j)
  \defeq
  \sum_{e:\,a(e)=j}1_{\F}.
  \label{eq:index-map-multiplicity}
\end{equation}
For a tagging challenge $\tau\in\F$, form
\begin{align}
  f_{a,u,\tau}(e)
  &\defeq u(e)+\tau\iota(a(e)),
  \label{eq:indexed-query-compression}\\
  t_{z,\tau}(j)
  &\defeq z(j)+\tau\iota(j).
  \label{eq:indexed-table-compression}
\end{align}
The corresponding commitment words are virtual affine combinations:
\begin{equation}
  w_{f,\tau}=w_u+\tau w_{\iota\circ a},
  \qquad
  w_{t,\tau}=w_z+\tau w_{\iota},
  \label{eq:indexed-virtual-compressions}
\end{equation}
where the two index words are public preprocessing data.

\begin{lemma}[Correct indexed selection gives a lookup]
\label{lem:indexed-selection-gives-lookup}
If \eqref{eq:indexed-selection-relation} holds, then for every $\tau\in\F$ the triple
\begin{equation}
  \bigl(f_{a,u,\tau},t_{z,\tau},m_a\bigr)
  \label{eq:indexed-selection-lookup-triple}
\end{equation}
satisfies the algebraic lookup relation of \cref{def:lookup-residue-multiplicities}.
\end{lemma}

\begin{proof}
Fix $\lambda\in\F$.  Every query row $e$ contributes the compressed pair value
\[
  u(e)+\tau\iota(a(e))
  =z(a(e))+\tau\iota(a(e)).
\]
For a fixed table row $j$, this value occurs once for every $e$ with $a(e)=j$, hence exactly with field-valued multiplicity $m_a(j)$.  Summing over all $j$ gives equality of the query multiplicity and the multiplicity-weighted table multiplicity at every $\lambda$, which is precisely \eqref{eq:algebraic-lookup-relation}.
\end{proof}

The converse can fail only if the random scalar compression collides two distinct index--value pairs.  The next statement is the indexed-lookup analogue of \cref{lem:tagged-compression}.

\begin{lemma}[Soundness of indexed compression]
\label{lem:indexed-selection-compression}
Assume \eqref{eq:sparse-index-characteristic} and suppose \eqref{eq:indexed-selection-relation} is false.  Then there is a set $T_{a,u,z}\subseteq\F$ with
\begin{equation}
  |T_{a,u,z}|
  \le\binom{2N}{2}
  \label{eq:indexed-selection-collision-count}
\end{equation}
such that for every $\tau\notin T_{a,u,z}$, the triple \eqref{eq:indexed-selection-lookup-triple} does not satisfy the algebraic lookup relation.
\end{lemma}

\begin{proof}
Consider the two tagged multisets
\begin{align}
  \mathcal A
  &\defeq
  \{\!\{(\iota(a(e)),u(e)):e\in\B^n\}\!\},
  \label{eq:indexed-left-pairs}\\
  \mathcal B
  &\defeq
  \biguplus_{j\in\B^n}
  |a^{-1}(j)|\cdot\{\!\{(\iota(j),z(j))\}\!\}.
  \label{eq:indexed-right-pairs}
\end{align}
Here the multiplicity on the right is interpreted as the ordinary integer number of preimages of $j$ under $a$; since $p>N$, its field image is exactly \eqref{eq:index-map-multiplicity}.

If $\mathcal A=\mathcal B$, then for each $j$ the right multiset contains exactly $|a^{-1}(j)|$ copies of the unique pair $(\iota(j),z(j))$.  The left multiset contains exactly the same number of pairs whose first coordinate is $\iota(j)$, namely $(\iota(j),u(e))$ for $e\in a^{-1}(j)$.  Equality therefore forces $u(e)=z(j)$ for every such $e$, which is \eqref{eq:indexed-selection-relation}.  Hence $\mathcal A\ne\mathcal B$.

The union of the supports of \eqref{eq:indexed-left-pairs} and \eqref{eq:indexed-right-pairs} contains at most $2N$ distinct pairs.  For two distinct pairs $(i,x)$ and $(j,y)$ in this union, the scalar compressions $x+\tau i$ and $y+\tau j$ can coincide for at most one $\tau$: there is no solution when $i=j$, and exactly one when $i\ne j$.  Let $T_{a,u,z}$ be the union of those collision values.  Then \eqref{eq:indexed-selection-collision-count} follows exactly as in \cref{lem:tagged-compression}.

For $\tau\notin T_{a,u,z}$, scalar compression is injective on the union of the two supports.  If the compressed triple satisfied the algebraic lookup relation, the compressed query multiset would equal the multiplicity-weighted compressed table multiset.  Injectivity would imply $\mathcal A=\mathcal B$, a contradiction.
\end{proof}

\begin{remark}[Collision-free extension tagging]
\label{rem:collision-free-index-tagging}
The random tag is not essential.  If both coordinates of every indexed pair lie in a proper subfield $\F_0\subsetneq\F$ and one fixes $\vartheta\in\F\setminus\F_0$, then
\[
  (i,x)\longmapsto x+\vartheta i
\]
is injective on $\F_0^2$.  In that configuration the term \eqref{eq:indexed-selection-collision-count} disappears.  We retain random compression because it works over one field without adding a subfield hypothesis.
\end{remark}

For later use we spell out the exact primitive batch used by an indexed selection.  After a lookup shift $\zeta\in\F$, let $q,h:\B^n\to\F$ denote the claimed inverse tables
\begin{equation}
  q(e)=\frac{1}{\zeta-f_{a,u,\tau}(e)},
  \qquad
  h(j)=\frac{1}{\zeta-t_{z,\tau}(j)}.
  \label{eq:indexed-selection-inverses}
\end{equation}

\begin{definition}[Indexed-selection sub-batch]
\label{def:indexed-selection-sub-batch}
Fix $a,\tau,\zeta$, commitment words $w_u,w_z,w_q,w_h$, and a scalar $S\in\F$.  The batch $\mathfrak B_{\mathrm{Sel}}[a]$ consists of:
\begin{enumerate}[label=(\roman*),leftmargin=*]
  \item the Hadamard relation
  \[
    q\had(\zeta\mathbf1-f_{a,u,\tau})=\mathbf1;
  \]
  \item the Hadamard relation
  \[
    h\had(\zeta\mathbf1-t_{z,\tau})=\mathbf1;
  \]
  \item the public linear-functional claim
  \[
    \ip{q}{\mathbf1}=S;
  \]
  \item the public linear-functional claim
  \[
    \ip{h}{m_a}=S.
  \]
\end{enumerate}
The compressed source words are the virtual words in \eqref{eq:indexed-virtual-compressions}, and $m_a$ is public.  Hence this is exactly the logarithmic-derivative lookup reduction of \cref{sec:lookup}, specialized so that its multiplicity table is public.
\end{definition}

Thus an indexed selection contributes two inverse Hadamard relations and two public linear-functional claims; the committed inner-product claim of the general lookup disappears because $m_a$ is public.

\subsection{Geometric row fingerprint}
\label{sec:sparse-row-fingerprint}

Let $\alpha\in\F$.  Define the geometric table
\begin{equation}
  g_{\alpha}(i)
  \defeq
  \alpha^{\iota(i)},
  \qquad i\in\B^n.
  \label{eq:geometric-row-table}
\end{equation}
and its table polynomial
\begin{equation}
  G_{\alpha}(X)
  \defeq
  \sum_{j=0}^{N-1}\alpha^jX^j.
  \label{eq:geometric-row-polynomial}
\end{equation}
For $\alpha X\ne1$,
\begin{equation}
  G_{\alpha}(X)
  =\frac{1-(\alpha X)^N}{1-\alpha X},
  \label{eq:geometric-series-formula}
\end{equation}
while $G_{\alpha}(X)=N\,1_{\F}$ when $\alpha X=1$.  Since the characteristic is odd and $N$ is a power of two, $N\,1_{\F}\ne0$.  Hence the public codeword $w_{g_\alpha}=\Enc_T(g_\alpha)$ and the reversed kernel used for the public linear functional $\ip{g_\alpha}{y}$ are evaluable at any verifier query point with $O(\log N)$ field operations by binary exponentiation.

For a claimed sparse-entry row-weight table $s:\B^n\to\F$, the relation
\begin{equation}
  s(e)=\alpha^{\iota(r(e))}
  =g_\alpha(r(e))
  \label{eq:row-weight-selection}
\end{equation}
is exactly the indexed-selection relation of \eqref{eq:indexed-selection-relation} with source table $g_\alpha$ and public map $r$.  Therefore it is authenticated by $\mathfrak B_{\mathrm{Sel}}[r]$ without materializing a dense matrix.

The final fingerprint identity is the following elementary fact.

\begin{lemma}[Sparse contraction identity]
\label{lem:sparse-contraction}
Let $u,s,v:\B^n\to\F$ satisfy
\begin{equation}
  u(e)=z(c(e)),
  \qquad
  s(e)=\alpha^{\iota(r(e))},
  \qquad
  v(e)=\mu(e)u(e).
  \label{eq:sparse-auxiliary-relations}
\end{equation}
Then
\begin{equation}
  \ip{s}{v}
  =
  \sum_{i\in\B^n}\alpha^{\iota(i)}(Mz)(i).
  \label{eq:sparse-contraction-identity}
\end{equation}
\end{lemma}

\begin{proof}
Using \eqref{eq:sparse-auxiliary-relations} and grouping sparse-entry slots by their public row index,
\begin{align*}
  \ip{s}{v}
  &=\sum_{e\in\B^n}
    \alpha^{\iota(r(e))}\mu(e)z(c(e))\\
  &=\sum_{i\in\B^n}\alpha^{\iota(i)}
    \sum_{e:\,r(e)=i}\mu(e)z(c(e))\\
  &=\sum_{i\in\B^n}\alpha^{\iota(i)}(Mz)(i),
\end{align*}
where the last equality is \eqref{eq:sparse-matvec-coordinate}.
\end{proof}

\begin{lemma}[Geometric fingerprint detects a false product]
\label{lem:sparse-geometric-fingerprint}
Suppose $y\ne Mz$.  Define
\begin{equation}
  D_{M,z,y}(T)
  \defeq
  \sum_{i\in\B^n}
  \bigl(y(i)-(Mz)(i)\bigr)T^{\iota(i)}.
  \label{eq:sparse-error-polynomial}
\end{equation}
Then $D_{M,z,y}$ is a nonzero polynomial of degree at most $N-1$.  Consequently
\begin{equation}
  \left|
    \left\{
      \alpha\in\F:
      \ip{g_\alpha}{y}
      =\sum_i\alpha^{\iota(i)}(Mz)(i)
    \right\}
  \right|
  \le N-1.
  \label{eq:sparse-fingerprint-bad-alpha}
\end{equation}
\end{lemma}

\begin{proof}
Because $y\ne Mz$, at least one coefficient of \eqref{eq:sparse-error-polynomial} is nonzero, so $D_{M,z,y}\ne0$.  Its largest possible exponent is $N-1$.  The equality in \eqref{eq:sparse-fingerprint-bad-alpha} is exactly $D_{M,z,y}(\alpha)=0$.  Apply the univariate root bound \cref{lem:root-bound}.
\end{proof}

\subsection{The sparse matrix--vector protocol}
\label{sec:sparse-matvec-protocol}

We can now state the complete reduction.  Besides the witness words $w_z,w_y$, the verification key contains public commitments to the length-$N$ sparse metadata and index-derived tables required below.  These are fixed once for a public matrix and can therefore be preprocessed.

\begin{definition}[Sparse matrix--vector sub-batch]
\label{def:sparse-matvec-sub-batch}
Fix challenges $\alpha,\tau_c,\tau_r,\zeta_c,\zeta_r\in\F$ and prover words $w_u,w_s,w_v$ together with the inverse words and scalar messages required by the two indexed lookups.  The batch $\mathfrak B_{\mathrm{SMV}}$ contains:
\begin{enumerate}[label=(\roman*),leftmargin=*]
  \item the indexed-selection sub-batch $\mathfrak B_{\mathrm{Sel}}[c]$ proving
  \[
    u(e)=z(c(e));
  \]
  \item the indexed-selection sub-batch $\mathfrak B_{\mathrm{Sel}}[r]$ proving
  \[
    s(e)=g_\alpha(r(e));
  \]
  \item the Hadamard claim
  \begin{equation}
    v=\mu\had u;
    \label{eq:sparse-contribution-hadamard}
  \end{equation}
  \item a committed inner-product claim
  \begin{equation}
    \ip{s}{v}=S;
    \label{eq:sparse-entry-inner-product}
  \end{equation}
  \item the public linear-functional claim
  \begin{equation}
    \ip{g_\alpha}{y}=S.
    \label{eq:sparse-output-linear-functional}
  \end{equation}
\end{enumerate}
Every primitive claim above is appended to one heterogeneous batch and is protected by one execution of $\Pi_{\mathrm{Het}}$.
\end{definition}

\begin{definition}[Sparse matrix--vector protocol $\Pi_{\mathrm{SMV}}$]
\label{def:sparse-matvec-protocol}
The instance is $(w_z,w_y)$ and the public sparse matrix metadata $(r,c,\mu)$.
\begin{enumerate}[label=\textnormal{\arabic*.},leftmargin=*]
  \item The verifier samples
  \begin{equation}
    \alpha,\tau_c,\tau_r\leftarrow\F.
    \label{eq:sparse-first-challenges}
  \end{equation}
  \item The prover sends commitment words $w_u,w_s,w_v$.  Honestly,
  \begin{equation}
    u(e)=z(c(e)),
    \qquad
    s(e)=\alpha^{\iota(r(e))},
    \qquad
    v(e)=\mu(e)u(e).
    \label{eq:sparse-honest-auxiliaries}
  \end{equation}
  \item The verifier forms the virtual compressed lookup words for the column selection and row-weight selection using \eqref{eq:indexed-virtual-compressions}, with source tables $z$ and $g_\alpha$, respectively, and samples independent lookup shifts
  \begin{equation}
    \zeta_c,\zeta_r\leftarrow\F.
    \label{eq:sparse-lookup-shifts}
  \end{equation}
  \item The prover sends the inverse words and scalar sum messages required by the two lookup sub-batches, and a scalar $S$.  Honestly,
  \begin{equation}
    S=\sum_{e\in\B^n}s(e)v(e).
    \label{eq:sparse-honest-S}
  \end{equation}
  \item The parties execute a single
  \begin{equation}
    \Pi_{\mathrm{Het}}\bigl(\mathfrak B_{\mathrm{SMV}}\bigr).
    \label{eq:sparse-run-one-batch}
  \end{equation}
\end{enumerate}
\end{definition}

\begin{theorem}[Sparse matrix--vector completeness]
\label{thm:sparse-matvec-completeness}
Suppose $(w_z,w_y)$ are honest commitments to tables satisfying $y=Mz$, and the prover follows \cref{def:sparse-matvec-protocol}.  Conditioned on the two lookup shifts and the shared Hadamard DEEP challenges avoiding their pole sets, the verifier accepts with probability $1$.
\end{theorem}

\begin{proof}
The honest $u$ satisfies the column indexed-selection relation, so \cref{lem:indexed-selection-gives-lookup} makes its compressed lookup true for every $\tau_c$.  The honest $s$ satisfies \eqref{eq:row-weight-selection}, so the row-weight lookup is true for every $\tau_r$.  The relation \eqref{eq:sparse-contribution-hadamard} holds by definition of $v$.

By \cref{lem:sparse-contraction},
\[
  \ip{s}{v}
  =\sum_i\alpha^{\iota(i)}(Mz)(i)
  =\sum_i\alpha^{\iota(i)}y(i)
  =\ip{g_\alpha}{y},
\]
so the two scalar claims \eqref{eq:sparse-entry-inner-product} and \eqref{eq:sparse-output-linear-functional} hold with the same $S$.  Hence every primitive claim in \cref{def:sparse-matvec-sub-batch} has a joint honest witness.  Conditional perfect completeness follows from \cref{thm:heterogeneous-completeness} and the lookup completeness theorem \cref{thm:lookup-completeness}.
\end{proof}

\subsection{Soundness}
\label{sec:sparse-matvec-soundness}

Let $s_{\mathrm{SMV}}$ be the number of distinct protected low-degree words in the heterogeneous batch of \cref{def:sparse-matvec-sub-batch} after exact deduplication.  We leave it symbolic because public matrix words, the constant-one word, inverse tables, and repeated source words can be shared across the two indexed lookups and with surrounding constraints.

For each indexed selection, define the outer bad-event bound
\begin{equation}
  \varepsilon_{\mathrm{Sel,out}}
  \defeq
  \frac{\binom{2N}{2}}{|\F|}
  +\frac{2N-1}{|\F|}.
  \label{eq:indexed-selection-outer-error}
\end{equation}
The first term is the tag-collision bound of \cref{lem:indexed-selection-compression}; the second is the bad logarithmic-derivative shift bound of \cref{lem:lookup-bad-challenges}.  As in \cref{rem:collision-free-index-tagging}, the first term can be removed by collision-free extension-field tagging.

\begin{lemma}[A joint sparse sub-batch witness implies the compressed matrix identity]
\label{lem:sparse-joint-witness}
Fix the outer challenges $\alpha,\tau_c,\tau_r,\zeta_c,\zeta_r$.  Suppose the heterogeneous batch $\mathfrak B_{\mathrm{SMV}}$ has a joint witness, and suppose neither indexed compression challenge lies in its collision set and neither lookup shift lies in the corresponding bad-shift set.  Then the unique decoded tables $z,y$ satisfy
\begin{equation}
  \ip{g_\alpha}{y}
  =
  \sum_{i\in\B^n}\alpha^{\iota(i)}(Mz)(i).
  \label{eq:sparse-joint-compressed-equality}
\end{equation}
\end{lemma}

\begin{proof}
Because the column compression challenge is collision-free and its logarithmic-derivative shift is not bad, a joint witness for the first indexed-selection sub-batch implies the underlying indexed relation
\[
  u(e)=z(c(e))
\]
by \cref{lem:indexed-selection-compression,lem:lookup-joint-witness}.  The same argument for the row-weight lookup gives
\[
  s(e)=g_\alpha(r(e))=\alpha^{\iota(r(e))}.
\]
The Hadamard subclaim in the same joint witness gives $v=\mu\had u$.  The committed inner-product and public linear-functional claims share the scalar $S$, hence
\[
  \ip{s}{v}=S=\ip{g_\alpha}{y}.
\]
Apply \cref{lem:sparse-contraction} to the three decoded auxiliary relations to obtain \eqref{eq:sparse-joint-compressed-equality}.
\end{proof}

\begin{theorem}[Sparse matrix--vector soundness]
\label{thm:sparse-matvec-soundness}
Assume
\begin{equation}
  p>N,
  \qquad
  0<\delta\le\frac{1-\rho}{2},
  \qquad
  2N<(1-\delta)M.
  \label{eq:sparse-matvec-parameters}
\end{equation}
If the instance $(w_z,w_y)$ has no witness in $\mathcal R^{\delta}_{\mathrm{SMV}}[r,c,\mu]$, then every deterministic prover is accepted by $\Pi_{\mathrm{SMV}}$ with probability at most
\begin{equation}
\begin{split}
  \varepsilon_{\mathrm{SMV}}
  \le{}&
  \frac{N-1}{|\F|}
  +2\varepsilon_{\mathrm{Sel,out}}
  +\frac{2(N-1)}{|\F|}
  +\frac{(s_{\mathrm{SMV}}-1)M}{|\F|}
  \\
  &+\varepsilon_{\mathrm{fold}}
  +(1-\delta)^\kappa.
\end{split}
\label{eq:sparse-matvec-soundness-bound}
\end{equation}
The same bound holds for randomized provers by averaging.  With collision-free extension tagging, delete the two $\binom{2N}{2}/|\F|$ terms contained in $2\varepsilon_{\mathrm{Sel,out}}$.
\end{theorem}

\begin{proof}
By \cref{prop:unique-table-near-word}, each of $w_z,w_y$ has at most one table within fibre distance $\delta$.  If one of them has none, the heterogeneous batch cannot have a compatible joint witness.  Indeed, a decoded compressed source word in either indexed lookup would, after subtracting its public affine tag, give a decoding of $w_z$ at the same fibre distance; similarly the public-linear-functional subclaim on $y$ requires a decoding of $w_y$.  Therefore \cref{thm:heterogeneous-soundness} applies directly, giving the shared Hadamard fingerprint term and the final batching/folding/query terms in \eqref{eq:sparse-matvec-soundness-bound}.

Otherwise let $z,y$ be the unique decoded tables.  Since the sparse matrix--vector relation is false, $y\ne Mz$.  By \cref{lem:sparse-geometric-fingerprint}, at most $N-1$ values of $\alpha$ satisfy the compressed equality \eqref{eq:sparse-joint-compressed-equality}.  This contributes the first term of \eqref{eq:sparse-matvec-soundness-bound}.

Fix an $\alpha$ outside that root set.  For each of the two indexed selections, \cref{lem:indexed-selection-compression} excludes all but at most $\binom{2N}{2}$ tagging challenges from faithfully representing the underlying indexed relation.  Conditional on a faithful tag, \cref{lem:lookup-bad-challenges} excludes all but at most $2N-1$ lookup shifts from detecting a false compressed lookup.  The union bound over the column and row indexed selections therefore contributes at most $2\varepsilon_{\mathrm{Sel,out}}$.

Outside all these outer bad events, \cref{lem:sparse-joint-witness} shows that the heterogeneous batch cannot have a joint witness, because such a witness would imply \eqref{eq:sparse-joint-compressed-equality}, contradicting the chosen $\alpha$.  Apply \cref{thm:heterogeneous-soundness}.  The batch contains Hadamard claims, so its remaining error is
\[
  \frac{2(N-1)}{|\F|}
  +\frac{(s_{\mathrm{SMV}}-1)M}{|\F|}
  +\varepsilon_{\mathrm{fold}}
  +(1-\delta)^\kappa.
\]
Adding the outer bad-event probabilities gives \eqref{eq:sparse-matvec-soundness-bound}.  Randomized provers follow by conditioning on their internal randomness and averaging.
\end{proof}

\begin{corollary}[Sparse matrix--vector knowledge by decoding]
\label{cor:sparse-matvec-knowledge}
Under the hypotheses of \cref{thm:sparse-matvec-soundness}, acceptance with probability larger than the right-hand side of \eqref{eq:sparse-matvec-soundness-bound} implies that the unique decodings of $w_z,w_y$ exist and satisfy $y=Mz$.  They are recoverable in polynomial time by Reed--Solomon unique decoding.
\end{corollary}

\begin{proof}
This is the contrapositive of \cref{thm:sparse-matvec-soundness}, together with uniqueness from \cref{prop:unique-table-near-word} and the standard decoding algorithms cited in \cref{cor:lookup-knowledge}.
\end{proof}

\subsection{Cost and consequence}
\label{sec:sparse-matvec-cost}

The reduction has a simple algebraic decomposition:
\begin{equation}
\begin{aligned}
  \text{sparse }y=Mz
  ={}&2\text{ indexed selections}\\
     &+1\text{ Hadamard relation}\\
     &+1\text{ committed inner product}\\
     &+1\text{ public geometric linear functional}.
\end{aligned}
\label{eq:sparse-matvec-decomposition}
\end{equation}
Each indexed selection is itself two inverse Hadamard relations plus two public linear functionals.  Hence, after expanding the selections, the sparse matrix--vector relation contributes five Hadamard claims, five public linear-functional claims, and one committed inner-product claim to the master batch.  They share one folding chain and one final query test.

The public matrix is stored by $O(N)$ sparse-entry metadata rather than $N^2$ dense entries.  If the original matrix has $K$ nonzero entries and vector dimension $d$, choose the ambient power of two
\begin{equation}
  N=2^{\lceil\log_2\max\{K,d\}\rceil}
  \label{eq:sparse-ambient-choice}
\end{equation}
and pad the remaining slots.  Public preprocessing therefore scales with the sparse representation.  The online prover constructs $u,s,v$ in $O(N)$ field operations apart from the Reed--Solomon encoding and the common folding backend.  At verification time the row/column maps, coefficient table, index words, and multiplicity tables are fixed preprocessing data; their Merkle roots can be included in a circuit-specific verification key, so the verifier does not scan the sparse list online.

\begin{remark}[Direct compression without preprocessing]
\label{rem:direct-sparse-compression}
There is a shorter but non-succinct reduction.  For a random multilinear point $r_0\in\F^n$, define
\[
  p_{r_0}(j)=\sum_i M_{i,j}\eq(i,r_0).
\]
Then
\[
  \widetilde y(r_0)=\ip{p_{r_0}}{z}
\]
is a public linear-functional claim, and a false matrix product is detected by a random multilinear evaluation.  However computing the full public kernel $p_{r_0}$ online requires scanning the sparse matrix.  The indexed-selection construction above trades this online scan for public sparse preprocessing and relations already native to the coefficient-extraction framework.
\end{remark}

The main consequence is that generic public sparse matrix--vector multiplication no longer requires an independent Sumcheck layer.  Its witness-dependent algebra has been reduced to the same lookup, Hadamard, inner-product, and public-kernel relations already handled by the Reed--Solomon/FRI backend.  The next section applies this result three times to the sparse R1CS matrices and combines the outputs with one additional Hadamard relation.

\section{Sparse R1CS from One Folding Batch}
\label{sec:r1cs}

The previous sections provide all algebraic primitives needed for rank-one constraint systems.  Let
\[
  A,B,C\in\F^{N\times N}
\]
be public sparse matrices and let $z:\B^n\to\F$ be the witness vector, identified with its length-$N$ table in the convention of \cref{sec:preliminaries}.  The R1CS relation is
\begin{equation}
  (Az)\had(Bz)=Cz.
  \label{eq:r1cs-standard-relation}
\end{equation}
The standard Sumcheck route proves matrix evaluations and the coordinatewise product through multilinear reductions.  Here we instead introduce committed tables
\[
  a=Az,
  \qquad
  b=Bz,
  \qquad
  c=Cz,
\]
prove each sparse matrix--vector relation by \cref{sec:sparse-matvec}, and prove $a\had b=c$ by \cref{sec:hadamard}.  Crucially, these are not run as four independent protocols: all protected low-degree words are appended to one heterogeneous batch and are tested by one folding chain and one final set of queries.

\subsection{The committed R1CS relation}
\label{sec:r1cs-relation}

For each $M\in\{A,B,C\}$, fix a sparse-entry representation
\begin{equation}
  (r_M,c_M,\mu_M):\B^n\longrightarrow \B^n\times\B^n\times\F
  \label{eq:r1cs-sparse-representations}
\end{equation}
as in \cref{sec:sparse-representation}.  Thus
\begin{equation}
  (Mz)(i)
  =\sum_{e:\,r_M(e)=i}\mu_M(e)z(c_M(e)).
  \label{eq:r1cs-sparse-coordinate}
\end{equation}
Padding slots have coefficient $0$, and the ambient power of two is chosen large enough for the witness dimension and every sparse list.

\begin{definition}[Committed sparse R1CS relation]
\label{def:r1cs-relation}
Fix $0<\delta\le(1-\rho)/2$.  The relation
$\mathcal R^{\delta}_{\mathrm{R1CS}}[A,B,C]$ consists of pairs
\[
  ((w_z,w_a,w_b,w_c),(z,a,b,c))
\]
with $w_z,w_a,w_b,w_c\in\F^L$ and tables $z,a,b,c:\B^n\to\F$ such that
\begin{align}
  \Delta_{\mathrm{fib}}(w_z,\Enc_T(z))&\le\delta,
  \label{eq:r1cs-distance-z}\\
  \Delta_{\mathrm{fib}}(w_a,\Enc_T(a))&\le\delta,
  \label{eq:r1cs-distance-a}\\
  \Delta_{\mathrm{fib}}(w_b,\Enc_T(b))&\le\delta,
  \label{eq:r1cs-distance-b}\\
  \Delta_{\mathrm{fib}}(w_c,\Enc_T(c))&\le\delta,
  \label{eq:r1cs-distance-c}\\
  a&=Az,
  \label{eq:r1cs-a}\\
  b&=Bz,
  \label{eq:r1cs-b}\\
  c&=Cz,
  \label{eq:r1cs-c}\\
  a\had b&=c.
  \label{eq:r1cs-had}
\end{align}
The four words form the oracle part of the instance and $(z,a,b,c)$ is the decoded witness.
\end{definition}

\begin{proposition}[Equivalence with the usual R1CS equation]
\label{prop:r1cs-equivalence}
For a table $z:\B^n\to\F$, the following are equivalent.
\begin{enumerate}[label=(\roman*),leftmargin=*]
  \item $z$ satisfies \eqref{eq:r1cs-standard-relation}.
  \item There exist tables $a,b,c:\B^n\to\F$ satisfying \eqref{eq:r1cs-a}--\eqref{eq:r1cs-had}.
\end{enumerate}
Moreover, when they exist the three tables are uniquely determined by $z$:
\[
  a=Az,\qquad b=Bz,\qquad c=Cz.
\]
\end{proposition}

\begin{proof}
If \eqref{eq:r1cs-standard-relation} holds, set $a=Az$, $b=Bz$, and $c=Cz$.  Then \eqref{eq:r1cs-a}--\eqref{eq:r1cs-c} hold by definition and \eqref{eq:r1cs-had} is exactly \eqref{eq:r1cs-standard-relation}.  Conversely, substituting \eqref{eq:r1cs-a}--\eqref{eq:r1cs-c} into \eqref{eq:r1cs-had} gives \eqref{eq:r1cs-standard-relation}.  The three matrix products are deterministic functions of $z$, which gives uniqueness.
\end{proof}

The point of the auxiliary tables is therefore not to weaken the relation.  They expose the R1CS equation as three sparse linear maps followed by one coordinatewise multiplication, exactly the two relation types already supported by the coefficient-extraction framework.

\subsection{One master batch for the whole R1CS instance}
\label{sec:r1cs-master-batch}

For $M\in\{A,B,C\}$, write $y_A=a$, $y_B=b$, and $y_C=c$.  Instantiate the sparse matrix--vector sub-batch of \cref{def:sparse-matvec-sub-batch} for the relation
\begin{equation}
  y_M=Mz.
  \label{eq:r1cs-one-matvec}
\end{equation}
We use one shared geometric row-fingerprint challenge $\alpha$ for all three matrices.  The indexed-selection tagging challenges and logarithmic-derivative shifts may be independent for the six selections.  Sharing $\alpha$ is sound because $w_z,w_a,w_b,w_c$ are fixed before it is sampled; \cref{thm:r1cs-soundness} shows that a false system then has at most $N-1$ masking values of $\alpha$, not three separate root-error terms.

\begin{definition}[R1CS heterogeneous sub-batch]
\label{def:r1cs-sub-batch}
Fix the outer challenges and all prover messages required by the three sparse matrix--vector reductions.  The batch $\mathfrak B_{\mathrm{R1CS}}$ is obtained by concatenating
\begin{equation}
  \mathfrak B_{\mathrm{SMV}}[A],
  \qquad
  \mathfrak B_{\mathrm{SMV}}[B],
  \qquad
  \mathfrak B_{\mathrm{SMV}}[C],
  \label{eq:r1cs-three-smv-batches}
\end{equation}
with the final committed Hadamard claim
\begin{equation}
  a\had b=c.
  \label{eq:r1cs-final-hadamard-claim}
\end{equation}
Exact repetitions of original commitments, public words, and virtual words are deduplicated before the master challenge of \cref{def:heterogeneous-batching-protocol} is sampled.
\end{definition}

Expanding \eqref{eq:r1cs-three-smv-batches} using \eqref{eq:sparse-matvec-decomposition}, the R1CS batch contains six indexed selections, three sparse contribution Hadamards, three sparse committed inner products, three geometric output linear functionals, and the final Hadamard relation.  After expanding the indexed selections themselves, there are sixteen Hadamard claims in total.  They all share the one Hadamard fingerprint pair $(\gamma,\theta)$ of \cref{sec:heterogeneous-batching}; therefore the Hadamard fingerprint error is paid once.

\begin{definition}[Sparse R1CS protocol $\Pi_{\mathrm{R1CS}}$]
\label{def:r1cs-protocol}
The public data are the sparse representations of $A,B,C$ and their preprocessed verification-key data.  The oracle instance is $(w_z,w_a,w_b,w_c)$, fixed before any verifier challenge.
\begin{enumerate}[label=\textnormal{\arabic*.},leftmargin=*]
  \item The verifier samples one shared row-fingerprint challenge
  \begin{equation}
    \alpha\leftarrow\F.
    \label{eq:r1cs-alpha}
  \end{equation}
  \item For each $M\in\{A,B,C\}$, the parties perform the outer, pre-batching part of \cref{def:sparse-matvec-protocol}: the prover supplies the auxiliary source, row-weight, and contribution words; the verifier samples the two indexed-selection tags and shifts; and the prover supplies the corresponding inverse words and scalar messages.  All these words are fixed before the master batch challenge.
  \item The final relation $a\had b=c$ and all Hadamard relations arising inside the three sparse reductions are prepared using the shared $(\gamma,\theta)$ required by \cref{def:heterogeneous-batching-protocol}.
  \item The parties execute exactly one
  \begin{equation}
    \Pi_{\mathrm{Het}}\bigl(\mathfrak B_{\mathrm{R1CS}}\bigr).
    \label{eq:r1cs-one-heterogeneous-batch}
  \end{equation}
\end{enumerate}
\end{definition}

Thus there is one master batching challenge, one folding chain, and one final $\kappa$-query test for the complete R1CS algebra.  The outer challenges exist only to reduce the matrix, lookup, and coordinatewise relations to the low-degree words protected by that single batch.

\subsection{Completeness}
\label{sec:r1cs-completeness}

\begin{theorem}[Sparse R1CS completeness]
\label{thm:r1cs-completeness}
Suppose $(w_z,w_a,w_b,w_c)$ are honest table commitments and $z$ satisfies \eqref{eq:r1cs-standard-relation}, with $a=Az$, $b=Bz$, and $c=Cz$.  If every lookup shift and every DEEP point used by the batch avoids its pole set, then the honest prover in \cref{def:r1cs-protocol} is accepted with probability $1$.
\end{theorem}

\begin{proof}
By \cref{prop:r1cs-equivalence}, the honest decoded tables satisfy the three sparse matrix--vector relations and the final Hadamard relation.  For each matrix $M\in\{A,B,C\}$, \cref{thm:sparse-matvec-completeness} supplies honest auxiliary words for $\mathfrak B_{\mathrm{SMV}}[M]$ and shows that every primitive relation in that sub-batch is true.  The final Hadamard relation is true by \eqref{eq:r1cs-standard-relation}, so \cref{thm:had-completeness} supplies its honest protected words.  Hence the concatenated batch $\mathfrak B_{\mathrm{R1CS}}$ has a joint witness.  Conditional perfect completeness follows from \cref{thm:heterogeneous-completeness}.
\end{proof}

\begin{remark}[Perfect completeness by rejection sampling]
\label{rem:r1cs-perfect-completeness}
The pole events are an artifact of writing DEEP and logarithmic-derivative virtual words as rational functions.  As in \cref{rem:had-perfect-completeness}, the verifier can rejection-sample the relevant extension-field challenges outside the explicitly known forbidden sets.  This yields perfect completeness without changing any algebraic relation.  If unrestricted challenges are preferred, a union bound gives an explicit completeness loss from the six lookup pole sets and the three shared Hadamard DEEP pole sets.
\end{remark}

\subsection{Soundness}
\label{sec:r1cs-soundness}

Let $s_{\mathrm{R1CS}}$ denote the number of distinct protected degree-$<N$ words in $\mathfrak B_{\mathrm{R1CS}}$ after exact deduplication.  Recall from \eqref{eq:indexed-selection-outer-error} that one indexed selection has outer bad-event bound
\begin{equation}
  \varepsilon_{\mathrm{Sel,out}}
  =\frac{\binom{2N}{2}}{|\F|}
   +\frac{2N-1}{|\F|}.
  \label{eq:r1cs-selection-outer-error}
\end{equation}
With collision-free extension tagging, the first term is absent.

The next lemma is the reason for sharing one matrix fingerprint challenge across $A,B,C$.

\begin{lemma}[One geometric challenge detects a false matrix equation]
\label{lem:r1cs-shared-alpha}
Fix decoded tables $z,a,b,c$.  If at least one of
\begin{equation}
  a=Az,
  \qquad
  b=Bz,
  \qquad
  c=Cz
  \label{eq:r1cs-three-matvec-equations}
\end{equation}
is false, then there exists a nonzero polynomial $D\in\F[T]$ of degree at most $N-1$ such that, for every $\alpha\in\F$ for which all three compressed sparse matrix identities hold, one has $D(\alpha)=0$.  Consequently the set of such $\alpha$ has cardinality at most $N-1$.
\end{lemma}

\begin{proof}
Choose any one false equation in \eqref{eq:r1cs-three-matvec-equations}; for definiteness, choose the first in the order $A,B,C$.  If the chosen equation is $a\ne Az$, take the error polynomial $D_{A,z,a}$ of \eqref{eq:sparse-error-polynomial}; if it is $b\ne Bz$ or $c\ne Cz$, use the analogous polynomial for $B$ or $C$.  By \cref{lem:sparse-geometric-fingerprint}, this polynomial is nonzero, has degree at most $N-1$, and vanishes at every $\alpha$ for which that chosen false matrix equation passes its compressed geometric identity.  If all three compressed identities hold, the chosen one holds in particular, so $D(\alpha)=0$.  The root bound gives at most $N-1$ such challenges.
\end{proof}

\begin{lemma}[A joint R1CS batch witness gives an R1CS witness]
\label{lem:r1cs-joint-witness}
Fix all outer challenges.  Suppose $\mathfrak B_{\mathrm{R1CS}}$ has a joint witness, every indexed-selection tag used by one selected false sparse matrix equation is collision-free, and its two logarithmic-derivative shifts lie outside the corresponding bad-shift sets.  If the shared row challenge $\alpha$ is outside the root set of \cref{lem:r1cs-shared-alpha}, then the decoded tables satisfy
\begin{equation}
  a=Az,
  \qquad b=Bz,
  \qquad c=Cz,
  \qquad a\had b=c.
  \label{eq:r1cs-joint-all-relations}
\end{equation}
In particular, if any matrix equation is false or the final Hadamard equation is false, then outside the stated outer bad events the batch has no joint witness.
\end{lemma}

\begin{proof}
If a matrix equation is false, choose the first false one as in \cref{lem:r1cs-shared-alpha}.  The corresponding sparse sub-batch is part of $\mathfrak B_{\mathrm{R1CS}}$.  Outside its two indexed-selection outer bad events, \cref{lem:sparse-joint-witness} implies the compressed matrix identity for that equation.  Since $\alpha$ is outside the root set from \cref{lem:r1cs-shared-alpha}, this is impossible.  Therefore a joint batch witness can exist only if all three matrix equations hold.

Independently, the final Hadamard relation \eqref{eq:r1cs-final-hadamard-claim} is itself a component of $\mathfrak B_{\mathrm{R1CS}}$.  A joint witness for the heterogeneous batch contains decoded tables satisfying every primitive claim, hence in particular $a\had b=c$.  Combining the four relations gives \eqref{eq:r1cs-joint-all-relations}.
\end{proof}

The previous lemma deliberately charges indexed-selection errors only for one false matrix equation.  We do not need a union bound over all six selections to prove soundness: if the decoded R1CS instance is false because of a matrix equation, one fixed false matrix already suffices to rule out a joint witness.  If all three matrix equations are true but $a\had b\ne c$, the final Hadamard claim alone suffices.

\begin{theorem}[Sparse R1CS soundness]
\label{thm:r1cs-soundness}
Assume
\begin{equation}
  p>N,
  \qquad
  0<\delta\le\frac{1-\rho}{2},
  \qquad
  2N<(1-\delta)M.
  \label{eq:r1cs-parameters}
\end{equation}
If $(w_z,w_a,w_b,w_c)$ has no witness in $\mathcal R^{\delta}_{\mathrm{R1CS}}[A,B,C]$, then every deterministic prover is accepted by $\Pi_{\mathrm{R1CS}}$ with probability at most
\begin{equation}
\begin{split}
  \varepsilon_{\mathrm{R1CS}}
  \le{}&
  \frac{N-1}{|\F|}
  +2\varepsilon_{\mathrm{Sel,out}}
  +\frac{2(N-1)}{|\F|}
  +\frac{(s_{\mathrm{R1CS}}-1)M}{|\F|}
  \\
  &+\varepsilon_{\mathrm{fold}}
  +(1-\delta)^\kappa.
\end{split}
\label{eq:r1cs-soundness-bound}
\end{equation}
The same bound holds for randomized provers by averaging.  Under collision-free extension tagging, delete the two $\binom{2N}{2}/|\F|$ contributions contained in $2\varepsilon_{\mathrm{Sel,out}}$.
\end{theorem}

\begin{proof}
By \cref{prop:unique-table-near-word}, each of the four original commitment words has at most one table within fibre distance $\delta$.  If one has no such table, the heterogeneous batch cannot have a compatible joint witness for all occurrences of that commitment; \cref{thm:heterogeneous-soundness} then gives the batching, folding, and query terms in \eqref{eq:r1cs-soundness-bound}.

Assume therefore that $w_z,w_a,w_b,w_c$ uniquely decode to $z,a,b,c$.  Since there is no R1CS witness, by \cref{prop:r1cs-equivalence} at least one of the three matrix equations or the final Hadamard equation is false.

First suppose at least one matrix equation is false.  By \cref{lem:r1cs-shared-alpha}, all three compressed matrix identities can hold for at most $N-1$ values of the shared challenge $\alpha$, contributing the first term of \eqref{eq:r1cs-soundness-bound}.  Fix $\alpha$ outside this set and choose the first false matrix equation.  Only the two indexed selections belonging to that sparse reduction need to be faithful.  Their tag-collision and bad-shift events contribute at most $2\varepsilon_{\mathrm{Sel,out}}$.  Outside those events, \cref{lem:r1cs-joint-witness} shows that the master heterogeneous batch has no joint witness.

If instead all three matrix equations hold, then $a\had b\ne c$.  In this case neither the $\alpha$ term nor the indexed-selection terms are needed, but we retain them in the uniform upper bound.  The final Hadamard relation is false, so the batch has no joint witness.

In either case, outside the outer events already charged, apply \cref{thm:heterogeneous-soundness} to $\mathfrak B_{\mathrm{R1CS}}$.  Since the batch contains Hadamard relations, its remaining error is
\[
  \frac{2(N-1)}{|\F|}
  +\frac{(s_{\mathrm{R1CS}}-1)M}{|\F|}
  +\varepsilon_{\mathrm{fold}}
  +(1-\delta)^\kappa.
\]
The first of these terms is the one shared Hadamard fingerprint error for all sixteen Hadamard claims, including the final R1CS multiplication constraint.  Adding the outer probabilities proves \eqref{eq:r1cs-soundness-bound}.  Randomized provers follow by conditioning on their internal randomness and averaging.
\end{proof}

\begin{corollary}[R1CS knowledge by decoding]
\label{cor:r1cs-knowledge}
Under the hypotheses of \cref{thm:r1cs-soundness}, acceptance with probability larger than the right-hand side of \eqref{eq:r1cs-soundness-bound} implies that the unique decodings of $w_z,w_a,w_b,w_c$ exist and satisfy
\[
  a=Az,
  \qquad
  b=Bz,
  \qquad
  c=Cz,
  \qquad
  a\had b=c.
\]
Hence the decoded $z$ satisfies the ordinary R1CS equation \eqref{eq:r1cs-standard-relation}.  The witness is recoverable in polynomial time by Reed--Solomon unique decoding.
\end{corollary}

\begin{proof}
This is the contrapositive of \cref{thm:r1cs-soundness}, followed by \cref{prop:r1cs-equivalence}.  Polynomial-time recoverability is the same unique-decoding statement used in \cref{cor:sparse-matvec-knowledge}.
\end{proof}

\subsection{Public inputs and the constant wire}
\label{sec:r1cs-public-inputs}

An NP statement normally fixes part of $z$ publicly, for example
\[
  z(0)=1
\]
and $z(j_t)=x_t$ for public input coordinates $j_1,\ldots,j_m$.  These conditions do not require a new protocol primitive.

\begin{proposition}[Public coordinate constraints are public linear functionals]
\label{prop:r1cs-public-coordinates}
Let $J=\{j_1,\ldots,j_m\}\subseteq\{0,\ldots,N-1\}$ be public and let $x_1,\ldots,x_m\in\F$.  For a challenge $\lambda\in\F$, define the public table
\begin{equation}
  p_{J,\lambda}(j)
  =\begin{cases}
    \lambda^{t-1},&j=j_t,\\
    0,&j\notin J.
  \end{cases}
  \label{eq:r1cs-public-selector-table}
\end{equation}
Then the public-input constraints $z(j_t)=x_t$ for all $t$ imply
\begin{equation}
  \ip{p_{J,\lambda}}{z}
  =\sum_{t=1}^{m}\lambda^{t-1}x_t.
  \label{eq:r1cs-public-input-fingerprint}
\end{equation}
If at least one coordinate constraint is false, \eqref{eq:r1cs-public-input-fingerprint} holds for at most $m-1$ values of $\lambda$.
\end{proposition}

\begin{proof}
The displayed equality expands to
\[
  \sum_{t=1}^{m}\lambda^{t-1}z(j_t)
  =\sum_{t=1}^{m}\lambda^{t-1}x_t.
\]
If some coordinate is wrong, their difference is a nonzero polynomial in $\lambda$ of degree at most $m-1$, so the root bound applies.
\end{proof}

The claim \eqref{eq:r1cs-public-input-fingerprint} is a public linear-functional claim of \cref{sec:linear-functionals}; its protected words can simply be appended to $\mathfrak B_{\mathrm{R1CS}}$ before the master challenge.  Thus the constant wire and public-input binding also share the same folding test.  If this optional check is used, add $(m-1)/|\F|$ to the outer soundness error.

\subsection{Cost and structural consequence}
\label{sec:r1cs-cost}

At the relation level the reduction is
\begin{equation}
\boxed{
  \mathrm{R1CS}(A,B,C;z)
  =3\times\mathrm{SparseMatVec}
   +1\times\mathrm{Hadamard}.
}
\label{eq:r1cs-decomposition-box}
\end{equation}
Using \eqref{eq:sparse-matvec-decomposition}, this becomes six indexed selections, three sparse contribution Hadamards, three committed inner products, three public geometric linear functionals, and one final Hadamard relation.  Every object is a table-form Kronecker commitment, a public succinct kernel, or a virtual word derived locally from such commitments.

If the matrices have $K_A,K_B,K_C$ nonzero entries and the witness dimension is $d$, each sparse representation uses an ambient power of two
\[
  N_M=2^{\lceil\log_2\max\{K_M,d\}\rceil}
\]
(or a common power of two large enough for all three, as assumed in this section).  The public matrix data are therefore linear in the sparse input size rather than quadratic in $d$.  The online prover builds the sparse auxiliary tables in linear time in the padded sparse lengths, before Reed--Solomon encoding and the common folding backend.

The cryptographic distinction from a Sumcheck-based R1CS reduction is now precise.  We do not modify Sumcheck or claim that Sumcheck itself requires post-quantum repair.  Instead, the multilinear algebra normally delegated to Sumcheck has been decomposed into coefficient-extraction, lookup, and Hadamard relations whose low-degree validity is enforced directly by one Reed--Solomon/FRI proximity argument.  The remaining questions for a complete SNARK concern compilation, zero knowledge, and concrete efficiency, not an unresolved R1CS algebraic reduction.

\section{Memory Consistency Arguments}
\label{sec:memory}

Memory consistency is a natural test of whether the preceding primitives form a sufficiently expressive proof-system backend.  A memory trace contains a sequence of accesses in execution order; correctness is a global temporal property, because the value returned by a read depends on the most recent preceding access to the same address.  Standard offline memory-checking arguments therefore introduce a second copy of the accesses sorted by address and time, prove that the two copies are permutations of one another, and check local transition rules on the sorted copy.  This section gives a self-contained reduction of that strategy to the primitives developed above.  The resulting argument uses only multiset equality, lookup, public linear functionals, and Hadamard relations, and hence all protected low-degree words can be placed in one heterogeneous folding batch.

The goal of the section is algebraic.  We prove equivalence between an ordinary finite-memory execution and the committed relation below, and we give an explicit interactive soundness bound.  We do not claim here that this construction is already the most efficient memory argument; the implementation comparison with specialized constructions such as Jolt and Twist--Shout is deferred to the experimental section \cite{Jolt23,TwistShout26}.

\subsection{Operational memory semantics}
\label{sec:memory-operational}

Fix an address-space size $A$ with
\begin{equation}
  1\le A\le N,
  \qquad
  p=\operatorname{char}(\F)>2\max\{A,N\}.
  \label{eq:memory-characteristic}
\end{equation}
We identify the integer intervals $\{0,\ldots,A-1\}$ and $\{0,\ldots,N-1\}$ with their canonical images in $\Fq\subseteq\F$.  Let
\begin{equation}
  \mu_0:\{0,\ldots,A-1\}\longrightarrow\F
  \label{eq:memory-initial-state}
\end{equation}
be a public initial memory state.

An execution trace of length $N$ consists of tables
\begin{equation}
  a,\omega,r,d:\B^n\longrightarrow\F,
  \label{eq:memory-execution-tables}
\end{equation}
where, for access index $i\in\{0,\ldots,N-1\}$,
\begin{itemize}[leftmargin=*]
  \item $a(i)$ is the address;
  \item $\omega(i)\in\{0,1\}$ is the write flag;
  \item $r(i)$ is the value observed immediately before the access;
  \item $d(i)$ is the proposed new value when $\omega(i)=1$.
\end{itemize}
The execution time is the public index $i$ itself.  A read has $\omega(i)=0$ and leaves memory unchanged; a write has $\omega(i)=1$ and changes the addressed cell to $d(i)$.  It is convenient to define the post-access value
\begin{equation}
  p(i)
  \defeq
  r(i)+\omega(i)\bigl(d(i)-r(i)\bigr).
  \label{eq:memory-post-value}
\end{equation}
Thus $p(i)=r(i)$ for a read and $p(i)=d(i)$ for a write.

\begin{definition}[Operationally consistent memory trace]
\label{def:memory-operational-consistency}
The trace $(a,\omega,r,d)$ is \emph{operationally consistent with $\mu_0$} if $a(i)\in\{0,\ldots,A-1\}$ and $\omega(i)\in\{0,1\}$ for every $i$, and there exist memory states
\[
  \mu_i:\{0,\ldots,A-1\}\to\F,
  \qquad 0\le i\le N,
\]
with initial state \eqref{eq:memory-initial-state} such that, for every access $i$,
\begin{align}
  r(i)&=\mu_i(a(i)),
  \label{eq:memory-read-semantics}\\
  \mu_{i+1}(x)&=\mu_i(x)
       \quad\text{for }x\ne a(i),
  \label{eq:memory-other-cells}\\
  \mu_{i+1}(a(i))&=p(i).
  \label{eq:memory-address-update}
\end{align}
\end{definition}

A shorter trace may be padded to a power of two with read-only accesses to a distinguished valid address.  Since such accesses do not change the state, we work with length exactly $N$ without loss of generality at the algebraic level.

\subsection{The sorted-trace characterization}
\label{sec:memory-sorted-characterization}

The prover supplies a second access table
\begin{equation}
  (\bar a,\bar t,\bar\omega,\bar r,\bar d):\B^n\longrightarrow\F^5,
  \label{eq:memory-sorted-trace}
\end{equation}
intended to contain exactly the execution records
\begin{equation}
  (a(i),i,\omega(i),r(i),d(i))
  \label{eq:memory-execution-record}
\end{equation}
ordered lexicographically by address and then by execution time.  Put
\begin{equation}
  \bar p(k)
  \defeq
  \bar r(k)+\bar\omega(k)\bigl(\bar d(k)-\bar r(k)\bigr).
  \label{eq:memory-sorted-post}
\end{equation}

For $0\le k<N-1$, define the adjacent differences
\begin{align}
  \Delta_a(k)&\defeq \bar a(k+1)-\bar a(k),
  \label{eq:memory-address-difference}\\
  \Delta_t(k)&\defeq \bar t(k+1)-\bar t(k).
  \label{eq:memory-time-difference}
\end{align}
We set $\Delta_a(N-1)=\Delta_t(N-1)=0$ only to obtain length-$N$ tables.

Let $s:\B^n\to\F$ be the claimed same-address indicator, with $s(N-1)=0$, and let $u:\B^n\to\F$ be an auxiliary inverse table.  On $k<N-1$ impose
\begin{align}
  s(k)\Delta_a(k)&=0,
  \label{eq:memory-same-address-zero}\\
  \Delta_a(k)u(k)&=1-s(k).
  \label{eq:memory-same-address-inverse}
\end{align}
These two equations force $s(k)=1$ exactly when the adjacent addresses are equal.

\begin{lemma}[Exact same-address indicator]
\label{lem:memory-same-address-indicator}
Assume \eqref{eq:memory-same-address-zero} and \eqref{eq:memory-same-address-inverse}.  Then for every $k<N-1$,
\[
  s(k)=1 \iff \Delta_a(k)=0,
  \qquad
  s(k)=0 \iff \Delta_a(k)\ne0.
\]
In particular $s(k)\in\{0,1\}$ without a separate Booleanity constraint.
\end{lemma}

\begin{proof}
If $\Delta_a(k)=0$, equation \eqref{eq:memory-same-address-inverse} gives $0=1-s(k)$, hence $s(k)=1$.  If $\Delta_a(k)\ne0$, equation \eqref{eq:memory-same-address-zero} gives $s(k)=0$ because $\F$ is a field; then \eqref{eq:memory-same-address-inverse} says $u(k)=\Delta_a(k)^{-1}$.  The converses follow immediately.
\end{proof}

To express strict time order only inside an address block, define
\begin{equation}
  q_t(k)
  \defeq
  s(k)\bigl(\Delta_t(k)-1\bigr)
  \qquad (k<N-1),
  \label{eq:memory-time-range-table}
\end{equation}
with $q_t(N-1)=0$.  We will require
\begin{equation}
  \Delta_a(k)\in\{0,\ldots,A-1\},
  \qquad
  q_t(k)\in\{0,\ldots,N-2\}.
  \label{eq:memory-range-conditions}
\end{equation}
The first range condition makes the address sequence nondecreasing; the second makes timestamps strictly increasing when $s(k)=1$.

\begin{lemma}[Range conditions imply lexicographic order]
\label{lem:memory-lexicographic-order}
Assume every $\bar a(k)$ lies in $\{0,\ldots,A-1\}$, every $\bar t(k)$ lies in $\{0,\ldots,N-1\}$, \eqref{eq:memory-characteristic}, \eqref{eq:memory-same-address-zero}--\eqref{eq:memory-same-address-inverse}, and \eqref{eq:memory-range-conditions}.  Then
\begin{equation}
  (\bar a(k),\bar t(k))
  <_{\mathrm{lex}}
  (\bar a(k+1),\bar t(k+1))
  \label{eq:memory-lex-order}
\end{equation}
for every $k<N-1$ for which the two records are distinct.  More precisely, either $\bar a(k+1)>\bar a(k)$ as ordinary integers, or the addresses are equal and $\bar t(k+1)>\bar t(k)$.
\end{lemma}

\begin{proof}
Because $\bar a(k),\bar a(k+1)\in\{0,\ldots,A-1\}$, the ordinary difference lies in $[-(A-1),A-1]$.  If it were negative, its field representative would be at least $p-(A-1)>A-1$ by \eqref{eq:memory-characteristic}, contradicting the first condition of \eqref{eq:memory-range-conditions}.  Hence $\bar a(k+1)\ge\bar a(k)$.

If the inequality is strict there is nothing more to prove.  If the addresses are equal, \cref{lem:memory-same-address-indicator} gives $s(k)=1$, so $q_t(k)=\Delta_t(k)-1$.  The second range condition says this field element is one of $0,\ldots,N-2$.  Since both timestamps lie in $\{0,\ldots,N-1\}$ and $p>2N$, a negative ordinary difference cannot have such a representative.  Therefore $\Delta_t(k)-1\ge0$ as an integer and $\bar t(k+1)>\bar t(k)$.
\end{proof}

Next define the block-start table $b:\B^n\to\F$ by
\begin{equation}
  b(0)=1,
  \qquad
  b(k)=1-s(k-1)\quad(1\le k<N).
  \label{eq:memory-block-start}
\end{equation}
Let $\iota_0:\B^n\to\F$ be a claimed initial-value table satisfying
\begin{equation}
  \iota_0(k)=\mu_0(\bar a(k)).
  \label{eq:memory-initial-selection}
\end{equation}
Finally define a continuity-difference table $e:\B^n\to\F$ by
\begin{equation}
  e(k)=\bar r(k+1)-\bar p(k)
  \quad(k<N-1),
  \qquad
  e(N-1)=0.
  \label{eq:memory-continuity-difference}
\end{equation}
The local state constraints are
\begin{align}
  \bar\omega\had(\bar\omega-\mathbf1)&=0,
  \label{eq:memory-write-boolean}\\
  \bar\omega\had(\bar d-\bar r)&=\bar p-\bar r,
  \label{eq:memory-post-hadamard}\\
  s\had e&=0,
  \label{eq:memory-continuity-hadamard}\\
  b\had(\bar r-\iota_0)&=0.
  \label{eq:memory-initial-hadamard}
\end{align}
Equation \eqref{eq:memory-continuity-hadamard} carries the post-access value to the next access of the same address; equation \eqref{eq:memory-initial-hadamard} anchors the first access of each address block to the public initial memory.

\begin{definition}[Locally valid sorted memory trace]
\label{def:memory-locally-valid-sorted}
A collection
\[
  (\bar a,\bar t,\bar\omega,\bar r,\bar d,
    \bar p,\Delta_a,\Delta_t,s,u,q_t,b,\iota_0,e)
\]
is \emph{locally valid} if:
\begin{enumerate}[label=(\roman*),leftmargin=*]
  \item $\bar a(k)\in\{0,\ldots,A-1\}$ and $\bar t(k)\in\{0,\ldots,N-1\}$;
  \item \eqref{eq:memory-address-difference}, \eqref{eq:memory-time-difference}, \eqref{eq:memory-block-start}, \eqref{eq:memory-continuity-difference} hold;
  \item \eqref{eq:memory-same-address-zero}, \eqref{eq:memory-same-address-inverse}, \eqref{eq:memory-time-range-table}, and \eqref{eq:memory-range-conditions} hold;
  \item \eqref{eq:memory-initial-selection} holds;
  \item all four relations \eqref{eq:memory-write-boolean}--\eqref{eq:memory-initial-hadamard} hold.
\end{enumerate}
\end{definition}

\begin{theorem}[Sorted-trace characterization of memory consistency]
\label{thm:memory-sorted-characterization}
Assume \eqref{eq:memory-characteristic}.  An execution trace $(a,\omega,r,d)$ is operationally consistent with $\mu_0$ if and only if there exists a locally valid sorted trace in the sense of \cref{def:memory-locally-valid-sorted} whose record multiset is exactly
\begin{equation}
  \bigl\{\!\!\bigl\{
    (a(i),i,\omega(i),r(i),d(i)):
    0\le i<N
  \bigr\}\!\!\bigr\}.
  \label{eq:memory-record-multiset}
\end{equation}
\end{theorem}

\begin{proof}
Suppose first that the execution trace is operationally consistent.  Sort the $N$ execution records lexicographically by the ordinary integer pair $(a(i),i)$ and write the resulting tables as in \eqref{eq:memory-sorted-trace}.  Since the execution times are distinct, this order is strict within every address block.  Define $\bar p$ by \eqref{eq:memory-sorted-post}, define the adjacent differences by \eqref{eq:memory-address-difference}--\eqref{eq:memory-time-difference}, let $s$ be the actual same-address indicator, set $u(k)=\Delta_a(k)^{-1}$ when $s(k)=0$ and arbitrarily when $s(k)=1$, and define $q_t,b,\iota_0,e$ by the displayed equations above.  Ordinary sortedness gives the two range conditions.  Equations \eqref{eq:memory-same-address-zero} and \eqref{eq:memory-same-address-inverse} hold by construction.  The write-flag and post-value equations follow from the operational update rule.  If two adjacent sorted records have the same address, there is no access to that address at an intermediate execution time, so the memory value immediately before the later record is exactly the post-access value of the earlier record; hence \eqref{eq:memory-continuity-hadamard}.  At the first record of an address block there is no earlier access to that address, so its read value is $\mu_0$ at that address; hence \eqref{eq:memory-initial-hadamard}.  Thus the sorted trace is locally valid and has record multiset \eqref{eq:memory-record-multiset}.

Conversely, suppose a locally valid sorted trace with multiset \eqref{eq:memory-record-multiset} exists.  By \cref{lem:memory-lexicographic-order}, the sorted records are nondecreasing by address and strictly increasing in execution time within every fixed address.  Because their multiset equals \eqref{eq:memory-record-multiset}, each execution time $i\in\{0,\ldots,N-1\}$ occurs exactly once, attached to the original access record at time $i$.

Fix an address $x$.  Its records therefore form one contiguous block
\[
  k_1<k_2<\cdots<k_m
\]
in sorted order, and their timestamps satisfy
\[
  \bar t(k_1)<\bar t(k_2)<\cdots<\bar t(k_m).
\]
By \eqref{eq:memory-block-start}, $b(k_1)=1$, so \eqref{eq:memory-initial-hadamard} gives
\[
  \bar r(k_1)=\mu_0(x).
\]
For every $j<m$, one has $s(k_j)=1$, and \eqref{eq:memory-continuity-hadamard} together with \eqref{eq:memory-continuity-difference} gives
\[
  \bar r(k_{j+1})=\bar p(k_j).
\]
Equation \eqref{eq:memory-write-boolean} makes every write flag Boolean, and \eqref{eq:memory-post-hadamard} gives exactly the read/write transition \eqref{eq:memory-post-value}.  Hence, along the strictly increasing execution times at address $x$, each access reads the state left by the preceding access to $x$, or $\mu_0(x)$ if there is none.

Define $\mu_i(x)$ to be $\mu_0(x)$ if no access to $x$ occurs before time $i$, and otherwise the post-value of the latest access to $x$ with time $<i$.  The preceding paragraph shows that the access at time $i$ reads $\mu_i(a(i))$ and leaves exactly $p(i)$ at that address.  All other addresses retain their latest previous post-value.  Therefore \eqref{eq:memory-read-semantics}--\eqref{eq:memory-address-update} hold for every $i$, proving operational consistency.
\end{proof}

\subsection{Authenticating the sorted trace}
\label{sec:memory-authentication}

We now reduce the conditions of \cref{thm:memory-sorted-characterization} to the committed primitives of the preceding sections.

\paragraph{Tuple permutation.}
For a challenge $\eta\in\F$, compress a five-component record by
\begin{equation}
  \chi_{\eta}(x_0,x_1,x_2,x_3,x_4)
  \defeq
  x_0+\eta x_1+\eta^2x_2+\eta^3x_3+\eta^4x_4.
  \label{eq:memory-tuple-compression}
\end{equation}
Define the execution and sorted compressed tables
\begin{align}
  c_{\mathrm{exec},\eta}(i)
  &\defeq
  \chi_\eta(a(i),i,\omega(i),r(i),d(i)),
  \label{eq:memory-exec-compression}\\
  c_{\mathrm{sort},\eta}(k)
  &\defeq
  \chi_\eta(\bar a(k),\bar t(k),\bar\omega(k),\bar r(k),\bar d(k)).
  \label{eq:memory-sort-compression}
\end{align}
Both commitment words are virtual linear combinations of the component commitments and the public timestamp table.

\begin{lemma}[Tuple compression detects a false record permutation]
\label{lem:memory-tuple-compression}
If the two multisets of five-component records differ, there is a set $T_{\mathrm{tuple}}\subseteq\F$ with
\begin{equation}
  |T_{\mathrm{tuple}}|
  \le
  4\binom{2N}{2}
  \label{eq:memory-tuple-collision-bound}
\end{equation}
such that for every $\eta\notin T_{\mathrm{tuple}}$ the scalar multisets
$\{\!\{c_{\mathrm{exec},\eta}(i)\}\!\}$ and
$\{\!\{c_{\mathrm{sort},\eta}(k)\}\!\}$ differ.
\end{lemma}

\begin{proof}
Consider the union of the supports of the two record multisets; it contains at most $2N$ distinct tuples.  For two distinct tuples $x\ne y$, the equality $\chi_\eta(x)=\chi_\eta(y)$ is the vanishing of a nonzero polynomial in $\eta$ of degree at most $4$, hence has at most four solutions.  Let $T_{\mathrm{tuple}}$ be the union of these collision sets over unordered pairs of distinct tuples.  This gives \eqref{eq:memory-tuple-collision-bound}.  Outside $T_{\mathrm{tuple}}$, $\chi_\eta$ is injective on the union of supports; equality of the compressed scalar multisets would then imply equality of the original tuple multisets, contrary to hypothesis.
\end{proof}

Thus the record permutation is reduced to the scalar multiset argument of \cref{sec:committed-multiset} applied to \eqref{eq:memory-exec-compression} and \eqref{eq:memory-sort-compression}.

\paragraph{Range lookups.}
Let $R_A$ be a public length-$N$ table containing every element of $\{0,\ldots,A-1\}$ at least once and arbitrary repetitions thereafter.  Let $R_{\Delta A}$ be the same range table, and let $R_T$ contain every element of $\{0,\ldots,N-2\}$ at least once.  Membership of the execution addresses $a$, adjacent differences $\Delta_a$, and conditional time differences $q_t$ in their respective integer ranges is proved by three instances of the lookup relation of \cref{sec:lookup}.  Multiplicity tables are prover supplied and decoded exactly as in \cref{def:lookup-relation}.  The tuple permutation then transfers address validity from the execution trace to the sorted trace and transfers the public timestamp multiset $\{0,\ldots,N-1\}$ to $\bar t$.

\paragraph{Initial-memory selection.}
The address map in \eqref{eq:memory-initial-selection} is witness dependent, so the public-index-map specialization of \cref{sec:indexed-selection} does not apply directly.  Instead, sample a tag $\tau_0\in\F$ and form
\begin{align}
  f_0(k)&\defeq \iota_0(k)+\tau_0\bar a(k),
  \label{eq:memory-initial-query-tag}\\
  t_0(j)&\defeq \mu_0(j)+\tau_0 j
  \qquad(0\le j<A),
  \label{eq:memory-initial-table-tag}
\end{align}
with the public table padded to length $N$.  A logarithmic-derivative lookup from $f_0$ into $t_0$ proves \eqref{eq:memory-initial-selection}, except when two distinct address--value pairs collide under the random tag.

\begin{lemma}[Soundness of public initial-memory selection]
\label{lem:memory-initial-selection}
Suppose \eqref{eq:memory-initial-selection} is false.  Then there is a set $T_{\mathrm{init}}\subseteq\F$ with
\begin{equation}
  |T_{\mathrm{init}}|
  \le\binom{N+A}{2}
  \label{eq:memory-initial-tag-bound}
\end{equation}
such that for every $\tau_0\notin T_{\mathrm{init}}$, the tagged query table \eqref{eq:memory-initial-query-tag} cannot satisfy the algebraic lookup relation against \eqref{eq:memory-initial-table-tag} for any multiplicity table.
\end{lemma}

\begin{proof}
Consider the multiset of query pairs $(\bar a(k),\iota_0(k))$ and the set of public pairs $(j,\mu_0(j))$, $0\le j<A$.  Their union contains at most $N+A$ distinct pairs.  Two distinct pairs $(i,x)$ and $(j,y)$ collide under $x+\tau_0 i$ and $y+\tau_0 j$ for at most one value of $\tau_0$, exactly as in \cref{lem:indexed-selection-compression}.  Outside the union $T_{\mathrm{init}}$ of all such values, scalar compression is injective on the union of supports.  If the compressed lookup relation held, every query pair would therefore equal a public pair with the same first coordinate, forcing $\iota_0(k)=\mu_0(\bar a(k))$ for every $k$, contrary to hypothesis.
\end{proof}

As in \cref{rem:collision-free-index-tagging}, extension-field tagging may remove the collision term \eqref{eq:memory-initial-tag-bound} entirely when addresses and values lie in a proper subfield.

\subsection{Linear adjacency identities}
\label{sec:memory-linear-identities}

The definitions of $\Delta_a,\Delta_t,b,e$ are linear but involve adjacent rows.  They can be checked without a new primitive by one random linear fingerprint.

For a table $x$ and a scalar $\lambda\in\F$, write
\begin{equation}
  \mathcal G_\lambda(x)
  \defeq
  \sum_{k=0}^{N-1}\lambda^k x(k).
  \label{eq:memory-geometric-functional}
\end{equation}
This is a public linear functional of \cref{sec:linear-functionals}.  Define five residual tables by
\begin{align}
  \ell_a(k)
  &\defeq
  \begin{cases}
    \Delta_a(k)-\bar a(k+1)+\bar a(k),&k<N-1,\\
    \Delta_a(N-1),&k=N-1,
  \end{cases}
  \label{eq:memory-residual-a}\\
  \ell_t(k)
  &\defeq
  \begin{cases}
    \Delta_t(k)-\bar t(k+1)+\bar t(k),&k<N-1,\\
    \Delta_t(N-1),&k=N-1,
  \end{cases}
  \label{eq:memory-residual-t}\\
  \ell_e(k)
  &\defeq
  \begin{cases}
    e(k)-\bar r(k+1)+\bar p(k),&k<N-1,\\
    e(N-1),&k=N-1,
  \end{cases}
  \label{eq:memory-residual-e}\\
  \ell_b(k)
  &\defeq
  \begin{cases}
    b(0)-1,&k=0,\\
    b(k)+s(k-1)-1,&1\le k<N,
  \end{cases}
  \label{eq:memory-residual-b}\\
  \ell_q(k)
  &\defeq
  \begin{cases}
    0,&k<N-1,\\
    q_t(N-1),&k=N-1.
  \end{cases}
  \label{eq:memory-residual-q}
\end{align}
Every geometric sum $\mathcal G_\lambda(\ell_\bullet)$ expands into public linear functionals of the committed component tables; the shifted terms merely change the public weight from $\lambda^k$ to $\lambda^{k-1}$ or $\lambda^{k+1}$ on the appropriate coordinates.

\begin{lemma}[One fingerprint for all adjacency identities]
\label{lem:memory-linear-fingerprint}
Suppose at least one of the five residual tables in \eqref{eq:memory-residual-a}--\eqref{eq:memory-residual-q} is nonzero.  Sample $\chi,\lambda\leftarrow\F$ and test
\begin{equation}
  \sum_{j=0}^{4}\chi^j\mathcal G_\lambda(\ell_j)=0,
  \qquad
  (\ell_0,\ell_1,\ell_2,\ell_3,\ell_4)
  =(\ell_a,\ell_t,\ell_e,\ell_b,\ell_q).
  \label{eq:memory-linear-fingerprint}
\end{equation}
Then \eqref{eq:memory-linear-fingerprint} holds with probability at most
\begin{equation}
  \frac{4}{|\F|}+\frac{N-1}{|\F|}.
  \label{eq:memory-linear-fingerprint-error}
\end{equation}
\end{lemma}

\begin{proof}
Choose a coordinate $k_0$ at which not all five residual values vanish.  The polynomial
\[
  R_{k_0}(Z)=\sum_{j=0}^{4}Z^j\ell_j(k_0)
\]
is nonzero of degree at most $4$, so for all but at most four values of $\chi$ the combined residual table
\[
  \ell_\chi(k)=\sum_{j=0}^{4}\chi^j\ell_j(k)
\]
is nonzero.  Fix such a $\chi$.  Then $\mathcal G_\lambda(\ell_\chi)$ is a nonzero polynomial in $\lambda$ of degree at most $N-1$, and therefore vanishes for at most $N-1$ values of $\lambda$.  A union bound proves \eqref{eq:memory-linear-fingerprint-error}.
\end{proof}

Thus all adjacency bookkeeping, including the terminal condition $q_t(N-1)=0$, costs one scalar public-linear-functional equality and the outer error \eqref{eq:memory-linear-fingerprint-error}.

\subsection{The memory sub-batch and protocol}
\label{sec:memory-protocol}

All remaining identities are pointwise multiplications and therefore Hadamard relations.  Besides \eqref{eq:memory-write-boolean}--\eqref{eq:memory-initial-hadamard}, we include
\begin{align}
  s\had\Delta_a&=0,
  \label{eq:memory-batch-same-zero}\\
  \Delta_a\had u&=\mathbf1-s,
  \label{eq:memory-batch-same-inverse}\\
  s\had(\Delta_t-\mathbf1)&=q_t.
  \label{eq:memory-batch-time}
\end{align}
Affine combinations such as $\bar\omega-\mathbf1$, $\bar d-\bar r$, and $\mathbf1-s$ are virtual linear words and require no new commitment.

\begin{definition}[Memory consistency sub-batch]
\label{def:memory-sub-batch}
Fix all outer compression, lookup-shift, and fingerprint challenges.  The heterogeneous batch $\mathfrak B_{\mathrm{Mem}}$ consists of:
\begin{enumerate}[label=(\roman*),leftmargin=*]
  \item the scalar multiset sub-batch proving equality of \eqref{eq:memory-exec-compression} and \eqref{eq:memory-sort-compression};
  \item four lookup sub-batches: execution-address range, $\Delta_a$ range, $q_t$ range, and public initial-memory selection \eqref{eq:memory-initial-query-tag}--\eqref{eq:memory-initial-table-tag};
  \item the seven Hadamard relations \eqref{eq:memory-write-boolean}, \eqref{eq:memory-post-hadamard}, \eqref{eq:memory-continuity-hadamard}, \eqref{eq:memory-initial-hadamard}, and \eqref{eq:memory-batch-same-zero}--\eqref{eq:memory-batch-time};
  \item the public linear-functional equality \eqref{eq:memory-linear-fingerprint}.
\end{enumerate}
Exact repetitions and affine virtual words are deduplicated before the master batching challenge is sampled.
\end{definition}

\begin{definition}[Memory protocol $\Pi_{\mathrm{Mem}}$]
\label{def:memory-protocol}
The public data are $A$, $\mu_0$, the range tables, and the public execution-time table $i\mapsto i$.  The original execution commitments $(w_a,w_\omega,w_r,w_d)$ are fixed before any challenge.
\begin{enumerate}[label=\textnormal{\arabic*.},leftmargin=*]
  \item The prover commits to the sorted-trace and auxiliary tables appearing in \cref{def:memory-locally-valid-sorted}.
  \item The verifier samples the tuple-compression challenge $\eta$ and the initial-memory tag $\tau_0$.
  \item The verifier samples the logarithmic-derivative shifts required by the multiset and four lookup reductions; the prover supplies the corresponding inverse tables and scalar messages.
  \item The verifier samples the linear-relation challenges $(\chi,\lambda)$ and the shared Hadamard fingerprint challenges $(\gamma,\theta)$ of \cref{sec:heterogeneous-batching}.
  \item The parties execute exactly one
  \begin{equation}
    \Pi_{\mathrm{Het}}(\mathfrak B_{\mathrm{Mem}}).
    \label{eq:memory-one-folding-batch}
  \end{equation}
\end{enumerate}
\end{definition}

The protocol therefore has one master batching challenge, one folding chain, and one final set of $\kappa$ proximity queries.  The earlier challenges only reduce tuple equality, range membership, and adjacency semantics to the low-degree relations protected by that one batch.

\subsection{Completeness and soundness}
\label{sec:memory-soundness}

\begin{theorem}[Memory completeness]
\label{thm:memory-completeness}
If $(a,\omega,r,d)$ is operationally consistent with $\mu_0$ and the original words are honest table-form commitments, then there are honest prover messages for \cref{def:memory-protocol}.  Conditional on all logarithmic-derivative shifts avoiding their pole sets and all DEEP points avoiding their poles, the verifier accepts with probability $1$.
\end{theorem}

\begin{proof}
By \cref{thm:memory-sorted-characterization}, sort the honest execution records and construct the locally valid auxiliary tables.  The tuple multisets are exactly equal, so their scalar compressions are equal for every $\eta$.  Each range lookup is true.  The initial-memory tagged lookup is true for every $\tau_0$ because every query pair is an actual public pair $(\bar a(k),\mu_0(\bar a(k)))$.  The five residual tables are zero, and all seven Hadamard relations hold by local validity.  Therefore every primitive claim in $\mathfrak B_{\mathrm{Mem}}$ is true.  Conditional perfect completeness follows from \cref{thm:heterogeneous-completeness} and the completeness theorems for the multiset and lookup reductions.
\end{proof}

Let $s_{\mathrm{Mem}}$ denote the number of distinct protected degree-$<N$ words in the deduplicated batch.  Let
\begin{equation}
  \varepsilon_{\mathrm{LD}}
  \defeq
  \frac{2N-1}{|\F|}
  \label{eq:memory-one-lookup-bad-shift}
\end{equation}
be the uniform bad-shift bound for each length-$N$ logarithmic-derivative multiset/lookup instance from \cref{lem:lookup-bad-challenges,lem:multiset-bad-shifts}.

\begin{definition}[Committed memory relation]
\label{def:memory-committed-relation}
Fix $0<\delta\le(1-\rho)/2$.  The relation $\mathcal R^{\delta}_{\mathrm{Mem}}[\mu_0]$ consists of pairs
\[
  ((w_a,w_\omega,w_r,w_d),(a,\omega,r,d))
\]
such that each original word is within fibre distance $\delta$ of the corresponding table-form commitment and $(a,\omega,r,d)$ is operationally consistent with $\mu_0$ in the sense of \cref{def:memory-operational-consistency}.
\end{definition}

\begin{lemma}[A joint memory-batch witness gives a valid sorted trace]
\label{lem:memory-joint-witness}
Fix all outer challenges.  Suppose $\mathfrak B_{\mathrm{Mem}}$ has a joint witness, the tuple-compression challenge is outside $T_{\mathrm{tuple}}$, the initial-memory tag is outside $T_{\mathrm{init}}$, every logarithmic-derivative shift lies outside the bad-shift set of its decoded false relation, and $(\chi,\lambda)$ lies outside the bad set of \cref{lem:memory-linear-fingerprint}.  Then the decoded original trace is operationally consistent with $\mu_0$.
\end{lemma}

\begin{proof}
A joint witness satisfies every primitive relation in \cref{def:memory-sub-batch}.  The scalar multiset relation, together with $\eta\notin T_{\mathrm{tuple}}$ and \cref{lem:memory-tuple-compression}, gives equality of the original and sorted record multisets.  The three range lookups give the address, $\Delta_a$, and $q_t$ range conditions.  The tagged initial-memory lookup and $\tau_0\notin T_{\mathrm{init}}$ give \eqref{eq:memory-initial-selection} by \cref{lem:memory-initial-selection}.  The linear fingerprint outside its bad set forces all five residual tables to vanish, hence \eqref{eq:memory-address-difference}, \eqref{eq:memory-time-difference}, \eqref{eq:memory-block-start}, \eqref{eq:memory-continuity-difference}, and $q_t(N-1)=0$.  Together with \eqref{eq:memory-batch-time} and $\Delta_t(N-1)=0$, this also forces $s(N-1)=0$.  The seven Hadamard components give the remaining local relations, including \eqref{eq:memory-time-range-table}.  Thus the decoded auxiliary tables form a locally valid sorted trace.  Operational consistency follows from \cref{thm:memory-sorted-characterization}.
\end{proof}

\begin{theorem}[Memory soundness]
\label{thm:memory-soundness}
Assume
\begin{equation}
  0<\delta\le\frac{1-\rho}{2},
  \qquad
  2N<(1-\delta)M,
  \qquad
  p>2\max\{A,N\}.
  \label{eq:memory-soundness-parameters}
\end{equation}
If $(w_a,w_\omega,w_r,w_d)$ has no witness in $\mathcal R^{\delta}_{\mathrm{Mem}}[\mu_0]$, then every deterministic prover is accepted by $\Pi_{\mathrm{Mem}}$ with probability at most
\begin{equation}
\begin{split}
  \varepsilon_{\mathrm{Mem}}
  \le{}&
  \frac{4\binom{2N}{2}}{|\F|}
  +\frac{\binom{N+A}{2}}{|\F|}
  +5\varepsilon_{\mathrm{LD}}
  +\frac{N+3}{|\F|}
  \\
  &+\frac{2(N-1)}{|\F|}
  +\frac{(s_{\mathrm{Mem}}-1)M}{|\F|}
  +\varepsilon_{\mathrm{fold}}
  +(1-\delta)^\kappa.
\end{split}
\label{eq:memory-soundness-bound}
\end{equation}
The same bound holds for randomized provers by averaging.  Collision-free extension tagging removes the two explicit compression-collision terms.  If lookup shifts are sampled from a larger extension field, replace the five occurrences of $|\F|$ in $\varepsilon_{\mathrm{LD}}$ by that challenge-field size.
\end{theorem}

\begin{proof}
By unique decoding, either one original commitment has no table within fibre distance $\delta$, in which case the heterogeneous batch has no compatible joint witness, or the four original words uniquely decode to a trace $(a,\omega,r,d)$.  Assume the latter.  Since the committed relation is false, this decoded trace is not operationally consistent.

The tuple-compression collision event contributes at most the first term of \eqref{eq:memory-soundness-bound} by \cref{lem:memory-tuple-compression}.  The initial-memory tag collision contributes the second term by \cref{lem:memory-initial-selection}.  There are five logarithmic-derivative instances before the master batch: one scalar multiset check for the record permutation and four lookups.  If any decoded one of these relations is false, \cref{lem:lookup-bad-challenges,lem:multiset-bad-shifts} charge at most $\varepsilon_{\mathrm{LD}}$ for its shift; a union bound gives $5\varepsilon_{\mathrm{LD}}$.  The five linear adjacency identities contribute at most $(4+N-1)/|\F|=(N+3)/|\F|$ by \cref{lem:memory-linear-fingerprint}.

Outside all these outer bad events, if the master batch had a joint witness, \cref{lem:memory-joint-witness} would imply that the decoded trace is operationally consistent, a contradiction.  Hence the heterogeneous batch has no joint witness.  Since it contains Hadamard relations, \cref{thm:heterogeneous-soundness} contributes the remaining terms
\[
  \frac{2(N-1)}{|\F|}
  +\frac{(s_{\mathrm{Mem}}-1)M}{|\F|}
  +\varepsilon_{\mathrm{fold}}
  +(1-\delta)^\kappa.
\]
Adding the outer events proves \eqref{eq:memory-soundness-bound}.  Randomized provers follow by conditioning on their internal randomness and averaging.
\end{proof}

\begin{corollary}[Memory knowledge by decoding]
\label{cor:memory-knowledge}
Under the hypotheses of \cref{thm:memory-soundness}, acceptance with probability larger than the right-hand side of \eqref{eq:memory-soundness-bound} implies that the unique decodings of the original access commitments exist and form an operationally consistent memory trace.  The decoded trace is recoverable in polynomial time by Reed--Solomon unique decoding.
\end{corollary}

\begin{proof}
This is the contrapositive of \cref{thm:memory-soundness} together with polynomial-time unique decoding.
\end{proof}

\subsection{Decomposition and cost}
\label{sec:memory-cost}

At the relation level the construction is
\begin{equation}
\boxed{
\begin{aligned}
  \mathrm{MemoryConsistency}
  ={}&\mathrm{TupleMultiset}
  +4\times\mathrm{Lookup}
  +7\times\mathrm{Hadamard}\\
  &+1\times\mathrm{LinearAdjacencyFingerprint}.
\end{aligned}}
\label{eq:memory-decomposition-box}
\end{equation}
The four lookups are the execution-address range check, the nondecreasing-address difference check, the conditional positive-time-difference check, and the selection of the public initial memory value.  All seven Hadamard relations share the one $(\gamma,\theta)$ fingerprint pair of \cref{sec:heterogeneous-batching}, and every primitive shares the same final folding chain and query set.

The prover constructs a sorted copy of $N$ records and a constant number of length-$N$ auxiliary tables.  Sorting therefore dominates the purely algebraic preprocessing at $O(N\log N)$ comparison work in a conventional implementation; once the order is known, all auxiliary tables are produced in $O(N)$ field operations before Reed--Solomon encoding.  The verifier does not scan the trace.  Its algebraic work consists of public challenge generation, scalar messages, local virtual-word arithmetic at the sampled positions, and the one FRI-style verifier.

The main structural conclusion is that offline memory consistency introduces no proof primitive beyond those already developed for coefficient extraction.  Global temporal correctness is converted into (i) a permutation of access records, (ii) range and initial-state lookups, and (iii) local pointwise transition relations.  Each of those is already expressible through logarithmic derivatives, public linear functionals, and Hadamard products, and therefore inherits the same transparent Reed--Solomon/FRI backend.

\section{Global Soundness and Decoding Knowledge}
\label{sec:soundness}

The preceding sections establish soundness separately for public linear functionals, committed inner products, Hadamard relations, lookup, multiset equality, sparse matrix--vector products, sparse R1CS, and memory consistency.  Every proof has the same shape.  First, relation-specific public challenges reduce the original statement to a heterogeneous batch of low-degree claims.  Second, one master batching challenge protects all low-degree component words simultaneously.  Third, one FRI-style folding chain certifies proximity of the master word to the Reed--Solomon code.  Finally, independent oracle queries test the accepting witness set.  This section isolates that common structure and gives a single composition theorem.

The purpose is twofold.  First, it prevents the application sections from hiding soundness losses inside informal reductions: every source of error is assigned to one explicit layer.  Second, it identifies the decoding state that will be used in the non-interactive compilation developed in the next section.  The result below is still an interactive, information-theoretic statement; no random-oracle assumption is used in this section.

\subsection{Committed relations and unique decoded states}
\label{sec:global-decoded-state}

Throughout this section assume
\begin{equation}
  0<\delta\le \frac{1-\rho}{2},
  \qquad
  2N<(1-\delta)M.
  \label{eq:global-soundness-parameters}
\end{equation}
The first inequality is the unique-decoding regime used throughout the paper, and the second is the root-counting margin required by the coefficient-extraction and Hadamard identity arguments.

Let $\mathbf w=(w_1,\ldots,w_t)$ be the tuple of \emph{original commitment words} of an application.  Auxiliary words sent after verifier challenges are not included in $\mathbf w$.  Each original word is intended to encode one table through \cref{def:table-commitment}.

\begin{definition}[Partial table decoder]
\label{def:partial-table-decoder}
For a word $w\in\F^L$, define
\[
  \operatorname{Dec}_\delta(w)=
  \begin{cases}
    f,&\text{if there is }f:\bits\to\F\text{ with }
    \Delta_{\mathrm{fib}}(w,\Enc_T(f))\le\delta,\\
    \bot,&\text{otherwise.}
  \end{cases}
\]
For a tuple $\mathbf w=(w_1,\ldots,w_t)$ put
\[
  \operatorname{Dec}_\delta(\mathbf w)
  =
  (\operatorname{Dec}_\delta(w_1),\ldots,\operatorname{Dec}_\delta(w_t)).
\]
\end{definition}

\begin{lemma}[Well-definedness and efficient recovery]
\label{lem:partial-decoder-well-defined}
Under \eqref{eq:global-soundness-parameters}, $\operatorname{Dec}_\delta$ is well defined.  Whenever $\operatorname{Dec}_\delta(w)\ne\bot$, it can be recovered in time polynomial in $M$ by Reed--Solomon unique decoding followed by the inverse table-encoding bijection.
\end{lemma}

\begin{proof}
Uniqueness is \cref{prop:unique-table-near-word}.  Since
\[
  \delta\le\frac{1-\rho}{2}
  <\frac12\left(1-\frac{N-1}{M}\right),
\]
standard Reed--Solomon unique decoding recovers the unique degree-$<N$ polynomial whenever it exists \cite{WelchBerlekamp86,Gao03}.  The table is then its coefficient vector by \cref{prop:table-isomorphism}.
\end{proof}

Thus an original commitment tuple determines at most one decoded application state.  This point is fundamental: all later occurrences of one commitment in different subprotocols refer to the same decoded table.  Auxiliary messages may depend on verifier challenges, but they cannot change the decoded state of the original commitments.

\subsection{Front-end reductions}
\label{sec:front-end-reductions}

We now formalize the part of an application that precedes the master heterogeneous batch.  This includes, for example, the logarithmic-derivative shift of a lookup, the tagging challenge of a prescribed permutation, the geometric fingerprint of a sparse matrix equation, and the tuple-compression and ordering challenges of the memory argument.

Let $\mathcal R^\delta$ be a committed relation whose instances contain an original commitment tuple $\mathbf w$.  Let
\[
  \omega=(\omega_1,\ldots,\omega_r)
\]
be the sequence of public \emph{outer challenges}, drawn before the master batching challenge $\beta$.  The prover may send messages between these challenges.  Once the outer transcript is fixed, the application produces a heterogeneous batch instance $\mathfrak B_\omega$ in the sense of \cref{def:heterogeneous-batch-instance}.

\begin{definition}[Sound front-end reduction]
\label{def:sound-front-end-reduction}
A front-end reduction from $\mathcal R^\delta$ to heterogeneous batching has error $\varepsilon_{\mathrm{red}}$ if the following hold for every deterministic prover and every fixed original commitment tuple $\mathbf w$.
\begin{enumerate}[label=(\arabic*),leftmargin=*]
  \item \textbf{Completeness.} If $\mathbf w$ has a witness in $\mathcal R^\delta$, then the honest prover can choose the auxiliary messages so that, for every outer challenge outside the explicitly declared pole sets, $\mathfrak B_\omega$ has a joint witness.
  \item \textbf{Witness lifting.} If $\mathbf w$ has no witness in $\mathcal R^\delta$, there is a set $\mathsf{Bad}_{\mathrm{red}}$ of outer transcripts such that
  \begin{equation}
    \Pr[\omega\in\mathsf{Bad}_{\mathrm{red}}]
    \le \varepsilon_{\mathrm{red}},
    \label{eq:front-end-bad-probability}
  \end{equation}
  and, for every $\omega\notin\mathsf{Bad}_{\mathrm{red}}$, the batch $\mathfrak B_\omega$ has no joint witness.
\end{enumerate}
\end{definition}

The probability in \eqref{eq:front-end-bad-probability} is over the sequential verifier challenges, with the prover's message at each round fixed by the preceding transcript.  Therefore the bad set need not be a Cartesian product.  The sequential union bound of \cref{lem:sequential-bad-challenges} is exactly the tool used in the application sections.

\begin{lemma}[Sequential root-event accounting]
\label{lem:global-sequential-accounting}
Suppose the outer challenges are drawn from finite sets $\Omega_1,\ldots,\Omega_r$.  Assume that, after every fixed prefix $(\omega_1,\ldots,\omega_{i-1})$, at most $b_i$ values of $\omega_i\in\Omega_i$ can cause the $i$-th reduction step to lose witness lifting.  Then the front-end reduction has
\begin{equation}
  \varepsilon_{\mathrm{red}}
  \le
  \sum_{i=1}^r\frac{b_i}{|\Omega_i|}.
  \label{eq:global-sequential-reduction-error}
\end{equation}
\end{lemma}

\begin{proof}
For every $i$, let $E_i$ be the event that the first failure of witness lifting occurs at challenge $i$.  Once the preceding prefix is fixed, at most $b_i$ values of the next challenge belong to $E_i$.  Apply \cref{lem:sequential-bad-challenges} and sum the resulting conditional fractions.
\end{proof}

This lemma covers all algebraic outer errors appearing earlier in the paper: Schwartz--Zippel fingerprints, logarithmic-derivative roots, tuple-compression collisions, tagged-index collisions, and the linear adjacency fingerprint of the memory argument.

\subsection{The master composition theorem}
\label{sec:master-soundness-composition}

For a heterogeneous batch $\mathfrak B$, let $s(\mathfrak B)$ be its number of distinct protected degree-$<N$ words after deduplication, as in \cref{def:protected-component-list}.  Let
\[
  h(\mathfrak B)
  =
  \begin{cases}
    1,&\text{if }\mathfrak B\text{ contains at least one Hadamard claim},\\
    0,&\text{otherwise.}
  \end{cases}
\]
The single bit $h(\mathfrak B)$ is enough because all Hadamard relations in one batch share the same fingerprint pair $(\gamma,\theta)$.

Define the backend error
\begin{equation}
  \varepsilon_{\mathrm{back}}(s,h)
  \defeq
  h\frac{2(N-1)}{|\F|}
  +\frac{(s-1)M}{|\F|}
  +\varepsilon_{\mathrm{fold}}
  +(1-\delta)^\kappa.
  \label{eq:global-backend-error}
\end{equation}

\begin{theorem}[Master interactive soundness theorem]
\label{thm:master-interactive-soundness}
Let $\mathcal R^\delta$ admit a sound front-end reduction with error $\varepsilon_{\mathrm{red}}$.  Suppose every resulting heterogeneous batch has at most $s$ protected words and Hadamard indicator at most $h\in\{0,1\}$.  Under \eqref{eq:global-soundness-parameters}, if an original commitment tuple $\mathbf w$ has no witness in $\mathcal R^\delta$, then every deterministic prover is accepted with probability at most
\begin{equation}
  \boxed{
  \varepsilon_{\mathrm{tot}}
  \le
  \varepsilon_{\mathrm{red}}
  +h\frac{2(N-1)}{|\F|}
  +\frac{(s-1)M}{|\F|}
  +\varepsilon_{\mathrm{fold}}
  +(1-\delta)^\kappa.}
  \label{eq:master-total-soundness}
\end{equation}
The same bound holds for randomized provers whose internal randomness is independent of the verifier challenges.
\end{theorem}

\begin{proof}
Let $E_{\mathrm{red}}$ be the event $\omega\in\mathsf{Bad}_{\mathrm{red}}$.  By \cref{def:sound-front-end-reduction},
\[
  \Pr[E_{\mathrm{red}}]\le\varepsilon_{\mathrm{red}}.
\]
Condition on any outer transcript outside $E_{\mathrm{red}}$.  Since the original relation is false, witness lifting implies that the resulting heterogeneous batch $\mathfrak B_\omega$ has no joint witness.  Apply \cref{thm:heterogeneous-soundness}.  Its conditional acceptance probability is at most
\[
  h(\mathfrak B_\omega)\frac{2(N-1)}{|\F|}
  +\frac{(s(\mathfrak B_\omega)-1)M}{|\F|}
  +\varepsilon_{\mathrm{fold}}
  +(1-\delta)^\kappa,
\]
which is bounded by \eqref{eq:global-backend-error} with the uniform values $s,h$.  Adding $\Pr[E_{\mathrm{red}}]$ proves \eqref{eq:master-total-soundness}.

For a randomized prover, condition on its internal random tape.  The deterministic argument applies to every fixed tape; averaging preserves the same bound.
\end{proof}

\begin{remark}[Where every error term comes from]
\label{rem:global-error-provenance}
The five summands of \eqref{eq:master-total-soundness} correspond to five logically distinct failures.
\begin{enumerate}[label=(\arabic*),leftmargin=*]
  \item $\varepsilon_{\mathrm{red}}$: an outer algebraic reduction challenge hits a root/collision set.
  \item $2(N-1)/|\F|$: the shared Hadamard fingerprint fails to detect at least one false pointwise product relation.
  \item $(s-1)M/|\F|$: the master batching challenge lies in the correlated-agreement exceptional set.
  \item $\varepsilon_{\mathrm{fold}}$: at least one FRI folding challenge is bad.
  \item $(1-\delta)^\kappa$: every final query lands in the accepting witness set despite its density being $<1-\delta$.
\end{enumerate}
No other algebraic soundness loss is hidden in the theorem.
\end{remark}

\subsection{A single decoding extractor}
\label{sec:global-decoding-extractor}

The application sections state ``knowledge by decoding'' separately.  The following theorem gives the common reason.

\begin{definition}[Application decoder]
\label{def:application-decoder}
For a committed relation $\mathcal R^\delta$ with original commitment tuple $\mathbf w=(w_1,\ldots,w_t)$, define
\[
  \operatorname{Ext}_{\mathcal R}(\mathbf w)
  =
  \begin{cases}
    \operatorname{Dec}_\delta(\mathbf w),&\text{if every component decodes and the decoded tuple is a witness for }\mathcal R^\delta,\\
    \bot,&\text{otherwise.}
  \end{cases}
\]
The relation-specific deterministic reconstruction of public or auxiliary data (for example $Mz$ in a sparse matrix relation, or the sorted trace in memory consistency) is applied after decoding the original tables.
\end{definition}

\begin{theorem}[Global decoding knowledge]
\label{thm:global-decoding-knowledge}
Under the hypotheses of \cref{thm:master-interactive-soundness}, let $\varepsilon_{\mathrm{tot}}$ denote the right-hand side of \eqref{eq:master-total-soundness}.  If a prover is accepted with probability strictly larger than $\varepsilon_{\mathrm{tot}}$, then
\[
  \operatorname{Ext}_{\mathcal R}(\mathbf w)\ne\bot.
\]
In particular, the original commitments have unique decoded tables forming a witness for $\mathcal R^\delta$, and this witness is recoverable in polynomial time.
\end{theorem}

\begin{proof}
If $\operatorname{Ext}_{\mathcal R}(\mathbf w)=\bot$, then by definition the original commitment tuple has no witness in $\mathcal R^\delta$.  The master soundness theorem would then bound the acceptance probability by $\varepsilon_{\mathrm{tot}}$, contradicting the hypothesis.  Polynomial-time extraction follows from \cref{lem:partial-decoder-well-defined} and the deterministic relation-specific reconstruction.
\end{proof}

\begin{remark}[What is and is not extracted]
\label{rem:global-extractor-scope}
The extractor is deliberately defined on the \emph{original} commitments.  Auxiliary inverse tables, DEEP quotient words, sparse-entry selections, sorted memory copies, and opening polynomials are proof messages rather than semantic witnesses.  Soundness shows that, outside the charged bad events, a successful low-degree explanation of those messages forces the uniquely decoded original tables to satisfy the intended application relation.  The extractor therefore need not reproduce the prover's auxiliary messages.
\end{remark}

\subsection{Challenge-by-challenge failure profile}
\label{sec:global-round-profile}

For later Fiat--Shamir compilation we need more than a final scalar soundness bound: we need to know at which public challenge a false execution can escape the no-witness state.  The previous proofs already provide such a decomposition.

Fix a deterministic prover and a false original instance.  Call a transcript prefix \emph{doomed} when, after excluding all bad events already charged at earlier rounds, the current prefix cannot be extended to a joint witness for the intended decoded relation.  We use the following challenge classes.

\begin{enumerate}[label=(\roman*),leftmargin=*]
  \item \textbf{Outer algebraic challenges.}  At the $i$-th reduction challenge, at most $b_i$ values can move a doomed prefix into a prefix for which witness lifting is no longer guaranteed.
  \item \textbf{Shared Hadamard fingerprint.}  If a decoded Hadamard relation is false, at most $2(N-1)$ pairs' worth of field probability are charged in total, as proved in \cref{lem:shared-hadamard-fingerprint}; all Hadamard claims share this one event.
  \item \textbf{Master batching challenge.}  Once the heterogeneous batch has no joint witness, at most $(s-1)M$ values of $\beta$ can yield the large common agreement required by \cref{lem:master-correlated-agreement} and \cref{prop:heterogeneous-deterministic-recovery}.
  \item \textbf{Folding challenges.}  At folding level $j$, at most $M_j$ values of $r_j$ are bad by \cref{lem:few-bad-folding-challenges}.
  \item \textbf{Final queries.}  If all preceding challenges are good and the initial master word has no valid large agreement, the final accepting witness set has density $<1-\delta$ by \cref{prop:codeword-chain}.  Therefore $\kappa$ independent queries all accept with probability at most $(1-\delta)^\kappa$.
\end{enumerate}

This yields an explicit sequential profile rather than only a union bound.

\begin{proposition}[Sequential knowledge-error profile]
\label{prop:sequential-knowledge-profile}
Assume the front-end reduction has conditional bad-set bounds $b_i/|\Omega_i|$ as in \cref{lem:global-sequential-accounting}.  For a false original instance, every accepting transcript belongs to the union of the following events:
\begin{equation}
\begin{aligned}
  &E_{\mathrm{red},i}
  &&(1\le i\le r),\\
  &E_{\mathrm{Had}},\\
  &E_{\beta},\\
  &E_{\mathrm{fold},j}
  &&(1\le j\le\ell),\\
  &E_{\mathrm{query}},
\end{aligned}
\label{eq:global-failure-events}
\end{equation}
with conditional probability bounds
\begin{equation}
\begin{aligned}
  \Pr[E_{\mathrm{red},i}\mid\text{previous prefix}]&\le \frac{b_i}{|\Omega_i|},\\
  \Pr[E_{\mathrm{Had}}\mid\text{previous prefix}]&\le h\frac{2(N-1)}{|\F|},\\
  \Pr[E_{\beta}\mid\text{previous prefix}]&\le \frac{(s-1)M}{|\F|},\\
  \Pr[E_{\mathrm{fold},j}\mid\text{previous prefix}]&\le \frac{M_j}{|\F|},\\
  \Pr[E_{\mathrm{query}}\mid\text{all previous challenges good}]&\le(1-\delta)^\kappa.
\end{aligned}
\label{eq:global-failure-profile}
\end{equation}
Consequently the sum of the first four families of conditional errors plus the final query error is exactly bounded by \eqref{eq:master-total-soundness}.
\end{proposition}

\begin{proof}
The outer bounds are the hypothesis and \cref{lem:global-sequential-accounting}.  The shared Hadamard bound is \cref{lem:shared-hadamard-fingerprint}.  Outside the outer and Hadamard bad events, witness lifting says that the heterogeneous batch has no joint witness.  The proof of \cref{thm:heterogeneous-soundness} then shows that more than $(s-1)M$ master challenges with large common agreement would, by correlated agreement and deterministic recovery, produce a joint witness, a contradiction.  This gives the $E_\beta$ bound.

Fix a master challenge outside $E_\beta$.  The initial master word has no fibre-closed agreement set of relative size $1-\delta$ with a codeword.  At each folding round, \cref{lem:few-bad-folding-challenges} excludes at most $M_j$ challenges.  If all folding challenges are good, the codeword-chain argument \cref{prop:codeword-chain} shows that the final witness-set density is strictly less than $1-\delta$; otherwise the chain would reconstruct a large initial agreement.  Finally \cref{lem:acceptance-witness-set,lem:independent-uniform-queries} gives the $(1-\delta)^\kappa$ query term.  The sum
\[
  \sum_{j=1}^{\ell}\frac{M_j}{|\F|}
  =\varepsilon_{\mathrm{fold}}
\]
is \eqref{eq:epsilon-fold}, so summing the profile yields \eqref{eq:master-total-soundness}.
\end{proof}

\begin{corollary}[Round-indexed decoding state]
\label{cor:round-indexed-decoding-state}
For every application proved through the front-end/heterogeneous-batch framework, the unique decoding of the original commitments is a valid knowledge state for the interactive protocol: outside the sequential bad events of \cref{prop:sequential-knowledge-profile}, no public challenge can turn a false decoded state into an accepting transcript.  If an accepting transcript avoids all those events, the decoded original commitments form a witness for the application relation.
\end{corollary}

\begin{proof}
The first statement is precisely the doomed-prefix implication in the proof of \cref{prop:sequential-knowledge-profile}.  If an accepting transcript avoids all listed events, the final-query event cannot occur either, so the assumption that the decoded state is false is impossible.  Hence the original commitments decode to a valid witness.
\end{proof}

The corollary is the form needed later: the semantic state is fixed by the original Reed--Solomon commitments, while each public challenge has an explicit bounded exceptional set.  No rewinding is used in the interactive extractor.

\subsection{Instantiation for the protocols of this paper}
\label{sec:global-instantiations}

The application theorems are instances of \cref{thm:master-interactive-soundness}.  It is useful to record where their front-end error comes from.

\begin{center}
\begin{tabular}{>{\raggedright\arraybackslash}p{0.23\linewidth} >{\raggedright\arraybackslash}p{0.49\linewidth} >{\raggedright\arraybackslash}p{0.18\linewidth}}
\toprule
Relation & Front-end bad events before the heterogeneous backend & Backend contains Hadamard \\
\midrule
Public linear functional & none & no \\
Committed inner product & none & no \\
Hadamard & none beyond the shared Hadamard fingerprint & yes \\
Lookup & logarithmic-derivative shift root & yes \\
Multiset equality & logarithmic-derivative shift root & yes \\
Tagged permutation & tag-compression collision plus multiset shift root & yes \\
Sparse matrix--vector & geometric row fingerprint plus two indexed-selection tag/shift events & yes \\
Sparse R1CS & shared matrix fingerprint plus indexed-selection events for one false matrix equation & yes \\
Memory consistency & tuple compression, initial-memory tag, five logarithmic-derivative shifts, and the linear adjacency fingerprint & yes \\
\bottomrule
\end{tabular}
\end{center}

The exact application-specific bounds were proved in \cref{thm:lookup-soundness,thm:multiset-soundness,thm:tagged-permutation-soundness,thm:sparse-matvec-soundness,thm:r1cs-soundness,thm:memory-soundness}.  The point of the master theorem is not to replace those sharper calculations, but to show that all of them have one common soundness architecture.  In particular, once the relation-specific front end has ruled out a joint witness, \emph{every application pays the folding backend only once}.

\subsection{Consequences for the architecture}
\label{sec:global-soundness-consequences}

The global theorem gives three structural conclusions.

First, coefficient extraction is compositional.  A new relation does not require a new proximity proof: it requires only a reduction whose good outer transcripts lift every joint low-degree witness to a witness of the original relation.  Once that lemma is proved, \cref{thm:master-interactive-soundness} supplies the remaining master-batch, folding, and query terms.

Second, the semantic extractor is uniform.  Evaluation, inner product, lookup, R1CS, and memory consistency all use the same operation on the original commitments: unique Reed--Solomon decoding followed by interpretation of the coefficient vector as a table.  The auxiliary proof words certify that this fixed decoded state satisfies increasingly expressive relations; they do not define a second semantic state.

Third, the soundness proof is already organized by public-coin rounds.  Outer algebraic challenges, the shared Hadamard fingerprint, the master challenge, the FRI folding challenges, and the final queries each have their own bounded failure event.  This is the structure required to discuss Fiat--Shamir and post-quantum random-oracle compilation precisely.  We turn to that question next.

\section{Non-Interactive Compilation and Post-Quantum Knowledge Soundness}
\label{sec:post-quantum}

The protocols developed so far are public-coin interactive oracle protocols.  Their algebraic soundness is information theoretic: no discrete-logarithm, pairing, or other number-theoretic assumption is used.  To obtain a succinct non-interactive argument, however, one must replace oracle messages by commitments and public challenges by hash outputs.  This step is cryptographic, and in the presence of a quantum adversary it must be analysed in the quantum random-oracle model (QROM), where the adversary may query the random oracle in superposition \cite{BonehEtAl11}.

This section gives that compilation precisely.  We first encode the protocols of \cref{sec:soundness} as interactive oracle reductions (IORs) with fibre symbols.  We then prove a relaxed round-by-round knowledge statement for partially opened oracle strings.  Finally we invoke, as an external compiler theorem, the post-quantum BCS theorem for IORs of Chiesa, Di, Hu and Zheng \cite{CDHZ25}.  All statements specific to the present protocols are proved here; the only external cryptographic result used in this section is the compiler theorem itself.  This is analogous to the use of correlated agreement in \cref{sec:preliminaries}: the external theorem is quoted with its hypotheses, and the work of this section is to verify those hypotheses for the coefficient-extraction framework.

\subsection{Fibre strings and the committed application relation}
\label{sec:pq-fibre-strings}

Recall that $L=L_0$ is the smooth multiplicative domain of order $M=2^{n+R}$ and that $L_j=L^{2^j}$ has order $M_j=M/2^j$.  For each $j$ with $M_j\ge 2$ and each $\eta\in L_{j+1}$, fix once and for all one square root $\xi_\eta\in L_j$; the other is $-\xi_\eta$.

\begin{definition}[Fibre string]
\label{def:pq-fibre-string}
For a word $u\in\F^{L_j}$ define
\[
  \operatorname{fib}_j(u)
  \in (\F^2)^{L_{j+1}},
  \qquad
  \operatorname{fib}_j(u)(\eta)
  =\bigl(u(\xi_\eta),u(-\xi_\eta)\bigr).
\]
For a tuple $\mathbf u=(u_1,\ldots,u_t)$ of words on the same domain, define the packed fibre string
\[
  \operatorname{Fib}_j(\mathbf u)(\eta)
  =\bigl(\operatorname{fib}_j(u_1)(\eta),\ldots,
         \operatorname{fib}_j(u_t)(\eta)\bigr).
\]
\end{definition}

The map $\operatorname{fib}_j$ is a bijection between words on $L_j$ and strings of $M_{j+1}$ fibre symbols.  Moreover, the ordinary Hamming distance between fibre strings is exactly the fibre distance of the underlying words:
\begin{equation}
  \Delta\bigl(\operatorname{fib}_j(u),\operatorname{fib}_j(v)\bigr)
  =\Delta_{\mathrm{fib},j}(u,v).
  \label{eq:pq-fibre-distance}
\end{equation}
Indeed, one fibre symbol differs precisely when the two words differ at at least one of the two points of that fibre.

Packing several original table commitments into one fibre symbol is useful for the non-interactive scheme.  It gives a single semantic Merkle root and, more importantly, makes the relaxed relation below express \emph{common} agreement of all original commitments on the same set of fibres.  This is exactly the agreement structure produced by the correlated-agreement arguments of \cref{sec:heterogeneous-batching,sec:soundness}.

Let $\mathcal R$ be any one of the application relations proved earlier: a linear-functional relation, committed inner product, Hadamard relation, lookup relation, multiset/permutation relation, sparse matrix--vector relation, sparse R1CS relation, or memory-consistency relation.  Let the semantic witness be a tuple of tables
\[
  \mathbf f=(f_1,\ldots,f_t),
  \qquad f_i:\bits\to\F,
\]
and let
\[
  \mathbf c(\mathbf f)
  =\bigl(\Enc_T(f_1),\ldots,\Enc_T(f_t)\bigr).
\]

\begin{definition}[Packed committed relation]
\label{def:pq-packed-relation}
The relation $R_{\mathcal R}$ has explicit instance $x$, implicit instance
\[
  y\in\Sigma^{M/2},
  \qquad
  \Sigma=(\F^2)^t,
\]
and witness $\mathbf f$.  It contains $(x,y,\mathbf f)$ exactly when
\begin{equation}
  y=\operatorname{Fib}_0\bigl(\mathbf c(\mathbf f)\bigr)
  \qquad\text{and}\qquad
  (x,\mathbf f)\in\mathcal R.
  \label{eq:pq-exact-relation}
\end{equation}
For $\delta\in[0,1]$, the relaxed relation $(R_{\mathcal R})_{\le\delta}$ consists of triples $(x,y,\mathbf f)$ with a partial string
\[
  y\in(\Sigma\sqcup\{\bot\})^{M/2}
\]
such that
\begin{equation}
  \Delta\left(
      y,
      \operatorname{Fib}_0\bigl(\mathbf c(\mathbf f)\bigr)
  \right)
  \le\delta
  \qquad\text{and}\qquad
  (x,\mathbf f)\in\mathcal R,
  \label{eq:pq-relaxed-relation}
\end{equation}
where $\bot$ differs from every symbol of $\Sigma$.
\end{definition}

Thus a relaxed witness gives one set of at least $(1-\delta)M/2$ fibres on which \emph{all} original table commitments agree simultaneously with their extracted Reed--Solomon codewords.  In particular, every individual commitment is within fibre distance $\delta$ of its decoded table, so \cref{eq:pq-relaxed-relation} implies every application relation used in the interactive sections.

\begin{remark}[Separate reusable roots]
\label{rem:pq-separate-roots}
The packed commitment is the canonical form used in this section because it minimizes authentication paths and gives the cleanest relaxed relation.  If an application requires separately reusable roots for the individual tables, the same construction can be stated with several implicit input strings; the compiler theorem of \cite{CDHZ25} allows multiple implicit instances.  No algebraic argument changes.  The packed form is sufficient for the theorem proved here and is the form we recommend for benchmarks of a complete proof.
\end{remark}

\subsection{Bit-string challenges}
\label{sec:pq-bit-challenges}

The interactive protocols were written with field challenges.  The Fiat--Shamir compiler derives bit strings from a random oracle.  We therefore fix an explicit conversion.

Let $\lambda_{\mathrm{ch}}\ge \lceil\log_2|\F|\rceil$.  Fix a bijection
\[
  \operatorname{num}:\{0,\ldots,|\F|-1\}\longrightarrow\F
\]
and define
\[
  \phi:\{0,1\}^{\lambda_{\mathrm{ch}}}\longrightarrow\F,
  \qquad
  \phi(s)=\operatorname{num}\bigl(\operatorname{int}(s)\bmod|\F|\bigr).
\]
Also fix a bijection
\[
  \operatorname{pos}:\{0,1\}^{n+R}\longrightarrow L.
\]

\begin{lemma}[Sampling from hash bits]
\label{lem:pq-sampling}
For every set $S\subseteq\F$,
\begin{equation}
  \Pr_{s\leftarrow\{0,1\}^{\lambda_{\mathrm{ch}}}}
  [\phi(s)\in S]
  \le
  |S|\left(\frac1{|\F|}+2^{-\lambda_{\mathrm{ch}}}\right).
  \label{eq:pq-field-sampling}
\end{equation}
If $q\leftarrow\{0,1\}^{\kappa(n+R)}$ is divided into $\kappa$ consecutive blocks of $n+R$ bits and each block is mapped through $\operatorname{pos}$, the resulting $\kappa$ points are independent and uniform in $L$.
\end{lemma}

\begin{proof}
Every $a\in\F$ has at most
\[
  \left\lceil\frac{2^{\lambda_{\mathrm{ch}}}}{|\F|}\right\rceil
  <\frac{2^{\lambda_{\mathrm{ch}}}}{|\F|}+1
\]
preimages under $\phi$.  Divide by $2^{\lambda_{\mathrm{ch}}}$ and sum over $a\in S$ to obtain \eqref{eq:pq-field-sampling}.  The second statement follows because applying the bijection $\operatorname{pos}$ independently to each block gives a bijection
\[
  \{0,1\}^{\kappa(n+R)}\simeq L^\kappa.
\]
\end{proof}

Consequently every root-counting error $b/|\F|$ in \cref{prop:sequential-knowledge-profile} becomes at most
\begin{equation}
  b\left(\frac1{|\F|}+2^{-\lambda_{\mathrm{ch}}}\right)
  \label{eq:pq-lifted-field-error}
\end{equation}
when its field challenge is generated from $\lambda_{\mathrm{ch}}$ hash bits.  The final query term $(1-\delta)^\kappa$ is unchanged.

\subsection{The protocol as an interactive oracle reduction}
\label{sec:pq-ior}

Fix one application protocol after all reductions of the preceding sections have been expanded.  Put every public challenge in chronological order: relation-specific outer challenges, the shared Hadamard challenges when present, the master heterogeneous batching challenge, the FRI folding challenges, and finally the query string.  A challenge that occurs before any substantive prover message is preceded by a fixed one-symbol dummy message.  This does not alter the interactive protocol but puts it in the standard ``message, then public challenge'' form.

Every oracle message is written as a fibre string.  Words on $L_j$ are encoded by $\operatorname{fib}_j$.  Several words sent in the same round are packed componentwise into one symbol.  A fully read scalar message, such as a claimed out-of-domain evaluation or the coefficients of the final polynomial, is represented by a short string that the verifier reads in full.  By padding fixed-length tuples, all messages can be embedded injectively in one finite alphabet; the choice of this encoding affects only byte counts, not the protocol or its soundness.

\begin{definition}[Application IOR]
\label{def:pq-application-ior}
Let $\Pi_{\mathcal R}^{\Sigma}$ denote the IOR whose explicit and implicit input is $(x;y)$ for the relation $R_{\mathcal R}$ of \cref{def:pq-packed-relation}.  Its messages are the packed fibre strings and fully read scalar strings of the corresponding interactive protocol.  Every ordinary field challenge is obtained as $\phi(vr_i)$ from a fresh
\[
  vr_i\in\{0,1\}^{\lambda_{\mathrm{ch}}},
\]
and the final query challenge is a string of length $\kappa(n+R)$ mapped to $L^\kappa$ through $\operatorname{pos}$.  On full strings, the verifier performs exactly the checks of the interactive protocol.  On partial strings it rejects whenever it attempts to read an undefined symbol $\bot$.
\end{definition}

The output relation is the trivial accepting relation
\[
  R_{\mathrm{acc}}=\{(\top,(),w):w\text{ arbitrary}\}.
\]
The IOR outputs $(\top,())$ exactly when the interactive verifier accepts.  Hence $\Pi_{\mathcal R}^{\Sigma}$ is an IOR from $R_{\mathcal R}$ to $R_{\mathrm{acc}}$ in the terminology of \cite{CDHZ25}.

\subsection{Partial strings and erasure-stable witness sets}
\label{sec:pq-erasures}

We now verify the hypothesis needed for quantum Fiat--Shamir extraction.  The compiler does not, in general, recover complete oracle messages from a quantum adversary.  It recovers only those positions that the proof opens.  Therefore the knowledge state must remain meaningful for partial strings.

For a partial fibre string $s$, write $\overline{s}$ for the full string obtained by replacing every $\bot$ by the zero symbol.  For a collection of partial oracle strings in a transcript, define $Z_j$ to be the set of points at level $j$ whose entire dependency cone contains no undefined symbol.  Concretely, a point belongs to $Z_j$ if every fibre symbol required to reconstruct and check that point from the original and auxiliary oracle messages is defined.  Because folding dependencies are unions of fibres, each $Z_j$ is fibre closed in the corresponding domain.

For a candidate codeword $c\in C_j$, let $W_j(c)$ be its full-word witness set from \cref{sec:folding-test}.  Define the erasure-restricted witness set
\begin{equation}
  W_j^{\bot}(c)=W_j(c)\cap Z_j,
  \qquad
  \operatorname{dens}_j^{\bot}(c)
  =\frac{|W_j^{\bot}(c)|}{|L_j|}.
  \label{eq:pq-erasure-witness-set}
\end{equation}
Undefined fibres therefore count against the density, exactly as $\bot$ counts as a disagreement in \cref{eq:pq-relaxed-relation}.

\begin{lemma}[Erasure stability]
\label{lem:pq-erasure-stability}
The challenge-by-challenge arguments of \cref{prop:sequential-knowledge-profile} remain valid for partial strings when every witness set is replaced by \eqref{eq:pq-erasure-witness-set}.  More precisely:
\begin{enumerate}[label=(\arabic*),leftmargin=*]
  \item Every outer reduction or shared-Hadamard root-counting argument has the same exceptional set of field challenges as in the full-word proof.
  \item Outside the master batching exceptional set, no candidate codeword has erasure-restricted initial density at least $1-\delta$ unless the partial semantic input together with the decoded tables lies in $(R_{\mathcal R})_{\le\delta}$.
  \item If a folding challenge is good in the sense of \cref{lem:few-bad-folding-challenges}, then
  \[
     W_j^{\bot}(c')\subseteq W_{j-1}^{\bot}(c)
  \]
  for the predecessor codeword $c$ reconstructed from a level-$j$ candidate $c'$ by the codeword-chain argument.
  \item If the verifier accepts a partial transcript after all good challenges, every final query point belongs to $W_\ell^{\bot}(c_\ell)$ for the final codeword $c_\ell$.
\end{enumerate}
\end{lemma}

\begin{proof}
Fill every undefined entry with zero and run the full-word algebraic proof on the filled words.  Root-counting and correlated-agreement arguments depend only on polynomial degree and on positions at which the relevant words are simultaneously defined and agree; replacing undefined positions cannot create a new defined agreement.  This proves (1).

For (2), correlated agreement applied to a master challenge outside its exceptional set produces one common fibre-closed agreement set $T$ for all protected component words.  Intersect $T$ with the set of fibres on which the semantic input is defined.  If the resulting set has size at least $(1-\delta)M/2$, the decoded original tables satisfy the application relation by the same deterministic witness-lifting argument as in \cref{thm:master-interactive-soundness}, while the packed partial input agrees with their encodings on those fibres.  Hence \eqref{eq:pq-relaxed-relation} holds.  Contrapositively, outside the relaxed relation the erasure-restricted initial density is $<1-\delta$.

For (3), the full-word codeword-chain proof gives $W_j(c')\subseteq W_{j-1}(c)$.  Every symbol needed to check a level-$j$ point is computed from symbols in its predecessor dependency cone, so $Z_j\subseteq Z_{j-1}$ under the natural projection of the fold.  Intersect the full-word inclusion with $Z_j$.

For (4), the partial verifier rejects any query that reads $\bot$.  Thus every accepted query lies in $Z_\ell$.  On its defined values it performs the same checks as the full-word verifier, so the query also lies in $W_\ell(c_\ell)$.
\end{proof}

\subsection{Relaxed round-by-round knowledge soundness}
\label{sec:pq-relaxed-rbr}

We use the relaxed round-by-round knowledge notion of Chiesa, Di, Hu and Zheng \cite{CDHZ25}.  Informally, a transcript prefix is assigned a binary knowledge state.  A state equal to one means that the partial implicit input already has an extracted witness for the relaxed relation, or that the transcript has escaped the doomed region of the folding proof.  A public challenge may change a zero state to a one state only with the bounded probability charged to that round.

Let the field-challenge rounds of an application have full-word exceptional-set bounds
\[
  b_1,b_2,\ldots,b_m,
\]
meaning that after every fixed prior transcript at most $b_i$ field values can cause the $i$-th bad event.  This list includes all outer algebraic reductions, the shared Hadamard fingerprint (represented sequentially by its challenge coordinates), the master batching challenge, and the $\ell$ folding challenges.  Let the last challenge be the $\kappa$ query points.

Define
\begin{equation}
  \epsilon_i^{\mathrm{rr}}
  =b_i\left(\frac1{|\F|}+2^{-\lambda_{\mathrm{ch}}}\right)
  \quad(1\le i\le m),
  \qquad
  \epsilon_{m+1}^{\mathrm{rr}}=(1-\delta)^\kappa.
  \label{eq:pq-rbr-errors}
\end{equation}

\begin{theorem}[Relaxed round-by-round decoding knowledge]
\label{thm:pq-relaxed-rbr}
Fix $0<\delta\le(1-\rho)/2$ and assume $2N<(1-\delta)M$.  For every application relation $\mathcal R$ covered by \cref{thm:master-interactive-soundness}, the IOR $\Pi_{\mathcal R}^{\Sigma}$ has relaxed round-by-round knowledge soundness for $R_{\mathcal R}$ with per-round errors bounded by \eqref{eq:pq-rbr-errors}.  The extractor associated with the knowledge state is the unique Reed--Solomon decoder of \cref{def:application-decoder}, applied to the packed semantic input after filling undefined positions arbitrarily; whenever the state certifies knowledge, the resulting tables form a witness for $(R_{\mathcal R})_{\le\delta}$.
\end{theorem}

\begin{proof}
Use the doomed-prefix state of \cref{cor:round-indexed-decoding-state}, replacing every full witness set by its erasure-restricted version from \cref{eq:pq-erasure-witness-set}.  By \cref{lem:pq-erasure-stability}, a false relaxed semantic state cannot leave the doomed region except through one of the same bounded exceptional sets as in the full-word argument.  For a field challenge, \cref{lem:pq-sampling} converts a set of at most $b_i$ bad field values into the first bound of \eqref{eq:pq-rbr-errors}.

If all field and folding challenges remain good, the final erasure-restricted witness-set density is strictly below $1-\delta$.  By \cref{lem:pq-erasure-stability}(4), acceptance requires every one of the $\kappa$ independent uniform query points to lie in that set.  The probability is therefore at most $(1-\delta)^\kappa$ by \cref{lem:independent-uniform-queries}.

If the state leaves the doomed region outside these events, the common agreement set obtained by the master batch has relative size at least $1-\delta$ and consists only of defined fibres.  Unique decoding gives the semantic tables; the relation-specific witness-lifting theorem gives $(x,\mathbf f)\in\mathcal R$; and the common agreement gives \eqref{eq:pq-relaxed-relation}.  This is exactly a witness for $(R_{\mathcal R})_{\le\delta}$.
\end{proof}

Define the single round-by-round parameter
\begin{equation}
  \epsilon_{\mathrm{rr}}(\delta)
  =\max\left\{
      \max_{1\le i\le m}
      b_i\left(\frac1{|\F|}+2^{-\lambda_{\mathrm{ch}}}\right),
      (1-\delta)^\kappa
    \right\}.
  \label{eq:pq-epsilon-rr}
\end{equation}
Unlike the overall interactive error of \cref{eq:master-total-soundness}, the quantity \eqref{eq:pq-epsilon-rr} is a \emph{maximum} over rounds.  This is the parameter consumed by the post-quantum Fiat--Shamir compiler.

\subsection{Salted Merkle commitments and the BCS compiler}
\label{sec:pq-bcs}

We now recall the exact external theorem used for compilation.  A salted Merkle commitment to a string $a=(a_1,\ldots,a_T)$ samples an independent salt for every leaf, hashes the pair $(a_i,\mathrm{salt}_i)$ at each leaf, and hashes pairs of child nodes up the tree.  The root is the commitment.  An opening contains the value, its salt, and the authentication path.  The salts are essential in the post-quantum multi-extraction theorem of \cite{CDHZ25}.

The BCS transformation \cite{BenSassonChiesaSpooner16} replaces every oracle message of an IOR by such a Merkle root.  Fiat--Shamir challenges are obtained by hashing the explicit instance, the semantic input root, all message roots produced so far, and fresh round salts.  The non-interactive proof contains the roots, the salts required by the compiler, all fully read messages, and authenticated openings of every oracle position queried by the IOR verifier.

For completeness, we state the specialization of the compiler theorem that we use.  The theorem itself is quoted, not reproved.

\begin{theorem}[Post-quantum BCS compiler for IORs, specialised \cite{CDHZ25}]
\label{thm:pq-cdhz}
Let $\Pi$ be a $\nu$-round IOR from a relation $R$ with implicit input to $R_{\mathrm{acc}}$.  Suppose $\Pi$ has relaxed round-by-round knowledge error at most $\epsilon_{\mathrm{rr}}(\delta)$, all public challenge strings have length at least $\lambda_{\min}$, and every implicit or oracle string has length at most $l_{\max}$.  Let $q_V$ be the number of random-oracle queries made by the compiled verifier, $q_s$ the number of string positions read by the IOR verifier, and let $\lambda_H$ be the Merkle-hash output length.  For a quantum adversary making at most $Q$ random-oracle queries, the salted-BCS compilation admits a polynomial-time straightline quantum extractor.  Its simulation error is zero, its commitment-commutativity error satisfies
\[
  \epsilon_{\mathrm{com}}(q)
  \le \frac{240(q+1)}{2^{\lambda_H}},
\]
and its extraction error is bounded by
\begin{align}
 \epsilon_{\mathrm{ext}}(Q)
 \le{}&320(Q+\nu+1)(Q+\nu)\epsilon_{\mathrm{rr}}(\delta)
      +\frac{4(Q+\nu+1)\nu}{2^{\lambda_{\min}}}
 \notag\\
 &+\frac{32l_{\max}
      +1280Q(2Q+1)^2
      +128q_s\lceil\log_2 l_{\max}\rceil}{2^{\lambda_H}}
 \notag\\
 &+\frac{960(Q+1)T_1^2
      +480q_V^2(Q+\nu+q_V+2)}{2^{\lambda_H}},
 \label{eq:pq-cdhz-bound}
\end{align}
where
\[
  T_1=Q+1+q_V+\nu+q_{\mathrm{CR}}.
\]
Here $q_{\mathrm{CR}}$ is the random-oracle query cost of deciding the relaxed committed relation.  Except with probability \eqref{eq:pq-cdhz-bound}, whenever the compiled verifier accepts, the extractor outputs a partial committed witness belonging to the relaxed committed relation $\mathrm{CR}^c[R,\delta]$.
\end{theorem}

The constants in \cref{thm:pq-cdhz} are those of the 2 March 2026 version of \cite{CDHZ25}.  We use the theorem only through the displayed bound and its extraction conclusion.

\subsection{Post-quantum knowledge soundness of coefficient-extraction arguments}
\label{sec:pq-main-theorem}

We can now combine the paper-specific relaxed round-by-round theorem with the external compiler.

\begin{theorem}[Post-quantum straightline knowledge soundness]
\label{thm:pq-main}
Let $\mathcal R$ be any application relation covered by \cref{thm:master-interactive-soundness}, and let $\Pi_{\mathcal R}^{\Sigma}$ be the IOR of \cref{def:pq-application-ior}.  Fix
\[
  0<\delta\le\frac{1-\rho}{2},
  \qquad
  2N<(1-\delta)M.
\]
Let $\nu$, $\lambda_{\min}$, $l_{\max}$, $q_s$, $q_V$ and $q_{\mathrm{CR}}$ denote the parameters of its salted-BCS compilation, and let $\epsilon_{\mathrm{rr}}(\delta)$ be \eqref{eq:pq-epsilon-rr}.  Then, in the QROM, the compiled non-interactive coefficient-extraction argument has post-quantum straightline knowledge soundness with extraction error bounded by \eqref{eq:pq-cdhz-bound}.

In particular, for every quantum adversary making at most $Q$ random-oracle queries and outputting a semantic Merkle root, an explicit application instance, and an accepting proof, except with probability $\epsilon_{\mathrm{ext}}(Q)$ the extractor outputs
\begin{enumerate}[label=(\arabic*),leftmargin=*]
  \item a partial packed fibre string $y$ whose defined positions have accepting Merkle openings under that root;
  \item a tuple of tables $\mathbf f=(f_1,\ldots,f_t)$;
  \item auxiliary commitment trapdoor data required by the committed relation,
\end{enumerate}
such that
\begin{equation}
  \Delta\left(
      y,
      \operatorname{Fib}_0(\Enc_T(f_1),\ldots,\Enc_T(f_t))
  \right)
  \le\delta
  \label{eq:pq-extracted-distance}
\end{equation}
and $(x,\mathbf f)\in\mathcal R$.
\end{theorem}

\begin{proof}
By \cref{thm:pq-relaxed-rbr}, $\Pi_{\mathcal R}^{\Sigma}$ satisfies the relaxed round-by-round knowledge hypothesis of \cref{thm:pq-cdhz} with parameter \eqref{eq:pq-epsilon-rr}.  Apply \cref{thm:pq-cdhz}.  Its extractor outputs a witness in the relaxed committed relation associated with $R_{\mathcal R}$.  Unfolding \cref{def:pq-packed-relation} gives exactly \eqref{eq:pq-extracted-distance} and $(x,\mathbf f)\in\mathcal R$.
\end{proof}

This theorem is the precise sense in which the framework is post-quantum.  Sumcheck itself was never the source of a quantum assumption: it is an information-theoretic protocol.  What the coefficient-extraction framework changes is the algebraic reduction.  The relations of this paper are reduced directly to Reed--Solomon proximity and then compiled using only hash-based commitments and the QROM Fiat--Shamir theorem above.  No pairing, elliptic-curve discrete logarithm, or structured reference string appears in the construction.

\subsection{Binding of the extracted semantic state}
\label{sec:pq-binding}

The relaxed extractor works one accepting proof at a time.  We therefore record the deterministic binding consequence needed to interpret the extracted tables as the state committed by the semantic root.

\begin{lemma}[Uniqueness under one semantic root]
\label{lem:pq-root-binding}
Assume $\delta\le(1-\rho)/2$.  Let the same semantic Merkle root admit two relaxed extracted witnesses
\[
  (y,\mathbf f,\mathrm{td}),
  \qquad
  (y',\mathbf f',\mathrm{td}')
\]
for possibly different explicit claims, each satisfying \eqref{eq:pq-extracted-distance}.  If $\mathbf f\ne\mathbf f'$, then the two trapdoors contain accepting openings of one fibre position under the same root to different packed symbols.  Hence one can compute a collision in the Merkle random oracle.
\end{lemma}

\begin{proof}
Let $D$ and $D'$ be the sets of fibre positions at which $y$ and $y'$ respectively agree with the packed codewords of $\mathbf f$ and $\mathbf f'$.  Then
\[
  |D|,|D'|\ge(1-\delta)M/2,
\]
so
\[
  |D\cap D'|\ge(1-2\delta)M/2\ge \rho M/2=N/2.
\]
If the two packed symbols agreed at every fibre of $D\cap D'$, then every component pair of Reed--Solomon codewords would agree on the at least $N$ points contained in those fibres.  Degree-$<N$ uniqueness would imply $\Enc_T(f_i)=\Enc_T(f_i')$ for every $i$, hence $f_i=f_i'$ by \cref{prop:table-isomorphism}.  This contradicts $\mathbf f\ne\mathbf f'$.  Therefore there is a fibre $\eta\in D\cap D'$ with $y(\eta)\ne y'(\eta)$.  Both openings are accepted under the same Merkle root, so tracing their authentication paths upward yields a collision at the first node where the two paths merge.
\end{proof}

Thus, except for the hash-collision event already accounted for by the compiler, one semantic root determines at most one decoded application state within the unique-decoding radius.

\subsection{Parameter consequences}
\label{sec:pq-parameters}

The dominant QROM term in \eqref{eq:pq-cdhz-bound} is typically
\begin{equation}
  320(Q+\nu+1)(Q+\nu)\epsilon_{\mathrm{rr}}(\delta),
  \label{eq:pq-dominant-term}
\end{equation}
so the interactive query error must be set substantially below the desired non-interactive security level.  This is a compiler loss, not an additional algebraic failure of coefficient extraction.

For the recommended unique-decoding rate $\rho=1/4$ and $\delta=3/8$,
\[
  (1-\delta)^\kappa=(5/8)^\kappa.
\]
If the query term dominates \eqref{eq:pq-epsilon-rr}, increasing $\kappa$ by one multiplies it by $5/8$.  For a target extraction error $2^{-\lambda}$ against a fixed query budget $Q$, a conservative first choice is therefore to select $\kappa$ so that
\begin{equation}
  320(Q+\nu+1)(Q+\nu)(5/8)^\kappa
  \le 2^{-\lambda-1},
  \label{eq:pq-kappa-choice}
\end{equation}
and then choose $\lambda_{\mathrm{ch}}$ and $\lambda_H$ so that all remaining terms of \eqref{eq:pq-cdhz-bound} sum to at most $2^{-\lambda-1}$.

The field size must simultaneously make every algebraic round error in \eqref{eq:pq-epsilon-rr} negligible.  If
\[
  b_{\max}=\max_i b_i,
\]
it is enough to require
\begin{equation}
  b_{\max}
  \left(\frac1{|\F|}+2^{-\lambda_{\mathrm{ch}}}\right)
  \ll (1-\delta)^\kappa.
  \label{eq:pq-field-choice}
\end{equation}
For large heterogeneous applications $b_{\max}$ is often the master-batching quantity $(s-1)M$; this is why the challenge field, rather than only the base arithmetic field, may need to be a substantial extension.

We postpone concrete parameter tables until the implementation section, where $\nu$, $s$, $q_s$ and $q_V$ can be measured from the actual transcript layout.  The formulas above are the exact constraints that the Rust parameter tool must evaluate.

\subsection{What has and has not been proved}
\label{sec:pq-scope}

It is useful to separate three statements that are sometimes conflated.

\begin{enumerate}[label=(\arabic*),leftmargin=*]
  \item \textbf{Interactive algebraic soundness.}  This is proved internally in \cref{sec:soundness}; it is information theoretic apart from the quoted Reed--Solomon correlated-agreement theorem.
  \item \textbf{Relaxed round-by-round decoding knowledge.}  This is proved internally in \cref{thm:pq-relaxed-rbr}; it shows that the interactive protocol has exactly the state structure required by the non-interactive compiler, including partial oracle strings.
  \item \textbf{Post-quantum Fiat--Shamir extraction.}  This follows by applying the external QROM compiler theorem \cref{thm:pq-cdhz}.  We do not reprove that general quantum theorem.
\end{enumerate}

Accordingly, the post-quantum claim of this paper is not based on the informal slogan that ``hashes are post-quantum.''  It is a theorem obtained by verifying the hypotheses of a specific QROM knowledge-sound compiler.  Conversely, the theorem is only as strong as that compiler model: the construction analysed here is the salted-BCS transcript described above, with domain-separated random-oracle calls and the exact message/challenge ordering used by the IOR.  A Rust implementation claiming this theorem must conform to that transcript; conformance is therefore a testable implementation obligation rather than a presentation detail.

\section{Complexity and Concrete Protocol Costs}
\label{sec:complexity}

The preceding sections were organized around correctness and soundness.  This section records the computational objects that an implementation must actually construct and gives symbolic prover, verifier, communication, and preprocessing costs.  We deliberately separate these accounting identities from measured wall-clock performance.  The latter depends on the field implementation, NTT layout, Merkle hashing, memory traffic, batching strategy, and hardware, and will be reported only after the Rust implementation is complete.

The main structural point is simple.  Heterogeneous batching removes repeated \emph{proximity proofs}: once all relation-specific words have been fixed, there is one master Reed--Solomon word, one folding chain, and one final set of $\kappa$ queries.  Batching does not make the relation-specific auxiliary commitments disappear.  Consequently, a meaningful performance comparison must distinguish
\begin{equation}
  \text{front-end algebra and commitments}
  \qquad\text{from}\qquad
  \text{the shared folding backend}.
  \label{eq:complexity-two-layers}
\end{equation}
This distinction is particularly important for lookup, sparse matrix multiplication, R1CS, and memory consistency, where many algebraic relations share one folding test.

\subsection{Cost model}
\label{sec:complexity-cost-model}

Throughout this section,
\[
  N=2^n,
  \qquad
  M=2^{n+R},
  \qquad
  \rho=N/M,
\]
and the folding test performs $\ell$ binary folds and $\kappa$ independent base-domain queries.  We count arithmetic in the challenge field $\F$ unless stated otherwise.

Rather than build an implementation-dependent FFT constant into the mathematical statement, we use the following kernel costs.

\begin{definition}[Implementation kernels]
\label{def:complexity-kernels}
Let
\begin{align}
  \mathsf{Mul}_N
  &\defeq \text{cost of multiplying two degree-$<N$ polynomials},\\
  \mathsf{Enc}_{N,M}
  &\defeq \text{cost of evaluating a degree-$<N$ polynomial on $L$},\\
  \mathsf{Merkle}_{M}(d)
  &\defeq \text{cost of committing to $d$ field elements per one of $M$ leaves},\\
  \mathsf{Fold}_{M,\ell}
  &\defeq \text{cost of one honest binary folding chain of length $\ell$},\\
  \mathsf{Open}_{M}(q)
  &\defeq \text{cost of producing/verifying authenticated openings of $q$ leaves}.
  \label{eq:complexity-kernel-costs}
\end{align}
The parameter $d$ in the Merkle cost allows several words fixed in the same transcript round to be packed into one vector-valued leaf.
\end{definition}

On a multiplicative domain supporting radix-two FFT/NTT algorithms, the standard bounds are
\begin{equation}
  \mathsf{Mul}_N=O(N\log N),
  \qquad
  \mathsf{Enc}_{N,M}=O(M\log M),
  \label{eq:standard-fft-bounds}
\end{equation}
with the usual possibility of computing a degree-$<N$ evaluation on an $M$-point domain by a padded transform.  The paper does not assume these particular algorithms for correctness.

The folding backend itself is linear in the domain size.

\begin{proposition}[Arithmetic size of one folding chain]
\label{prop:folding-arithmetic-cost}
Let $M_j=M/2^j$.  An honest binary folding chain of length $\ell$ computes
\begin{equation}
  \sum_{j=1}^{\ell}M_j
  =M(1-2^{-\ell})
  <M
  \label{eq:fold-total-output-cells}
\end{equation}
folded field elements.  Hence, once the initial word is available, the arithmetic work of all folds is $O(M)$ field operations, and the total number of field elements in all explicitly materialized folded words is less than $M$.
\end{proposition}

\begin{proof}
At fold $j$, the output domain has $M_j=M/2^j$ points.  Summing the geometric series gives
\[
  \sum_{j=1}^{\ell}\frac{M}{2^j}
  =M\left(1-2^{-\ell}\right).
\]
The local fold formula of \cref{lem:fold-local-line} uses a constant number of field operations per output point, proving the arithmetic bound.
\end{proof}

If only the checkpoint set $\mathcal C\subseteq[1,\ell-1]$ of \cref{sec:folding-test} is committed, the arithmetic folds are still computed but only the checkpoint words require Merkle roots.  Thus higher arity and checkpointing alter hashing and proof size without changing the algebraic reduction developed in this paper.

\subsection{One master word rather than many proximity tests}
\label{sec:complexity-master-word}

Let
\[
  \mathcal U=(u_0,\ldots,u_{s-1})
\]
be the protected component list of \cref{def:protected-component-list}, and let $P_i\in\F[X]_{<N}$ denote the honest polynomial underlying $u_i$.  The master challenge $\beta$ defines
\begin{equation}
  u_{\beta}
  =\sum_{i=0}^{s-1}\beta^iu_i,
  \qquad
  P_{\beta}
  =\sum_{i=0}^{s-1}\beta^iP_i.
  \label{eq:complexity-master-poly}
\end{equation}
The naive implementation could materialize all $s$ words on $L$ and combine them pointwise, at cost $O(sM)$.  This is unnecessary for the honest prover.

\begin{proposition}[Coefficient-side construction of the master word]
\label{prop:master-coefficient-side}
Assume the honest prover has coefficient representations of the protected polynomials $P_0,\ldots,P_{s-1}$.  Then it can construct the master polynomial $P_{\beta}$ in
\begin{equation}
  O(sN)
  \label{eq:master-coeff-cost}
\end{equation}
field operations and obtain the initial folding word by one evaluation of $P_{\beta}$ on $L$, at cost $\mathsf{Enc}_{N,M}$.  In particular, heterogeneous batching need not incur an $O(sM)$ pointwise master-word construction.
\end{proposition}

\begin{proof}
Use Horner accumulation coefficientwise:
\[
  P_{\beta}
  =P_0+\beta(P_1+\beta(P_2+\cdots+\beta P_{s-1})\cdots).
\]
There are $N$ coefficient positions and $s-1$ scalar multiplication/addition steps per position, giving $O(sN)$ field operations.  Evaluation of the resulting degree-$<N$ polynomial on $L$ is exactly one call to the encoding kernel.
\end{proof}

The hypothesis of \cref{prop:master-coefficient-side} is natural for the present protocols.  Reversal is a coefficient permutation; affine virtual words are coefficientwise linear combinations; DEEP quotients are obtained by synthetic division in $O(N)$ operations; and the high-part polynomial $H$ is obtained while splitting a product already computed by the prover.  Thus the implementation should keep coefficient representations of protected words whenever possible and regard full-domain evaluations as commitment objects, not as the primary representation for batching.

\begin{corollary}[Backend sharing]
\label{cor:backend-sharing}
For a heterogeneous batch with $s$ protected words, the final low-degree backend requires, beyond relation-specific commitments,
\begin{equation}
  O(sN)+\mathsf{Enc}_{N,M}+\mathsf{Fold}_{M,\ell}
  \label{eq:shared-backend-cost}
\end{equation}
arithmetic work, one folding transcript, and $\kappa$ final proximity queries.  It does not require $s$ independent Reed--Solomon proximity proofs.
\end{corollary}

\begin{proof}
By \cref{prop:master-coefficient-side}, the $s$ protected degree-$<N$ polynomials can be combined in coefficient form in $O(sN)$ field operations and evaluated on $L$ with one call to the encoding kernel.  The heterogeneous protocol of \cref{def:heterogeneous-batching-protocol} then invokes one folding proximity test on that master word, so the remaining work is one $\mathsf{Fold}_{M,\ell}$ execution and the common set of $\kappa$ final queries.  No component word starts an independent folding transcript.  This gives \eqref{eq:shared-backend-cost}.
\end{proof}

\subsection{Primitive relation costs}
\label{sec:complexity-primitives}

We next count the relation-specific objects.  A \emph{new oracle word} means a length-$M$ word that was not already part of the application instance and must therefore be committed by the prover.  Public and virtual words do not count as new oracles.  The counts below are before packing several words into one vector-valued Merkle leaf.

\begin{proposition}[Primitive word counts]
\label{prop:primitive-word-counts}
Ignoring exact deduplication between distinct claims, the primitive protocols have the costs in \cref{tab:primitive-costs}.
\end{proposition}

\begin{table}[ht]
\centering
\small
\begin{tabular}{@{}p{3.1cm}ccc p{4.2cm}@{}}
\toprule
Relation & $s$ & new oracles & products & extra coefficient work \\
\midrule
Public linear functional & $3$ & $1$ & $1$ & $O(N)$ \\
Committed inner product & $4$ & $1$ & $1$ & $O(N)$ reversal and split \\
Hadamard $a\had b=c$ & $9$ & $2$ & $1$ & $O(N)$ scaling and quotients \\
Lookup & $\le25$ & $8$ & $3$ & $O(N)$ inverse and quotient work \\
Multiset equality & $\le24$ & $8$ & $2$ & $O(N)$ inverse and quotient work \\
\bottomrule
\end{tabular}
\caption{Naive per-relation accounting before deduplication, packing, fusion, and sharing with surrounding claims.  ``Polynomial products'' are products of two degree-$<N$ polynomials used by coefficient extraction.}
\label{tab:primitive-costs}
\end{table}

\begin{proof}
A public linear functional protects $(w,w_A,h)$ and introduces only $w_A$; its universal split requires one product $UK_p$.  An inner product protects $(w_a,w_b^\star,w_A,h)$, introduces only $w_A$, and requires one product $V_aV_b^\star$.  A Hadamard claim protects the nine words listed in \eqref{eq:had-protected-words}; its prover sends $w_Q,w_A$ and computes the single product $Q_\gamma V_b^\star$.

For lookup, \cref{def:lookup-sub-batch} contains two Hadamard claims, one public linear-functional claim, and one committed inner-product claim.  Hence the protected-word count is at most
\[
  2\cdot9+3+4=25.
\]
The two inverse tables $w_q,w_h$ are new source oracles.  The two Hadamard claims contribute four opening/scaled oracles, while the linear-functional and inner-product claims contribute one opening oracle each, for $2+4+1+1=8$ new words.  The same decomposition gives three coefficient products: one for each Hadamard and one for the inner product.

For multiset equality, the sub-batch contains two Hadamard claims and two public linear functionals, giving
\[
  2\cdot9+2\cdot3=24
\]
protected words.  The inverse tables contribute two source oracles, the two Hadamard claims four new oracles, and the two linear claims two opening oracles, again totaling eight.  Only the two Hadamard coefficient products are required.  All remaining operations are coefficient permutations, affine combinations, synthetic divisions, or tablewise inversions and therefore require $O(N)$ field operations apart from the cost of producing their commitments.
\end{proof}

The table intentionally counts calls rather than pretending that all calls have the same constant.  For example, $Q_\gamma(X)=V_a(\gamma X)$ is obtained by $N$ scalar multiplications on coefficients, whereas a degree-$<N$ polynomial product normally uses an NTT or another multiplication algorithm.

\subsection{Verifier locality}
\label{sec:complexity-verifier}

The verifier's algebra outside Merkle authentication is small because virtual words are locally computable.

\begin{proposition}[Local arithmetic of the primitive relations]
\label{prop:primitive-verifier-locality}
For one base-domain query point, excluding the common folding checks and Merkle authentication:
\begin{enumerate}[label=(\roman*),leftmargin=*]
  \item a public linear functional requires one source value and one opening value, plus the cost of evaluating its public kernel;
  \item an inner product requires values of $w_a,w_A$ at the query point and $w_b$ at the inverse point, followed by $O(1)$ field operations;
  \item a Hadamard claim requires values of $w_a,w_c,w_Q,w_A$ at the query point and $w_b$ at the inverse point, followed by $O(1)$ field operations and three quotient divisions;
  \item lookup and multiset relations add no asymptotically new local algebra beyond their constituent Hadamard, linear-functional, and inner-product claims.
\end{enumerate}
All quotient divisions over a batch of query points may be implemented with one batched inversion per denominator family.
\end{proposition}

\begin{proof}
Items (i)--(iii) are \cref{prop:ip-local-verifier-work,prop:had-local-verifier-work} and the public-kernel formula of \cref{sec:linear-functionals}.  Item (iv) follows because lookup and multiset protocols are defined by heterogeneous composition of those primitives.  Standard batch inversion replaces $q$ inversions of nonzero field elements by one inversion and $O(q)$ multiplications.
\end{proof}

For multilinear evaluation, the public kernel
\[
  E_z^\star(\xi)
  =\prod_{i=1}^{n}\bigl(z_i+(1-z_i)\xi^{2^{i-1}}\bigr)
\]
requires $O(n)$ field operations per queried base point if powers are produced by repeated squaring.  The arbitrary inner-product primitive avoids this $O(n)$ kernel computation because its second factor is read from the inverse fibre instead.

\subsection{Sparse matrix--vector and R1CS costs}
\label{sec:complexity-r1cs}

The sparse protocol is designed so that no dense $N\times N$ object is formed.  Let $K=\operatorname{nnz}(M)$ before padding.  Public preprocessing stores $O(N)$ padded sparse-entry metadata, or $O(K)$ meaningful nonzero entries plus padding descriptors in an implementation that generates the dummy tail implicitly.

\begin{proposition}[Primitive decomposition of sparse matrix--vector multiplication]
\label{prop:smv-primitive-count}
The sparse matrix--vector batch of \cref{def:sparse-matvec-sub-batch} contains
\begin{equation}
  5\ \text{Hadamard claims},
  \qquad
  5\ \text{public linear functionals},
  \qquad
  1\ \text{committed inner product}.
  \label{eq:smv-primitive-count}
\end{equation}
Consequently, before exact deduplication,
\begin{equation}
  s_{\mathrm{SMV}}\le5\cdot9+5\cdot3+4=64.
  \label{eq:smv-naive-s}
\end{equation}
\end{proposition}

\begin{proof}
Each indexed-selection sub-batch of \cref{def:indexed-selection-sub-batch} consists of two Hadamard relations and two public linear functionals.  There are two indexed selections, contributing four of each.  The sparse contribution relation $v=\mu\had u$ contributes one further Hadamard; \eqref{eq:sparse-entry-inner-product} contributes one committed inner product; and \eqref{eq:sparse-output-linear-functional} contributes one public linear functional.  The protected-word bound follows from \eqref{eq:naive-protected-count}.
\end{proof}

A naive, unfused sparse-matrix prover introduces seven semantic/outer auxiliary words before the primitive opening words: $w_u,w_s,w_v$ and two inverse words for each of the two indexed selections.  The primitive coefficient-extraction layer then introduces
\[
  2\cdot5+5+1=16
\]
additional opening/scaled words.  Thus a standalone sparse-matrix proof can contain up to $23$ new length-$M$ oracle words before packing and sharing.  This number is an implementation target, not a lower bound: vector-valued leaves, reuse of inverse tables, direct generation of public words, and fusion of several coefficient products can reduce both roots and memory traffic.

For R1CS the primitive decomposition is completely explicit.

\begin{proposition}[Primitive decomposition of sparse R1CS]
\label{prop:r1cs-primitive-count}
The R1CS batch of \cref{def:r1cs-sub-batch} contains
\begin{equation}
  16\ \text{Hadamard claims},
  \qquad
  15\ \text{public linear functionals},
  \qquad
  3\ \text{committed inner products}.
  \label{eq:r1cs-primitive-count}
\end{equation}
Hence its naive protected-word count is
\begin{equation}
  s_{\mathrm{R1CS}}
  \le16\cdot9+15\cdot3+3\cdot4
  =201.
  \label{eq:r1cs-protected-upper}
\end{equation}
An optional public-input fingerprint adds one public linear-functional claim and raises this naive bound to $204$.
\end{proposition}

\begin{proof}
Each sparse matrix--vector relation contributes the counts of \cref{prop:smv-primitive-count}.  Three such relations therefore contribute $15$ Hadamards, $15$ public linear functionals, and $3$ inner products.  The final R1CS multiplication relation $a\had b=c$ contributes one additional Hadamard.  Substitution into \eqref{eq:naive-protected-count} gives \eqref{eq:r1cs-protected-upper}.
\end{proof}

The important point is that \cref{prop:r1cs-primitive-count} does \emph{not} imply 34 independent folding proofs.  All $34$ primitive relations contribute to one list of protected words and hence to one master fold.  Their cost is paid in auxiliary word construction and commitment, whereas the final proximity backend is shared.

\subsection{Memory-checking costs}
\label{sec:complexity-memory}

The memory construction is larger algebraically but has the same backend shape.

\begin{proposition}[Primitive decomposition of memory consistency]
\label{prop:memory-primitive-count}
The memory sub-batch of \cref{def:memory-sub-batch} expands to
\begin{equation}
  17\ \text{Hadamard claims},
  \qquad
  7\ \text{public linear functionals},
  \qquad
  4\ \text{committed inner products}.
  \label{eq:memory-primitive-count}
\end{equation}
Therefore, before exact deduplication,
\begin{equation}
  s_{\mathrm{Mem}}
  \le17\cdot9+7\cdot3+4\cdot4
  =190.
  \label{eq:memory-protected-upper}
\end{equation}
\end{proposition}

\begin{proof}
The scalar multiset sub-batch contributes two Hadamards and two public linear functionals.  Each of the four lookup sub-batches contributes two Hadamards, one public linear functional, and one committed inner product, for a total of eight Hadamards, four linear functionals, and four inner products.  The local memory semantics add the seven explicit Hadamard relations in \cref{def:memory-sub-batch}, and the adjacency residual check adds one public linear functional.  Summing gives \eqref{eq:memory-primitive-count}; \eqref{eq:memory-protected-upper} follows from \eqref{eq:naive-protected-count}.
\end{proof}

As for R1CS, the protected-word count controls the field term in the current unique-decoding batching theorem, but it is not the number of Merkle roots and it is not the number of folding proofs.  Many protected words are virtual, and many committed words can share leaves because they are fixed in the same round and queried at the same positions.

\subsection{Merkle roots, opened values, and proof size}
\label{sec:complexity-proof-size}

A byte-level proof-size formula depends on commitment layout.  We therefore state a layout-independent accounting identity and then the simple packed-leaf upper bound that the Rust implementation will instantiate.

Let $R_{\mathrm{root}}$ be the number of Merkle roots appearing in the non-interactive transcript, let $Q_{\mathrm{leaf}}$ be the total number of distinct authenticated leaves opened across all roots after deduplicating repeated positions, and let $H_{\mathrm{auth}}$ be the number of sibling hashes carried by the chosen Merkle opening or multiproof representation.  Let $B_F$ and $B_H$ denote the serialized byte lengths of one field element and one hash output.

\begin{proposition}[Generic proof-size accounting]
\label{prop:generic-proof-size}
Apart from fixed-size scalar challenges that are recomputed by Fiat--Shamir, a compiled proof has serialized size
\begin{equation}
  B_F\,Q_F
  +B_H\bigl(R_{\mathrm{root}}+H_{\mathrm{auth}}+Q_{\mathrm{salt}}\bigr)
  +O(1),
  \label{eq:proof-size-accounting}
\end{equation}
where $Q_F$ is the number of transmitted field elements in opened leaves, scalar prover messages, and the final explicit polynomial, and $Q_{\mathrm{salt}}$ is the number of opened salted leaves.  In particular, vector-valued leaves change $Q_F$ but can reduce $R_{\mathrm{root}}$ and $H_{\mathrm{auth}}$ substantially.
\end{proposition}

\begin{proof}
The salted-BCS transcript of \cref{sec:pq-bcs} contains Merkle roots, prover scalar messages, the final explicit polynomial, and authenticated openings.  Each opened leaf contributes its serialized field values and one salt; each authentication path or multiproof contributes sibling hashes.  Challenges are recomputed from the transcript and therefore need not be transmitted.  Summing those serialized objects gives \eqref{eq:proof-size-accounting}.
\end{proof}

For a binary folding chain with checkpoint set $\mathcal C$, the final polynomial contributes exactly
\begin{equation}
  N_\ell=N/2^\ell
  \label{eq:final-poly-elements}
\end{equation}
field elements before any padding convention.  The query portion scales linearly with $\kappa$ times the number of committed roots that must be opened at the corresponding fibre positions.  Consequently, reducing the number of queries, packing same-round words, and using Merkle multiproofs can improve proof size independently of the algebraic primitive count.

This also explains why the protected-word count $s$ should never be converted directly into a byte estimate.  A protected word can be public, virtual, or already committed; only actual oracle roots and opened committed leaves enter \eqref{eq:proof-size-accounting}.

\subsection{Asymptotic comparison with Sumcheck}
\label{sec:complexity-vs-sumcheck}

The purpose of this paper is not to claim an asymptotic prover improvement over Sumcheck.  In fact, for several isolated relations the opposite comparison is immediate.

\begin{proposition}[Arithmetic comparison for an isolated degree-two sum]
\label{prop:ip-vs-sumcheck-complexity}
For
\[
  S=\sum_{w\in\B^n}a(w)b(w),
\]
a standard multilinear degree-two Sumcheck prover can be implemented in $O(N)$ field operations once the two length-$N$ evaluation tables are available, whereas the coefficient-extraction inner-product prover performs one degree-$<N$ polynomial product and one low-part commitment.  With FFT/NTT multiplication and encoding this is, in the generic model,
\begin{equation}
  O(N\log N)+\mathsf{Enc}_{N,M}
  \label{eq:ip-ce-generic-cost}
\end{equation}
before the shared folding backend.
\end{proposition}

\begin{proof}
The usual multilinear Sumcheck recurrence folds two length-$N$ tables once per Boolean coordinate; across all rounds the table lengths form $N+N/2+\cdots<2N$, giving linear arithmetic.  The coefficient-extraction protocol instead computes $V_aV_b^\star$ and commits to the resulting low part, as shown in \cref{sec:ip-costs}.  Applying \eqref{eq:standard-fft-bounds} yields \eqref{eq:ip-ce-generic-cost}.
\end{proof}

Thus the potential advantages of coefficient extraction are architectural rather than an automatic arithmetic speedup:
\begin{enumerate}[label=(\roman*),leftmargin=*]
  \item the same Reed--Solomon/FRI backend handles evaluation, inner products, Hadamard relations, lookup, permutation, sparse R1CS, and memory consistency;
  \item heterogeneous relations share one folding chain and one query phase;
  \item there is no recursive Sumcheck transcript whose algebraic rounds must be interleaved with a separate polynomial-commitment opening protocol;
  \item the operations are regular polynomial multiplication, scaling, encoding, hashing, and folding, which are natural targets for parallel and hardware acceleration;
  \item the cryptographic backend remains transparent and hash based.
\end{enumerate}
Whether these structural advantages compensate for extra polynomial products and auxiliary commitments is precisely what must be measured.

For a pointwise arithmetic circuit with $T_\times$ multiplication gates, the present construction can materialize one auxiliary table per multiplication gate and prove each multiplication by a Hadamard relation.  Without fusion this incurs $T_\times$ relation-specific products/commitments, whereas an optimized Sumcheck prover may traverse the original witness tables in $O(T_\times N)$ arithmetic.  Consequently, a benchmark that reports only verifier time or only the final folding time would be incomplete.

\subsection{Benchmark obligations}
\label{sec:complexity-benchmark-plan}

The Rust implementation will report the following quantities separately.

\begin{enumerate}[label=(\arabic*),leftmargin=*]
  \item \textbf{Front-end arithmetic.}  Time for table construction, inverses, coefficient scaling, polynomial products, synthetic divisions, and coefficient-side master batching.
  \item \textbf{Reed--Solomon encoding.}  Time and memory for every newly committed degree-$<N$ polynomial, reported both individually and after any fusion.
  \item \textbf{Merkle commitments.}  Hash time, root count, packed leaf width, and peak memory.
  \item \textbf{Folding backend.}  Time for the one shared master encoding/folding chain, final polynomial, and query generation.
  \item \textbf{Verifier.}  Field operations, inversions, hashes, bytes read, and total wall-clock time.
  \item \textbf{Proof size.}  Field payload, roots, salts, authentication hashes, and total serialized bytes, so that improvements from packing or multiproofs remain visible.
\end{enumerate}

The primary controlled comparison must use the same machine, compiler, field arithmetic, hash function, security target, and serialization policy.  In particular, the first baselines will be:
\begin{equation}
\begin{aligned}
  &\text{tree-based multilinear Sumcheck},\\
  &\text{Sumcheck + the same RS/Merkle opening backend}.
\end{aligned}
  \label{eq:benchmark-primary-baselines}
\end{equation}
so that the effect of replacing Sumcheck algebra by coefficient extraction is isolated from unrelated implementation choices.  External systems such as Spartan, HyperPlonk, BaseFold, DeepFold, WHIR, and Jolt/Lasso should be reported in a separate table unless their public implementations can be run under the same environment and security parameters.

\begin{remark}[No benchmark claim in the current manuscript]
\label{rem:no-benchmark-claim}
The theorems of this section are symbolic accounting statements.  They do not imply that coefficient extraction is faster, smaller, or more memory efficient than an optimized Sumcheck implementation.  The present mathematical contribution is that all relations considered in this paper can be reduced to one transparent Reed--Solomon/FRI proof architecture with explicit soundness and post-quantum compilation.  Performance superiority, if any, is an empirical claim and will be made only from reproducible measurements.
\end{remark}

\subsection{Summary of the implementation target}
\label{sec:complexity-summary}

The complete prover pipeline suggested by the theory is therefore
\begin{equation}
\begin{split}
  &\text{construct relation-specific coefficient tables and auxiliary words}
  \\
  &\quad\longrightarrow\
  \text{commit the semantic and auxiliary source words}
  \\
  &\quad\longrightarrow\
  \text{form all protected polynomials in coefficient form}
  \\
  &\quad\longrightarrow\
  \text{combine them into one master polynomial in $O(sN)$ work}
  \\
  &\quad\longrightarrow\
  \text{one RS encoding + one FRI-style folding chain + $\kappa$ queries}.
\end{split}
\label{eq:implementation-pipeline}
\end{equation}
The corresponding verifier performs relation-specific local computations only at the queried positions and then checks the common folding path.  The comparison and conclusion sections will use this accounting to position the construction against existing proof-system architectures and to state the remaining optimization problems precisely.

\section{Comparison with Existing Proof-System Architectures}
\label{sec:comparison}

This section positions the coefficient-extraction framework relative to the principal proof-system architectures that motivate it.  The comparison is deliberately structural.  The manuscript does not yet contain a complete Rust implementation of every relation, and therefore we do not compare wall-clock times obtained on different machines, with different fields, hash functions, security levels, or commitment layouts.  The controlled benchmarking methodology is fixed in \cref{sec:complexity-benchmark-plan}; numerical claims from other papers are used here only to identify the design points that those systems target.

The main distinction is the location of the algebraic reduction.  In a conventional multilinear architecture, a relation such as
\[
  \sum_{w\in\B^n} a(w)b(w)=S
\]
is represented directly as a multivariate claim and reduced by Sumcheck to one or more random evaluations.  A polynomial commitment scheme is then responsible for authenticating those terminal evaluations.  In the present framework, the relation is first rewritten as a coefficient or pointwise identity among univariate degree-$<N$ objects; the resulting words are batched before one Reed--Solomon proximity test.  Thus the two generic pipelines are
\begin{equation}
\begin{aligned}
  \text{multivariate relation}
  &\longrightarrow \text{Sumcheck}
   \longrightarrow \text{terminal evaluations}
  \\
  &\longrightarrow \text{PCS},
  \\[0.4ex]
  \text{multivariate relation}
  &\longrightarrow \text{coefficient/Hadamard reduction}
  \\
  &\longrightarrow \text{low-degree word family}
   \longrightarrow \text{RS proximity}.
\end{aligned}
\label{eq:comparison-two-pipelines}
\end{equation}
Neither pipeline is universally preferable.  Their arithmetic costs, communication patterns, and opportunities for batching are different, and \cref{rem:no-benchmark-claim} remains in force.

\subsection{Comparison axes}
\label{sec:comparison-axes}

To avoid conflating unrelated properties, we compare systems along the following axes.

\begin{enumerate}[label=(\roman*),leftmargin=*]
  \item \textbf{Algebraic front end.}  What transforms the original relation into low-degree claims: Sumcheck, a grand product, a logarithmic derivative, a coefficient identity, or a constrained-code condition?
  \item \textbf{Commitment/proximity backend.}  Is consistency enforced by a pairing-based polynomial commitment, a linear-code commitment, Reed--Solomon proximity, or another oracle proof?
  \item \textbf{Setup and post-quantum orientation.}  Is the backend transparent and hash based, or does it rely on a structured reference string and pairing/discrete-log assumptions?  We distinguish this cryptographic architecture from the information-theoretic security of the algebraic reduction itself.
  \item \textbf{Relation coverage.}  Is the primitive designed only for multilinear evaluation, or does the surrounding system also address R1CS, permutation, lookup, or memory consistency?
  \item \textbf{Composition.}  Are heterogeneous claims combined by repeated Sumchecks/openings, by a dedicated accumulation protocol, or by one low-degree batch?
  \item \textbf{Concrete cost.}  Prover time, verifier time, proof size, memory, and preprocessing must be compared only under compatible parameters or clearly identified as paper-reported external measurements.
\end{enumerate}

These axes explain why a comparison of polynomial-commitment proof sizes alone would be incomplete.  The present construction moves work from the recursive Sumcheck transcript into polynomial products, auxiliary low-degree words, and one common proximity backend.  A PCS-only benchmark therefore does not measure the intended R1CS/lookup/memory architecture.

\subsection{Sumcheck-centered multilinear systems}
\label{sec:comparison-sumcheck-systems}

\paragraph{Spartan.}
Spartan is a canonical example of a Sumcheck-centered R1CS proof system~\cite{Setty20}.  Its algebraic reduction exploits the multilinear extension of the R1CS matrices and witness, while a sparse-polynomial mechanism handles the public matrices.  The relation
\[
  (Az)\had(Bz)=Cz
\]
is reduced through multilinear polynomial identities and Sumcheck.  Our \cref{sec:sparse-matvec,sec:r1cs} targets the same algebraic bottleneck with a different decomposition:
\[
  \mathrm{R1CS}
  \longrightarrow
  3\times\mathrm{SparseMatVec}
  +1\times\mathrm{Hadamard}
  \longrightarrow
  \text{one heterogeneous RS batch}.
\]
The difference is therefore not the statement proved, but the mechanism that transports the statement to the commitment layer.  Spartan's recursive reduction keeps the witness in multilinear form; our construction materializes sparse-entry and Hadamard auxiliary tables and pays the corresponding commitment cost.

\paragraph{HyperPlonk.}
HyperPlonk uses multilinear polynomial IOPs to obtain linear-time prover algorithms for Plonkish constraints and high-degree custom gates~\cite{HyperPlonk23}.  Its zero checks and permutation checks are expressed as Sumcheck instances over virtual multilinear polynomials; in particular, its permutation argument rewrites the wiring relation as a zero relation and applies Sumcheck.  By contrast, \cref{sec:permutation} uses tagged multiset equality and logarithmic derivatives, and reduces the result to inverse Hadamard relations and linear functionals.  Both approaches exploit randomization to compress many coordinate constraints, but the terminal proof objects are different: random-point multilinear evaluations in HyperPlonk versus low-degree univariate words in the present framework.

\paragraph{Lasso and Jolt.}
Lasso develops lookup arguments whose efficiency comes from exploiting structured tables and sparse multilinear polynomials, with Sumcheck as a central proving primitive~\cite{Lasso23}.  Jolt uses lookup arguments as the main mechanism for proving virtual-machine instruction semantics~\cite{Jolt23}.  Our lookup protocol of \cref{sec:lookup} is closer algebraically to the logarithmic-derivative family than to Lasso's original sparse-table reduction: after a challenge $\beta$, inverse tables enforce
\[
  q\had(\beta\mathbf 1-f)=\mathbf 1,
  \qquad
  h\had(\beta\mathbf 1-t)=\mathbf 1,
\]
and one inner-product equality enforces multiplicities.  The distinction is again the backend.  The present protocol converts these relations to the common Hadamard/inner-product instruction set before the RS proximity phase instead of invoking an independent Sumcheck transcript.

\paragraph{Twist and Shout.}
Twist and Shout revisit lookup and memory checking with the explicit goal of reducing prover bottlenecks in zkVM-style workloads~\cite{TwistShout26}.  Their results are an important benchmark target because \cref{sec:memory} addresses the same high-level object: consistency of an execution-order memory trace with an address/time-sorted trace.  Our reduction
\[
  \mathrm{Memory}
  =
  \mathrm{TupleMultiset}
  +\mathrm{Lookup}
  +\mathrm{Hadamard}
  +\mathrm{LinearAdjacencyFingerprint}
\]
should therefore be viewed as an alternative algebraic architecture, not as a claim of lower prover cost.  The comparison that matters is the full end-to-end memory argument under the same field and security target.

\paragraph{BinarySpartan and Celer.}
Two 2026 systems sharpen the practical baselines for the front-end relations studied here.  BinarySpartan instantiates Spartan over binary fields with a transparent hash-based multilinear commitment and incorporates several recent Sumcheck optimizations~\cite{BinarySpartan26}.  It is therefore a particularly relevant controlled baseline for the claim that coefficient extraction can replace, rather than merely reimplement, the R1CS Sumcheck layer.  Celer targets lookup workloads in which the number of queries is much larger than the table and reduces commitment overhead while keeping field work linear in the query size~\cite{Celer26}.  Our logarithmic-derivative reduction has a different objective---unifying lookup with the same RS backend as the other relations---so Celer is a natural specialized lookup baseline rather than an architecture that the present theory subsumes.

\subsection{Code-based multilinear commitments and folding}
\label{sec:comparison-code-mlpcs}

The closest conceptual neighbors of this work are multilinear commitments that combine a code-based commitment with folding or Reed--Solomon proximity.  They share our goal of obtaining transparent, hash-based commitment layers without pairings, but differ in how the multilinear relation is linked to the code.

\paragraph{BaseFold.}
BaseFold formalizes foldable codes and interleaves its folding IOPP with the classical Sumcheck protocol~\cite{BaseFold24}.  Its construction uses a random evaluation produced by the folding process to answer the final Sumcheck claim, with shared verifier randomness.  This is almost the opposite interface from \eqref{eq:comparison-two-pipelines}: BaseFold synchronizes Sumcheck and folding, whereas coefficient extraction tries to eliminate the separate Sumcheck layer for the relations treated in this paper before folding starts.

The distinction is particularly visible for a committed inner product.  BaseFold can treat the claim through the multilinear Sumcheck/PCS interface.  Here it becomes
\[
  \ip{a}{b}
  =\coeff{N-1}{V_aV_b^\star},
\]
followed by the universal split identity and one proximity test.  Whether the extra polynomial multiplication is worthwhile is a concrete question, not an algebraic one; \cref{prop:ip-vs-sumcheck-complexity} records the generic arithmetic disadvantage of the isolated coefficient-extraction route.

\paragraph{DeepFold.}
DeepFold instantiates a Reed--Solomon version of the foldable-code approach and moves the analysis toward the list-decoding regime, reducing the query burden of the unique-decoding approach~\cite{DeepFold25}.  It is therefore especially relevant to the open parameter question in \cref{sec:complexity}: our current batching theorems use unique-decoding correlated agreement, so the query count can dominate proof size.  DeepFold demonstrates that list-decoding analysis is not merely an asymptotic refinement; it materially changes the useful concrete regime of an RS-based multilinear PCS.  A list-decoded version of the heterogeneous coefficient-extraction batching lemma would be complementary to, rather than a replacement for, the algebraic reductions developed here.

\paragraph{FRI-Binius.}
Diamond and Posen adapt the BaseFold paradigm to binary towers and relate the resulting folding to the additive transform of Lin--Chung--Han~\cite{FRIBinius24}.  This line shows that the multilinear/folding interface is not tied to large prime fields.  The present manuscript instead assumes the multiplicative-domain structure of \cref{sec:preliminaries} because inversion closure gives the virtual reversal
\[
  w^\star(\xi)=\xi^{N-1}w(\xi^{-1})
\]
for free.  Extending the coefficient-extraction instruction set to additive/binary domains would require a replacement for that virtual inversion symmetry; the coefficient identities themselves are field-independent.

\paragraph{WHIR.}
WHIR approaches multilinear commitments through constrained Reed--Solomon proximity rather than by simply attaching a standard FRI opening to a multilinear reduction~\cite{WHIR24}.  Its design is therefore the closest comparison at the level of philosophy: encode the evaluation relation into the code/proximity statement itself.  The difference is the kind of constraint.  WHIR develops a dedicated constrained-RS proximity test, whereas our framework first rewrites a family of relations into ordinary low-degree words and then deliberately reuses a generic RS folding test.  WHIR's specialization can give a substantially more query-efficient verifier; our intended advantage is relation uniformity---evaluation, inner product, Hadamard, lookup, permutation, sparse R1CS, and memory all terminate in the same ordinary low-degree interface.

\paragraph{TensorSwitch.}
TensorSwitch is a 2026 hash-based multilinear PCS built from tensor codes; its analysis explicitly uses list decoding and correlated agreement and targets both prover time and proof size~\cite{TensorSwitch26}.  It strengthens the evidence that the unique-decoding regime used in this manuscript is a conservative backend choice rather than an inherent requirement of transparent multilinear commitments.  TensorSwitch is not a drop-in replacement for the coefficient-extraction algebra developed here, but it is a direct benchmark for the commitment layer and a concrete motivation for the list-decoding open problem in \cref{sec:limitations}.

\subsection{Field-agnostic linear-code systems}
\label{sec:comparison-brakedown}

Brakedown gives a linear-time, field-agnostic SNARK/PCS architecture based on linear codes rather than FFT-friendly Reed--Solomon domains~\cite{Brakedown23}.  It illustrates another important point in the design space: a transparent post-quantum-oriented proof system need not use FRI or smooth multiplicative subgroups at all.  Our construction makes the opposite engineering choice.  Smooth RS domains are used aggressively because they provide
\begin{enumerate}[label=(\roman*),leftmargin=*]
  \item fast polynomial products and encodings;
  \item inversion closure for virtual reversal;
  \item a simple degree filtration under even--odd folding;
  \item direct compatibility with the Kronecker kernels of \cref{sec:coefficient-calculus}.
\end{enumerate}
This choice sacrifices field agnosticism.  Brakedown and related tensor/linear-code commitments are therefore not simply slower or faster versions of the same primitive; they optimize a different axis.


\paragraph{Neo and SuperNeo.}
Neo and SuperNeo provide a different post-quantum route: lattice-based folding for general constraint systems over small fields, with pay-per-bit commitments under Module-SIS assumptions~\cite{NeoSuperNeo26}.  They still invoke Sumcheck as part of the folding architecture.  They are therefore useful in this comparison because they separate two choices that are sometimes conflated: eliminating group-based commitments does not require eliminating Sumcheck, and eliminating Sumcheck from the algebraic front end does not require lattice commitments.  The present work chooses hash-based RS proximity; Neo/SuperNeo choose lattice-based folding.

\subsection{Pairing-based multilinear commitments}
\label{sec:comparison-pairing}

The middle-coefficient identity used throughout this paper is classical and is not specific to hash-based commitments.  Pairing-based multilinear PCS constructions provide an instructive contrast because they can exploit algebraic homomorphism and structured reference strings to obtain much shorter openings.

\paragraph{Gemini and Mercury.}
Gemini reduces multilinear evaluation to a sequence of univariate commitments and restriction relations, typically instantiated with KZG~\cite{Gemini22}.  Mercury starts from the same table/univariate viewpoint and develops a pairing-based multilinear PCS with constant-size evaluation proofs and linear field work in the opening procedure~\cite{Mercury25}.  These systems demonstrate that univariate representations of multilinear tables are already powerful even without RS proximity.

The present work differs cryptographically and compositionally.  The commitment is a Merkle commitment to an RS codeword rather than a homomorphic group commitment; reversal and quotient relations are read or derived as virtual words; and arbitrary relation families are folded into one proximity test.  The cost is that proof size is governed by hash openings and query complexity rather than a constant number of group elements.  Conversely, the resulting backend does not rely on pairing-friendly groups or a structured reference string.

\paragraph{Samaritan.}
Samaritan combines a logarithmic-derivative-based multilinear PIOP with a new pairing-based multilinear PCS to obtain an R1CS SNARK with linear-time prover and logarithmic proof/verification complexity~\cite{Samaritan25}.  This is especially relevant because both Samaritan and the present paper use logarithmic-derivative ideas to avoid less favorable lookup/permutation mechanisms.  The important difference is where the logarithmic-derivative relations terminate.  Samaritan compiles them through its multilinear PIOP/PCS stack; here they are decomposed into inverse Hadamard and inner-product relations and then absorbed into the heterogeneous RS batch.  Consequently the two constructions share some front-end algebra while making opposite commitment assumptions.

\subsection{Architecture matrix}
\label{sec:comparison-matrix}

\Cref{tab:architecture-comparison} summarizes the comparison at the level relevant to this paper.  ``Sumcheck role'' refers to the conventional recursive Sumcheck protocol in the algebraic front end; it does not imply anything about whether a system is post-quantum secure.  ``Backend'' describes the mechanism that ultimately authenticates low-degree objects.  Entries marked ``PCS-dependent'' inherit the setup/security properties of the chosen PCS instantiation.

\clearpage
\begin{table}[!ht]
\centering
\footnotesize
\begin{tabular}{@{}>{\raggedright\arraybackslash}p{2.35cm}>{\raggedright\arraybackslash}p{3.1cm}>{\raggedright\arraybackslash}p{4.55cm}>{\raggedright\arraybackslash}p{4.05cm}@{}}
\toprule
Architecture & Main target & Algebraic front end & Backend/setup \\
\midrule
Spartan & R1CS & multilinear R1CS identities and sparse matrix evaluation & PCS-dependent \\
HyperPlonk & Plonkish constraints and permutation & zero/permutation checks over virtual MLPs & PCS-dependent \\
Lasso/Jolt & lookup and VM execution & sparse-table lookup and multilinear reductions & PCS-dependent \\
Twist--Shout & lookup and memory & dedicated multilinear memory/lookup arguments & PCS-dependent \\
BaseFold & multilinear PCS & interleaved MLE Sumcheck and code folding & transparent foldable-code commitment \\
DeepFold & multilinear PCS & BaseFold-style reduction with RS list-decoding analysis & transparent RS/hash backend \\
FRI-Binius & multilinear PCS over binary towers & BaseFold-style multilinear reduction & binary-FRI/additive-code backend \\
WHIR & multilinear PCS & constrained Reed--Solomon relation & transparent constrained-RS/hash backend \\
Brakedown & R1CS/PCS & multilinear/tensor IOP & transparent linear-code backend \\
Mercury & multilinear PCS & univariate/table reduction & KZG, structured reference string, pairings \\
Samaritan & R1CS & LogSpartan and logarithmic derivatives & updatable SRS, pairing-based PCS \\
This work & MLE, IP, Hadamard, lookup, permutation, sparse R1CS, memory & coefficient extraction and pointwise reductions & transparent RS/Merkle/FRI-style folding \\
\bottomrule
\end{tabular}
\caption{Structural comparison of the algebraic front end and commitment/proximity backend.}
\label{tab:architecture-comparison}
\end{table}

\begin{table}[!ht]
\centering
\footnotesize
\begin{tabular}{@{}>{\raggedright\arraybackslash}p{2.65cm}>{\raggedright\arraybackslash}p{4.15cm}>{\raggedright\arraybackslash}p{7.15cm}@{}}
\toprule
Architecture & Conventional Sumcheck role & Relation to this work \\
\midrule
Spartan & central & same R1CS target; different sparse and multiplication reduction \\
HyperPlonk & central & alternative permutation/Hadamard reduction \\
Lasso/Jolt & central & lookup reduced instead to IP/Hadamard/RS relations \\
Twist--Shout & central & same memory/lookup bottleneck; different instruction set \\
BaseFold & explicit and interleaved with folding & closest folding baseline \\
DeepFold & explicit and interleaved with folding & motivates a list-decoded version of heterogeneous batching \\
FRI-Binius & explicit in the BaseFold-style reduction & suggests an additive/binary-domain extension path \\
WHIR & no conventional separate Sumcheck in the PCS & similar goal of moving an evaluation relation into the proximity statement \\
Brakedown & architecture-dependent & field-agnostic alternative to smooth RS domains \\
Mercury & no FRI/Sumcheck coupling in the PCS & related univariate/table viewpoint with a different cryptographic backend \\
Samaritan & used in its multilinear PIOP stack & related logarithmic-derivative algebra with different commitment assumptions \\
This work & removed for the relations covered here & one heterogeneous ordinary-RS batch \\
\bottomrule
\end{tabular}
\caption{Role of conventional Sumcheck and the specific point of comparison with the coefficient-extraction framework.  Neither table ranks the systems or combines performance measurements obtained under incompatible parameters.}
\label{tab:sumcheck-role-comparison}
\end{table}
\clearpage

\subsection{What is new, and what is classical}
\label{sec:comparison-novelty}

The elementary identity
\[
  \ip{a}{b}=\coeff{N-1}{V_aV_b^\star}
\]
is classical, as are polynomial reversal, random linear fingerprints, logarithmic-derivative multiset tests, DEEP quotients, Merkle commitments, and FRI-style proximity testing.  The claim of this paper is therefore not novelty of any one of those ingredients in isolation.

The framework contribution is the following composition of facts.
\begin{enumerate}[label=(\arabic*),leftmargin=*]
  \item A single table-form Kronecker commitment supports public multilinear evaluation and arbitrary committed inner products through the same coefficient functional.
  \item Inversion closure of the RS domain makes reversal of a committed table a \emph{virtual} codeword rather than another committed oracle.
  \item The split identity
  \[
    XUK=A+vX^N+X^{N+1}H
  \]
  converts the selected coefficient into a local identity in which the high part $H$ is also virtual.
  \item Geometric weighting and DEEP quotients lift the coefficient primitive from scalar inner products to pointwise Hadamard relations.
  \item Logarithmic-derivative lookup and multiset equality reduce to those same IP/Hadamard primitives.
  \item Sparse matrix multiplication, R1CS, and offline memory consistency reduce to the same instruction set.
  \item The protected low-degree words from all these heterogeneous relations are combined before proximity testing, so one RS folding chain and one final query phase certify the batch.
  \item The resulting IOR admits the relaxed round-by-round decoding state required for the post-quantum QROM compilation of \cref{sec:post-quantum}.
\end{enumerate}
The novelty question for the complete paper is therefore whether this \emph{systematic transparent composition}---rather than the middle-coefficient identity itself---has appeared previously.  The related systems above each contain nearby ingredients, but they organize the algebra/commitment boundary differently.

\subsection{Claims deliberately not made}
\label{sec:comparison-no-claims}

Several conclusions do \emph{not} follow from the current theory.

\begin{enumerate}[label=(\roman*),leftmargin=*]
  \item We do not claim that coefficient extraction asymptotically improves the prover complexity of Sumcheck.  \Cref{prop:ip-vs-sumcheck-complexity} shows that an isolated degree-two inner product is generically cheaper arithmetically with an optimized multilinear Sumcheck prover.
  \item We do not claim smaller proofs or faster verification than DeepFold, WHIR, pairing-based PCS constructions, or other optimized commitments.  The current unique-decoding backend can require many queries.
  \item We do not claim field agnosticism.  The virtual-reversal construction uses the chosen smooth inversion-closed multiplicative domain.
  \item We do not yet claim zero knowledge.  The present manuscript establishes knowledge soundness and transparent commitment consistency; witness hiding requires a separate masking layer compatible with all virtual relations.
  \item We do not identify ordinary interactive Sumcheck as a quantum-vulnerable primitive.  The post-quantum distinction concerns the cryptographic commitment and non-interactive compilation layers, as explained in \cref{sec:post-quantum}.
  \item We do not use published benchmark numbers from different systems as evidence that this construction is faster.  Such figures are contextual only until the controlled Rust study of \cref{sec:complexity-benchmark-plan} is complete.
\end{enumerate}

These limitations are useful because they isolate the empirical hypothesis of the project.  The mathematical result is a common coefficient-extraction/RS representation for a broad set of relations.  The implementation question is whether a fused implementation of those relations, with shared NTTs, packed commitments, one master encoding, and one folding/query phase, offers a useful performance point relative to mature Sumcheck-centered systems.

\subsection{Benchmark matrix for the companion implementation}
\label{sec:comparison-benchmark-matrix}

The companion repository will maintain two classes of baselines.

\paragraph{Controlled baselines.}
These are implemented or run under the same field, hash, compiler, CPU, security target, commitment layout, and serialization format as the coefficient-extraction prover.  At minimum we require:
\begin{align*}
  &\text{degree-two Sumcheck}
    &&\text{vs.}\quad \Pi_{\mathrm{IP}},\\
  &\text{pointwise multiplication via Sumcheck/standard MLP reduction}
    &&\text{vs.}\quad \Pi_{\mathrm{Had}},\\
  &\text{log-derivative lookup with Sumcheck}
    &&\text{vs.}\quad \Pi_{\mathrm{Lookup}},\\
  &\text{sparse R1CS with the tree-based Sumcheck baseline}
    &&\text{vs.}\quad \Pi_{\mathrm{R1CS}},\\
  &\text{offline memory consistency baseline}
    &&\text{vs.}\quad \Pi_{\mathrm{Mem}}.
\end{align*}
Every comparison reports front-end arithmetic, encoding, hashing, folding, verifier work, proof size, and peak memory separately.

\paragraph{External reproducible baselines.}
When public implementations of BaseFold, DeepFold, WHIR, TensorSwitch, Spartan/BinarySpartan, HyperPlonk, Lasso/Jolt, Celer, Twist--Shout, Brakedown, Mercury, Samaritan, or Neo/SuperNeo can be built under compatible parameters, the repository will record the exact revision and command line and run them on the same machine.  Otherwise, their paper-reported measurements will appear in a separate table explicitly labelled \emph{reported by the authors}; they will not be mixed into speedup ratios against our own measurements.

\begin{remark}[Comparison criterion]
\label{rem:comparison-criterion}
The experiment that most directly tests the thesis of this paper is not a standalone PCS opening.  It is a heterogeneous workload in which evaluation, multiplication, lookup/permutation, and sparse constraints are all required simultaneously.  In that setting the coefficient-extraction architecture pays multiple relation-specific auxiliary constructions but amortizes the low-degree backend over the complete batch.  A fair benchmark must therefore report both sides of that tradeoff.
\end{remark}

\subsection{Summary}
\label{sec:comparison-summary}

Existing systems demonstrate several successful interfaces between multilinear algebra and cryptographic commitment: recursive Sumcheck followed by a PCS, Sumcheck interleaved with foldable-code proximity, constrained Reed--Solomon testing, field-agnostic linear-code commitments, and homomorphic pairing-based multilinear commitments.  The interface studied here is different:
\begin{equation}
  \boxed{
  \begin{gathered}
    \text{relation}
    \longrightarrow
    \text{coefficient/Hadamard instruction set}
    \\
    \longrightarrow
    \text{one ordinary RS low-degree batch}
    \longrightarrow
    \text{FRI-style proximity}
  \end{gathered}}
  \label{eq:comparison-framework-summary}
\end{equation}
The next section states the limitations and open problems that determine whether this architectural simplification can become a competitive proof system in practice.

\section{Limitations, Open Problems, and Research Agenda}
\label{sec:limitations}

The preceding sections establish a common coefficient-extraction framework for public linear functionals, committed inner products, pointwise Hadamard relations, logarithmic-derivative lookups and multiset checks, sparse matrix--vector multiplication, R1CS, and an offline memory-consistency relation.  The purpose of this section is to state precisely what these results do \emph{not} yet establish, and to isolate the mathematical and systems questions that determine whether the framework can become a competitive general-purpose proof architecture.

The central distinction is between three layers:
\begin{equation}
  \boxed{
  \begin{gathered}
    \text{algebraic expressibility}
    \quad\longrightarrow\quad
    \text{interactive proof efficiency}
    \\
    \quad\longrightarrow\quad
    \text{non-interactive cryptographic efficiency}.
  \end{gathered}}
  \label{eq:limitations-three-layers}
\end{equation}
The present manuscript resolves the first layer for the relations above and proves soundness and decoding knowledge for the resulting interactive protocols in the unique-decoding regime.  It also identifies a post-quantum QROM compilation route in \cref{sec:post-quantum}.  It does not yet establish that the complete construction is faster, smaller, zero knowledge, recursive, or optimal relative to mature Sumcheck-centered systems.

\subsection{The unique-decoding restriction}
\label{sec:limitations-unique-decoding}

All correlated-agreement arguments in this paper use the unique-decoding form of the Reed--Solomon theorem quoted in \cref{thm:correlated-agreement}, with
\[
  0<\delta\le \delta_\star\defeq\frac{1-\rho}{2}.
\]
This restriction propagates to the generic coefficient-extraction argument, the heterogeneous batch, and every application built from them.  In particular, the final query term has the form
\[
  (1-\delta)^\kappa,
\]
and the maximal admissible $\delta$ is only half the relative minimum-distance bound $1-\rho$.

At rate $\rho=1/4$, the largest unique-decoding radius is $\delta_\star=3/8$.  Thus obtaining a target query error $2^{-\lambda}$ requires
\begin{equation}
  \kappa
  \ge
  \left\lceil
    \frac{\lambda}{-\log_2(1-\delta_\star)}
  \right\rceil.
  \label{eq:limitations-query-count}
\end{equation}
The large number of authenticated queries is therefore not intrinsic to coefficient extraction; it is a consequence of the decoding regime used by the current proof.

\paragraph{Open problem 1: list-decoded heterogeneous batching.}
The first mathematical priority is to replace the unique-decoding use of correlated agreement by a list-decoding statement, ideally up to a Johnson-type radius.  The desired theorem has the following form.  Let
\[
  W_\beta=\sum_{i=0}^{t}\beta^i u_i
\]
be the master batching curve.  If $W_\beta$ is close to the target RS code for sufficiently many $\beta$, then the words $u_i$ should admit a \emph{small common list} of low-degree explanations on a large agreement set.  The protocol soundness proof must then show that the algebraic identities and the decoded commitment words isolate one consistent member of this list.

The obstruction is not the coefficient identity itself.  The identities in \cref{sec:coefficient-calculus} remain valid independently of the decoding radius.  The missing ingredient is a list-valued analogue of the common-agreement decoding state used in \cref{sec:heterogeneous-batching,sec:soundness}.  A successful theorem would reduce $\kappa$ without changing the algebraic front end.

\subsection{Rate and degree budget}
\label{sec:limitations-rate}

The universal split identity
\[
  XU(X)K(X)=A(X)+vX^N+X^{N+1}H(X)
\]
uses polynomials $U,K,A,H$ of degree $<N$.  Consequently the defect polynomial used in the soundness proof has degree at most $2N$.  To conclude that it is identically zero from its vanishing on an agreement set, the present analysis requires enough evaluation points beyond this $2N$ degree budget.  This is one reason the manuscript uses a constant-rate RS code with substantial redundancy.

\paragraph{Open problem 2: higher-rate coefficient extraction.}
A stronger opening identity could lower the degree of the defect polynomial or split the product into several lower-degree pieces.  Either improvement would permit a higher code rate, potentially approaching $1/2$, while preserving the same selected coefficient.  A useful target is an identity whose soundness polynomial has degree $N+O(1)$ rather than $2N+O(1)$.

This question is distinct from list decoding.  List decoding improves how much corruption the proximity proof tolerates; a higher-rate identity reduces the encoded domain size $M$ itself.  The two improvements are complementary.

\subsection{Higher-degree sums and auxiliary tables}
\label{sec:limitations-higher-degree}

For a degree-two Boolean sum,
\[
  S=\sum_{w\in\B^n}a(w)b(w),
\]
coefficient extraction gives a direct one-product identity.  For higher pointwise degree, the present framework introduces auxiliary tables and proves their multiplication gates with Hadamard arguments.  For example,
\[
  \sum_w a(w)b(w)c(w)
\]
is handled by committing to $p=a\had b$ and then proving $\ip{p}{c}$.

This reduction is general but may materialize many intermediate vectors.  If a pointwise arithmetic circuit has $g_\times$ multiplication gates, a direct vectorized realization creates up to $g_\times$ auxiliary degree-$<N$ tables before batching.

\paragraph{Open problem 3: direct multilinear coefficient identities for higher arity.}
One can express diagonal tensor contractions as selected coefficients of multivariate products, but naive Kronecker substitution can increase the univariate degree from $O(N)$ to $\Omega(N^2)$.  The open problem is to find a family of structured encodings for
\[
  \sum_w\prod_{j=1}^{d} f_j(w)
\]
that keeps both the degree and the number of committed auxiliary words linear in $N$ for fixed $d$, while retaining local verifier access compatible with an RS/FRI backend.  Such an identity could remove entire layers of Hadamard auxiliary tables.

\subsection{Sparse matrix preprocessing}
\label{sec:limitations-sparse-preprocessing}

The sparse matrix protocol of \cref{sec:sparse-matvec} represents a public matrix by its nonzero-entry stream
\[
  (\mathsf{row}_e,\mathsf{col}_e,\mathsf{val}_e)_{e<K},
  \qquad K=\operatorname{nnz}(M),
\]
and reduces matrix--vector multiplication to indexed selection, a Hadamard multiplication, and aggregation.  This preserves work proportional to the sparse representation rather than to $N^2$.

The construction nevertheless assumes that the sparse entry representation and the public indexing data are available in a form that can be preprocessed consistently.  The theorem does not by itself optimize repeated proofs for the same matrix, nor does it establish the best commitment layout for very large static R1CS matrices.

\paragraph{Open problem 4: direct sparse kernels.}
For fixed public $M$, define the random row-compression vector
\[
  p_r(j)=\sum_i M_{ij}\eq(i,r).
\]
Then
\[
  \widetilde y(r)=\ip{p_r}{z}
\]
for $y=Mz$.  If $p_r$ had a succinct generating kernel that could be evaluated without materializing all $N$ coordinates, sparse matrix multiplication would reduce directly to the inner-product primitive.  The question is therefore whether common sparse constraint matrices admit structured factorizations, low-rank decompositions, or preprocessing data from which
\[
  K_{p_r}(X)=\sum_j p_r(j)X^{N-1-j}
\]
can be evaluated in $\operatorname{polylog}N$ or in time proportional only to the local sparsity touched by a verifier query.

\subsection{Lookup and permutation costs}
\label{sec:limitations-lookup}

The lookup and permutation arguments of \cref{sec:lookup,sec:permutation} are algebraically compact: they replace grand products by logarithmic derivatives and then reduce the resulting equations to inverse Hadamard relations and linear functionals.  Their efficiency, however, depends on committed inverse tables such as
\[
  q_i=(\beta-f_i)^{-1}.
\]
Constructing these tables costs one inversion per coordinate naively, or one batched inversion plus linear additional arithmetic.  Moreover, each inverse table participates as another protected low-degree word in the heterogeneous batch.

\paragraph{Open problem 5: fused inverse-table construction.}
The implementation should determine whether several inverse tables can share denominator products, batch inversions, encoding passes, and Merkle leaves.  The relevant comparison is not only field-operation count but memory traffic and hashing.  A mathematical refinement would avoid explicit inverse tables altogether by proving a denominator-cleared rational identity with bounded degree and no poles; whether such a formulation preserves the simple one-batch soundness proof is open.

\subsection{Proof size and authenticated-query cost}
\label{sec:limitations-proof-size}

The heterogeneous batch saves the low-degree test: once the master word is formed, the folding chain and final query phase are paid once.  It does \emph{not} make authentication paths free.  A query may still need values from several underlying committed words, and those values must be authenticated unless they can be derived virtually from already opened data.

The proof size therefore depends on three independent quantities:
\begin{enumerate}[label=(\roman*),leftmargin=*]
  \item the number $\kappa$ of query points;
  \item the number of physically committed words whose values are required at each query;
  \item the amount of Merkle-path sharing provided by the leaf layout, caps, and multiproofs.
\end{enumerate}
The theory of \cref{sec:heterogeneous-batching} controls the first algebraic combination but does not optimize the latter two representation choices.

\paragraph{Open problem 6: commitment packing and shared authentication.}
The implementation should place vectors that are commonly opened together in the same leaf and exploit multiproofs across common paths.  For R1CS and memory checking, this may be as important as reducing field multiplications.  A complete cost model should optimize the leaf arity, Merkle cap size, folding arity, and query count jointly.

\subsection{Zero knowledge}
\label{sec:limitations-zk}

The present protocols are arguments of knowledge; they are not zero-knowledge protocols.  Table commitments and query openings may reveal witness-dependent field values.  The post-quantum compilation of \cref{sec:post-quantum} protects soundness and extraction in the QROM but does not provide witness privacy.

\paragraph{Open problem 7: masking compatible with virtual relations.}
A zero-knowledge transformation must preserve all identities that are evaluated locally by the verifier.  In particular, randomization of $V_a$ must interact correctly with reversal
\[
  V_a^*(X)=X^{N-1}V_a(X^{-1}),
\]
with DEEP quotient words, with Hadamard products, and with logarithmic-derivative inverse constraints.  Independent masks generally destroy multiplicative relations.  A useful construction should provide correlated masks whose induced error terms are either public, cancel algebraically, or are themselves included in the heterogeneous batch.

For a full SNARK application, this is a necessary cryptographic layer rather than an optional optimization.

\subsection{Concrete post-quantum security}
\label{sec:limitations-pq-concrete}

The QROM result of \cref{sec:post-quantum} derives security from the round-by-round decoding state and the compiler of Chiesa, Di, Hu, and Zheng~\cite{CDHZ25}.  The compiler is general, and its concrete extraction bound contains substantial losses in the adversary's number of quantum random-oracle queries.  Thus an interactive error near $2^{-128}$ does not automatically imply a $128$-bit bound for the compiled non-interactive argument.

\paragraph{Open problem 8: parameter optimization for a target QROM level.}
For a chosen quantum-query budget $Q$ and target extraction error $2^{-\lambda}$, one must optimize jointly
\[
  |\F|,\qquad \rho,\qquad \delta,\qquad \kappa,\qquad
  \lambda_{\rm ch},\qquad \lambda_H,
\]
together with the number of protocol rounds and oracle positions read.  The objective should be a \emph{compiled} security target rather than an interactive-only target.  A list-decoding improvement would help here twice: by reducing $\kappa$ and by decreasing the round-by-round proximity error that is amplified by the compiler.

\subsection{Domain and field restrictions}
\label{sec:limitations-domain}

Virtual reversal relies on
\[
  \xi\in L\Longrightarrow \xi^{-1}\in L,
\]
and the folding backend requires a smooth hierarchy of evaluation domains.  These conditions are natural for multiplicative-subgroup RS codes but are not available in every field or every code family.

The framework is therefore not field agnostic in its present form.  This limitation is architectural rather than merely implementational: the local formula
\[
  w^*(\xi)=\xi^{N-1}w(\xi^{-1})
\]
would cease to be an oracle-local transformation if $\xi^{-1}$ left the committed domain.

\paragraph{Open problem 9: coefficient extraction over other foldable domains.}
Possible directions include additive domains, affine/circle domains, binary towers, or general foldable linear codes.  The invariant to preserve is not necessarily multiplicative inversion itself, but an efficiently computable domain automorphism that realizes coefficient reversal or the appropriate adjoint transformation locally on committed words.

\subsection{Recursion and proof aggregation}
\label{sec:limitations-recursion}

Nothing in the current manuscript proves that the verifier can itself be represented cheaply inside the same constraint system.  A recursive implementation must arithmetize hashing, Merkle verification, extension-field operations, folding checks, and any DEEP inversions.  A protocol with an efficient native verifier may still be expensive recursively.

\paragraph{Open problem 10: recursion-friendly instantiation.}
The correct metric is the number and type of constraints needed to verify one proof inside a target field or VM.  This may favor a different hash, field tower, folding arity, or commitment packing than the native benchmark.  Recursive proof aggregation should therefore be treated as a separate parameter regime, not inferred from native verifier timings.

\subsection{Memory model limitations}
\label{sec:limitations-memory-model}

The memory argument of \cref{sec:memory} is an \emph{offline} consistency proof: the complete access trace is available, a sorted copy is authenticated, and consistency is reduced to permutation, range/ordering checks, and local state-transition relations.  This model covers an important class of zkVM memory arguments but is not an online authenticated data structure.

The current theorem also assumes bounded encodings for addresses and timestamps suitable for the range tables used in the lookup layer.  Very large address spaces, wraparound semantics, concurrency, atomic operations, or nondeterministic I/O require additional relations.

\paragraph{Open problem 11: streaming and stateful memory.}
A natural extension is an incremental version in which a commitment to memory state is carried between proof segments and only a bounded access batch is processed at once.  Such a construction would need a composable state commitment and a proof that segment boundaries preserve the same memory semantics as the global sorted-trace theorem.

\subsection{Backend modularity}
\label{sec:limitations-backend}

Coefficient extraction produces ordinary low-degree RS relations before the proximity test.  The paper instantiates the backend with the FRI-style folding test analysed in \cref{sec:generic-protocol}, but the algebraic front end is conceptually separable from that choice.

\paragraph{Open problem 12: stronger proximity backends.}
It is natural to ask whether the same master low-degree word can be certified by a more query-efficient RS proximity system, including list-decoding-oriented constructions such as WHIR~\cite{WHIR24} or other constrained-code techniques.  A replacement backend must support the commitment/oracle interface and the round-by-round knowledge state needed by the non-interactive compiler.  If it does, improvements in proximity testing could be inherited without changing the coefficient-extraction reductions.

\subsection{The principal empirical hypothesis}
\label{sec:limitations-empirical}

The paper's algebraic result does not imply an implementation speedup.  The empirical hypothesis is narrower:

\begin{quote}
When a workload simultaneously requires evaluations, inner products, multiplication constraints, lookups or permutations, and sparse linear relations, the cost of constructing relation-specific auxiliary words may be offset by fusing their encodings, commitments, and one shared low-degree proximity proof.
\end{quote}

The controlled benchmark plan of \cref{sec:complexity-benchmark-plan} must test this hypothesis by decomposing prover time into
\[
  \text{front-end arithmetic}
  +\text{encoding}
  +\text{hashing}
  +\text{folding}
  +\text{serialization},
\]
and by reporting peak memory, verifier time, and proof bytes separately.  Failure to beat a specialized Sumcheck implementation on one isolated primitive would not refute the architecture; conversely, a faster inner-product microbenchmark would not establish superiority on R1CS or memory checking.

\subsection{A prioritized research program}
\label{sec:limitations-roadmap}

The remaining work can be ordered by dependency rather than by apparent novelty.

\begin{enumerate}[label=\textbf{R\arabic*.},leftmargin=*]
  \item \textbf{Proof audit and formalization.}
  Formalize the algebraic reductions and the heterogeneous-batching soundness proof in Lean, preserving the theorem boundaries of this manuscript.  This is the prerequisite for treating the framework as a reusable proof-system layer rather than a collection of informal reductions.

  \item \textbf{Reference implementation and controlled baselines.}
  Implement one common RS/FRI backend and the IP, Hadamard, lookup, permutation, sparse-matrix, R1CS, and memory front ends.  Benchmark them against the tree-based Sumcheck baseline under identical field, hash, hardware, compiler, and security settings.

  \item \textbf{List-decoding soundness.}
  Generalize the batching theorem beyond unique decoding.  This is the highest-leverage mathematical improvement for proof size and verifier cost.

  \item \textbf{Packed commitment engineering.}
  Optimize leaf layout, Merkle multiproofs, NTT sharing, batching of inversions, and construction of the master word.  These optimizations determine whether heterogeneous batching is practically useful.

  \item \textbf{Zero knowledge.}
  Add a masking layer compatible with reversal, DEEP quotients, Hadamard products, and logarithmic derivatives, and prove simulation rather than only knowledge soundness.

  \item \textbf{Compiled PQ parameters.}
  Choose parameter sets from the QROM extraction bound rather than from interactive soundness alone, and publish executable scripts that reproduce every security number.

  \item \textbf{Higher-rate and direct higher-degree identities.}
  Explore improved split identities, structured multi-product kernels, and direct sparse kernels.  These are the directions most likely to change asymptotic or large-instance constants rather than only implementation overhead.

  \item \textbf{Recursion and alternative backends.}
  Study recursion-friendly hash/field choices and whether the same algebraic front end can profit from more query-efficient proximity tests.
\end{enumerate}

\subsection{Summary}
\label{sec:limitations-summary}

The framework developed in this paper establishes that a broad family of relations normally presented through Sumcheck can instead be expressed as coefficient, Hadamard, and logarithmic-derivative identities over one table-form Kronecker representation, and that these identities can be combined before one RS proximity proof.  The remaining questions are no longer primarily about expressibility.  They concern decoding radius, degree/rate tradeoffs, auxiliary-word cost, privacy, concrete QROM parameters, and implementation efficiency.

The most important conceptual boundary is therefore:
\begin{equation}
  \boxed{
  \begin{gathered}
    \text{proved here: a common transparent algebraic/proximity architecture},\\
    \text{not yet proved: that this architecture dominates optimized alternatives in practice.}
  \end{gathered}}
  \label{eq:limitations-final-boundary}
\end{equation}
The next section concludes by collecting the algebraic principle and the resulting proof-system architecture in its shortest form.

\section{Conclusion}
\label{sec:conclusion}

This work develops coefficient extraction as a common algebraic interface between multilinear proof relations and Reed--Solomon proximity testing.  The starting identity is elementary: for vectors $a,b\in\F^N$,
\[
  \ip{a}{b}
  =
  [X^{N-1}]\,V_a(X)V_b^\star(X),
\]
where $V_a(X)=\sum_{j=0}^{N-1}a_jX^j$ and $V_b^\star(X)=\sum_{j=0}^{N-1}b_jX^{N-1-j}$.  The contribution of the framework is to turn this coefficient identity into a reusable proof-system primitive.  The selected coefficient is isolated by the split
\begin{equation}
  XU(X)K(X)
  =
  A(X)+vX^N+X^{N+1}H(X),
  \label{eq:conclusion-master-identity}
\end{equation}
with all protected objects represented as degree-$<N$ Reed--Solomon words.  Reversal, high-part recovery, and DEEP quotient relations are exposed as virtual or locally derived words, and random linear batching reduces the resulting family of low-degree claims to a single proximity instance.

The first consequence is a uniform treatment of public linear functionals.  Table-form multilinear evaluation, Boolean-hypercube summation, and more general succinct public kernels all become instances of the same coefficient-extraction protocol.  Reversal then removes the restriction that the second vector be public: an arbitrary committed inner product is represented by the middle coefficient of $V_aV_b^\star$, with the reversed word obtained from the original commitment by the inversion map on the evaluation domain.  Thus the canonical degree-two Boolean sum
\[
  \sum_{w\in\B^n} a(w)b(w)
\]
can be certified without an $n$-round Sumcheck reduction.

The second step is pointwise multiplication.  The Hadamard relation
\[
  a\had b=c
\]
is compressed by a random geometric fingerprint and rewritten as a coefficient claim involving $V_a(\gamma X)V_b^\star(X)$, while DEEP quotient words bind the scaled and out-of-domain values to the original commitments.  This supplies the multiplication gate of the framework.  Together, linear relations and Hadamard relations are sufficient to represent bounded pointwise arithmetic computations on committed tables; higher-degree Boolean sums can therefore be reduced to auxiliary multiplication tables followed by a final linear functional.

The same instruction set extends beyond polynomial evaluation.  Logarithmic-derivative lookup reduces to reciprocal Hadamard relations and an inner-product equality.  Multiset equality and tagged permutation checks use the same reciprocal structure.  Sparse matrix--vector multiplication is represented by an $O(\operatorname{nnz}(M))$ sparse-entry stream, indexed selection, pointwise multiplication, and linear aggregation.  Consequently an R1CS relation
\[
  (Az)\had(Bz)=Cz
\]
reduces to three sparse matrix--vector relations and one Hadamard relation.  An offline memory-consistency argument further decomposes into tuple-multiset consistency, lookups, local Hadamard transition constraints, and linear adjacency fingerprints.  These reductions show that the coefficient-extraction view is not confined to a multilinear polynomial commitment: it defines a small algebraic proof language whose basic operations already cover several central components of modern proof systems.

A central technical point is that these heterogeneous relations need not incur independent proximity proofs.  Once each relation has been expressed as a family of degree-bounded words, the words are protected by one random batching curve and one folding chain.  The global analysis of \cref{sec:soundness} separates the error into an outer algebraic-reduction term and a common Reed--Solomon/proximity term.  In the notation of that section, the soundness error is bounded by
\[
  \varepsilon_{\mathrm{tot}}
  \le
  \varepsilon_{\mathrm{red}}
  +h\frac{2(N-1)}{|\F|}
  +\frac{(s-1)M}{|\F|}
  +\varepsilon_{\mathrm{fold}}
  +(1-\delta)^\kappa.
\]
This factorization is conceptually useful: lookup, permutation, sparse algebra, R1CS, and memory checking contribute relation-specific algebraic errors, while the low-degree certification is shared.

The post-quantum interpretation should be stated carefully.  Sumcheck is an information-theoretic interactive protocol and is not, by itself, a non-post-quantum primitive.  The result here is instead a transformation of the \emph{algebraic role} played by Sumcheck for the relations studied in this paper.  The recursive hypercube reduction is replaced by coefficient identities and pointwise constraints that feed directly into a transparent Reed--Solomon/FRI layer.  The non-interactive construction of \cref{sec:post-quantum} then commits to oracle words with salted Merkle trees and uses a QROM-secure IOR-to-NIR compiler.  Under the stated round-by-round decoding-knowledge conditions, this yields post-quantum straight-line knowledge soundness for the compiled framework.  Thus the architecture separates three concerns:
\[
  \boxed{
  \text{relation algebra}
  \;\longrightarrow\;
  \text{RS proximity}
  \;\longrightarrow\;
  \text{hash-based non-interactive compilation}.}
\]

The present results do not establish that this architecture is universally preferable to Sumcheck.  Sumcheck remains a direct and highly optimized method for general low-degree hypercube sums.  Coefficient extraction trades recursive reduction for univariate encoding, auxiliary tables, and proximity testing.  The practical question is therefore whether shared encoding, virtual words, and one heterogeneous proximity proof compensate for those costs on workloads containing several relation types at once.  \Cref{sec:complexity} makes the accounting explicit, and \cref{sec:comparison} separates structural comparisons from performance claims that require implementation.  The companion Rust implementation and controlled baselines are intended to answer this empirical question under identical fields, security parameters, commitment layouts, hardware, and compiler settings.

Several mathematical improvements could materially change that tradeoff.  The current soundness analysis is confined to unique decoding; extending heterogeneous batching to a list-decoding regime would reduce the number of authenticated queries.  A higher-rate split identity could reduce the Reed--Solomon domain size.  Direct higher-arity coefficient identities might eliminate intermediate Hadamard tables, while direct sparse kernels could reduce the auxiliary structure needed for matrix relations.  Zero knowledge, recursion-friendly instantiations, and tighter concrete QROM parameters remain separate design problems.  These limitations are isolated in \cref{sec:limitations} so that improvements to the algebraic front end, proximity backend, and cryptographic compiler can be studied independently.

The broader conclusion is that the classical middle-coefficient identity becomes significantly more expressive when it is combined with Kronecker encodings and code-based proximity.  Structured public kernels turn evaluation into coefficient extraction; reversal turns committed vectors into compatible kernels; geometric fingerprints turn pointwise multiplication into one coefficient claim; reciprocal tables turn lookup and permutation into Hadamard and inner-product relations; and sparse streams extend the same language to R1CS and memory consistency.  FRI is then used not to reproduce Sumcheck round by round, but as the common certification mechanism for the low-degree objects produced by these reductions.

This suggests a proof-system design principle that is independent of any one application:
\begin{quote}
First search for a coefficient or local polynomial identity that exposes the target relation as a small family of low-degree univariate objects.  Then exploit virtual transformations and random batching so that those objects can be certified by one transparent proximity argument.
\end{quote}
The framework developed in this paper gives a concrete realization of that principle for linear, bilinear, lookup, permutation, sparse-linear, R1CS, and memory relations.  Its next stage is no longer to establish algebraic expressibility, but to determine---through formal verification and controlled implementation---how far this common coefficient-extraction layer can simplify and accelerate complete transparent proof systems.


% The abstract will be audited once the complete theorem/reference audit is finished.

\printbibliography

\end{document}
